\documentclass[10pt]{article}
\usepackage[a4paper,landscape,margin=1.5cm]{geometry}
\usepackage{amsmath,amssymb,longtable,array}
\renewcommand{\thetable}{6.\arabic{table}}
\newcommand{\possprim}{\textbf{P}}
\newcommand{\nongcr}{\textbf{N}}
\newcommand{\newrow}{\textbf{new}}
\newcommand{\Alt}{\mathrm{Alt}}
\setlength{\LTleft}{0pt}
\begin{document}
\title{Feasible characters: updated tables}
\author{Recomputed with \texttt{feasgen.g} (FeasChar repository)}
\date{}
\maketitle
\noindent Table numbers and captions are those of Chapter 6 of A.~J.~Litterick, \emph{On non-generic
finite subgroups of exceptional algebraic groups}, Mem.\ Amer.\ Math.\ Soc.\ 1207 (2018).

\noindent The rows of each table are the feasible characters compatible with the adjoint module, all
end-node Weyl modules and, in characteristic $p$, the irreducible modules $L(\lambda)$ that are
smaller than the corresponding Weyl modules. All element orders are used, and the class
assignment must be compatible with every power map. Only $L(G)$ and $V_{\min}$ are shown, as in
the Memoir.

\noindent Rows are listed up to automorphisms of the character table of $H$ (and, for $F_4$
with $p = 2$, up to the graph automorphism of $F_4$). Columns are labelled by degree, with letters
in the order of the irreducible (Brauer) characters in GAP's character table library. These
labels need not agree with the Memoir's.

\noindent \emph{Provenance.} These tables were computed and typeset by code written by an AI
system (Claude, Anthropic), October 2026, under the direction of the Memoir's author. They are
computational results and have not been independently checked by hand. See the README.

\noindent \possprim, \nongcr: as in the Memoir, carried over from the Memoir row with the same
composition factors. \newrow: not in the Memoir. Differences from the Memoir are listed in
\texttt{DIFFERENCES.md}.
\newpage
% Updated tables for F4 (Memoir Chapter 6), generated by tools/fg_tables.g
\section*{$F_4$: feasible characters on $L(G) = V_{52}$ and $V_{\min} = V_{26}$}

\setcounter{table}{0}
\begin{longtable}{rc|cccc|ccc}
\caption{Alt$_{10} < F_4$, $p = 2$}\\
& & \multicolumn{4}{c|}{$V_{52}$} & \multicolumn{3}{c}{$V_{26}$} \\
& & $1$ & $8$ & $16$ & $26$ & $1$ & $8$ & $16$ \\ \hline
\endfirsthead
& & \multicolumn{4}{c|}{$V_{52}$} & \multicolumn{3}{c}{$V_{26}$} \\
& & $1$ & $8$ & $16$ & $26$ & $1$ & $8$ & $16$ \\ \hline
\endhead
1) &  & 2 & 1 & 1 & 1 & 2 & 1 & 1 \\
\end{longtable}
\noindent Modules tested: V52, L26, V26.\\

\setcounter{table}{1}
\begin{longtable}{rc|ccccc|cccc}
\caption{Alt$_{9} < F_4$, $p = 2$}\\
& & \multicolumn{5}{c|}{$V_{52}$} & \multicolumn{4}{c}{$V_{26}$} \\
& & $1$ & $8_{a}$ & $8_{b}$ & $8_{c}$ & $26$ & $1$ & $8_{a}$ & $8_{b}$ & $8_{c}$ \\ \hline
\endfirsthead
& & \multicolumn{5}{c|}{$V_{52}$} & \multicolumn{4}{c}{$V_{26}$} \\
& & $1$ & $8_{a}$ & $8_{b}$ & $8_{c}$ & $26$ & $1$ & $8_{a}$ & $8_{b}$ & $8_{c}$ \\ \hline
\endhead
1) &  & 2 & 1 & 1 & 1 & 1 & 2 & 1 & 1 & 1 \\
\end{longtable}
\noindent Modules tested: V52, L26, V26.\\

\setcounter{table}{2}
\begin{longtable}{rc|ccc|cc}
\caption{Alt$_{7} < F_4$, $p = 5$}\\
& & \multicolumn{3}{c|}{$V_{52}$} & \multicolumn{2}{c}{$V_{26}$} \\
& & $8$ & $10_{a}$ & $10_{b}$ & $1$ & $8$ \\ \hline
\endfirsthead
& & \multicolumn{3}{c|}{$V_{52}$} & \multicolumn{2}{c}{$V_{26}$} \\
& & $8$ & $10_{a}$ & $10_{b}$ & $1$ & $8$ \\ \hline
\endhead
1) &  & 4 & 1 & 1 & 2 & 3 \\
\end{longtable}
\noindent Modules tested: V52, V26.\\

\setcounter{table}{3}
\begin{longtable}{rc|cccccc|cccc}
\caption{Alt$_{7} < F_4$, $p = 2$}\\
& & \multicolumn{6}{c|}{$V_{52}$} & \multicolumn{4}{c}{$V_{26}$} \\
& & $1$ & $4_{a}$ & $4_{b}$ & $6$ & $14$ & $20$ & $1$ & $4_{a}$ & $4_{b}$ & $6$ \\ \hline
\endfirsthead
& & \multicolumn{6}{c|}{$V_{52}$} & \multicolumn{4}{c}{$V_{26}$} \\
& & $1$ & $4_{a}$ & $4_{b}$ & $6$ & $14$ & $20$ & $1$ & $4_{a}$ & $4_{b}$ & $6$ \\ \hline
\endhead
1) &  & 4 & 2 & 2 & 2 & 0 & 1 & 4 & 2 & 2 & 1 \\
2) &  & 4 & 2 & 2 & 3 & 1 & 0 & 4 & 2 & 2 & 1 \\
\end{longtable}
\noindent Memoir rows not reproduced: 5) (not compatible with L26).\\
Modules tested: V52, L26, V26.\\

\setcounter{table}{4}
\begin{longtable}{rc|ccccccc|ccccccc}
\caption{Alt$_{6} < F_4$, $p = 0$ or $p \nmid |H|$}\\
& & \multicolumn{7}{c|}{$V_{52}$} & \multicolumn{7}{c}{$V_{26}$} \\
& & $1$ & $5_{a}$ & $5_{b}$ & $8_{a}$ & $8_{b}$ & $9$ & $10$ & $1$ & $5_{a}$ & $5_{b}$ & $8_{a}$ & $8_{b}$ & $9$ & $10$ \\ \hline
\endfirsthead
& & \multicolumn{7}{c|}{$V_{52}$} & \multicolumn{7}{c}{$V_{26}$} \\
& & $1$ & $5_{a}$ & $5_{b}$ & $8_{a}$ & $8_{b}$ & $9$ & $10$ & $1$ & $5_{a}$ & $5_{b}$ & $8_{a}$ & $8_{b}$ & $9$ & $10$ \\ \hline
\endhead
1) &  & 0 & 0 & 0 & 4 & 0 & 0 & 2 & 2 & 0 & 0 & 3 & 0 & 0 & 0 \\
2) & \possprim & 0 & 0 & 0 & 3 & 1 & 0 & 2 & 0 & 0 & 0 & 0 & 1 & 2 & 0 \\
3) &  & 0 & 0 & 0 & 2 & 2 & 0 & 2 & 1 & 0 & 0 & 1 & 1 & 1 & 0 \\
4) &  & 2 & 1 & 1 & 0 & 0 & 0 & 4 & 1 & 0 & 0 & 1 & 1 & 1 & 0 \\
5) &  & 3 & 0 & 0 & 0 & 0 & 1 & 4 & 0 & 1 & 1 & 1 & 1 & 0 & 0 \\
6) &  & 3 & 0 & 0 & 0 & 0 & 1 & 4 & 1 & 3 & 0 & 0 & 0 & 0 & 1 \\
\end{longtable}
\noindent Modules tested: V52, V26.\\

\setcounter{table}{5}
\begin{longtable}{rc|ccccc|ccccc}
\caption{Alt$_{6} < F_4$, $p = 5$}\\
& & \multicolumn{5}{c|}{$V_{52}$} & \multicolumn{5}{c}{$V_{26}$} \\
& & $1$ & $5_{a}$ & $5_{b}$ & $8$ & $10$ & $1$ & $5_{a}$ & $5_{b}$ & $8$ & $10$ \\ \hline
\endfirsthead
& & \multicolumn{5}{c|}{$V_{52}$} & \multicolumn{5}{c}{$V_{26}$} \\
& & $1$ & $5_{a}$ & $5_{b}$ & $8$ & $10$ & $1$ & $5_{a}$ & $5_{b}$ & $8$ & $10$ \\ \hline
\endhead
1) & \possprim & 0 & 0 & 0 & 4 & 2 & 2 & 0 & 0 & 3 & 0 \\
2) &  & 2 & 1 & 1 & 0 & 4 & 2 & 0 & 0 & 3 & 0 \\
3) &  & 4 & 0 & 0 & 1 & 4 & 0 & 1 & 1 & 2 & 0 \\
4) &  & 4 & 0 & 0 & 1 & 4 & 1 & 3 & 0 & 0 & 1 \\
\end{longtable}
\noindent Modules tested: V52, V26.\\

\setcounter{table}{6}
\begin{longtable}{rc|ccccc|ccccc}
\caption{Alt$_{5} < F_4$, $p = 0$ or $p \nmid |H|$}\\
& & \multicolumn{5}{c|}{$V_{52}$} & \multicolumn{5}{c}{$V_{26}$} \\
& & $1$ & $3_{a}$ & $3_{b}$ & $4$ & $5$ & $1$ & $3_{a}$ & $3_{b}$ & $4$ & $5$ \\ \hline
\endfirsthead
& & \multicolumn{5}{c|}{$V_{52}$} & \multicolumn{5}{c}{$V_{26}$} \\
& & $1$ & $3_{a}$ & $3_{b}$ & $4$ & $5$ & $1$ & $3_{a}$ & $3_{b}$ & $4$ & $5$ \\ \hline
\endhead
1) & \possprim & 0 & 5 & 3 & 2 & 4 & 0 & 0 & 1 & 2 & 3 \\
2) &  & 3 & 5 & 3 & 5 & 1 & 0 & 0 & 1 & 2 & 3 \\
3) &  & 14 & 1 & 0 & 0 & 7 & 0 & 7 & 0 & 0 & 1 \\
4) &  & 0 & 4 & 4 & 2 & 4 & 1 & 1 & 1 & 1 & 3 \\
5) &  & 3 & 4 & 4 & 5 & 1 & 1 & 1 & 1 & 1 & 3 \\
6) &  & 0 & 6 & 2 & 2 & 4 & 2 & 3 & 0 & 0 & 3 \\
7) &  & 3 & 6 & 2 & 5 & 1 & 2 & 3 & 0 & 0 & 3 \\
8) &  & 3 & 5 & 3 & 5 & 1 & 3 & 0 & 1 & 5 & 0 \\
9) &  & 3 & 4 & 4 & 5 & 1 & 4 & 1 & 1 & 4 & 0 \\
10) &  & 3 & 6 & 2 & 5 & 1 & 5 & 3 & 0 & 3 & 0 \\
11) &  & 8 & 13 & 0 & 0 & 1 & 8 & 6 & 0 & 0 & 0 \\
\end{longtable}
\noindent Modules tested: V52, V26.\\

\setcounter{table}{7}
\begin{longtable}{rc|cccc|cccc}
\caption{Alt$_{5} < F_4$, $p = 3$}\\
& & \multicolumn{4}{c|}{$V_{52}$} & \multicolumn{4}{c}{$V_{26}$} \\
& & $1$ & $3_{a}$ & $3_{b}$ & $4$ & $1$ & $3_{a}$ & $3_{b}$ & $4$ \\ \hline
\endfirsthead
& & \multicolumn{4}{c|}{$V_{52}$} & \multicolumn{4}{c}{$V_{26}$} \\
& & $1$ & $3_{a}$ & $3_{b}$ & $4$ & $1$ & $3_{a}$ & $3_{b}$ & $4$ \\ \hline
\endhead
1) & \nongcr & 21 & 1 & 0 & 7 & 1 & 7 & 0 & 1 \\
2) & \possprim, \nongcr & 4 & 5 & 3 & 6 & 3 & 0 & 1 & 5 \\
3) & \possprim, \nongcr & 4 & 4 & 4 & 6 & 4 & 1 & 1 & 4 \\
4) & \nongcr & 4 & 6 & 2 & 6 & 5 & 3 & 0 & 3 \\
5) &  & 9 & 13 & 0 & 1 & 8 & 6 & 0 & 0 \\
\end{longtable}
\noindent Modules tested: V52, V26, L25.\\

\setcounter{table}{8}
\begin{longtable}{rc|ccc|ccc}
\caption{$M_{11} < F_4$, $p = 11$}\\
& & \multicolumn{3}{c|}{$V_{52}$} & \multicolumn{3}{c}{$V_{26}$} \\
& & $10_{a}$ & $10_{b}$ & $16$ & $1$ & $9$ & $16$ \\ \hline
\endfirsthead
& & \multicolumn{3}{c|}{$V_{52}$} & \multicolumn{3}{c}{$V_{26}$} \\
& & $10_{a}$ & $10_{b}$ & $16$ & $1$ & $9$ & $16$ \\ \hline
\endhead
1) &  & 1 & 1 & 2 & 1 & 1 & 1 \\
\end{longtable}
\noindent Modules tested: V52, V26.\\

\setcounter{table}{9}
\begin{longtable}{rc|ccc|cc}
\caption{$J_{1} < F_4$, $p = 11$}\\
& & \multicolumn{3}{c|}{$V_{52}$} & \multicolumn{2}{c}{$V_{26}$} \\
& & $1$ & $7$ & $14$ & $1$ & $7$ \\ \hline
\endfirsthead
& & \multicolumn{3}{c|}{$V_{52}$} & \multicolumn{2}{c}{$V_{26}$} \\
& & $1$ & $7$ & $14$ & $1$ & $7$ \\ \hline
\endhead
1) &  & 3 & 5 & 1 & 5 & 3 \\
\end{longtable}
\noindent Modules tested: V52, V26.\\

\setcounter{table}{10}
\begin{longtable}{rc|ccccc|cc}
\caption{$J_{2} < F_4$, $p = 2$}\\
& & \multicolumn{5}{c|}{$V_{52}$} & \multicolumn{2}{c}{$V_{26}$} \\
& & $1$ & $6_{a}$ & $6_{b}$ & $14_{a}$ & $14_{b}$ & $1$ & $6_{a}$ \\ \hline
\endfirsthead
& & \multicolumn{5}{c|}{$V_{52}$} & \multicolumn{2}{c}{$V_{26}$} \\
& & $1$ & $6_{a}$ & $6_{b}$ & $14_{a}$ & $14_{b}$ & $1$ & $6_{a}$ \\ \hline
\endhead
1) & \nongcr & 8 & 3 & 2 & 0 & 1 & 8 & 3 \\
2) & \nongcr & 8 & 5 & 0 & 1 & 0 & 8 & 3 \\
\end{longtable}
\noindent Modules tested: V52, L26, V26.\\

\setcounter{table}{11}
\begin{longtable}{rc|cccccc|cccccc}
\caption{$L_2(7) < F_4$, $p = 0$ or $p \nmid |H|$}\\
& & \multicolumn{6}{c|}{$V_{52}$} & \multicolumn{6}{c}{$V_{26}$} \\
& & $1$ & $3_{a}$ & $3_{b}$ & $6$ & $7$ & $8$ & $1$ & $3_{a}$ & $3_{b}$ & $6$ & $7$ & $8$ \\ \hline
\endfirsthead
& & \multicolumn{6}{c|}{$V_{52}$} & \multicolumn{6}{c}{$V_{26}$} \\
& & $1$ & $3_{a}$ & $3_{b}$ & $6$ & $7$ & $8$ & $1$ & $3_{a}$ & $3_{b}$ & $6$ & $7$ & $8$ \\ \hline
\endhead
1) & \possprim & 0 & 1 & 1 & 0 & 2 & 4 & 0 & 1 & 1 & 2 & 0 & 1 \\
2) &  & 0 & 1 & 1 & 0 & 2 & 4 & 2 & 0 & 0 & 0 & 0 & 3 \\
3) &  & 3 & 1 & 1 & 0 & 5 & 1 & 0 & 1 & 1 & 2 & 0 & 1 \\
4) &  & 3 & 1 & 1 & 0 & 5 & 1 & 2 & 0 & 0 & 0 & 0 & 3 \\
5) &  & 3 & 1 & 1 & 0 & 5 & 1 & 5 & 0 & 0 & 0 & 3 & 0 \\
6) &  & 8 & 0 & 0 & 6 & 0 & 1 & 0 & 3 & 3 & 0 & 0 & 1 \\
7) &  & 8 & 6 & 6 & 0 & 0 & 1 & 8 & 3 & 3 & 0 & 0 & 0 \\
\end{longtable}
\noindent Modules tested: V52, V26.\\

\setcounter{table}{12}
\begin{longtable}{rc|ccccc|ccccc}
\caption{$L_2(7) < F_4$, $p = 3$}\\
& & \multicolumn{5}{c|}{$V_{52}$} & \multicolumn{5}{c}{$V_{26}$} \\
& & $1$ & $3_{a}$ & $3_{b}$ & $6$ & $7$ & $1$ & $3_{a}$ & $3_{b}$ & $6$ & $7$ \\ \hline
\endfirsthead
& & \multicolumn{5}{c|}{$V_{52}$} & \multicolumn{5}{c}{$V_{26}$} \\
& & $1$ & $3_{a}$ & $3_{b}$ & $6$ & $7$ & $1$ & $3_{a}$ & $3_{b}$ & $6$ & $7$ \\ \hline
\endhead
1) & \possprim, \nongcr & 4 & 1 & 1 & 0 & 6 & 1 & 1 & 1 & 2 & 1 \\
2) & \nongcr & 4 & 1 & 1 & 0 & 6 & 5 & 0 & 0 & 0 & 3 \\
3) &  & 9 & 0 & 0 & 6 & 1 & 1 & 3 & 3 & 0 & 1 \\
4) &  & 9 & 6 & 6 & 0 & 1 & 8 & 3 & 3 & 0 & 0 \\
\end{longtable}
\noindent Memoir rows not reproduced: 3) (not compatible with L25).\\
Modules tested: V52, V26, L25.\\

\setcounter{table}{13}
\begin{longtable}{rc|ccccccc|cccccc}
\caption{$L_2(8) < F_4$, $p = 0$ or $p \nmid |H|$}\\
& & \multicolumn{7}{c|}{$V_{52}$} & \multicolumn{6}{c}{$V_{26}$} \\
& & $1$ & $7_{a}$ & $7_{b}$ & $7_{c}$ & $7_{d}$ & $8$ & $9_{a}$ & $1$ & $7_{a}$ & $7_{b}$ & $8$ & $9_{a}$ & $9_{c}$ \\ \hline
\endfirsthead
& & \multicolumn{7}{c|}{$V_{52}$} & \multicolumn{6}{c}{$V_{26}$} \\
& & $1$ & $7_{a}$ & $7_{b}$ & $7_{c}$ & $7_{d}$ & $8$ & $9_{a}$ & $1$ & $7_{a}$ & $7_{b}$ & $8$ & $9_{a}$ & $9_{c}$ \\ \hline
\endhead
1) & \possprim & 0 & 2 & 1 & 1 & 1 & 1 & 1 & 0 & 0 & 0 & 1 & 1 & 1 \\
2) &  & 0 & 2 & 1 & 1 & 1 & 1 & 1 & 1 & 1 & 0 & 0 & 1 & 1 \\
3) &  & 1 & 2 & 1 & 1 & 1 & 2 & 0 & 2 & 0 & 0 & 3 & 0 & 0 \\
4) &  & 1 & 2 & 1 & 1 & 1 & 2 & 0 & 3 & 1 & 0 & 2 & 0 & 0 \\
5) &  & 3 & 1 & 5 & 1 & 0 & 0 & 0 & 5 & 0 & 3 & 0 & 0 & 0 \\
\end{longtable}
\noindent Modules tested: V52, V26.\\

\setcounter{table}{14}
\begin{longtable}{rc|cccccc|cccc}
\caption{$L_2(8) < F_4$, $p = 7$}\\
& & \multicolumn{6}{c|}{$V_{52}$} & \multicolumn{4}{c}{$V_{26}$} \\
& & $1$ & $7_{a}$ & $7_{b}$ & $7_{c}$ & $7_{d}$ & $8$ & $1$ & $7_{a}$ & $7_{b}$ & $8$ \\ \hline
\endfirsthead
& & \multicolumn{6}{c|}{$V_{52}$} & \multicolumn{4}{c}{$V_{26}$} \\
& & $1$ & $7_{a}$ & $7_{b}$ & $7_{c}$ & $7_{d}$ & $8$ & $1$ & $7_{a}$ & $7_{b}$ & $8$ \\ \hline
\endhead
1) & \possprim, \nongcr & 1 & 2 & 1 & 1 & 1 & 2 & 2 & 0 & 0 & 3 \\
2) & \nongcr & 1 & 2 & 1 & 1 & 1 & 2 & 3 & 1 & 0 & 2 \\
3) &  & 3 & 1 & 5 & 1 & 0 & 0 & 5 & 0 & 3 & 0 \\
\end{longtable}
\noindent Modules tested: V52, V26.\\

\setcounter{table}{15}
\begin{longtable}{rc|ccc|cccc}
\caption{$L_2(8) < F_4$, $p = 3$}\\
& & \multicolumn{3}{c|}{$V_{52}$} & \multicolumn{4}{c}{$V_{26}$} \\
& & $1$ & $7$ & $9_{a}$ & $1$ & $7$ & $9_{a}$ & $9_{c}$ \\ \hline
\endfirsthead
& & \multicolumn{3}{c|}{$V_{52}$} & \multicolumn{4}{c}{$V_{26}$} \\
& & $1$ & $7$ & $9_{a}$ & $1$ & $7$ & $9_{a}$ & $9_{c}$ \\ \hline
\endhead
1) & \possprim, \nongcr & 1 & 6 & 1 & 1 & 1 & 1 & 1 \\
2) & \nongcr & 3 & 7 & 0 & 5 & 3 & 0 & 0 \\
\end{longtable}
\noindent Modules tested: V52, V26, L25.\\

\setcounter{table}{16}
\begin{longtable}{rc|ccccc|cccc}
\caption{$L_2(13) < F_4$, $p = 0$ or $p \nmid |H|$}\\
& & \multicolumn{5}{c|}{$V_{52}$} & \multicolumn{4}{c}{$V_{26}$} \\
& & $1$ & $7_{a}$ & $12_{a}$ & $12_{b}$ & $14_{b}$ & $1$ & $7_{a}$ & $12_{a}$ & $14_{a}$ \\ \hline
\endfirsthead
& & \multicolumn{5}{c|}{$V_{52}$} & \multicolumn{4}{c}{$V_{26}$} \\
& & $1$ & $7_{a}$ & $12_{a}$ & $12_{b}$ & $14_{b}$ & $1$ & $7_{a}$ & $12_{a}$ & $14_{a}$ \\ \hline
\endhead
1) & \possprim & 0 & 0 & 1 & 1 & 2 & 0 & 0 & 1 & 1 \\
2) &  & 3 & 5 & 0 & 0 & 1 & 5 & 3 & 0 & 0 \\
\end{longtable}
\noindent Modules tested: V52, V26.\\

\setcounter{table}{17}
\begin{longtable}{rc|cccc|cccc}
\caption{$L_2(13) < F_4$, $p = 7$}\\
& & \multicolumn{4}{c|}{$V_{52}$} & \multicolumn{4}{c}{$V_{26}$} \\
& & $1$ & $7_{a}$ & $12$ & $14_{b}$ & $1$ & $7_{a}$ & $12$ & $14_{a}$ \\ \hline
\endfirsthead
& & \multicolumn{4}{c|}{$V_{52}$} & \multicolumn{4}{c}{$V_{26}$} \\
& & $1$ & $7_{a}$ & $12$ & $14_{b}$ & $1$ & $7_{a}$ & $12$ & $14_{a}$ \\ \hline
\endhead
1) & \possprim & 0 & 0 & 2 & 2 & 0 & 0 & 1 & 1 \\
2) &  & 3 & 5 & 0 & 1 & 5 & 3 & 0 & 0 \\
\end{longtable}
\noindent Modules tested: V52, V26.\\

\setcounter{table}{18}
\begin{longtable}{rc|ccccc|cccc}
\caption{$L_2(13) < F_4$, $p = 3$}\\
& & \multicolumn{5}{c|}{$V_{52}$} & \multicolumn{4}{c}{$V_{26}$} \\
& & $1$ & $7_{a}$ & $7_{b}$ & $12_{a}$ & $12_{b}$ & $1$ & $7_{a}$ & $12_{a}$ & $13$ \\ \hline
\endfirsthead
& & \multicolumn{5}{c|}{$V_{52}$} & \multicolumn{4}{c}{$V_{26}$} \\
& & $1$ & $7_{a}$ & $7_{b}$ & $12_{a}$ & $12_{b}$ & $1$ & $7_{a}$ & $12_{a}$ & $13$ \\ \hline
\endhead
1) & \possprim & 0 & 2 & 2 & 1 & 1 & 1 & 0 & 1 & 1 \\
2) &  & 3 & 6 & 1 & 0 & 0 & 5 & 3 & 0 & 0 \\
\end{longtable}
\noindent Modules tested: V52, V26, L25.\\

\setcounter{table}{19}
\begin{longtable}{rc|cccccc|cccc}
\caption{$L_2(13) < F_4$, $p = 2$}\\
& & \multicolumn{6}{c|}{$V_{52}$} & \multicolumn{4}{c}{$V_{26}$} \\
& & $1$ & $6_{a}$ & $6_{b}$ & $12_{a}$ & $12_{b}$ & $14$ & $1$ & $6_{a}$ & $12_{a}$ & $14$ \\ \hline
\endfirsthead
& & \multicolumn{6}{c|}{$V_{52}$} & \multicolumn{4}{c}{$V_{26}$} \\
& & $1$ & $6_{a}$ & $6_{b}$ & $12_{a}$ & $12_{b}$ & $14$ & $1$ & $6_{a}$ & $12_{a}$ & $14$ \\ \hline
\endhead
1) & \possprim & 0 & 0 & 0 & 1 & 1 & 2 & 0 & 0 & 1 & 1 \\
2) & \nongcr & 8 & 5 & 0 & 0 & 0 & 1 & 8 & 3 & 0 & 0 \\
3) & \nongcr & 8 & 3 & 2 & 0 & 0 & 1 & 8 & 3 & 0 & 0 \\
\end{longtable}
\noindent Modules tested: V52, L26, V26.\\

\setcounter{table}{20}
\begin{longtable}{rc|ccc|cccc}
\caption{$L_2(17) < F_4$, $p = 0$ or $p \nmid |H|$}\\
& & \multicolumn{3}{c|}{$V_{52}$} & \multicolumn{4}{c}{$V_{26}$} \\
& & $16_{a}$ & $18_{b}$ & $18_{c}$ & $1$ & $9_{a}$ & $16_{a}$ & $17$ \\ \hline
\endfirsthead
& & \multicolumn{3}{c|}{$V_{52}$} & \multicolumn{4}{c}{$V_{26}$} \\
& & $16_{a}$ & $18_{b}$ & $18_{c}$ & $1$ & $9_{a}$ & $16_{a}$ & $17$ \\ \hline
\endhead
1) & \possprim & 1 & 1 & 1 & 0 & 1 & 0 & 1 \\
2) &  & 1 & 1 & 1 & 1 & 1 & 1 & 0 \\
\end{longtable}
\noindent Modules tested: V52, V26.\\

\setcounter{table}{21}
\begin{longtable}{rc|ccc|ccc}
\caption{$L_2(17) < F_4$, $p = 3$}\\
& & \multicolumn{3}{c|}{$V_{52}$} & \multicolumn{3}{c}{$V_{26}$} \\
& & $16$ & $18_{b}$ & $18_{c}$ & $1$ & $9_{a}$ & $16$ \\ \hline
\endfirsthead
& & \multicolumn{3}{c|}{$V_{52}$} & \multicolumn{3}{c}{$V_{26}$} \\
& & $16$ & $18_{b}$ & $18_{c}$ & $1$ & $9_{a}$ & $16$ \\ \hline
\endhead
1) & \possprim & 1 & 1 & 1 & 1 & 1 & 1 \\
\end{longtable}
\noindent Modules tested: V52, V26, L25.\\

\setcounter{table}{22}
\begin{longtable}{rc|cccc|ccc}
\caption{$L_2(17) < F_4$, $p = 2$}\\
& & \multicolumn{4}{c|}{$V_{52}$} & \multicolumn{3}{c}{$V_{26}$} \\
& & $1$ & $8_{a}$ & $8_{b}$ & $16_{a}$ & $1$ & $8_{a}$ & $8_{b}$ \\ \hline
\endfirsthead
& & \multicolumn{4}{c|}{$V_{52}$} & \multicolumn{3}{c}{$V_{26}$} \\
& & $1$ & $8_{a}$ & $8_{b}$ & $16_{a}$ & $1$ & $8_{a}$ & $8_{b}$ \\ \hline
\endhead
1) &  & 4 & 2 & 2 & 1 & 2 & 2 & 1 \\
\end{longtable}
\noindent Modules tested: V52, L26, V26.\\

\setcounter{table}{23}
\begin{longtable}{rc|cc|c}
\caption{$L_2(25) < F_4$, $p \neq 2,3,5$}\\
& & \multicolumn{2}{c|}{$V_{52}$} & \multicolumn{1}{c}{$V_{26}$} \\
& & $26_{d}$ & $26_{e}$ & $26_{a}$ \\ \hline
\endfirsthead
& & \multicolumn{2}{c|}{$V_{52}$} & \multicolumn{1}{c}{$V_{26}$} \\
& & $26_{d}$ & $26_{e}$ & $26_{a}$ \\ \hline
\endhead
1) & \possprim & 1 & 1 & 1 \\
\end{longtable}
\noindent Modules tested: V52, V26.\\

\setcounter{table}{24}
\begin{longtable}{rc|c|cc}
\caption{$L_2(25) < F_4$, $p = 3$}\\
& & \multicolumn{1}{c|}{$V_{52}$} & \multicolumn{2}{c}{$V_{26}$} \\
& & $26$ & $1$ & $25$ \\ \hline
\endfirsthead
& & \multicolumn{1}{c|}{$V_{52}$} & \multicolumn{2}{c}{$V_{26}$} \\
& & $26$ & $1$ & $25$ \\ \hline
\endhead
1) & \possprim & 2 & 1 & 1 \\
\end{longtable}
\noindent Modules tested: V52, V26, L25.\\

\setcounter{table}{25}
\begin{longtable}{rc|c|c}
\caption{$L_2(25) < F_4$, $p = 2$}\\
& & \multicolumn{1}{c|}{$V_{52}$} & \multicolumn{1}{c}{$V_{26}$} \\
& & $26$ & $26$ \\ \hline
\endfirsthead
& & \multicolumn{1}{c|}{$V_{52}$} & \multicolumn{1}{c}{$V_{26}$} \\
& & $26$ & $26$ \\ \hline
\endhead
1) & \possprim & 2 & 1 \\
\end{longtable}
\noindent Modules tested: V52, L26, V26.\\

\setcounter{table}{26}
\begin{longtable}{rc|cc|c}
\caption{$L_2(27) < F_4$, $p \neq 2,3,7$}\\
& & \multicolumn{2}{c|}{$V_{52}$} & \multicolumn{1}{c}{$V_{26}$} \\
& & $26_{a}$ & $26_{b}$ & $26_{e}$ \\ \hline
\endfirsthead
& & \multicolumn{2}{c|}{$V_{52}$} & \multicolumn{1}{c}{$V_{26}$} \\
& & $26_{a}$ & $26_{b}$ & $26_{e}$ \\ \hline
\endhead
1) & \possprim & 1 & 1 & 1 \\
\end{longtable}
\noindent Modules tested: V52, V26.\\

\setcounter{table}{27}
\begin{longtable}{rc|c|cc}
\caption{$L_2(27) < F_4$, $p = 7$}\\
& & \multicolumn{1}{c|}{$V_{52}$} & \multicolumn{2}{c}{$V_{26}$} \\
& & $26$ & $13_{a}$ & $13_{b}$ \\ \hline
\endfirsthead
& & \multicolumn{1}{c|}{$V_{52}$} & \multicolumn{2}{c}{$V_{26}$} \\
& & $26$ & $13_{a}$ & $13_{b}$ \\ \hline
\endhead
1) & \possprim & 2 & 1 & 1 \\
\end{longtable}
\noindent Modules tested: V52, V26.\\

\setcounter{table}{28}
\begin{longtable}{rc|cc|c}
\caption{$L_2(27) < F_4$, $p = 2$}\\
& & \multicolumn{2}{c|}{$V_{52}$} & \multicolumn{1}{c}{$V_{26}$} \\
& & $26_{a}$ & $26_{b}$ & $26_{a}$ \\ \hline
\endfirsthead
& & \multicolumn{2}{c|}{$V_{52}$} & \multicolumn{1}{c}{$V_{26}$} \\
& & $26_{a}$ & $26_{b}$ & $26_{a}$ \\ \hline
\endhead
1) & \possprim & 1 & 1 & 1 \\
\end{longtable}
\noindent Modules tested: V52, L26, V26.\\

\setcounter{table}{29}
\begin{longtable}{rc|cc|c}
\caption{$L_3(3) < F_4$, $p \neq 2, 3$}\\
& & \multicolumn{2}{c|}{$V_{52}$} & \multicolumn{1}{c}{$V_{26}$} \\
& & $26_{b}$ & $26_{c}$ & $26_{a}$ \\ \hline
\endfirsthead
& & \multicolumn{2}{c|}{$V_{52}$} & \multicolumn{1}{c}{$V_{26}$} \\
& & $26_{b}$ & $26_{c}$ & $26_{a}$ \\ \hline
\endhead
1) & \possprim & 1 & 1 & 1 \\
\end{longtable}
\noindent Modules tested: V52, V26.\\

\setcounter{table}{30}
\begin{longtable}{rc|c|c}
\caption{$L_3(3) < F_4$, $p = 2$}\\
& & \multicolumn{1}{c|}{$V_{52}$} & \multicolumn{1}{c}{$V_{26}$} \\
& & $26$ & $26$ \\ \hline
\endfirsthead
& & \multicolumn{1}{c|}{$V_{52}$} & \multicolumn{1}{c}{$V_{26}$} \\
& & $26$ & $26$ \\ \hline
\endhead
1) & \possprim & 2 & 1 \\
\end{longtable}
\noindent Modules tested: V52, L26, V26.\\

\setcounter{table}{31}
\begin{longtable}{rc|cc|c}
\caption{$L_4(3) < F_4$, $p = 2$}\\
& & \multicolumn{2}{c|}{$V_{52}$} & \multicolumn{1}{c}{$V_{26}$} \\
& & $26_{a}$ & $26_{b}$ & $26_{a}$ \\ \hline
\endfirsthead
& & \multicolumn{2}{c|}{$V_{52}$} & \multicolumn{1}{c}{$V_{26}$} \\
& & $26_{a}$ & $26_{b}$ & $26_{a}$ \\ \hline
\endhead
1) & \possprim & 1 & 1 & 1 \\
\end{longtable}
\noindent Modules tested: V52, L26, V26.\\

\setcounter{table}{32}
\begin{longtable}{rc|cccccc|cccc}
\caption{$U_3(3) < F_4$, $p = 0$ or $p \nmid |H|$}\\
& & \multicolumn{6}{c|}{$V_{52}$} & \multicolumn{4}{c}{$V_{26}$} \\
& & $1$ & $7_{a}$ & $7_{b}$ & $7_{c}$ & $14$ & $21_{a}$ & $1$ & $6$ & $7_{a}$ & $14$ \\ \hline
\endfirsthead
& & \multicolumn{6}{c|}{$V_{52}$} & \multicolumn{4}{c}{$V_{26}$} \\
& & $1$ & $7_{a}$ & $7_{b}$ & $7_{c}$ & $14$ & $21_{a}$ & $1$ & $6$ & $7_{a}$ & $14$ \\ \hline
\endhead
1) &  & 3 & 5 & 0 & 0 & 1 & 0 & 5 & 0 & 3 & 0 \\
2) &  & 3 & 0 & 2 & 2 & 0 & 1 & 0 & 2 & 0 & 1 \\
\end{longtable}
\noindent Modules tested: V52, V26.\\

\setcounter{table}{33}
\begin{longtable}{rc|cccccccc|ccccc}
\caption{$U_3(3) < F_4$, $p = 7$}\\
& & \multicolumn{8}{c|}{$V_{52}$} & \multicolumn{5}{c}{$V_{26}$} \\
& & $1$ & $6$ & $7_{a}$ & $7_{b}$ & $7_{c}$ & $14$ & $21_{a}$ & $26$ & $1$ & $6$ & $7_{a}$ & $14$ & $26$ \\ \hline
\endfirsthead
& & \multicolumn{8}{c|}{$V_{52}$} & \multicolumn{5}{c}{$V_{26}$} \\
& & $1$ & $6$ & $7_{a}$ & $7_{b}$ & $7_{c}$ & $14$ & $21_{a}$ & $26$ & $1$ & $6$ & $7_{a}$ & $14$ & $26$ \\ \hline
\endhead
1) & \possprim & 0 & 2 & 0 & 0 & 0 & 1 & 0 & 1 & 0 & 0 & 0 & 0 & 1 \\
2) &  & 3 & 0 & 0 & 2 & 2 & 0 & 1 & 0 & 0 & 2 & 0 & 1 & 0 \\
3) &  & 3 & 0 & 5 & 0 & 0 & 1 & 0 & 0 & 5 & 0 & 3 & 0 & 0 \\
\end{longtable}
\noindent Modules tested: V52, V26.\\

\setcounter{table}{34}
\begin{longtable}{rc|c|c}
\caption{$^{3}D_4(2) < F_4$, $p \neq 2, 3$ (irreducible on both)}\\
& & \multicolumn{1}{c|}{$V_{52}$} & \multicolumn{1}{c}{$V_{26}$} \\
& & $52$ & $26$ \\ \hline
\endfirsthead
& & \multicolumn{1}{c|}{$V_{52}$} & \multicolumn{1}{c}{$V_{26}$} \\
& & $52$ & $26$ \\ \hline
\endhead
1) & \possprim & 1 & 1 \\
\end{longtable}
\noindent Modules tested: V52, V26.\\

\setcounter{table}{35}
\begin{longtable}{rc|c|cc}
\caption{$^{3}D_4(2) < F_4$, $p = 3$}\\
& & \multicolumn{1}{c|}{$V_{52}$} & \multicolumn{2}{c}{$V_{26}$} \\
& & $52$ & $1$ & $25$ \\ \hline
\endfirsthead
& & \multicolumn{1}{c|}{$V_{52}$} & \multicolumn{2}{c}{$V_{26}$} \\
& & $52$ & $1$ & $25$ \\ \hline
\endhead
1) & \possprim & 1 & 1 & 1 \\
\end{longtable}
\noindent Modules tested: V52, V26, L25.\\


\newpage
% Updated tables for E6 (Memoir Chapter 6), generated by tools/fg_tables.g
\section*{$E_6$: feasible characters on $L(G) = V_{78}$ and $V_{\min} = V_{27}$}

\setcounter{table}{36}
\begin{longtable}{rc|cccc|ccc}
\caption{Alt$_{12} < E_6$, $p = 2$}\\
& & \multicolumn{4}{c|}{$V_{78}$} & \multicolumn{3}{c}{$V_{27}$} \\
& & $1$ & $16_{a}$ & $16_{b}$ & $44$ & $1$ & $10$ & $16_{a}$ \\ \hline
\endfirsthead
& & \multicolumn{4}{c|}{$V_{78}$} & \multicolumn{3}{c}{$V_{27}$} \\
& & $1$ & $16_{a}$ & $16_{b}$ & $44$ & $1$ & $10$ & $16_{a}$ \\ \hline
\endhead
1) &  & 2 & 1 & 1 & 1 & 1 & 1 & 1 \\
\end{longtable}
\noindent Modules tested: V78, V27.\\

\setcounter{table}{37}
\begin{longtable}{rc|cccc|ccc}
\caption{Alt$_{11} < E_6$, $p = 2$}\\
& & \multicolumn{4}{c|}{$V_{78}$} & \multicolumn{3}{c}{$V_{27}$} \\
& & $1$ & $16_{a}$ & $16_{b}$ & $44$ & $1$ & $10$ & $16_{a}$ \\ \hline
\endfirsthead
& & \multicolumn{4}{c|}{$V_{78}$} & \multicolumn{3}{c}{$V_{27}$} \\
& & $1$ & $16_{a}$ & $16_{b}$ & $44$ & $1$ & $10$ & $16_{a}$ \\ \hline
\endhead
1) &  & 2 & 1 & 1 & 1 & 1 & 1 & 1 \\
\end{longtable}
\noindent Modules tested: V78, V27.\\

\setcounter{table}{38}
\begin{longtable}{rc|cccc|cccc}
\caption{Alt$_{10} < E_6$, $p = 2$}\\
& & \multicolumn{4}{c|}{$V_{78}$} & \multicolumn{4}{c}{$V_{27}$} \\
& & $1$ & $8$ & $16$ & $26$ & $1$ & $8$ & $16$ & $26$ \\ \hline
\endfirsthead
& & \multicolumn{4}{c|}{$V_{78}$} & \multicolumn{4}{c}{$V_{27}$} \\
& & $1$ & $8$ & $16$ & $26$ & $1$ & $8$ & $16$ & $26$ \\ \hline
\endhead
1) &  & 2 & 1 & 1 & 2 & 1 & 0 & 0 & 1 \\
2) &  & 4 & 2 & 2 & 1 & 3 & 1 & 1 & 0 \\
\end{longtable}
\noindent Modules tested: V78, V27.\\

\setcounter{table}{39}
\begin{longtable}{rc|ccccc|ccccc}
\caption{Alt$_{9} < E_6$, $p = 2$}\\
& & \multicolumn{5}{c|}{$V_{78}$} & \multicolumn{5}{c}{$V_{27}$} \\
& & $1$ & $8_{a}$ & $8_{b}$ & $8_{c}$ & $26$ & $1$ & $8_{a}$ & $8_{b}$ & $8_{c}$ & $26$ \\ \hline
\endfirsthead
& & \multicolumn{5}{c|}{$V_{78}$} & \multicolumn{5}{c}{$V_{27}$} \\
& & $1$ & $8_{a}$ & $8_{b}$ & $8_{c}$ & $26$ & $1$ & $8_{a}$ & $8_{b}$ & $8_{c}$ & $26$ \\ \hline
\endhead
1) & \possprim & 2 & 1 & 1 & 1 & 2 & 1 & 0 & 0 & 0 & 1 \\
2) &  & 4 & 2 & 2 & 2 & 1 & 3 & 1 & 1 & 1 & 0 \\
\end{longtable}
\noindent Modules tested: V78, V27.\\

\setcounter{table}{40}
\begin{longtable}{rc|cccccc|cc}
\caption{Alt$_{7} < E_6$, $p = 0$ or $p \nmid |H|$}\\
& & \multicolumn{6}{c|}{$V_{78}$} & \multicolumn{2}{c}{$V_{27}$} \\
& & $1$ & $6$ & $10_{a}$ & $10_{b}$ & $14_{a}$ & $15$ & $6$ & $15$ \\ \hline
\endfirsthead
& & \multicolumn{6}{c|}{$V_{78}$} & \multicolumn{2}{c}{$V_{27}$} \\
& & $1$ & $6$ & $10_{a}$ & $10_{b}$ & $14_{a}$ & $15$ & $6$ & $15$ \\ \hline
\endhead
1) &  & 3 & 1 & 2 & 2 & 1 & 1 & 2 & 1 \\
\end{longtable}
\noindent Modules tested: V78, V27.\\

\setcounter{table}{41}
\begin{longtable}{rc|ccccc|cc}
\caption{$3\cdot$Alt$_{7} < E_6$, $p = 0$ or $p \nmid |H|$}\\
& & \multicolumn{5}{c|}{$V_{78}$} & \multicolumn{2}{c}{$V_{27}$} \\
& & $1$ & $10_{a}$ & $10_{b}$ & $14_{a}$ & $21_{a}$ & $6_{b}$ & $15_{b}$ \\ \hline
\endfirsthead
& & \multicolumn{5}{c|}{$V_{78}$} & \multicolumn{2}{c}{$V_{27}$} \\
& & $1$ & $10_{a}$ & $10_{b}$ & $14_{a}$ & $21_{a}$ & $6_{b}$ & $15_{b}$ \\ \hline
\endhead
1) &  & 3 & 2 & 2 & 1 & 1 & 2 & 1 \\
\end{longtable}
\noindent Modules tested: V78, V27.\\

\setcounter{table}{42}
\begin{longtable}{rc|cccc|ccc}
\caption{Alt$_{7} < E_6$, $p = 7$}\\
& & \multicolumn{4}{c|}{$V_{78}$} & \multicolumn{3}{c}{$V_{27}$} \\
& & $1$ & $5$ & $10$ & $14_{a}$ & $1$ & $5$ & $10$ \\ \hline
\endfirsthead
& & \multicolumn{4}{c|}{$V_{78}$} & \multicolumn{3}{c}{$V_{27}$} \\
& & $1$ & $5$ & $10$ & $14_{a}$ & $1$ & $5$ & $10$ \\ \hline
\endhead
1) & \nongcr & 4 & 2 & 5 & 1 & 2 & 3 & 1 \\
\end{longtable}
\noindent Modules tested: V78, V27.\\

\setcounter{table}{43}
\begin{longtable}{rc|cccc|cc}
\caption{$3\cdot$Alt$_{7} < E_6$, $p = 7$}\\
& & \multicolumn{4}{c|}{$V_{78}$} & \multicolumn{2}{c}{$V_{27}$} \\
& & $1$ & $10$ & $14_{a}$ & $21_{a}$ & $6_{a}$ & $15_{a}$ \\ \hline
\endfirsthead
& & \multicolumn{4}{c|}{$V_{78}$} & \multicolumn{2}{c}{$V_{27}$} \\
& & $1$ & $10$ & $14_{a}$ & $21_{a}$ & $6_{a}$ & $15_{a}$ \\ \hline
\endhead
1) &  & 3 & 4 & 1 & 1 & 2 & 1 \\
\end{longtable}
\noindent Modules tested: V78, V27.\\

\setcounter{table}{44}
\begin{longtable}{rc|ccccccc|ccccc}
\caption{Alt$_{7} < E_6$, $p = 5$}\\
& & \multicolumn{7}{c|}{$V_{78}$} & \multicolumn{5}{c}{$V_{27}$} \\
& & $1$ & $6$ & $8$ & $10_{a}$ & $10_{b}$ & $15$ & $35$ & $1$ & $6$ & $8$ & $13$ & $15$ \\ \hline
\endfirsthead
& & \multicolumn{7}{c|}{$V_{78}$} & \multicolumn{5}{c}{$V_{27}$} \\
& & $1$ & $6$ & $8$ & $10_{a}$ & $10_{b}$ & $15$ & $35$ & $1$ & $6$ & $8$ & $13$ & $15$ \\ \hline
\endhead
1) & \possprim & 0 & 0 & 1 & 0 & 0 & 0 & 2 & 0 & 1 & 1 & 1 & 0 \\
2) &  & 3 & 2 & 1 & 2 & 2 & 1 & 0 & 0 & 2 & 0 & 0 & 1 \\
3) &  & 2 & 0 & 7 & 1 & 1 & 0 & 0 & 3 & 0 & 3 & 0 & 0 \\
\end{longtable}
\noindent Modules tested: V78, V27.\\

\setcounter{table}{45}
\begin{longtable}{rc|ccccccc|cccc}
\caption{$3\cdot$Alt$_{7} < E_6$, $p = 5$}\\
& & \multicolumn{7}{c|}{$V_{78}$} & \multicolumn{4}{c}{$V_{27}$} \\
& & $1$ & $6_{a}$ & $8$ & $10_{a}$ & $10_{b}$ & $13$ & $35$ & $3_{a}$ & $6_{b}$ & $15_{b}$ & $21_{a}$ \\ \hline
\endfirsthead
& & \multicolumn{7}{c|}{$V_{78}$} & \multicolumn{4}{c}{$V_{27}$} \\
& & $1$ & $6_{a}$ & $8$ & $10_{a}$ & $10_{b}$ & $13$ & $35$ & $3_{a}$ & $6_{b}$ & $15_{b}$ & $21_{a}$ \\ \hline
\endhead
1) & \possprim & 0 & 0 & 1 & 0 & 0 & 0 & 2 & 0 & 1 & 0 & 1 \\
2) &  & 3 & 1 & 2 & 2 & 2 & 1 & 0 & 0 & 2 & 1 & 0 \\
3) &  & 14 & 0 & 8 & 0 & 0 & 0 & 0 & 7 & 1 & 0 & 0 \\
\end{longtable}
\noindent Modules tested: V78, V27.\\

\setcounter{table}{46}
\begin{longtable}{rc|cccccc|cc}
\caption{Alt$_{7} < E_6$, $p = 3$}\\
& & \multicolumn{6}{c|}{$V_{78}$} & \multicolumn{2}{c}{$V_{27}$} \\
& & $1$ & $6$ & $10_{a}$ & $10_{b}$ & $13$ & $15$ & $6$ & $15$ \\ \hline
\endfirsthead
& & \multicolumn{6}{c|}{$V_{78}$} & \multicolumn{2}{c}{$V_{27}$} \\
& & $1$ & $6$ & $10_{a}$ & $10_{b}$ & $13$ & $15$ & $6$ & $15$ \\ \hline
\endhead
1) & \possprim, \nongcr & 4 & 1 & 2 & 2 & 1 & 1 & 2 & 1 \\
\end{longtable}
\noindent Modules tested: V78, L77, V27.\\

\setcounter{table}{47}
\begin{longtable}{rc|cccccc|cccccc}
\caption{Alt$_{7} < E_6$, $p = 2$}\\
& & \multicolumn{6}{c|}{$V_{78}$} & \multicolumn{6}{c}{$V_{27}$} \\
& & $1$ & $4_{a}$ & $4_{b}$ & $6$ & $14$ & $20$ & $1$ & $4_{a}$ & $4_{b}$ & $6$ & $14$ & $20$ \\ \hline
\endfirsthead
& & \multicolumn{6}{c|}{$V_{78}$} & \multicolumn{6}{c}{$V_{27}$} \\
& & $1$ & $4_{a}$ & $4_{b}$ & $6$ & $14$ & $20$ & $1$ & $4_{a}$ & $4_{b}$ & $6$ & $14$ & $20$ \\ \hline
\endhead
1) & \nongcr & 4 & 2 & 2 & 3 & 0 & 2 & 1 & 0 & 0 & 1 & 0 & 1 \\
2) &  & 4 & 2 & 2 & 4 & 1 & 1 & 1 & 0 & 0 & 1 & 0 & 1 \\
3) & \nongcr & 4 & 2 & 2 & 5 & 2 & 0 & 1 & 0 & 0 & 2 & 1 & 0 \\
4) &  & 8 & 4 & 4 & 3 & 0 & 1 & 5 & 2 & 2 & 1 & 0 & 0 \\
5) &  & 8 & 4 & 4 & 4 & 1 & 0 & 5 & 2 & 2 & 1 & 0 & 0 \\
\end{longtable}
\noindent Modules tested: V78, V27.\\

\setcounter{table}{48}
\begin{longtable}{rc|cccccc|cc}
\caption{$3\cdot$Alt$_{7} < E_6$, $p = 2$}\\
& & \multicolumn{6}{c|}{$V_{78}$} & \multicolumn{2}{c}{$V_{27}$} \\
& & $1$ & $4_{a}$ & $4_{b}$ & $6_{a}$ & $14$ & $20$ & $6_{b}$ & $15_{a}$ \\ \hline
\endfirsthead
& & \multicolumn{6}{c|}{$V_{78}$} & \multicolumn{2}{c}{$V_{27}$} \\
& & $1$ & $4_{a}$ & $4_{b}$ & $6_{a}$ & $14$ & $20$ & $6_{b}$ & $15_{a}$ \\ \hline
\endhead
1) & \nongcr & 4 & 2 & 2 & 3 & 0 & 2 & 2 & 1 \\
2) &  & 4 & 2 & 2 & 4 & 1 & 1 & 2 & 1 \\
\end{longtable}
\noindent Modules tested: V78, V27.\\

\setcounter{table}{49}
\begin{longtable}{rc|ccccccc|ccccccc}
\caption{Alt$_{6} < E_6$, $p = 0$ or $p \nmid |H|$}\\
& & \multicolumn{7}{c|}{$V_{78}$} & \multicolumn{7}{c}{$V_{27}$} \\
& & $1$ & $5_{a}$ & $5_{b}$ & $8_{a}$ & $8_{b}$ & $9$ & $10$ & $1$ & $5_{a}$ & $5_{b}$ & $8_{a}$ & $8_{b}$ & $9$ & $10$ \\ \hline
\endfirsthead
& & \multicolumn{7}{c|}{$V_{78}$} & \multicolumn{7}{c}{$V_{27}$} \\
& & $1$ & $5_{a}$ & $5_{b}$ & $8_{a}$ & $8_{b}$ & $9$ & $10$ & $1$ & $5_{a}$ & $5_{b}$ & $8_{a}$ & $8_{b}$ & $9$ & $10$ \\ \hline
\endhead
1) & \possprim & 0 & 0 & 0 & 3 & 2 & 2 & 2 & 0 & 1 & 1 & 0 & 1 & 1 & 0 \\
2) &  & 2 & 1 & 1 & 1 & 0 & 2 & 4 & 0 & 1 & 1 & 0 & 1 & 1 & 0 \\
3) &  & 0 & 0 & 0 & 3 & 2 & 2 & 2 & 1 & 0 & 0 & 0 & 1 & 2 & 0 \\
4) &  & 2 & 1 & 1 & 1 & 0 & 2 & 4 & 1 & 0 & 0 & 0 & 1 & 2 & 0 \\
5) &  & 1 & 0 & 0 & 3 & 3 & 1 & 2 & 1 & 1 & 1 & 1 & 1 & 0 & 0 \\
6) &  & 3 & 1 & 1 & 1 & 1 & 1 & 4 & 1 & 1 & 1 & 1 & 1 & 0 & 0 \\
7) &  & 1 & 0 & 0 & 3 & 3 & 1 & 2 & 2 & 0 & 0 & 1 & 1 & 1 & 0 \\
8) &  & 3 & 1 & 1 & 1 & 1 & 1 & 4 & 2 & 0 & 0 & 1 & 1 & 1 & 0 \\
9) &  & 4 & 3 & 0 & 0 & 0 & 1 & 5 & 2 & 3 & 0 & 0 & 0 & 0 & 1 \\
10) &  & 2 & 0 & 0 & 7 & 0 & 0 & 2 & 3 & 0 & 0 & 3 & 0 & 0 & 0 \\
\end{longtable}
\noindent Modules tested: V78, V27.\\

\setcounter{table}{50}
\begin{longtable}{rc|ccccccc|ccccc}
\caption{$3\cdot$Alt$_{6} < E_6$, $p = 0$ or $p \nmid |H|$}\\
& & \multicolumn{7}{c|}{$V_{78}$} & \multicolumn{5}{c}{$V_{27}$} \\
& & $1$ & $5_{a}$ & $5_{b}$ & $8_{a}$ & $8_{b}$ & $9_{a}$ & $10$ & $3_{a}$ & $3_{c}$ & $6_{a}$ & $9_{b}$ & $15_{a}$ \\ \hline
\endfirsthead
& & \multicolumn{7}{c|}{$V_{78}$} & \multicolumn{5}{c}{$V_{27}$} \\
& & $1$ & $5_{a}$ & $5_{b}$ & $8_{a}$ & $8_{b}$ & $9_{a}$ & $10$ & $3_{a}$ & $3_{c}$ & $6_{a}$ & $9_{b}$ & $15_{a}$ \\ \hline
\endhead
1) &  & 1 & 0 & 0 & 3 & 3 & 1 & 2 & 0 & 0 & 2 & 0 & 1 \\
2) &  & 3 & 1 & 1 & 1 & 1 & 1 & 4 & 0 & 0 & 2 & 0 & 1 \\
3) & \possprim & 0 & 0 & 0 & 3 & 2 & 2 & 2 & 0 & 1 & 1 & 2 & 0 \\
4) &  & 2 & 1 & 1 & 1 & 0 & 2 & 4 & 0 & 1 & 1 & 2 & 0 \\
5) &  & 2 & 0 & 0 & 7 & 0 & 0 & 2 & 3 & 0 & 3 & 0 & 0 \\
6) &  & 14 & 0 & 0 & 8 & 0 & 0 & 0 & 7 & 0 & 1 & 0 & 0 \\
7) &  & 1 & 0 & 0 & 3 & 3 & 1 & 2 & 1 & 1 & 2 & 1 & 0 \\
8) &  & 3 & 1 & 1 & 1 & 1 & 1 & 4 & 1 & 1 & 2 & 1 & 0 \\
9) &  & 8 & 0 & 0 & 1 & 1 & 6 & 0 & 3 & 3 & 0 & 1 & 0 \\
\end{longtable}
\noindent Modules tested: V78, V27.\\

\setcounter{table}{51}
\begin{longtable}{rc|ccccc|ccccc}
\caption{Alt$_{6} < E_6$, $p = 5$}\\
& & \multicolumn{5}{c|}{$V_{78}$} & \multicolumn{5}{c}{$V_{27}$} \\
& & $1$ & $5_{a}$ & $5_{b}$ & $8$ & $10$ & $1$ & $5_{a}$ & $5_{b}$ & $8$ & $10$ \\ \hline
\endfirsthead
& & \multicolumn{5}{c|}{$V_{78}$} & \multicolumn{5}{c}{$V_{27}$} \\
& & $1$ & $5_{a}$ & $5_{b}$ & $8$ & $10$ & $1$ & $5_{a}$ & $5_{b}$ & $8$ & $10$ \\ \hline
\endhead
1) & \possprim, \nongcr & 2 & 0 & 0 & 7 & 2 & 1 & 1 & 1 & 2 & 0 \\
2) & \nongcr & 4 & 1 & 1 & 3 & 4 & 1 & 1 & 1 & 2 & 0 \\
3) &  & 5 & 3 & 0 & 1 & 5 & 2 & 3 & 0 & 0 & 1 \\
4) & \nongcr & 2 & 0 & 0 & 7 & 2 & 3 & 0 & 0 & 3 & 0 \\
5) & \nongcr & 4 & 1 & 1 & 3 & 4 & 3 & 0 & 0 & 3 & 0 \\
\end{longtable}
\noindent Modules tested: V78, V27.\\

\setcounter{table}{52}
\begin{longtable}{rc|ccccc|ccc}
\caption{$3\cdot$Alt$_{6} < E_6$, $p = 5$}\\
& & \multicolumn{5}{c|}{$V_{78}$} & \multicolumn{3}{c}{$V_{27}$} \\
& & $1$ & $5_{a}$ & $5_{b}$ & $8$ & $10$ & $3_{a}$ & $6_{a}$ & $15_{a}$ \\ \hline
\endfirsthead
& & \multicolumn{5}{c|}{$V_{78}$} & \multicolumn{3}{c}{$V_{27}$} \\
& & $1$ & $5_{a}$ & $5_{b}$ & $8$ & $10$ & $3_{a}$ & $6_{a}$ & $15_{a}$ \\ \hline
\endhead
1) & \possprim, \nongcr & 2 & 0 & 0 & 7 & 2 & 0 & 2 & 1 \\
2) & \nongcr & 4 & 1 & 1 & 3 & 4 & 0 & 2 & 1 \\
3) & \possprim, \nongcr & 2 & 0 & 0 & 7 & 2 & 3 & 3 & 0 \\
4) & \nongcr & 4 & 1 & 1 & 3 & 4 & 3 & 3 & 0 \\
5) & \nongcr & 14 & 0 & 0 & 8 & 0 & 7 & 1 & 0 \\
\end{longtable}
\noindent Modules tested: V78, V27.\\

\setcounter{table}{53}
\begin{longtable}{rc|ccccc|ccccc}
\caption{Alt$_{5} < E_6$, $p = 0$ or $p \nmid |H|$}\\
& & \multicolumn{5}{c|}{$V_{78}$} & \multicolumn{5}{c}{$V_{27}$} \\
& & $1$ & $3_{a}$ & $3_{b}$ & $4$ & $5$ & $1$ & $3_{a}$ & $3_{b}$ & $4$ & $5$ \\ \hline
\endfirsthead
& & \multicolumn{5}{c|}{$V_{78}$} & \multicolumn{5}{c}{$V_{27}$} \\
& & $1$ & $3_{a}$ & $3_{b}$ & $4$ & $5$ & $1$ & $3_{a}$ & $3_{b}$ & $4$ & $5$ \\ \hline
\endhead
1) &  & 5 & 1 & 1 & 3 & 11 & 0 & 3 & 3 & 1 & 1 \\
2) &  & 8 & 1 & 1 & 6 & 8 & 0 & 3 & 3 & 1 & 1 \\
3) &  & 0 & 5 & 4 & 4 & 7 & 1 & 0 & 1 & 2 & 3 \\
4) &  & 3 & 5 & 4 & 7 & 4 & 1 & 0 & 1 & 2 & 3 \\
5) &  & 14 & 8 & 0 & 0 & 8 & 1 & 7 & 0 & 0 & 1 \\
6) &  & 1 & 5 & 5 & 3 & 7 & 2 & 1 & 1 & 1 & 3 \\
7) &  & 4 & 5 & 5 & 6 & 4 & 2 & 1 & 1 & 1 & 3 \\
8) &  & 2 & 9 & 2 & 2 & 7 & 3 & 3 & 0 & 0 & 3 \\
9) &  & 5 & 9 & 2 & 5 & 4 & 3 & 3 & 0 & 0 & 3 \\
10) &  & 6 & 5 & 4 & 10 & 1 & 4 & 0 & 1 & 5 & 0 \\
11) &  & 7 & 5 & 5 & 9 & 1 & 5 & 1 & 1 & 4 & 0 \\
12) &  & 8 & 9 & 2 & 8 & 1 & 6 & 3 & 0 & 3 & 0 \\
13) &  & 16 & 19 & 0 & 0 & 1 & 9 & 6 & 0 & 0 & 0 \\
\end{longtable}
\noindent Modules tested: V78, V27.\\

\setcounter{table}{54}
\begin{longtable}{rc|cccc|cccc}
\caption{Alt$_{5} < E_6$, $p = 3$}\\
& & \multicolumn{4}{c|}{$V_{78}$} & \multicolumn{4}{c}{$V_{27}$} \\
& & $1$ & $3_{a}$ & $3_{b}$ & $4$ & $1$ & $3_{a}$ & $3_{b}$ & $4$ \\ \hline
\endfirsthead
& & \multicolumn{4}{c|}{$V_{78}$} & \multicolumn{4}{c}{$V_{27}$} \\
& & $1$ & $3_{a}$ & $3_{b}$ & $4$ & $1$ & $3_{a}$ & $3_{b}$ & $4$ \\ \hline
\endhead
1) & \possprim, \nongcr & 15 & 1 & 0 & 15 & 0 & 2 & 3 & 3 \\
2) & \nongcr & 16 & 1 & 1 & 14 & 1 & 3 & 3 & 2 \\
3) & \nongcr & 22 & 8 & 0 & 8 & 2 & 7 & 0 & 1 \\
4) & \possprim, \nongcr & 7 & 5 & 4 & 11 & 4 & 0 & 1 & 5 \\
5) & \nongcr & 8 & 5 & 5 & 10 & 5 & 1 & 1 & 4 \\
6) & \nongcr & 9 & 9 & 2 & 9 & 6 & 3 & 0 & 3 \\
7) &  & 17 & 19 & 0 & 1 & 9 & 6 & 0 & 0 \\
\end{longtable}
\noindent Modules tested: V78, L77, V27.\\

\setcounter{table}{55}
\begin{longtable}{rc|cccc|cccc}
\caption{$M_{11} < E_6$, $p = 0$ or $p \nmid |H|$}\\
& & \multicolumn{4}{c|}{$V_{78}$} & \multicolumn{4}{c}{$V_{27}$} \\
& & $1$ & $16_{a}$ & $16_{b}$ & $45$ & $1$ & $10_{a}$ & $11$ & $16_{a}$ \\ \hline
\endfirsthead
& & \multicolumn{4}{c|}{$V_{78}$} & \multicolumn{4}{c}{$V_{27}$} \\
& & $1$ & $16_{a}$ & $16_{b}$ & $45$ & $1$ & $10_{a}$ & $11$ & $16_{a}$ \\ \hline
\endhead
1) &  & 1 & 1 & 1 & 1 & 0 & 0 & 1 & 1 \\
2) &  & 1 & 1 & 1 & 1 & 1 & 1 & 0 & 1 \\
\end{longtable}
\noindent Modules tested: V78, V27.\\

\setcounter{table}{56}
\begin{longtable}{rc|ccccc|cccc}
\caption{$M_{11} < E_6$, $p = 11$}\\
& & \multicolumn{5}{c|}{$V_{78}$} & \multicolumn{4}{c}{$V_{27}$} \\
& & $1$ & $9$ & $10_{a}$ & $10_{b}$ & $16$ & $1$ & $9$ & $11$ & $16$ \\ \hline
\endfirsthead
& & \multicolumn{5}{c|}{$V_{78}$} & \multicolumn{4}{c}{$V_{27}$} \\
& & $1$ & $9$ & $10_{a}$ & $10_{b}$ & $16$ & $1$ & $9$ & $11$ & $16$ \\ \hline
\endhead
1) &  & 1 & 1 & 1 & 1 & 3 & 0 & 0 & 1 & 1 \\
2) &  & 1 & 1 & 1 & 1 & 3 & 2 & 1 & 0 & 1 \\
\end{longtable}
\noindent Modules tested: V78, V27.\\

\setcounter{table}{57}
\begin{longtable}{rc|cccc|cccc}
\caption{$M_{11} < E_6$, $p = 5$}\\
& & \multicolumn{4}{c|}{$V_{78}$} & \multicolumn{4}{c}{$V_{27}$} \\
& & $1$ & $16_{a}$ & $16_{b}$ & $45$ & $1$ & $10_{a}$ & $11$ & $16_{a}$ \\ \hline
\endfirsthead
& & \multicolumn{4}{c|}{$V_{78}$} & \multicolumn{4}{c}{$V_{27}$} \\
& & $1$ & $16_{a}$ & $16_{b}$ & $45$ & $1$ & $10_{a}$ & $11$ & $16_{a}$ \\ \hline
\endhead
1) & \possprim, \nongcr & 1 & 1 & 1 & 1 & 0 & 0 & 1 & 1 \\
2) & \nongcr & 1 & 1 & 1 & 1 & 1 & 1 & 0 & 1 \\
\end{longtable}
\noindent Modules tested: V78, V27.\\

\setcounter{table}{58}
\begin{longtable}{rc|ccccccc|cccccc}
\caption{$M_{11} < E_6$, $p = 3$}\\
& & \multicolumn{7}{c|}{$V_{78}$} & \multicolumn{6}{c}{$V_{27}$} \\
& & $1$ & $5_{a}$ & $5_{b}$ & $10_{b}$ & $10_{c}$ & $24$ & $45$ & $1$ & $5_{a}$ & $5_{b}$ & $10_{a}$ & $10_{b}$ & $10_{c}$ \\ \hline
\endfirsthead
& & \multicolumn{7}{c|}{$V_{78}$} & \multicolumn{6}{c}{$V_{27}$} \\
& & $1$ & $5_{a}$ & $5_{b}$ & $10_{b}$ & $10_{c}$ & $24$ & $45$ & $1$ & $5_{a}$ & $5_{b}$ & $10_{a}$ & $10_{b}$ & $10_{c}$ \\ \hline
\endhead
1) & \possprim, \nongcr & 3 & 1 & 1 & 1 & 1 & 0 & 1 & 2 & 1 & 0 & 1 & 0 & 1 \\
2) & \possprim, \nongcr & 3 & 1 & 1 & 1 & 1 & 0 & 1 & 2 & 2 & 1 & 0 & 0 & 1 \\
3) & \possprim, \nongcr & 4 & 1 & 1 & 2 & 2 & 1 & 0 & 2 & 2 & 1 & 0 & 1 & 0 \\
\end{longtable}
\noindent Modules tested: V78, L77, V27.\\

\setcounter{table}{59}
\begin{longtable}{rc|ccccc|ccc}
\caption{$M_{11} < E_6$, $p = 2$}\\
& & \multicolumn{5}{c|}{$V_{78}$} & \multicolumn{3}{c}{$V_{27}$} \\
& & $1$ & $10$ & $16_{a}$ & $16_{b}$ & $44$ & $1$ & $10$ & $16_{a}$ \\ \hline
\endfirsthead
& & \multicolumn{5}{c|}{$V_{78}$} & \multicolumn{3}{c}{$V_{27}$} \\
& & $1$ & $10$ & $16_{a}$ & $16_{b}$ & $44$ & $1$ & $10$ & $16_{a}$ \\ \hline
\endhead
1) &  & 2 & 0 & 1 & 1 & 1 & 1 & 1 & 1 \\
2) & \nongcr & 4 & 3 & 0 & 0 & 1 & 1 & 1 & 1 \\
\end{longtable}
\noindent Modules tested: V78, V27.\\

\setcounter{table}{60}
\begin{longtable}{rc|c|cc}
\caption{$M_{12} < E_6$, $p = 5$}\\
& & \multicolumn{1}{c|}{$V_{78}$} & \multicolumn{2}{c}{$V_{27}$} \\
& & $78$ & $11_{a}$ & $16_{a}$ \\ \hline
\endfirsthead
& & \multicolumn{1}{c|}{$V_{78}$} & \multicolumn{2}{c}{$V_{27}$} \\
& & $78$ & $11_{a}$ & $16_{a}$ \\ \hline
\endhead
1) & \possprim & 1 & 1 & 1 \\
\end{longtable}
\noindent Modules tested: V78, V27.\\

\setcounter{table}{61}
\begin{longtable}{rc|cccc|ccc}
\caption{$M_{12} < E_6$, $p = 2$}\\
& & \multicolumn{4}{c|}{$V_{78}$} & \multicolumn{3}{c}{$V_{27}$} \\
& & $1$ & $16_{a}$ & $16_{b}$ & $44$ & $1$ & $10$ & $16_{a}$ \\ \hline
\endfirsthead
& & \multicolumn{4}{c|}{$V_{78}$} & \multicolumn{3}{c}{$V_{27}$} \\
& & $1$ & $16_{a}$ & $16_{b}$ & $44$ & $1$ & $10$ & $16_{a}$ \\ \hline
\endhead
1) &  & 2 & 1 & 1 & 1 & 1 & 1 & 1 \\
\end{longtable}
\noindent Modules tested: V78, V27.\\

\setcounter{table}{62}
\begin{longtable}{rc|cccc|cc}
\caption{$3\cdot M_{22} < E_6$, $p = 2$}\\
& & \multicolumn{4}{c|}{$V_{78}$} & \multicolumn{2}{c}{$V_{27}$} \\
& & $1$ & $10_{a}$ & $10_{b}$ & $34$ & $6_{a}$ & $15_{a}$ \\ \hline
\endfirsthead
& & \multicolumn{4}{c|}{$V_{78}$} & \multicolumn{2}{c}{$V_{27}$} \\
& & $1$ & $10_{a}$ & $10_{b}$ & $34$ & $6_{a}$ & $15_{a}$ \\ \hline
\endhead
1) & \nongcr & 4 & 2 & 2 & 1 & 2 & 1 \\
\end{longtable}
\noindent Modules tested: V78, V27.\\

\setcounter{table}{63}
\begin{longtable}{rc|cccc|ccc}
\caption{$J_{1} < E_6$, $p = 11$}\\
& & \multicolumn{4}{c|}{$V_{78}$} & \multicolumn{3}{c}{$V_{27}$} \\
& & $1$ & $7$ & $14$ & $64$ & $1$ & $7$ & $27$ \\ \hline
\endfirsthead
& & \multicolumn{4}{c|}{$V_{78}$} & \multicolumn{3}{c}{$V_{27}$} \\
& & $1$ & $7$ & $14$ & $64$ & $1$ & $7$ & $27$ \\ \hline
\endhead
1) & \possprim & 0 & 0 & 1 & 1 & 0 & 0 & 1 \\
2) &  & 8 & 8 & 1 & 0 & 6 & 3 & 0 \\
\end{longtable}
\noindent Modules tested: V78, V27.\\

\setcounter{table}{64}
\begin{longtable}{rc|ccccc|cccc}
\caption{$J_{2} < E_6$, $p = 2$}\\
& & \multicolumn{5}{c|}{$V_{78}$} & \multicolumn{4}{c}{$V_{27}$} \\
& & $1$ & $6_{a}$ & $6_{b}$ & $14_{a}$ & $64_{b}$ & $1$ & $6_{a}$ & $6_{b}$ & $14_{a}$ \\ \hline
\endfirsthead
& & \multicolumn{5}{c|}{$V_{78}$} & \multicolumn{4}{c}{$V_{27}$} \\
& & $1$ & $6_{a}$ & $6_{b}$ & $14_{a}$ & $64_{b}$ & $1$ & $6_{a}$ & $6_{b}$ & $14_{a}$ \\ \hline
\endhead
1) & \nongcr & 8 & 4 & 3 & 2 & 0 & 1 & 2 & 0 & 1 \\
2) & \possprim & 0 & 0 & 0 & 1 & 1 & 1 & 1 & 1 & 1 \\
3) & \nongcr & 16 & 8 & 0 & 1 & 0 & 9 & 3 & 0 & 0 \\
\end{longtable}
\noindent Modules tested: V78, V27.\\

\setcounter{table}{65}
\begin{longtable}{rc|c|cc}
\caption{$3 \cdot J_{3} < E_6$, $p = 2$}\\
& & \multicolumn{1}{c|}{$V_{78}$} & \multicolumn{2}{c}{$V_{27}$} \\
& & $78_{a}$ & $9_{a}$ & $18_{c}$ \\ \hline
\endfirsthead
& & \multicolumn{1}{c|}{$V_{78}$} & \multicolumn{2}{c}{$V_{27}$} \\
& & $78_{a}$ & $9_{a}$ & $18_{c}$ \\ \hline
\endhead
1) & \possprim & 1 & 1 & 1 \\
\end{longtable}
\noindent Modules tested: V78, V27.\\

\setcounter{table}{66}
\begin{longtable}{rc|c|c}
\caption{$3 \cdot Fi_{22} < E_6$, $p = 2$ (irreducible on both)}\\
& & \multicolumn{1}{c|}{$V_{78}$} & \multicolumn{1}{c}{$V_{27}$} \\
& & $78$ & $27_{a}$ \\ \hline
\endfirsthead
& & \multicolumn{1}{c|}{$V_{78}$} & \multicolumn{1}{c}{$V_{27}$} \\
& & $78$ & $27_{a}$ \\ \hline
\endhead
1) & \possprim & 1 & 1 \\
\end{longtable}
\noindent Modules tested: V78, V27.\\

\setcounter{table}{67}
\begin{longtable}{rc|cccccc|cccccc}
\caption{$L_2(7) < E_6$, $p = 0$ or $p \nmid |H|$}\\
& & \multicolumn{6}{c|}{$V_{78}$} & \multicolumn{6}{c}{$V_{27}$} \\
& & $1$ & $3_{a}$ & $3_{b}$ & $6$ & $7$ & $8$ & $1$ & $3_{a}$ & $3_{b}$ & $6$ & $7$ & $8$ \\ \hline
\endfirsthead
& & \multicolumn{6}{c|}{$V_{78}$} & \multicolumn{6}{c}{$V_{27}$} \\
& & $1$ & $3_{a}$ & $3_{b}$ & $6$ & $7$ & $8$ & $1$ & $3_{a}$ & $3_{b}$ & $6$ & $7$ & $8$ \\ \hline
\endhead
1) & \possprim & 0 & 2 & 2 & 2 & 2 & 5 & 0 & 0 & 0 & 2 & 1 & 1 \\
2) &  & 0 & 2 & 2 & 2 & 2 & 5 & 1 & 1 & 1 & 2 & 0 & 1 \\
3) &  & 2 & 1 & 1 & 0 & 2 & 7 & 3 & 0 & 0 & 0 & 0 & 3 \\
4) &  & 3 & 2 & 2 & 2 & 5 & 2 & 0 & 0 & 0 & 2 & 1 & 1 \\
5) &  & 3 & 2 & 2 & 2 & 5 & 2 & 1 & 1 & 1 & 2 & 0 & 1 \\
6) &  & 5 & 1 & 1 & 0 & 5 & 4 & 3 & 0 & 0 & 0 & 0 & 3 \\
7) &  & 8 & 1 & 1 & 0 & 8 & 1 & 6 & 0 & 0 & 0 & 3 & 0 \\
8) &  & 8 & 3 & 3 & 6 & 0 & 2 & 1 & 3 & 3 & 0 & 0 & 1 \\
9) &  & 14 & 0 & 0 & 0 & 0 & 8 & 0 & 7 & 0 & 1 & 0 & 0 \\
10) &  & 16 & 9 & 9 & 0 & 0 & 1 & 9 & 3 & 3 & 0 & 0 & 0 \\
\end{longtable}
\noindent Modules tested: V78, V27.\\

\setcounter{table}{68}
\begin{longtable}{rc|ccccc|ccccc}
\caption{$L_2(7) < E_6$, $p = 3$}\\
& & \multicolumn{5}{c|}{$V_{78}$} & \multicolumn{5}{c}{$V_{27}$} \\
& & $1$ & $3_{a}$ & $3_{b}$ & $6$ & $7$ & $1$ & $3_{a}$ & $3_{b}$ & $6$ & $7$ \\ \hline
\endfirsthead
& & \multicolumn{5}{c|}{$V_{78}$} & \multicolumn{5}{c}{$V_{27}$} \\
& & $1$ & $3_{a}$ & $3_{b}$ & $6$ & $7$ & $1$ & $3_{a}$ & $3_{b}$ & $6$ & $7$ \\ \hline
\endhead
1) & \possprim, \nongcr & 5 & 2 & 2 & 2 & 7 & 1 & 0 & 0 & 2 & 2 \\
2) & \nongcr & 5 & 2 & 2 & 2 & 7 & 2 & 1 & 1 & 2 & 1 \\
3) & \nongcr & 9 & 1 & 1 & 0 & 9 & 6 & 0 & 0 & 0 & 3 \\
4) & \nongcr & 10 & 3 & 3 & 6 & 2 & 2 & 3 & 3 & 0 & 1 \\
5) &  & 17 & 9 & 9 & 0 & 1 & 9 & 3 & 3 & 0 & 0 \\
6) & \nongcr & 22 & 0 & 0 & 0 & 8 & 0 & 7 & 0 & 1 & 0 \\
\end{longtable}
\noindent Modules tested: V78, L77, V27.\\

\setcounter{table}{69}
\begin{longtable}{rc|ccccccccc|ccccccc}
\caption{$L_2(8) < E_6$, $p = 0$ or $p \nmid |H|$}\\
& & \multicolumn{9}{c|}{$V_{78}$} & \multicolumn{7}{c}{$V_{27}$} \\
& & $1$ & $7_{a}$ & $7_{b}$ & $7_{c}$ & $7_{d}$ & $8$ & $9_{a}$ & $9_{b}$ & $9_{c}$ & $1$ & $7_{a}$ & $7_{b}$ & $8$ & $9_{a}$ & $9_{b}$ & $9_{c}$ \\ \hline
\endfirsthead
& & \multicolumn{9}{c|}{$V_{78}$} & \multicolumn{7}{c}{$V_{27}$} \\
& & $1$ & $7_{a}$ & $7_{b}$ & $7_{c}$ & $7_{d}$ & $8$ & $9_{a}$ & $9_{b}$ & $9_{c}$ & $1$ & $7_{a}$ & $7_{b}$ & $8$ & $9_{a}$ & $9_{b}$ & $9_{c}$ \\ \hline
\endhead
1) &  & 0 & 2 & 1 & 1 & 1 & 2 & 2 & 0 & 1 & 1 & 0 & 0 & 1 & 1 & 0 & 1 \\
2) & \possprim & 0 & 2 & 1 & 1 & 1 & 2 & 1 & 1 & 1 & 0 & 0 & 0 & 0 & 1 & 1 & 1 \\
3) &  & 1 & 3 & 1 & 1 & 1 & 1 & 2 & 0 & 1 & 2 & 1 & 0 & 0 & 1 & 0 & 1 \\
4) &  & 3 & 2 & 1 & 1 & 1 & 5 & 0 & 0 & 0 & 3 & 0 & 0 & 3 & 0 & 0 & 0 \\
5) &  & 4 & 3 & 1 & 1 & 1 & 4 & 0 & 0 & 0 & 4 & 1 & 0 & 2 & 0 & 0 & 0 \\
6) &  & 8 & 1 & 8 & 1 & 0 & 0 & 0 & 0 & 0 & 6 & 0 & 3 & 0 & 0 & 0 & 0 \\
\end{longtable}
\noindent Modules tested: V78, V27.\\

\setcounter{table}{70}
\begin{longtable}{rc|cccccc|cccc}
\caption{$L_2(8) < E_6$, $p = 7$}\\
& & \multicolumn{6}{c|}{$V_{78}$} & \multicolumn{4}{c}{$V_{27}$} \\
& & $1$ & $7_{a}$ & $7_{b}$ & $7_{c}$ & $7_{d}$ & $8$ & $1$ & $7_{a}$ & $7_{b}$ & $8$ \\ \hline
\endfirsthead
& & \multicolumn{6}{c|}{$V_{78}$} & \multicolumn{4}{c}{$V_{27}$} \\
& & $1$ & $7_{a}$ & $7_{b}$ & $7_{c}$ & $7_{d}$ & $8$ & $1$ & $7_{a}$ & $7_{b}$ & $8$ \\ \hline
\endhead
1) & \nongcr & 3 & 2 & 1 & 1 & 1 & 5 & 3 & 0 & 0 & 3 \\
2) & \nongcr & 4 & 3 & 1 & 1 & 1 & 4 & 4 & 1 & 0 & 2 \\
3) &  & 8 & 1 & 8 & 1 & 0 & 0 & 6 & 0 & 3 & 0 \\
\end{longtable}
\noindent Modules tested: V78, V27.\\

\setcounter{table}{71}
\begin{longtable}{rc|ccccc|ccccc}
\caption{$L_2(8) < E_6$, $p = 3$}\\
& & \multicolumn{5}{c|}{$V_{78}$} & \multicolumn{5}{c}{$V_{27}$} \\
& & $1$ & $7$ & $9_{a}$ & $9_{b}$ & $9_{c}$ & $1$ & $7$ & $9_{a}$ & $9_{b}$ & $9_{c}$ \\ \hline
\endfirsthead
& & \multicolumn{5}{c|}{$V_{78}$} & \multicolumn{5}{c}{$V_{27}$} \\
& & $1$ & $7$ & $9_{a}$ & $9_{b}$ & $9_{c}$ & $1$ & $7$ & $9_{a}$ & $9_{b}$ & $9_{c}$ \\ \hline
\endhead
1) & \nongcr & 2 & 7 & 2 & 0 & 1 & 2 & 1 & 1 & 0 & 1 \\
2) & \possprim, \nongcr & 2 & 7 & 1 & 1 & 1 & 0 & 0 & 1 & 1 & 1 \\
3) & \nongcr & 8 & 10 & 0 & 0 & 0 & 6 & 3 & 0 & 0 & 0 \\
\end{longtable}
\noindent Modules tested: V78, L77, V27.\\

\setcounter{table}{72}
\begin{longtable}{rc|ccccccc|ccccccc}
\caption{$L_2(11) < E_6$, $p = 0$ or $p \nmid |H|$}\\
& & \multicolumn{7}{c|}{$V_{78}$} & \multicolumn{7}{c}{$V_{27}$} \\
& & $1$ & $5_{a}$ & $5_{b}$ & $10_{a}$ & $11$ & $12_{a}$ & $12_{b}$ & $1$ & $5_{a}$ & $5_{b}$ & $10_{a}$ & $10_{b}$ & $11$ & $12_{b}$ \\ \hline
\endfirsthead
& & \multicolumn{7}{c|}{$V_{78}$} & \multicolumn{7}{c}{$V_{27}$} \\
& & $1$ & $5_{a}$ & $5_{b}$ & $10_{a}$ & $11$ & $12_{a}$ & $12_{b}$ & $1$ & $5_{a}$ & $5_{b}$ & $10_{a}$ & $10_{b}$ & $11$ & $12_{b}$ \\ \hline
\endhead
1) & \possprim & 0 & 1 & 1 & 1 & 2 & 2 & 1 & 0 & 1 & 0 & 0 & 1 & 0 & 1 \\
2) &  & 1 & 1 & 1 & 1 & 3 & 1 & 1 & 1 & 1 & 0 & 0 & 1 & 1 & 0 \\
3) &  & 4 & 1 & 1 & 4 & 0 & 1 & 1 & 2 & 2 & 1 & 1 & 0 & 0 & 0 \\
\end{longtable}
\noindent Modules tested: V78, V27.\\

\setcounter{table}{73}
\begin{longtable}{rc|ccccc|cccccc}
\caption{$L_2(11) < E_6$, $p = 5$}\\
& & \multicolumn{5}{c|}{$V_{78}$} & \multicolumn{6}{c}{$V_{27}$} \\
& & $1$ & $5_{a}$ & $5_{b}$ & $10_{a}$ & $11$ & $1$ & $5_{a}$ & $5_{b}$ & $10_{a}$ & $10_{b}$ & $11$ \\ \hline
\endfirsthead
& & \multicolumn{5}{c|}{$V_{78}$} & \multicolumn{6}{c}{$V_{27}$} \\
& & $1$ & $5_{a}$ & $5_{b}$ & $10_{a}$ & $11$ & $1$ & $5_{a}$ & $5_{b}$ & $10_{a}$ & $10_{b}$ & $11$ \\ \hline
\endhead
1) & \nongcr & 3 & 1 & 1 & 1 & 5 & 1 & 1 & 0 & 0 & 1 & 1 \\
2) & \nongcr & 6 & 1 & 1 & 4 & 2 & 2 & 2 & 1 & 1 & 0 & 0 \\
\end{longtable}
\noindent Modules tested: V78, V27.\\

\setcounter{table}{74}
\begin{longtable}{rc|cccccc|ccccc}
\caption{$L_2(11) < E_6$, $p = 3$}\\
& & \multicolumn{6}{c|}{$V_{78}$} & \multicolumn{5}{c}{$V_{27}$} \\
& & $1$ & $5_{a}$ & $5_{b}$ & $10$ & $12_{a}$ & $12_{b}$ & $1$ & $5_{a}$ & $5_{b}$ & $10$ & $12_{b}$ \\ \hline
\endfirsthead
& & \multicolumn{6}{c|}{$V_{78}$} & \multicolumn{5}{c}{$V_{27}$} \\
& & $1$ & $5_{a}$ & $5_{b}$ & $10$ & $12_{a}$ & $12_{b}$ & $1$ & $5_{a}$ & $5_{b}$ & $10$ & $12_{b}$ \\ \hline
\endhead
1) & \possprim, \nongcr & 2 & 1 & 1 & 3 & 2 & 1 & 0 & 2 & 1 & 0 & 1 \\
2) & \nongcr & 4 & 1 & 1 & 4 & 1 & 1 & 2 & 2 & 1 & 1 & 0 \\
\end{longtable}
\noindent Modules tested: V78, L77, V27.\\

\setcounter{table}{75}
\begin{longtable}{rc|cccccc|ccccc}
\caption{$L_2(11) < E_6$, $p = 2$}\\
& & \multicolumn{6}{c|}{$V_{78}$} & \multicolumn{5}{c}{$V_{27}$} \\
& & $1$ & $5_{a}$ & $5_{b}$ & $10$ & $12_{a}$ & $12_{b}$ & $1$ & $5_{a}$ & $5_{b}$ & $10$ & $12_{b}$ \\ \hline
\endfirsthead
& & \multicolumn{6}{c|}{$V_{78}$} & \multicolumn{5}{c}{$V_{27}$} \\
& & $1$ & $5_{a}$ & $5_{b}$ & $10$ & $12_{a}$ & $12_{b}$ & $1$ & $5_{a}$ & $5_{b}$ & $10$ & $12_{b}$ \\ \hline
\endhead
1) &  & 2 & 0 & 0 & 4 & 2 & 1 & 0 & 1 & 0 & 1 & 1 \\
2) & \possprim, \nongcr & 2 & 3 & 3 & 1 & 2 & 1 & 0 & 1 & 0 & 1 & 1 \\
3) & \nongcr & 4 & 1 & 1 & 4 & 1 & 1 & 2 & 2 & 1 & 1 & 0 \\
4) & \possprim, \nongcr & 4 & 4 & 4 & 1 & 1 & 1 & 2 & 2 & 1 & 1 & 0 \\
\end{longtable}
\noindent Modules tested: V78, V27.\\

\setcounter{table}{76}
\begin{longtable}{rc|ccccccccc|ccccc}
\caption{$L_2(13) < E_6$, $p = 0$ or $p \nmid |H|$}\\
& & \multicolumn{9}{c|}{$V_{78}$} & \multicolumn{5}{c}{$V_{27}$} \\
& & $1$ & $7_{a}$ & $7_{b}$ & $12_{a}$ & $12_{b}$ & $12_{c}$ & $13$ & $14_{a}$ & $14_{b}$ & $1$ & $7_{a}$ & $12_{a}$ & $13$ & $14_{a}$ \\ \hline
\endfirsthead
& & \multicolumn{9}{c|}{$V_{78}$} & \multicolumn{5}{c}{$V_{27}$} \\
& & $1$ & $7_{a}$ & $7_{b}$ & $12_{a}$ & $12_{b}$ & $12_{c}$ & $13$ & $14_{a}$ & $14_{b}$ & $1$ & $7_{a}$ & $12_{a}$ & $13$ & $14_{a}$ \\ \hline
\endhead
1) & \possprim & 0 & 0 & 0 & 1 & 1 & 1 & 0 & 1 & 2 & 0 & 0 & 0 & 1 & 1 \\
2) &  & 0 & 0 & 0 & 2 & 1 & 0 & 0 & 1 & 2 & 1 & 0 & 1 & 0 & 1 \\
3) &  & 1 & 2 & 2 & 1 & 1 & 1 & 1 & 0 & 0 & 0 & 0 & 0 & 1 & 1 \\
4) &  & 1 & 2 & 2 & 2 & 1 & 0 & 1 & 0 & 0 & 1 & 0 & 1 & 0 & 1 \\
5) &  & 8 & 8 & 0 & 0 & 0 & 0 & 0 & 0 & 1 & 6 & 3 & 0 & 0 & 0 \\
\end{longtable}
\noindent Modules tested: V78, V27.\\

\setcounter{table}{77}
\begin{longtable}{rc|cccccc|cccc}
\caption{$L_2(13) < E_6$, $p = 7$}\\
& & \multicolumn{6}{c|}{$V_{78}$} & \multicolumn{4}{c}{$V_{27}$} \\
& & $1$ & $7_{a}$ & $7_{b}$ & $12$ & $14_{a}$ & $14_{b}$ & $1$ & $7_{a}$ & $12$ & $14_{a}$ \\ \hline
\endfirsthead
& & \multicolumn{6}{c|}{$V_{78}$} & \multicolumn{4}{c}{$V_{27}$} \\
& & $1$ & $7_{a}$ & $7_{b}$ & $12$ & $14_{a}$ & $14_{b}$ & $1$ & $7_{a}$ & $12$ & $14_{a}$ \\ \hline
\endhead
1) &  & 0 & 0 & 0 & 3 & 1 & 2 & 1 & 0 & 1 & 1 \\
2) & \nongcr & 2 & 2 & 2 & 4 & 0 & 0 & 1 & 0 & 1 & 1 \\
3) & \nongcr & 8 & 8 & 0 & 0 & 0 & 1 & 6 & 3 & 0 & 0 \\
\end{longtable}
\noindent Modules tested: V78, V27.\\

\setcounter{table}{78}
\begin{longtable}{rc|ccccccc|cccc}
\caption{$L_2(13) < E_6$, $p = 3$}\\
& & \multicolumn{7}{c|}{$V_{78}$} & \multicolumn{4}{c}{$V_{27}$} \\
& & $1$ & $7_{a}$ & $7_{b}$ & $12_{a}$ & $12_{b}$ & $12_{c}$ & $13$ & $1$ & $7_{a}$ & $12_{a}$ & $13$ \\ \hline
\endfirsthead
& & \multicolumn{7}{c|}{$V_{78}$} & \multicolumn{4}{c}{$V_{27}$} \\
& & $1$ & $7_{a}$ & $7_{b}$ & $12_{a}$ & $12_{b}$ & $12_{c}$ & $13$ & $1$ & $7_{a}$ & $12_{a}$ & $13$ \\ \hline
\endhead
1) & \possprim & 1 & 2 & 2 & 1 & 1 & 1 & 1 & 1 & 0 & 0 & 2 \\
2) &  & 1 & 2 & 2 & 2 & 1 & 0 & 1 & 2 & 0 & 1 & 1 \\
3) &  & 8 & 9 & 1 & 0 & 0 & 0 & 0 & 6 & 3 & 0 & 0 \\
\end{longtable}
\noindent Modules tested: V78, L77, V27.\\

\setcounter{table}{79}
\begin{longtable}{rc|ccccccc|ccccc}
\caption{$L_2(13) < E_6$, $p = 2$}\\
& & \multicolumn{7}{c|}{$V_{78}$} & \multicolumn{5}{c}{$V_{27}$} \\
& & $1$ & $6_{a}$ & $6_{b}$ & $12_{a}$ & $12_{b}$ & $12_{c}$ & $14$ & $1$ & $6_{a}$ & $6_{b}$ & $12_{a}$ & $14$ \\ \hline
\endfirsthead
& & \multicolumn{7}{c|}{$V_{78}$} & \multicolumn{5}{c}{$V_{27}$} \\
& & $1$ & $6_{a}$ & $6_{b}$ & $12_{a}$ & $12_{b}$ & $12_{c}$ & $14$ & $1$ & $6_{a}$ & $6_{b}$ & $12_{a}$ & $14$ \\ \hline
\endhead
1) & \possprim & 0 & 0 & 0 & 1 & 1 & 1 & 3 & 1 & 1 & 1 & 0 & 1 \\
2) &  & 0 & 0 & 0 & 2 & 1 & 0 & 3 & 1 & 0 & 0 & 1 & 1 \\
3) &  & 2 & 1 & 0 & 0 & 0 & 0 & 5 & 1 & 2 & 0 & 0 & 1 \\
4) & \nongcr & 6 & 3 & 3 & 1 & 1 & 1 & 0 & 1 & 1 & 1 & 0 & 1 \\
5) & \nongcr & 6 & 3 & 3 & 2 & 1 & 0 & 0 & 1 & 0 & 0 & 1 & 1 \\
6) & \nongcr & 8 & 4 & 3 & 0 & 0 & 0 & 2 & 1 & 2 & 0 & 0 & 1 \\
7) & \nongcr & 16 & 8 & 0 & 0 & 0 & 0 & 1 & 9 & 3 & 0 & 0 & 0 \\
\end{longtable}
\noindent Modules tested: V78, V27.\\

\setcounter{table}{80}
\begin{longtable}{rc|cccccc|ccccc}
\caption{$L_2(17) < E_6$, $p = 0$ or $p \nmid |H|$}\\
& & \multicolumn{6}{c|}{$V_{78}$} & \multicolumn{5}{c}{$V_{27}$} \\
& & $1$ & $9_{a}$ & $16_{a}$ & $17$ & $18_{b}$ & $18_{c}$ & $1$ & $9_{a}$ & $16_{a}$ & $17$ & $18_{a}$ \\ \hline
\endfirsthead
& & \multicolumn{6}{c|}{$V_{78}$} & \multicolumn{5}{c}{$V_{27}$} \\
& & $1$ & $9_{a}$ & $16_{a}$ & $17$ & $18_{b}$ & $18_{c}$ & $1$ & $9_{a}$ & $16_{a}$ & $17$ & $18_{a}$ \\ \hline
\endhead
1) & \possprim & 0 & 1 & 1 & 1 & 1 & 1 & 0 & 1 & 0 & 0 & 1 \\
2) &  & 0 & 1 & 1 & 1 & 1 & 1 & 1 & 1 & 0 & 1 & 0 \\
3) &  & 1 & 1 & 2 & 0 & 1 & 1 & 2 & 1 & 1 & 0 & 0 \\
\end{longtable}
\noindent Modules tested: V78, V27.\\

\setcounter{table}{81}
\begin{longtable}{rc|ccccc|cccc}
\caption{$L_2(17) < E_6$, $p = 3$}\\
& & \multicolumn{5}{c|}{$V_{78}$} & \multicolumn{4}{c}{$V_{27}$} \\
& & $1$ & $9_{a}$ & $16$ & $18_{b}$ & $18_{c}$ & $1$ & $9_{a}$ & $16$ & $18_{a}$ \\ \hline
\endfirsthead
& & \multicolumn{5}{c|}{$V_{78}$} & \multicolumn{4}{c}{$V_{27}$} \\
& & $1$ & $9_{a}$ & $16$ & $18_{b}$ & $18_{c}$ & $1$ & $9_{a}$ & $16$ & $18_{a}$ \\ \hline
\endhead
1) & \possprim, \nongcr & 1 & 1 & 2 & 1 & 1 & 0 & 1 & 0 & 1 \\
2) & \nongcr & 1 & 1 & 2 & 1 & 1 & 2 & 1 & 1 & 0 \\
\end{longtable}
\noindent Modules tested: V78, L77, V27.\\

\setcounter{table}{82}
\begin{longtable}{rc|cccc|cccc}
\caption{$L_2(17) < E_6$, $p = 2$}\\
& & \multicolumn{4}{c|}{$V_{78}$} & \multicolumn{4}{c}{$V_{27}$} \\
& & $1$ & $8_{a}$ & $8_{b}$ & $16_{a}$ & $1$ & $8_{a}$ & $8_{b}$ & $16_{a}$ \\ \hline
\endfirsthead
& & \multicolumn{4}{c|}{$V_{78}$} & \multicolumn{4}{c}{$V_{27}$} \\
& & $1$ & $8_{a}$ & $8_{b}$ & $16_{a}$ & $1$ & $8_{a}$ & $8_{b}$ & $16_{a}$ \\ \hline
\endhead
1) & \nongcr & 6 & 3 & 2 & 2 & 3 & 1 & 0 & 1 \\
2) & \nongcr & 6 & 4 & 3 & 1 & 3 & 2 & 1 & 0 \\
\end{longtable}
\noindent Modules tested: V78, V27.\\

\setcounter{table}{83}
\begin{longtable}{rc|cccc|cc}
\caption{$L_2(19) < E_6$, $p = 0$ or $p \nmid |H|$}\\
& & \multicolumn{4}{c|}{$V_{78}$} & \multicolumn{2}{c}{$V_{27}$} \\
& & $18_{a}$ & $20_{b}$ & $20_{c}$ & $20_{d}$ & $9_{a}$ & $18_{d}$ \\ \hline
\endfirsthead
& & \multicolumn{4}{c|}{$V_{78}$} & \multicolumn{2}{c}{$V_{27}$} \\
& & $18_{a}$ & $20_{b}$ & $20_{c}$ & $20_{d}$ & $9_{a}$ & $18_{d}$ \\ \hline
\endhead
1) & \possprim & 1 & 1 & 1 & 1 & 1 & 1 \\
\end{longtable}
\noindent Modules tested: V78, V27.\\

\setcounter{table}{84}
\begin{longtable}{rc|cccc|cc}
\caption{$L_2(19) < E_6$, $p = 5$}\\
& & \multicolumn{4}{c|}{$V_{78}$} & \multicolumn{2}{c}{$V_{27}$} \\
& & $18$ & $20_{b}$ & $20_{c}$ & $20_{d}$ & $9_{a}$ & $9_{b}$ \\ \hline
\endfirsthead
& & \multicolumn{4}{c|}{$V_{78}$} & \multicolumn{2}{c}{$V_{27}$} \\
& & $18$ & $20_{b}$ & $20_{c}$ & $20_{d}$ & $9_{a}$ & $9_{b}$ \\ \hline
\endhead
1) & \possprim & 1 & 1 & 1 & 1 & 2 & 1 \\
\end{longtable}
\noindent Modules tested: V78, V27.\\

\setcounter{table}{85}
\begin{longtable}{rc|ccc|cc}
\caption{$L_2(19) < E_6$, $p = 3$}\\
& & \multicolumn{3}{c|}{$V_{78}$} & \multicolumn{2}{c}{$V_{27}$} \\
& & $1$ & $18_{a}$ & $19$ & $9_{a}$ & $18_{d}$ \\ \hline
\endfirsthead
& & \multicolumn{3}{c|}{$V_{78}$} & \multicolumn{2}{c}{$V_{27}$} \\
& & $1$ & $18_{a}$ & $19$ & $9_{a}$ & $18_{d}$ \\ \hline
\endhead
1) & \possprim & 3 & 1 & 3 & 1 & 1 \\
\end{longtable}
\noindent Modules tested: V78, L77, V27.\\

\setcounter{table}{86}
\begin{longtable}{rc|cccc|cc}
\caption{$L_2(19) < E_6$, $p = 2$}\\
& & \multicolumn{4}{c|}{$V_{78}$} & \multicolumn{2}{c}{$V_{27}$} \\
& & $18_{a}$ & $20_{b}$ & $20_{c}$ & $20_{d}$ & $9_{a}$ & $18_{b}$ \\ \hline
\endfirsthead
& & \multicolumn{4}{c|}{$V_{78}$} & \multicolumn{2}{c}{$V_{27}$} \\
& & $18_{a}$ & $20_{b}$ & $20_{c}$ & $20_{d}$ & $9_{a}$ & $18_{b}$ \\ \hline
\endhead
1) & \possprim & 1 & 1 & 1 & 1 & 1 & 1 \\
\end{longtable}
\noindent Modules tested: V78, V27.\\

\setcounter{table}{87}
\begin{longtable}{rc|cccccc|ccc}
\caption{$L_2(25) < E_6$, $p = 0$ or $p \nmid |H|$}\\
& & \multicolumn{6}{c|}{$V_{78}$} & \multicolumn{3}{c}{$V_{27}$} \\
& & $1$ & $25$ & $26_{a}$ & $26_{b}$ & $26_{d}$ & $26_{e}$ & $1$ & $26_{a}$ & $26_{c}$ \\ \hline
\endfirsthead
& & \multicolumn{6}{c|}{$V_{78}$} & \multicolumn{3}{c}{$V_{27}$} \\
& & $1$ & $25$ & $26_{a}$ & $26_{b}$ & $26_{d}$ & $26_{e}$ & $1$ & $26_{a}$ & $26_{c}$ \\ \hline
\endhead
1) &  & 0 & 0 & 1 & 0 & 1 & 1 & 1 & 1 & 0 \\
2) &  & 1 & 1 & 0 & 2 & 0 & 0 & 1 & 1 & 0 \\
3) &  & 1 & 1 & 0 & 2 & 0 & 0 & 1 & 0 & 1 \\
\end{longtable}
\noindent Modules tested: V78, V27.\\

\setcounter{table}{88}
\begin{longtable}{rc|cccccc|ccc}
\caption{$L_2(25) < E_6$, $p = 13$}\\
& & \multicolumn{6}{c|}{$V_{78}$} & \multicolumn{3}{c}{$V_{27}$} \\
& & $1$ & $24$ & $26_{a}$ & $26_{b}$ & $26_{d}$ & $26_{e}$ & $1$ & $26_{a}$ & $26_{c}$ \\ \hline
\endfirsthead
& & \multicolumn{6}{c|}{$V_{78}$} & \multicolumn{3}{c}{$V_{27}$} \\
& & $1$ & $24$ & $26_{a}$ & $26_{b}$ & $26_{d}$ & $26_{e}$ & $1$ & $26_{a}$ & $26_{c}$ \\ \hline
\endhead
1) &  & 0 & 0 & 1 & 0 & 1 & 1 & 1 & 1 & 0 \\
2) &  & 2 & 1 & 0 & 2 & 0 & 0 & 1 & 1 & 0 \\
3) &  & 2 & 1 & 0 & 2 & 0 & 0 & 1 & 0 & 1 \\
\end{longtable}
\noindent Modules tested: V78, V27.\\

\setcounter{table}{89}
\begin{longtable}{rc|ccc|cccc}
\caption{$L_2(25) < E_6$, $p = 3$}\\
& & \multicolumn{3}{c|}{$V_{78}$} & \multicolumn{4}{c}{$V_{27}$} \\
& & $1$ & $25$ & $26$ & $1$ & $13_{a}$ & $13_{b}$ & $25$ \\ \hline
\endfirsthead
& & \multicolumn{3}{c|}{$V_{78}$} & \multicolumn{4}{c}{$V_{27}$} \\
& & $1$ & $25$ & $26$ & $1$ & $13_{a}$ & $13_{b}$ & $25$ \\ \hline
\endhead
1) &  & 1 & 1 & 2 & 1 & 1 & 1 & 0 \\
2) &  & 1 & 1 & 2 & 2 & 0 & 0 & 1 \\
\end{longtable}
\noindent Modules tested: V78, L77, V27.\\

\setcounter{table}{90}
\begin{longtable}{rc|cccc|cc}
\caption{$L_2(25) < E_6$, $p = 2$}\\
& & \multicolumn{4}{c|}{$V_{78}$} & \multicolumn{2}{c}{$V_{27}$} \\
& & $1$ & $12_{a}$ & $12_{b}$ & $26$ & $1$ & $26$ \\ \hline
\endfirsthead
& & \multicolumn{4}{c|}{$V_{78}$} & \multicolumn{2}{c}{$V_{27}$} \\
& & $1$ & $12_{a}$ & $12_{b}$ & $26$ & $1$ & $26$ \\ \hline
\endhead
1) &  & 0 & 0 & 0 & 3 & 1 & 1 \\
2) &  & 6 & 3 & 3 & 0 & 1 & 1 \\
\end{longtable}
\noindent Modules tested: V78, V27.\\

\setcounter{table}{91}
\begin{longtable}{rc|ccc|cc}
\caption{$L_2(27) < E_6$, $p = 0$ or $p \nmid |H|$}\\
& & \multicolumn{3}{c|}{$V_{78}$} & \multicolumn{2}{c}{$V_{27}$} \\
& & $26_{a}$ & $26_{b}$ & $26_{e}$ & $1$ & $26_{e}$ \\ \hline
\endfirsthead
& & \multicolumn{3}{c|}{$V_{78}$} & \multicolumn{2}{c}{$V_{27}$} \\
& & $26_{a}$ & $26_{b}$ & $26_{e}$ & $1$ & $26_{e}$ \\ \hline
\endhead
1) &  & 1 & 1 & 1 & 1 & 1 \\
\end{longtable}
\noindent Modules tested: V78, V27.\\

\setcounter{table}{92}
\begin{longtable}{rc|ccc|cc}
\caption{$L_2(27) < E_6$, $p = 13$}\\
& & \multicolumn{3}{c|}{$V_{78}$} & \multicolumn{2}{c}{$V_{27}$} \\
& & $26_{a}$ & $26_{b}$ & $26_{e}$ & $1$ & $26_{e}$ \\ \hline
\endfirsthead
& & \multicolumn{3}{c|}{$V_{78}$} & \multicolumn{2}{c}{$V_{27}$} \\
& & $26_{a}$ & $26_{b}$ & $26_{e}$ & $1$ & $26_{e}$ \\ \hline
\endhead
1) &  & 1 & 1 & 1 & 1 & 1 \\
\end{longtable}
\noindent Modules tested: V78, V27.\\

\setcounter{table}{93}
\begin{longtable}{rc|ccc|ccc}
\caption{$L_2(27) < E_6$, $p = 7$}\\
& & \multicolumn{3}{c|}{$V_{78}$} & \multicolumn{3}{c}{$V_{27}$} \\
& & $13_{a}$ & $13_{b}$ & $26$ & $1$ & $13_{a}$ & $13_{b}$ \\ \hline
\endfirsthead
& & \multicolumn{3}{c|}{$V_{78}$} & \multicolumn{3}{c}{$V_{27}$} \\
& & $13_{a}$ & $13_{b}$ & $26$ & $1$ & $13_{a}$ & $13_{b}$ \\ \hline
\endhead
1) &  & 1 & 1 & 2 & 1 & 1 & 1 \\
\end{longtable}
\noindent Modules tested: V78, V27.\\

\setcounter{table}{94}
\begin{longtable}{rc|ccc|cccc}
\caption{$L_2(27) < E_6$, $p = 2$}\\
& & \multicolumn{3}{c|}{$V_{78}$} & \multicolumn{4}{c}{$V_{27}$} \\
& & $26_{a}$ & $26_{b}$ & $26_{c}$ & $1$ & $13_{a}$ & $13_{b}$ & $26_{a}$ \\ \hline
\endfirsthead
& & \multicolumn{3}{c|}{$V_{78}$} & \multicolumn{4}{c}{$V_{27}$} \\
& & $26_{a}$ & $26_{b}$ & $26_{c}$ & $1$ & $13_{a}$ & $13_{b}$ & $26_{a}$ \\ \hline
\endhead
1) &  & 1 & 1 & 1 & 1 & 1 & 1 & 0 \\
2) &  & 2 & 0 & 1 & 1 & 0 & 0 & 1 \\
\end{longtable}
\noindent Modules tested: V78, V27.\\

\setcounter{table}{95}
\begin{longtable}{rc|ccccccccc|ccc}
\caption{$L_3(3) < E_6$, $p = 0$ or $p \nmid |H|$}\\
& & \multicolumn{9}{c|}{$V_{78}$} & \multicolumn{3}{c}{$V_{27}$} \\
& & $1$ & $13$ & $16_{a}$ & $16_{b}$ & $16_{c}$ & $16_{d}$ & $26_{a}$ & $26_{b}$ & $26_{c}$ & $1$ & $26_{a}$ & $27$ \\ \hline
\endfirsthead
& & \multicolumn{9}{c|}{$V_{78}$} & \multicolumn{3}{c}{$V_{27}$} \\
& & $1$ & $13$ & $16_{a}$ & $16_{b}$ & $16_{c}$ & $16_{d}$ & $26_{a}$ & $26_{b}$ & $26_{c}$ & $1$ & $26_{a}$ & $27$ \\ \hline
\endhead
1) & \possprim & 0 & 0 & 0 & 0 & 0 & 0 & 1 & 1 & 1 & 0 & 0 & 1 \\
2) &  & 0 & 0 & 0 & 0 & 0 & 0 & 1 & 1 & 1 & 1 & 1 & 0 \\
3) &  & 1 & 1 & 1 & 1 & 1 & 1 & 0 & 0 & 0 & 0 & 0 & 1 \\
4) &  & 1 & 1 & 1 & 1 & 1 & 1 & 0 & 0 & 0 & 1 & 1 & 0 \\
\end{longtable}
\noindent Modules tested: V78, V27.\\

\setcounter{table}{96}
\begin{longtable}{rc|cccccc|cccc}
\caption{$L_3(3) < E_6$, $p = 13$}\\
& & \multicolumn{6}{c|}{$V_{78}$} & \multicolumn{4}{c}{$V_{27}$} \\
& & $1$ & $13$ & $16$ & $26_{a}$ & $26_{b}$ & $26_{c}$ & $1$ & $11$ & $16$ & $26_{a}$ \\ \hline
\endfirsthead
& & \multicolumn{6}{c|}{$V_{78}$} & \multicolumn{4}{c}{$V_{27}$} \\
& & $1$ & $13$ & $16$ & $26_{a}$ & $26_{b}$ & $26_{c}$ & $1$ & $11$ & $16$ & $26_{a}$ \\ \hline
\endhead
1) & \possprim & 0 & 0 & 0 & 1 & 1 & 1 & 0 & 1 & 1 & 0 \\
2) &  & 0 & 0 & 0 & 1 & 1 & 1 & 1 & 0 & 0 & 1 \\
3) &  & 1 & 1 & 4 & 0 & 0 & 0 & 0 & 1 & 1 & 0 \\
4) &  & 1 & 1 & 4 & 0 & 0 & 0 & 1 & 0 & 0 & 1 \\
\end{longtable}
\noindent Modules tested: V78, V27.\\

\setcounter{table}{97}
\begin{longtable}{rc|ccccccc|cc}
\caption{$L_3(3) < E_6$, $p = 2$}\\
& & \multicolumn{7}{c|}{$V_{78}$} & \multicolumn{2}{c}{$V_{27}$} \\
& & $1$ & $12$ & $16_{a}$ & $16_{b}$ & $16_{c}$ & $16_{d}$ & $26$ & $1$ & $26$ \\ \hline
\endfirsthead
& & \multicolumn{7}{c|}{$V_{78}$} & \multicolumn{2}{c}{$V_{27}$} \\
& & $1$ & $12$ & $16_{a}$ & $16_{b}$ & $16_{c}$ & $16_{d}$ & $26$ & $1$ & $26$ \\ \hline
\endhead
1) &  & 0 & 0 & 0 & 0 & 0 & 0 & 3 & 1 & 1 \\
2) &  & 2 & 1 & 1 & 1 & 1 & 1 & 0 & 1 & 1 \\
\end{longtable}
\noindent Modules tested: V78, V27.\\

\setcounter{table}{98}
\begin{longtable}{rc|cc|cc}
\caption{$L_4(3) < E_6$, $p = 2$}\\
& & \multicolumn{2}{c|}{$V_{78}$} & \multicolumn{2}{c}{$V_{27}$} \\
& & $26_{a}$ & $26_{b}$ & $1$ & $26_{a}$ \\ \hline
\endfirsthead
& & \multicolumn{2}{c|}{$V_{78}$} & \multicolumn{2}{c}{$V_{27}$} \\
& & $26_{a}$ & $26_{b}$ & $1$ & $26_{a}$ \\ \hline
\endhead
1) &  & 2 & 1 & 1 & 1 \\
\end{longtable}
\noindent Modules tested: V78, V27.\\

\setcounter{table}{99}
\begin{longtable}{rc|ccccccccccc|ccccc}
\caption{$U_3(3) < E_6$, $p = 0$ or $p \nmid |H|$}\\
& & \multicolumn{11}{c|}{$V_{78}$} & \multicolumn{5}{c}{$V_{27}$} \\
& & $1$ & $6$ & $7_{a}$ & $7_{b}$ & $7_{c}$ & $14$ & $21_{a}$ & $28_{a}$ & $28_{b}$ & $32_{a}$ & $32_{b}$ & $1$ & $6$ & $7_{a}$ & $14$ & $27$ \\ \hline
\endfirsthead
& & \multicolumn{11}{c|}{$V_{78}$} & \multicolumn{5}{c}{$V_{27}$} \\
& & $1$ & $6$ & $7_{a}$ & $7_{b}$ & $7_{c}$ & $14$ & $21_{a}$ & $28_{a}$ & $28_{b}$ & $32_{a}$ & $32_{b}$ & $1$ & $6$ & $7_{a}$ & $14$ & $27$ \\ \hline
\endhead
1) & \possprim & 0 & 0 & 0 & 0 & 0 & 1 & 0 & 0 & 0 & 1 & 1 & 0 & 0 & 0 & 0 & 1 \\
2) &  & 1 & 0 & 1 & 1 & 1 & 0 & 0 & 1 & 1 & 0 & 0 & 0 & 0 & 0 & 0 & 1 \\
3) &  & 3 & 2 & 0 & 2 & 2 & 1 & 1 & 0 & 0 & 0 & 0 & 1 & 2 & 0 & 1 & 0 \\
4) &  & 8 & 0 & 8 & 0 & 0 & 1 & 0 & 0 & 0 & 0 & 0 & 6 & 0 & 3 & 0 & 0 \\
\end{longtable}
\noindent Modules tested: V78, V27.\\

\setcounter{table}{100}
\begin{longtable}{rc|cccccccccc|ccccc}
\caption{$U_3(3) < E_6$, $p = 7$}\\
& & \multicolumn{10}{c|}{$V_{78}$} & \multicolumn{5}{c}{$V_{27}$} \\
& & $1$ & $6$ & $7_{a}$ & $7_{b}$ & $7_{c}$ & $14$ & $21_{a}$ & $26$ & $28_{a}$ & $28_{b}$ & $1$ & $6$ & $7_{a}$ & $14$ & $26$ \\ \hline
\endfirsthead
& & \multicolumn{10}{c|}{$V_{78}$} & \multicolumn{5}{c}{$V_{27}$} \\
& & $1$ & $6$ & $7_{a}$ & $7_{b}$ & $7_{c}$ & $14$ & $21_{a}$ & $26$ & $28_{a}$ & $28_{b}$ & $1$ & $6$ & $7_{a}$ & $14$ & $26$ \\ \hline
\endhead
1) &  & 0 & 2 & 0 & 0 & 0 & 1 & 0 & 2 & 0 & 0 & 1 & 0 & 0 & 0 & 1 \\
2) &  & 1 & 0 & 1 & 1 & 1 & 0 & 0 & 0 & 1 & 1 & 1 & 0 & 0 & 0 & 1 \\
3) &  & 3 & 2 & 0 & 2 & 2 & 1 & 1 & 0 & 0 & 0 & 1 & 2 & 0 & 1 & 0 \\
4) &  & 8 & 0 & 8 & 0 & 0 & 1 & 0 & 0 & 0 & 0 & 6 & 0 & 3 & 0 & 0 \\
\end{longtable}
\noindent Modules tested: V78, V27.\\

\setcounter{table}{101}
\begin{longtable}{rc|cccccccc|cccccc}
\caption{$U_4(2) < E_6$, $p = 0$ or $p \nmid |H|$}\\
& & \multicolumn{8}{c|}{$V_{78}$} & \multicolumn{6}{c}{$V_{27}$} \\
& & $1$ & $5_{a}$ & $5_{b}$ & $10_{a}$ & $10_{b}$ & $15_{a}$ & $20$ & $24$ & $1$ & $5_{a}$ & $5_{b}$ & $6$ & $10_{b}$ & $15_{a}$ \\ \hline
\endfirsthead
& & \multicolumn{8}{c|}{$V_{78}$} & \multicolumn{6}{c}{$V_{27}$} \\
& & $1$ & $5_{a}$ & $5_{b}$ & $10_{a}$ & $10_{b}$ & $15_{a}$ & $20$ & $24$ & $1$ & $5_{a}$ & $5_{b}$ & $6$ & $10_{b}$ & $15_{a}$ \\ \hline
\endhead
1) &  & 3 & 0 & 0 & 2 & 2 & 1 & 1 & 0 & 0 & 0 & 0 & 2 & 0 & 1 \\
2) &  & 4 & 1 & 1 & 2 & 2 & 0 & 0 & 1 & 2 & 2 & 1 & 0 & 1 & 0 \\
\end{longtable}
\noindent Modules tested: V78, V27.\\

\setcounter{table}{102}
\begin{longtable}{rc|cccccccc|cccccc}
\caption{$U_4(2) < E_6$, $p = 5$}\\
& & \multicolumn{8}{c|}{$V_{78}$} & \multicolumn{6}{c}{$V_{27}$} \\
& & $1$ & $5_{a}$ & $5_{b}$ & $10_{a}$ & $10_{b}$ & $15_{a}$ & $20$ & $23$ & $1$ & $5_{a}$ & $5_{b}$ & $6$ & $10_{b}$ & $15_{a}$ \\ \hline
\endfirsthead
& & \multicolumn{8}{c|}{$V_{78}$} & \multicolumn{6}{c}{$V_{27}$} \\
& & $1$ & $5_{a}$ & $5_{b}$ & $10_{a}$ & $10_{b}$ & $15_{a}$ & $20$ & $23$ & $1$ & $5_{a}$ & $5_{b}$ & $6$ & $10_{b}$ & $15_{a}$ \\ \hline
\endhead
1) &  & 3 & 0 & 0 & 2 & 2 & 1 & 1 & 0 & 0 & 0 & 0 & 2 & 0 & 1 \\
2) &  & 5 & 1 & 1 & 2 & 2 & 0 & 0 & 1 & 2 & 2 & 1 & 0 & 1 & 0 \\
\end{longtable}
\noindent Modules tested: V78, V27.\\

\setcounter{table}{103}
\begin{longtable}{rc|ccc|cc}
\caption{$3 \cdot U_4(3) < E_6$, $p = 2$}\\
& & \multicolumn{3}{c|}{$V_{78}$} & \multicolumn{2}{c}{$V_{27}$} \\
& & $1$ & $20$ & $34_{a}$ & $6_{a}$ & $15_{a}$ \\ \hline
\endfirsthead
& & \multicolumn{3}{c|}{$V_{78}$} & \multicolumn{2}{c}{$V_{27}$} \\
& & $1$ & $20$ & $34_{a}$ & $6_{a}$ & $15_{a}$ \\ \hline
\endhead
1) & \nongcr & 4 & 2 & 1 & 2 & 1 \\
\end{longtable}
\noindent Modules tested: V78, V27.\\

\setcounter{table}{104}
\begin{longtable}{rc|c|c}
\caption{$3 \cdot \Omega_7(3) < E_6$, $p = 2$ (irreducible on both)}\\
& & \multicolumn{1}{c|}{$V_{78}$} & \multicolumn{1}{c}{$V_{27}$} \\
& & $78$ & $27_{a}$ \\ \hline
\endfirsthead
& & \multicolumn{1}{c|}{$V_{78}$} & \multicolumn{1}{c}{$V_{27}$} \\
& & $78$ & $27_{a}$ \\ \hline
\endhead
1) & \possprim & 1 & 1 \\
\end{longtable}
\noindent Modules tested: V78, V27.\\

\setcounter{table}{105}
\begin{longtable}{rc|c|c}
\caption{$3 \cdot G_2(3) < E_6$, $p = 2$ (irreducible on both)}\\
& & \multicolumn{1}{c|}{$V_{78}$} & \multicolumn{1}{c}{$V_{27}$} \\
& & $78$ & $27_{a}$ \\ \hline
\endfirsthead
& & \multicolumn{1}{c|}{$V_{78}$} & \multicolumn{1}{c}{$V_{27}$} \\
& & $78$ & $27_{a}$ \\ \hline
\endhead
1) & \possprim & 1 & 1 \\
\end{longtable}
\noindent Modules tested: V78, V27.\\

\setcounter{table}{106}
\begin{longtable}{rc|cc|cc}
\caption{$^{3}D_4(2) < E_6$, $p \neq 2,3$}\\
& & \multicolumn{2}{c|}{$V_{78}$} & \multicolumn{2}{c}{$V_{27}$} \\
& & $26$ & $52$ & $1$ & $26$ \\ \hline
\endfirsthead
& & \multicolumn{2}{c|}{$V_{78}$} & \multicolumn{2}{c}{$V_{27}$} \\
& & $26$ & $52$ & $1$ & $26$ \\ \hline
\endhead
1) &  & 1 & 1 & 1 & 1 \\
\end{longtable}
\noindent Modules tested: V78, V27.\\

\setcounter{table}{107}
\begin{longtable}{rc|ccc|cc}
\caption{$^{3}D_4(2) < E_6$, $p = 3$}\\
& & \multicolumn{3}{c|}{$V_{78}$} & \multicolumn{2}{c}{$V_{27}$} \\
& & $1$ & $25$ & $52$ & $1$ & $25$ \\ \hline
\endfirsthead
& & \multicolumn{3}{c|}{$V_{78}$} & \multicolumn{2}{c}{$V_{27}$} \\
& & $1$ & $25$ & $52$ & $1$ & $25$ \\ \hline
\endhead
1) &  & 1 & 1 & 1 & 2 & 1 \\
\end{longtable}
\noindent Modules tested: V78, L77, V27.\\

\setcounter{table}{108}
\begin{longtable}{rc|c|c}
\caption{$^{2}F_4(2)' < E_6$, $p \neq 2$, $3$ (irreducible on both)}\\
& & \multicolumn{1}{c|}{$V_{78}$} & \multicolumn{1}{c}{$V_{27}$} \\
& & $78$ & $27_{a}$ \\ \hline
\endfirsthead
& & \multicolumn{1}{c|}{$V_{78}$} & \multicolumn{1}{c}{$V_{27}$} \\
& & $78$ & $27_{a}$ \\ \hline
\endhead
1) & \possprim & 1 & 1 \\
\end{longtable}
\noindent Modules tested: V78, V27.\\

\setcounter{table}{109}
\begin{longtable}{rc|cc|c}
\caption{$^{2}F_4(2)' < E_6$, $p = 3$}\\
& & \multicolumn{2}{c|}{$V_{78}$} & \multicolumn{1}{c}{$V_{27}$} \\
& & $1$ & $77$ & $27_{a}$ \\ \hline
\endfirsthead
& & \multicolumn{2}{c|}{$V_{78}$} & \multicolumn{1}{c}{$V_{27}$} \\
& & $1$ & $77$ & $27_{a}$ \\ \hline
\endhead
1) & \possprim & 1 & 1 & 1 \\
\end{longtable}
\noindent Modules tested: V78, L77, V27.\\


\newpage
% Updated tables for E7 (Memoir Chapter 6), generated by tools/fg_tables.g
\section*{$E_7$: feasible characters on $L(G) = V_{133}$ and $V_{\min} = V_{56}$}

\setcounter{table}{110}
\begin{longtable}{rc|ccc|cc}
\caption{Alt$_{13} < E_7$, $p = 2$}\\
& & \multicolumn{3}{c|}{$V_{133}$} & \multicolumn{2}{c}{$V_{56}$} \\
& & $1$ & $32_{a}$ & $64$ & $12$ & $32_{b}$ \\ \hline
\endfirsthead
& & \multicolumn{3}{c|}{$V_{133}$} & \multicolumn{2}{c}{$V_{56}$} \\
& & $1$ & $32_{a}$ & $64$ & $12$ & $32_{b}$ \\ \hline
\endhead
1) &  & 5 & 2 & 1 & 2 & 1 \\
\end{longtable}
\noindent Modules tested: V133, L132, V912, V56.\\

\setcounter{table}{111}
\begin{longtable}{rc|ccccc|cccc}
\caption{Alt$_{12} < E_7$, $p = 2$}\\
& & \multicolumn{5}{c|}{$V_{133}$} & \multicolumn{4}{c}{$V_{56}$} \\
& & $1$ & $10$ & $16_{a}$ & $16_{b}$ & $44$ & $1$ & $10$ & $16_{a}$ & $16_{b}$ \\ \hline
\endfirsthead
& & \multicolumn{5}{c|}{$V_{133}$} & \multicolumn{4}{c}{$V_{56}$} \\
& & $1$ & $10$ & $16_{a}$ & $16_{b}$ & $44$ & $1$ & $10$ & $16_{a}$ & $16_{b}$ \\ \hline
\endhead
1) & \nongcr & 5 & 2 & 2 & 2 & 1 & 4 & 2 & 1 & 1 \\
\end{longtable}
\noindent Modules tested: V133, L132, V912, V56.\\

\setcounter{table}{112}
\begin{longtable}{rc|ccccc|cccc}
\caption{Alt$_{11} < E_7$, $p = 2$}\\
& & \multicolumn{5}{c|}{$V_{133}$} & \multicolumn{4}{c}{$V_{56}$} \\
& & $1$ & $10$ & $16_{a}$ & $16_{b}$ & $44$ & $1$ & $10$ & $16_{a}$ & $16_{b}$ \\ \hline
\endfirsthead
& & \multicolumn{5}{c|}{$V_{133}$} & \multicolumn{4}{c}{$V_{56}$} \\
& & $1$ & $10$ & $16_{a}$ & $16_{b}$ & $44$ & $1$ & $10$ & $16_{a}$ & $16_{b}$ \\ \hline
\endhead
1) &  & 5 & 2 & 2 & 2 & 1 & 4 & 2 & 1 & 1 \\
\end{longtable}
\noindent Modules tested: V133, L132, V912, V56.\\

\setcounter{table}{113}
\begin{longtable}{rc|cccc|c}
\caption{Alt$_{10} < E_7$, $p = 5$}\\
& & \multicolumn{4}{c|}{$V_{133}$} & \multicolumn{1}{c}{$V_{56}$} \\
& & $28$ & $35_{a}$ & $35_{b}$ & $35_{c}$ & $28$ \\ \hline
\endfirsthead
& & \multicolumn{4}{c|}{$V_{133}$} & \multicolumn{1}{c}{$V_{56}$} \\
& & $28$ & $35_{a}$ & $35_{b}$ & $35_{c}$ & $28$ \\ \hline
\endhead
1) & \possprim & 1 & 1 & 1 & 1 & 2 \\
\end{longtable}
\noindent Modules tested: V133, V912, V56.\\

\setcounter{table}{114}
\begin{longtable}{rc|cccc|cccc}
\caption{Alt$_{10} < E_7$, $p = 2$}\\
& & \multicolumn{4}{c|}{$V_{133}$} & \multicolumn{4}{c}{$V_{56}$} \\
& & $1$ & $8$ & $16$ & $26$ & $1$ & $8$ & $16$ & $26$ \\ \hline
\endfirsthead
& & \multicolumn{4}{c|}{$V_{133}$} & \multicolumn{4}{c}{$V_{56}$} \\
& & $1$ & $8$ & $16$ & $26$ & $1$ & $8$ & $16$ & $26$ \\ \hline
\endhead
1) & \nongcr & 5 & 1 & 1 & 4 & 4 & 0 & 0 & 2 \\
2) & \nongcr & 11 & 4 & 4 & 1 & 8 & 2 & 2 & 0 \\
\end{longtable}
\noindent Modules tested: V133, L132, V912, V56.\\

\setcounter{table}{115}
\begin{longtable}{rc|ccccc|c}
\caption{Alt$_{9} < E_7$, $p = 0$ or $p \nmid |H|$}\\
& & \multicolumn{5}{c|}{$V_{133}$} & \multicolumn{1}{c}{$V_{56}$} \\
& & $8$ & $27$ & $28$ & $35_{a}$ & $35_{b}$ & $28$ \\ \hline
\endfirsthead
& & \multicolumn{5}{c|}{$V_{133}$} & \multicolumn{1}{c}{$V_{56}$} \\
& & $8$ & $27$ & $28$ & $35_{a}$ & $35_{b}$ & $28$ \\ \hline
\endhead
1) & \possprim & 1 & 1 & 1 & 1 & 1 & 2 \\
\end{longtable}
\noindent Modules tested: V133, V912, V56.\\

\setcounter{table}{116}
\begin{longtable}{rc|ccccc|c}
\caption{Alt$_{9} < E_7$, $p = 7$}\\
& & \multicolumn{5}{c|}{$V_{133}$} & \multicolumn{1}{c}{$V_{56}$} \\
& & $8$ & $19$ & $28$ & $35_{a}$ & $35_{b}$ & $28$ \\ \hline
\endfirsthead
& & \multicolumn{5}{c|}{$V_{133}$} & \multicolumn{1}{c}{$V_{56}$} \\
& & $8$ & $19$ & $28$ & $35_{a}$ & $35_{b}$ & $28$ \\ \hline
\endhead
1) & \possprim & 2 & 1 & 1 & 1 & 1 & 2 \\
\end{longtable}
\noindent Modules tested: V133, V912, V56.\\

\setcounter{table}{117}
\begin{longtable}{rc|ccccc|c}
\caption{Alt$_{9} < E_7$, $p = 5$}\\
& & \multicolumn{5}{c|}{$V_{133}$} & \multicolumn{1}{c}{$V_{56}$} \\
& & $8$ & $27$ & $28$ & $35_{a}$ & $35_{b}$ & $28$ \\ \hline
\endfirsthead
& & \multicolumn{5}{c|}{$V_{133}$} & \multicolumn{1}{c}{$V_{56}$} \\
& & $8$ & $27$ & $28$ & $35_{a}$ & $35_{b}$ & $28$ \\ \hline
\endhead
1) & \possprim & 1 & 1 & 1 & 1 & 1 & 2 \\
\end{longtable}
\noindent Modules tested: V133, V912, V56.\\

\setcounter{table}{118}
\begin{longtable}{rc|ccccc|cc}
\caption{Alt$_{9} < E_7$, $p = 3$}\\
& & \multicolumn{5}{c|}{$V_{133}$} & \multicolumn{2}{c}{$V_{56}$} \\
& & $1$ & $7$ & $21$ & $27$ & $35$ & $7$ & $21$ \\ \hline
\endfirsthead
& & \multicolumn{5}{c|}{$V_{133}$} & \multicolumn{2}{c}{$V_{56}$} \\
& & $1$ & $7$ & $21$ & $27$ & $35$ & $7$ & $21$ \\ \hline
\endhead
1) & \possprim, \nongcr & 1 & 2 & 1 & 1 & 2 & 2 & 2 \\
\end{longtable}
\noindent Modules tested: V133, V912, L856, V56.\\

\setcounter{table}{119}
\begin{longtable}{rc|ccccc|ccccc}
\caption{Alt$_{9} < E_7$, $p = 2$}\\
& & \multicolumn{5}{c|}{$V_{133}$} & \multicolumn{5}{c}{$V_{56}$} \\
& & $1$ & $8_{a}$ & $8_{b}$ & $8_{c}$ & $26$ & $1$ & $8_{a}$ & $8_{b}$ & $8_{c}$ & $26$ \\ \hline
\endfirsthead
& & \multicolumn{5}{c|}{$V_{133}$} & \multicolumn{5}{c}{$V_{56}$} \\
& & $1$ & $8_{a}$ & $8_{b}$ & $8_{c}$ & $26$ & $1$ & $8_{a}$ & $8_{b}$ & $8_{c}$ & $26$ \\ \hline
\endhead
1) & \nongcr & 5 & 1 & 1 & 1 & 4 & 4 & 0 & 0 & 0 & 2 \\
2) &  & 11 & 4 & 4 & 4 & 1 & 8 & 2 & 2 & 2 & 0 \\
\end{longtable}
\noindent Modules tested: V133, L132, V912, V56.\\

\setcounter{table}{120}
\begin{longtable}{rc|ccccc|cc}
\caption{Alt$_{8} < E_7$, $p = 0$ or $p \nmid |H|$}\\
& & \multicolumn{5}{c|}{$V_{133}$} & \multicolumn{2}{c}{$V_{56}$} \\
& & $1$ & $7$ & $20$ & $21_{a}$ & $35$ & $7$ & $21_{a}$ \\ \hline
\endfirsthead
& & \multicolumn{5}{c|}{$V_{133}$} & \multicolumn{2}{c}{$V_{56}$} \\
& & $1$ & $7$ & $20$ & $21_{a}$ & $35$ & $7$ & $21_{a}$ \\ \hline
\endhead
1) &  & 1 & 3 & 1 & 1 & 2 & 2 & 2 \\
\end{longtable}
\noindent Modules tested: V133, V912, V56.\\

\setcounter{table}{121}
\begin{longtable}{rc|ccccc|cc}
\caption{Alt$_{8} < E_7$, $p = 7$}\\
& & \multicolumn{5}{c|}{$V_{133}$} & \multicolumn{2}{c}{$V_{56}$} \\
& & $1$ & $7$ & $19$ & $21_{a}$ & $35$ & $7$ & $21_{a}$ \\ \hline
\endfirsthead
& & \multicolumn{5}{c|}{$V_{133}$} & \multicolumn{2}{c}{$V_{56}$} \\
& & $1$ & $7$ & $19$ & $21_{a}$ & $35$ & $7$ & $21_{a}$ \\ \hline
\endhead
1) &  & 2 & 3 & 1 & 1 & 2 & 2 & 2 \\
\end{longtable}
\noindent Modules tested: V133, V912, V56.\\

\setcounter{table}{122}
\begin{longtable}{rc|cccccccc|ccc}
\caption{Alt$_{8} < E_7$, $p = 5$}\\
& & \multicolumn{8}{c|}{$V_{133}$} & \multicolumn{3}{c}{$V_{56}$} \\
& & $1$ & $7$ & $13$ & $20$ & $21_{a}$ & $35$ & $43$ & $70$ & $7$ & $21_{a}$ & $21_{b}$ \\ \hline
\endfirsthead
& & \multicolumn{8}{c|}{$V_{133}$} & \multicolumn{3}{c}{$V_{56}$} \\
& & $1$ & $7$ & $13$ & $20$ & $21_{a}$ & $35$ & $43$ & $70$ & $7$ & $21_{a}$ & $21_{b}$ \\ \hline
\endhead
1) & \possprim & 0 & 0 & 0 & 1 & 0 & 0 & 1 & 1 & 2 & 0 & 2 \\
2) &  & 1 & 2 & 1 & 0 & 0 & 1 & 0 & 1 & 2 & 0 & 2 \\
3) &  & 1 & 3 & 0 & 1 & 1 & 2 & 0 & 0 & 2 & 2 & 0 \\
\end{longtable}
\noindent Modules tested: V133, V912, V56.\\

\setcounter{table}{123}
\begin{longtable}{rc|ccccc|cc}
\caption{Alt$_{8} < E_7$, $p = 3$}\\
& & \multicolumn{5}{c|}{$V_{133}$} & \multicolumn{2}{c}{$V_{56}$} \\
& & $1$ & $7$ & $13$ & $21$ & $35$ & $7$ & $21$ \\ \hline
\endfirsthead
& & \multicolumn{5}{c|}{$V_{133}$} & \multicolumn{2}{c}{$V_{56}$} \\
& & $1$ & $7$ & $13$ & $21$ & $35$ & $7$ & $21$ \\ \hline
\endhead
1) & \possprim, \nongcr & 1 & 4 & 1 & 1 & 2 & 2 & 2 \\
\end{longtable}
\noindent Modules tested: V133, V912, L856, V56.\\

\setcounter{table}{124}
\begin{longtable}{rc|cccccc|ccccc}
\caption{Alt$_{7} < E_7$, $p = 0$ or $p \nmid |H|$}\\
& & \multicolumn{6}{c|}{$V_{133}$} & \multicolumn{5}{c}{$V_{56}$} \\
& & $1$ & $6$ & $10_{a}$ & $10_{b}$ & $14_{a}$ & $15$ & $1$ & $6$ & $10_{a}$ & $10_{b}$ & $15$ \\ \hline
\endfirsthead
& & \multicolumn{6}{c|}{$V_{133}$} & \multicolumn{5}{c}{$V_{56}$} \\
& & $1$ & $6$ & $10_{a}$ & $10_{b}$ & $14_{a}$ & $15$ & $1$ & $6$ & $10_{a}$ & $10_{b}$ & $15$ \\ \hline
\endhead
1) &  & 8 & 1 & 0 & 0 & 1 & 7 & 0 & 6 & 1 & 1 & 0 \\
2) &  & 4 & 5 & 2 & 2 & 1 & 3 & 2 & 4 & 0 & 0 & 2 \\
\end{longtable}
\noindent Modules tested: V133, V912, V56.\\

\setcounter{table}{125}
\begin{longtable}{rc|cccc|ccc}
\caption{Alt$_{7} < E_7$, $p = 7$}\\
& & \multicolumn{4}{c|}{$V_{133}$} & \multicolumn{3}{c}{$V_{56}$} \\
& & $1$ & $5$ & $10$ & $14_{a}$ & $1$ & $5$ & $10$ \\ \hline
\endfirsthead
& & \multicolumn{4}{c|}{$V_{133}$} & \multicolumn{3}{c}{$V_{56}$} \\
& & $1$ & $5$ & $10$ & $14_{a}$ & $1$ & $5$ & $10$ \\ \hline
\endhead
1) & \nongcr & 9 & 8 & 7 & 1 & 6 & 6 & 2 \\
\end{longtable}
\noindent Modules tested: V133, V912, V56.\\

\setcounter{table}{126}
\begin{longtable}{rc|cccccccc|ccccccc}
\caption{Alt$_{7} < E_7$, $p = 5$}\\
& & \multicolumn{8}{c|}{$V_{133}$} & \multicolumn{7}{c}{$V_{56}$} \\
& & $1$ & $6$ & $8$ & $10_{a}$ & $10_{b}$ & $13$ & $15$ & $35$ & $1$ & $6$ & $8$ & $10_{a}$ & $10_{b}$ & $13$ & $15$ \\ \hline
\endfirsthead
& & \multicolumn{8}{c|}{$V_{133}$} & \multicolumn{7}{c}{$V_{56}$} \\
& & $1$ & $6$ & $8$ & $10_{a}$ & $10_{b}$ & $13$ & $15$ & $35$ & $1$ & $6$ & $8$ & $10_{a}$ & $10_{b}$ & $13$ & $15$ \\ \hline
\endhead
1) & \possprim & 0 & 3 & 7 & 1 & 1 & 3 & 0 & 0 & 0 & 0 & 2 & 2 & 2 & 0 & 0 \\
2) &  & 8 & 2 & 1 & 0 & 0 & 0 & 7 & 0 & 0 & 6 & 0 & 1 & 1 & 0 & 0 \\
3) & \nongcr & 1 & 2 & 3 & 0 & 0 & 2 & 0 & 2 & 2 & 2 & 2 & 0 & 0 & 2 & 0 \\
4) &  & 1 & 3 & 2 & 0 & 0 & 1 & 1 & 2 & 2 & 2 & 2 & 0 & 0 & 2 & 0 \\
5) & \nongcr & 4 & 1 & 2 & 1 & 1 & 4 & 0 & 1 & 2 & 2 & 2 & 0 & 0 & 2 & 0 \\
6) & \nongcr & 4 & 2 & 1 & 1 & 1 & 3 & 1 & 1 & 2 & 2 & 2 & 0 & 0 & 2 & 0 \\
7) &  & 4 & 6 & 1 & 2 & 2 & 0 & 3 & 0 & 2 & 4 & 0 & 0 & 0 & 0 & 2 \\
8) &  & 9 & 0 & 13 & 1 & 1 & 0 & 0 & 0 & 8 & 0 & 6 & 0 & 0 & 0 & 0 \\
\end{longtable}
\noindent Modules tested: V133, V912, V56.\\

\setcounter{table}{127}
\begin{longtable}{rc|cccc|cccc}
\caption{$2\cdot$Alt$_{7} < E_7$, $p = 5$}\\
& & \multicolumn{4}{c|}{$V_{133}$} & \multicolumn{4}{c}{$V_{56}$} \\
& & $8$ & $10_{a}$ & $10_{b}$ & $35$ & $4_{a}$ & $4_{b}$ & $14_{a}$ & $20_{b}$ \\ \hline
\endfirsthead
& & \multicolumn{4}{c|}{$V_{133}$} & \multicolumn{4}{c}{$V_{56}$} \\
& & $8$ & $10_{a}$ & $10_{b}$ & $35$ & $4_{a}$ & $4_{b}$ & $14_{a}$ & $20_{b}$ \\ \hline
\endhead
1) & \possprim & 1 & 1 & 1 & 3 & 1 & 1 & 2 & 1 \\
\end{longtable}
\noindent Modules tested: V133, V912, V56.\\

\setcounter{table}{128}
\begin{longtable}{rc|cccccc|ccccc}
\caption{Alt$_{7} < E_7$, $p = 3$}\\
& & \multicolumn{6}{c|}{$V_{133}$} & \multicolumn{5}{c}{$V_{56}$} \\
& & $1$ & $6$ & $10_{a}$ & $10_{b}$ & $13$ & $15$ & $1$ & $6$ & $10_{a}$ & $10_{b}$ & $15$ \\ \hline
\endfirsthead
& & \multicolumn{6}{c|}{$V_{133}$} & \multicolumn{5}{c}{$V_{56}$} \\
& & $1$ & $6$ & $10_{a}$ & $10_{b}$ & $13$ & $15$ & $1$ & $6$ & $10_{a}$ & $10_{b}$ & $15$ \\ \hline
\endhead
1) &  & 9 & 1 & 0 & 0 & 1 & 7 & 0 & 6 & 1 & 1 & 0 \\
2) & \nongcr & 5 & 5 & 2 & 2 & 1 & 3 & 2 & 4 & 0 & 0 & 2 \\
\end{longtable}
\noindent Modules tested: V133, V912, L856, V56.\\

\setcounter{table}{129}
\begin{longtable}{rc|cccc|cccc}
\caption{$2\cdot$Alt$_{7} < E_7$, $p = 3$}\\
& & \multicolumn{4}{c|}{$V_{133}$} & \multicolumn{4}{c}{$V_{56}$} \\
& & $1$ & $10_{a}$ & $10_{b}$ & $13$ & $4_{a}$ & $4_{b}$ & $6_{b}$ & $6_{c}$ \\ \hline
\endfirsthead
& & \multicolumn{4}{c|}{$V_{133}$} & \multicolumn{4}{c}{$V_{56}$} \\
& & $1$ & $10_{a}$ & $10_{b}$ & $13$ & $4_{a}$ & $4_{b}$ & $6_{b}$ & $6_{c}$ \\ \hline
\endhead
1) & \nongcr & 22 & 1 & 1 & 7 & 1 & 1 & 7 & 1 \\
\end{longtable}
\noindent Modules tested: V133, V912, L856, V56.\\

\setcounter{table}{130}
\begin{longtable}{rc|cccccc|cccccc}
\caption{Alt$_{7} < E_7$, $p = 2$}\\
& & \multicolumn{6}{c|}{$V_{133}$} & \multicolumn{6}{c}{$V_{56}$} \\
& & $1$ & $4_{a}$ & $4_{b}$ & $6$ & $14$ & $20$ & $1$ & $4_{a}$ & $4_{b}$ & $6$ & $14$ & $20$ \\ \hline
\endfirsthead
& & \multicolumn{6}{c|}{$V_{133}$} & \multicolumn{6}{c}{$V_{56}$} \\
& & $1$ & $4_{a}$ & $4_{b}$ & $6$ & $14$ & $20$ & $1$ & $4_{a}$ & $4_{b}$ & $6$ & $14$ & $20$ \\ \hline
\endhead
1) & \nongcr & 15 & 0 & 0 & 1 & 8 & 0 & 0 & 1 & 1 & 8 & 0 & 0 \\
2) & \nongcr & 7 & 2 & 2 & 5 & 0 & 4 & 4 & 0 & 0 & 2 & 0 & 2 \\
3) & \nongcr & 7 & 2 & 2 & 6 & 1 & 3 & 4 & 0 & 0 & 2 & 0 & 2 \\
4) & \nongcr & 7 & 2 & 2 & 9 & 4 & 0 & 4 & 0 & 0 & 4 & 2 & 0 \\
5) &  & 19 & 8 & 8 & 5 & 0 & 1 & 12 & 4 & 4 & 2 & 0 & 0 \\
6) &  & 19 & 8 & 8 & 6 & 1 & 0 & 12 & 4 & 4 & 2 & 0 & 0 \\
\end{longtable}
\noindent Modules tested: V133, L132, V912, V56.\\

\setcounter{table}{131}
\begin{longtable}{rc|ccccccc|ccccccc}
\caption{Alt$_{6} < E_7$, $p = 0$ or $p \nmid |H|$}\\
& & \multicolumn{7}{c|}{$V_{133}$} & \multicolumn{7}{c}{$V_{56}$} \\
& & $1$ & $5_{a}$ & $5_{b}$ & $8_{a}$ & $8_{b}$ & $9$ & $10$ & $1$ & $5_{a}$ & $5_{b}$ & $8_{a}$ & $8_{b}$ & $9$ & $10$ \\ \hline
\endfirsthead
& & \multicolumn{7}{c|}{$V_{133}$} & \multicolumn{7}{c}{$V_{56}$} \\
& & $1$ & $5_{a}$ & $5_{b}$ & $8_{a}$ & $8_{b}$ & $9$ & $10$ & $1$ & $5_{a}$ & $5_{b}$ & $8_{a}$ & $8_{b}$ & $9$ & $10$ \\ \hline
\endhead
1) & \possprim & 0 & 3 & 3 & 4 & 3 & 3 & 2 & 0 & 0 & 0 & 2 & 0 & 0 & 4 \\
2) &  & 1 & 5 & 2 & 3 & 2 & 3 & 3 & 0 & 0 & 0 & 2 & 0 & 0 & 4 \\
3) &  & 2 & 1 & 1 & 4 & 3 & 5 & 2 & 0 & 0 & 0 & 2 & 0 & 0 & 4 \\
4) &  & 2 & 4 & 4 & 2 & 1 & 3 & 4 & 0 & 0 & 0 & 2 & 0 & 0 & 4 \\
5) &  & 3 & 3 & 0 & 3 & 2 & 5 & 3 & 0 & 0 & 0 & 2 & 0 & 0 & 4 \\
6) &  & 4 & 2 & 2 & 2 & 1 & 5 & 4 & 0 & 0 & 0 & 2 & 0 & 0 & 4 \\
7) &  & 5 & 1 & 1 & 2 & 1 & 6 & 4 & 0 & 4 & 4 & 2 & 0 & 0 & 0 \\
8) &  & 3 & 4 & 1 & 4 & 2 & 3 & 3 & 1 & 2 & 2 & 0 & 1 & 3 & 0 \\
9) &  & 4 & 3 & 3 & 3 & 1 & 3 & 4 & 1 & 2 & 2 & 0 & 1 & 3 & 0 \\
10) &  & 1 & 2 & 2 & 4 & 3 & 4 & 2 & 2 & 2 & 2 & 2 & 0 & 2 & 0 \\
11) &  & 2 & 4 & 1 & 3 & 2 & 4 & 3 & 2 & 2 & 2 & 2 & 0 & 2 & 0 \\
12) &  & 3 & 3 & 3 & 2 & 1 & 4 & 4 & 2 & 2 & 2 & 2 & 0 & 2 & 0 \\
13) &  & 4 & 0 & 0 & 5 & 3 & 5 & 2 & 3 & 0 & 0 & 0 & 1 & 5 & 0 \\
14) &  & 6 & 1 & 1 & 3 & 1 & 5 & 4 & 3 & 0 & 0 & 0 & 1 & 5 & 0 \\
15) &  & 3 & 0 & 0 & 4 & 3 & 6 & 2 & 4 & 0 & 0 & 2 & 0 & 4 & 0 \\
16) &  & 5 & 1 & 1 & 2 & 1 & 6 & 4 & 4 & 0 & 0 & 2 & 0 & 4 & 0 \\
17) &  & 4 & 2 & 2 & 5 & 5 & 1 & 2 & 4 & 2 & 2 & 2 & 2 & 0 & 0 \\
18) &  & 5 & 4 & 1 & 4 & 4 & 1 & 3 & 4 & 2 & 2 & 2 & 2 & 0 & 0 \\
19) &  & 6 & 3 & 3 & 3 & 3 & 1 & 4 & 4 & 2 & 2 & 2 & 2 & 0 & 0 \\
20) &  & 6 & 0 & 0 & 5 & 5 & 3 & 2 & 6 & 0 & 0 & 2 & 2 & 2 & 0 \\
21) &  & 8 & 1 & 1 & 3 & 3 & 3 & 4 & 6 & 0 & 0 & 2 & 2 & 2 & 0 \\
22) &  & 9 & 9 & 0 & 0 & 0 & 1 & 7 & 6 & 6 & 0 & 0 & 0 & 0 & 2 \\
23) &  & 9 & 0 & 0 & 13 & 0 & 0 & 2 & 8 & 0 & 0 & 6 & 0 & 0 & 0 \\
\end{longtable}
\noindent Memoir rows not reproduced: 7) (not compatible with V912); 9) (not compatible with V912).\\
Modules tested: V133, V912, V56.\\

\setcounter{table}{132}
\begin{longtable}{rc|ccccccc|cccccc}
\caption{$2\cdot$Alt$_{6} < E_7$, $p = 0$ or $p \nmid |H|$}\\
& & \multicolumn{7}{c|}{$V_{133}$} & \multicolumn{6}{c}{$V_{56}$} \\
& & $1$ & $5_{a}$ & $5_{b}$ & $8_{a}$ & $8_{b}$ & $9$ & $10_{a}$ & $4_{a}$ & $4_{b}$ & $8_{c}$ & $8_{d}$ & $10_{b}$ & $10_{c}$ \\ \hline
\endfirsthead
& & \multicolumn{7}{c|}{$V_{133}$} & \multicolumn{6}{c}{$V_{56}$} \\
& & $1$ & $5_{a}$ & $5_{b}$ & $8_{a}$ & $8_{b}$ & $9$ & $10_{a}$ & $4_{a}$ & $4_{b}$ & $8_{c}$ & $8_{d}$ & $10_{b}$ & $10_{c}$ \\ \hline
\endhead
1) & \possprim & 0 & 0 & 0 & 4 & 3 & 3 & 5 & 0 & 0 & 2 & 0 & 3 & 1 \\
2) &  & 2 & 1 & 1 & 2 & 1 & 3 & 7 & 0 & 0 & 2 & 0 & 3 & 1 \\
3) &  & 3 & 1 & 1 & 3 & 1 & 2 & 7 & 3 & 0 & 1 & 2 & 2 & 0 \\
4) &  & 1 & 0 & 0 & 5 & 3 & 2 & 5 & 1 & 1 & 0 & 1 & 3 & 1 \\
5) &  & 3 & 1 & 1 & 3 & 1 & 2 & 7 & 1 & 1 & 0 & 1 & 3 & 1 \\
\end{longtable}
\noindent Modules tested: V133, V912, V56.\\

\setcounter{table}{133}
\begin{longtable}{rc|ccccc|ccccc}
\caption{Alt$_{6} < E_7$, $p = 5$}\\
& & \multicolumn{5}{c|}{$V_{133}$} & \multicolumn{5}{c}{$V_{56}$} \\
& & $1$ & $5_{a}$ & $5_{b}$ & $8$ & $10$ & $1$ & $5_{a}$ & $5_{b}$ & $8$ & $10$ \\ \hline
\endfirsthead
& & \multicolumn{5}{c|}{$V_{133}$} & \multicolumn{5}{c}{$V_{56}$} \\
& & $1$ & $5_{a}$ & $5_{b}$ & $8$ & $10$ & $1$ & $5_{a}$ & $5_{b}$ & $8$ & $10$ \\ \hline
\endhead
1) & \possprim, \nongcr & 3 & 3 & 3 & 10 & 2 & 0 & 0 & 0 & 2 & 4 \\
2) & \possprim, \nongcr & 4 & 5 & 2 & 8 & 3 & 0 & 0 & 0 & 2 & 4 \\
3) & \possprim, \nongcr & 5 & 4 & 4 & 6 & 4 & 0 & 0 & 0 & 2 & 4 \\
4) & \possprim, \nongcr & 7 & 1 & 1 & 12 & 2 & 0 & 0 & 0 & 2 & 4 \\
5) & \possprim, \nongcr & 8 & 3 & 0 & 10 & 3 & 0 & 0 & 0 & 2 & 4 \\
6) & \nongcr & 9 & 2 & 2 & 8 & 4 & 0 & 0 & 0 & 2 & 4 \\
7) & \possprim, \nongcr & 9 & 0 & 0 & 13 & 2 & 0 & 4 & 4 & 2 & 0 \\
8) & \nongcr & 11 & 1 & 1 & 9 & 4 & 0 & 4 & 4 & 2 & 0 \\
9) & \nongcr & 5 & 2 & 2 & 11 & 2 & 4 & 2 & 2 & 4 & 0 \\
10) & \nongcr & 6 & 4 & 1 & 9 & 3 & 4 & 2 & 2 & 4 & 0 \\
11) & \nongcr & 7 & 3 & 3 & 7 & 4 & 4 & 2 & 2 & 4 & 0 \\
12) &  & 10 & 9 & 0 & 1 & 7 & 6 & 6 & 0 & 0 & 2 \\
13) & \nongcr & 9 & 0 & 0 & 13 & 2 & 8 & 0 & 0 & 6 & 0 \\
14) & \nongcr & 11 & 1 & 1 & 9 & 4 & 8 & 0 & 0 & 6 & 0 \\
\end{longtable}
\noindent Modules tested: V133, V912, V56.\\

\setcounter{table}{134}
\begin{longtable}{rc|ccccc|cccc}
\caption{$2\cdot$Alt$_{6} < E_7$, $p = 5$}\\
& & \multicolumn{5}{c|}{$V_{133}$} & \multicolumn{4}{c}{$V_{56}$} \\
& & $1$ & $5_{a}$ & $5_{b}$ & $8$ & $10_{a}$ & $4_{a}$ & $4_{b}$ & $10_{b}$ & $10_{c}$ \\ \hline
\endfirsthead
& & \multicolumn{5}{c|}{$V_{133}$} & \multicolumn{4}{c}{$V_{56}$} \\
& & $1$ & $5_{a}$ & $5_{b}$ & $8$ & $10_{a}$ & $4_{a}$ & $4_{b}$ & $10_{b}$ & $10_{c}$ \\ \hline
\endhead
1) & \possprim, \nongcr & 3 & 0 & 0 & 10 & 5 & 2 & 2 & 3 & 1 \\
2) & \possprim, \nongcr & 5 & 1 & 1 & 6 & 7 & 2 & 2 & 3 & 1 \\
3) & \possprim, \nongcr & 5 & 1 & 1 & 6 & 7 & 6 & 3 & 2 & 0 \\
\end{longtable}
\noindent Modules tested: V133, V912, V56.\\

\setcounter{table}{135}
\begin{longtable}{rc|ccccc|ccccc}
\caption{Alt$_{5} < E_7$, $p = 0$ or $p \nmid |H|$}\\
& & \multicolumn{5}{c|}{$V_{133}$} & \multicolumn{5}{c}{$V_{56}$} \\
& & $1$ & $3_{a}$ & $3_{b}$ & $4$ & $5$ & $1$ & $3_{a}$ & $3_{b}$ & $4$ & $5$ \\ \hline
\endfirsthead
& & \multicolumn{5}{c|}{$V_{133}$} & \multicolumn{5}{c}{$V_{56}$} \\
& & $1$ & $3_{a}$ & $3_{b}$ & $4$ & $5$ & $1$ & $3_{a}$ & $3_{b}$ & $4$ & $5$ \\ \hline
\endhead
1) &  & 7 & 7 & 5 & 10 & 10 & 0 & 0 & 1 & 2 & 9 \\
2) &  & 8 & 10 & 6 & 3 & 13 & 0 & 9 & 1 & 4 & 2 \\
3) &  & 11 & 10 & 6 & 6 & 10 & 0 & 9 & 1 & 4 & 2 \\
4) &  & 3 & 6 & 5 & 8 & 13 & 0 & 6 & 4 & 4 & 2 \\
5) &  & 6 & 6 & 5 & 11 & 10 & 0 & 6 & 4 & 4 & 2 \\
6) &  & 6 & 6 & 5 & 11 & 10 & 1 & 2 & 0 & 1 & 9 \\
7) &  & 6 & 7 & 7 & 5 & 13 & 2 & 6 & 6 & 2 & 2 \\
8) &  & 9 & 7 & 7 & 8 & 10 & 2 & 6 & 6 & 2 & 2 \\
9) &  & 4 & 7 & 5 & 7 & 13 & 3 & 0 & 1 & 5 & 6 \\
10) &  & 7 & 7 & 5 & 10 & 10 & 3 & 0 & 1 & 5 & 6 \\
11) &  & 3 & 6 & 5 & 8 & 13 & 4 & 2 & 0 & 4 & 6 \\
12) &  & 6 & 6 & 5 & 11 & 10 & 4 & 2 & 0 & 4 & 6 \\
13) &  & 17 & 22 & 0 & 0 & 10 & 4 & 14 & 0 & 0 & 2 \\
14) &  & 9 & 15 & 2 & 2 & 13 & 4 & 10 & 4 & 0 & 2 \\
15) &  & 12 & 15 & 2 & 5 & 10 & 4 & 10 & 4 & 0 & 2 \\
16) &  & 6 & 7 & 7 & 5 & 13 & 6 & 2 & 2 & 2 & 6 \\
17) &  & 9 & 7 & 7 & 8 & 10 & 6 & 2 & 2 & 2 & 6 \\
18) &  & 9 & 15 & 2 & 2 & 13 & 8 & 6 & 0 & 0 & 6 \\
19) &  & 12 & 15 & 2 & 5 & 10 & 8 & 6 & 0 & 0 & 6 \\
20) &  & 16 & 7 & 5 & 19 & 1 & 9 & 0 & 1 & 11 & 0 \\
21) &  & 15 & 6 & 5 & 20 & 1 & 10 & 2 & 0 & 10 & 0 \\
22) &  & 18 & 7 & 7 & 17 & 1 & 12 & 2 & 2 & 8 & 0 \\
23) &  & 21 & 15 & 2 & 14 & 1 & 14 & 6 & 0 & 6 & 0 \\
24) &  & 35 & 31 & 0 & 0 & 1 & 20 & 12 & 0 & 0 & 0 \\
\end{longtable}
\noindent Modules tested: V133, V912, V56.\\

\setcounter{table}{136}
\begin{longtable}{rc|ccccc|cccc}
\caption{$2\cdot$Alt$_{5} < E_7$, $p = 0$ or $p \nmid |H|$}\\
& & \multicolumn{5}{c|}{$V_{133}$} & \multicolumn{4}{c}{$V_{56}$} \\
& & $1$ & $3_{a}$ & $3_{b}$ & $4_{a}$ & $5$ & $2_{a}$ & $2_{b}$ & $4_{b}$ & $6$ \\ \hline
\endfirsthead
& & \multicolumn{5}{c|}{$V_{133}$} & \multicolumn{4}{c}{$V_{56}$} \\
& & $1$ & $3_{a}$ & $3_{b}$ & $4_{a}$ & $5$ & $2_{a}$ & $2_{b}$ & $4_{b}$ & $6$ \\ \hline
\endhead
1) & \possprim & 0 & 9 & 8 & 8 & 10 & 0 & 2 & 4 & 6 \\
2) &  & 1 & 10 & 8 & 7 & 10 & 2 & 1 & 5 & 5 \\
3) &  & 3 & 15 & 5 & 5 & 10 & 5 & 0 & 7 & 3 \\
4) &  & 3 & 9 & 8 & 11 & 7 & 0 & 2 & 4 & 6 \\
5) &  & 3 & 9 & 8 & 11 & 7 & 3 & 5 & 1 & 6 \\
6) &  & 3 & 10 & 10 & 5 & 10 & 0 & 0 & 2 & 8 \\
7) &  & 4 & 10 & 8 & 10 & 7 & 2 & 1 & 5 & 5 \\
8) &  & 4 & 10 & 8 & 10 & 7 & 5 & 4 & 2 & 5 \\
9) &  & 6 & 15 & 5 & 8 & 7 & 5 & 0 & 7 & 3 \\
10) &  & 6 & 15 & 5 & 8 & 7 & 8 & 3 & 4 & 3 \\
11) &  & 6 & 10 & 10 & 8 & 7 & 0 & 0 & 2 & 8 \\
12) &  & 8 & 1 & 0 & 8 & 18 & 0 & 2 & 4 & 6 \\
13) &  & 8 & 13 & 9 & 6 & 7 & 0 & 8 & 1 & 6 \\
14) &  & 9 & 2 & 0 & 7 & 18 & 2 & 1 & 5 & 5 \\
15) &  & 11 & 1 & 0 & 11 & 15 & 0 & 2 & 4 & 6 \\
16) &  & 11 & 1 & 0 & 11 & 15 & 3 & 5 & 1 & 6 \\
17) &  & 11 & 2 & 2 & 5 & 18 & 0 & 0 & 2 & 8 \\
18) &  & 12 & 2 & 0 & 10 & 15 & 2 & 1 & 5 & 5 \\
19) &  & 12 & 2 & 0 & 10 & 15 & 5 & 4 & 2 & 5 \\
20) &  & 13 & 18 & 9 & 1 & 7 & 7 & 0 & 9 & 1 \\
21) &  & 13 & 18 & 9 & 1 & 7 & 10 & 3 & 6 & 1 \\
22) &  & 14 & 28 & 0 & 0 & 7 & 14 & 0 & 7 & 0 \\
23) &  & 14 & 2 & 2 & 8 & 15 & 0 & 0 & 2 & 8 \\
24) &  & 16 & 5 & 1 & 6 & 15 & 0 & 8 & 1 & 6 \\
25) &  & 21 & 10 & 1 & 1 & 15 & 7 & 0 & 9 & 1 \\
26) &  & 21 & 10 & 1 & 1 & 15 & 10 & 3 & 6 & 1 \\
27) &  & 36 & 10 & 1 & 16 & 0 & 16 & 9 & 0 & 1 \\
28) &  & 52 & 27 & 0 & 0 & 0 & 26 & 0 & 1 & 0 \\
\end{longtable}
\noindent Modules tested: V133, V912, V56.\\

\setcounter{table}{137}
\begin{longtable}{rc|cccc|cccc}
\caption{Alt$_{5} < E_7$, $p = 3$}\\
& & \multicolumn{4}{c|}{$V_{133}$} & \multicolumn{4}{c}{$V_{56}$} \\
& & $1$ & $3_{a}$ & $3_{b}$ & $4$ & $1$ & $3_{a}$ & $3_{b}$ & $4$ \\ \hline
\endfirsthead
& & \multicolumn{4}{c|}{$V_{133}$} & \multicolumn{4}{c}{$V_{56}$} \\
& & $1$ & $3_{a}$ & $3_{b}$ & $4$ & $1$ & $3_{a}$ & $3_{b}$ & $4$ \\ \hline
\endhead
1) & \possprim, \nongcr & 17 & 7 & 5 & 20 & 1 & 4 & 5 & 7 \\
2) & \nongcr & 21 & 10 & 6 & 16 & 2 & 9 & 1 & 6 \\
3) & \possprim, \nongcr & 16 & 6 & 5 & 21 & 2 & 6 & 4 & 6 \\
4) & \nongcr & 19 & 7 & 7 & 18 & 4 & 6 & 6 & 4 \\
5) & \nongcr & 27 & 22 & 0 & 10 & 6 & 14 & 0 & 2 \\
6) & \nongcr & 22 & 15 & 2 & 15 & 6 & 10 & 4 & 2 \\
7) & \possprim, \nongcr & 17 & 7 & 5 & 20 & 9 & 0 & 1 & 11 \\
8) & \nongcr & 16 & 6 & 5 & 21 & 10 & 2 & 0 & 10 \\
9) & \nongcr & 19 & 7 & 7 & 18 & 12 & 2 & 2 & 8 \\
10) & \nongcr & 22 & 15 & 2 & 15 & 14 & 6 & 0 & 6 \\
11) & \nongcr & 36 & 31 & 0 & 1 & 20 & 12 & 0 & 0 \\
\end{longtable}
\noindent Modules tested: V133, V912, L856, V56.\\

\setcounter{table}{138}
\begin{longtable}{rc|cccc|ccc}
\caption{$2\cdot$Alt$_{5} < E_7$, $p = 3$}\\
& & \multicolumn{4}{c|}{$V_{133}$} & \multicolumn{3}{c}{$V_{56}$} \\
& & $1$ & $3_{a}$ & $3_{b}$ & $4$ & $2_{a}$ & $2_{b}$ & $6$ \\ \hline
\endfirsthead
& & \multicolumn{4}{c|}{$V_{133}$} & \multicolumn{3}{c}{$V_{56}$} \\
& & $1$ & $3_{a}$ & $3_{b}$ & $4$ & $2_{a}$ & $2_{b}$ & $6$ \\ \hline
\endhead
1) & \nongcr & 15 & 13 & 9 & 13 & 1 & 9 & 6 \\
2) & \nongcr & 31 & 5 & 1 & 21 & 1 & 9 & 6 \\
3) &  & 52 & 27 & 0 & 0 & 27 & 1 & 0 \\
4) & \possprim, \nongcr & 13 & 10 & 10 & 15 & 2 & 2 & 8 \\
5) & \nongcr & 29 & 2 & 2 & 23 & 2 & 2 & 8 \\
6) & \possprim, \nongcr & 10 & 9 & 8 & 18 & 4 & 6 & 6 \\
7) & \nongcr & 26 & 1 & 0 & 26 & 4 & 6 & 6 \\
8) & \possprim, \nongcr & 11 & 10 & 8 & 17 & 7 & 6 & 5 \\
9) & \nongcr & 27 & 2 & 0 & 25 & 7 & 6 & 5 \\
10) & \possprim, \nongcr & 13 & 15 & 5 & 15 & 12 & 7 & 3 \\
11) & \nongcr & 20 & 18 & 9 & 8 & 16 & 9 & 1 \\
12) & \nongcr & 36 & 10 & 1 & 16 & 16 & 9 & 1 \\
13) & \nongcr & 21 & 28 & 0 & 7 & 21 & 7 & 0 \\
\end{longtable}
\noindent Modules tested: V133, V912, L856, V56.\\

\setcounter{table}{139}
\begin{longtable}{rc|ccccccc|ccccc}
\caption{$M_{11} < E_7$, $p = 0$ or $p \nmid |H|$}\\
& & \multicolumn{7}{c|}{$V_{133}$} & \multicolumn{5}{c}{$V_{56}$} \\
& & $1$ & $10_{a}$ & $11$ & $16_{a}$ & $16_{b}$ & $45$ & $55$ & $1$ & $10_{a}$ & $11$ & $16_{a}$ & $16_{b}$ \\ \hline
\endfirsthead
& & \multicolumn{7}{c|}{$V_{133}$} & \multicolumn{5}{c}{$V_{56}$} \\
& & $1$ & $10_{a}$ & $11$ & $16_{a}$ & $16_{b}$ & $45$ & $55$ & $1$ & $10_{a}$ & $11$ & $16_{a}$ & $16_{b}$ \\ \hline
\endhead
1) &  & 2 & 0 & 2 & 2 & 2 & 1 & 0 & 2 & 0 & 2 & 1 & 1 \\
2) &  & 3 & 0 & 1 & 2 & 2 & 0 & 1 & 2 & 0 & 2 & 1 & 1 \\
3) &  & 4 & 2 & 0 & 2 & 2 & 1 & 0 & 4 & 2 & 0 & 1 & 1 \\
\end{longtable}
\noindent Modules tested: V133, V912, V56.\\

\setcounter{table}{140}
\begin{longtable}{rc|ccccccc|cccc}
\caption{$M_{11} < E_7$, $p = 11$}\\
& & \multicolumn{7}{c|}{$V_{133}$} & \multicolumn{4}{c}{$V_{56}$} \\
& & $1$ & $9$ & $10_{a}$ & $10_{b}$ & $11$ & $16$ & $55$ & $1$ & $9$ & $11$ & $16$ \\ \hline
\endfirsthead
& & \multicolumn{7}{c|}{$V_{133}$} & \multicolumn{4}{c}{$V_{56}$} \\
& & $1$ & $9$ & $10_{a}$ & $10_{b}$ & $11$ & $16$ & $55$ & $1$ & $9$ & $11$ & $16$ \\ \hline
\endhead
1) &  & 2 & 1 & 1 & 1 & 2 & 5 & 0 & 2 & 0 & 2 & 2 \\
2) &  & 3 & 0 & 0 & 0 & 1 & 4 & 1 & 2 & 0 & 2 & 2 \\
3) & \nongcr & 6 & 3 & 1 & 1 & 0 & 5 & 0 & 6 & 2 & 0 & 2 \\
\end{longtable}
\noindent Modules tested: V133, V912, V56.\\

\setcounter{table}{141}
\begin{longtable}{rc|ccccccc|ccccc}
\caption{$M_{11} < E_7$, $p = 5$}\\
& & \multicolumn{7}{c|}{$V_{133}$} & \multicolumn{5}{c}{$V_{56}$} \\
& & $1$ & $10_{a}$ & $11$ & $16_{a}$ & $16_{b}$ & $45$ & $55$ & $1$ & $10_{a}$ & $11$ & $16_{a}$ & $16_{b}$ \\ \hline
\endfirsthead
& & \multicolumn{7}{c|}{$V_{133}$} & \multicolumn{5}{c}{$V_{56}$} \\
& & $1$ & $10_{a}$ & $11$ & $16_{a}$ & $16_{b}$ & $45$ & $55$ & $1$ & $10_{a}$ & $11$ & $16_{a}$ & $16_{b}$ \\ \hline
\endhead
1) &  & 1 & 1 & 2 & 0 & 0 & 1 & 1 & 2 & 0 & 2 & 1 & 1 \\
2) & \nongcr & 2 & 0 & 2 & 2 & 2 & 1 & 0 & 2 & 0 & 2 & 1 & 1 \\
3) &  & 2 & 1 & 1 & 0 & 0 & 0 & 2 & 2 & 0 & 2 & 1 & 1 \\
4) & \nongcr & 3 & 0 & 1 & 2 & 2 & 0 & 1 & 2 & 0 & 2 & 1 & 1 \\
5) &  & 3 & 3 & 0 & 0 & 0 & 1 & 1 & 4 & 2 & 0 & 1 & 1 \\
6) & \nongcr & 4 & 2 & 0 & 2 & 2 & 1 & 0 & 4 & 2 & 0 & 1 & 1 \\
\end{longtable}
\noindent Modules tested: V133, V912, V56.\\

\setcounter{table}{142}
\begin{longtable}{rc|cccccccc|cccccc}
\caption{$M_{11} < E_7$, $p = 3$}\\
& & \multicolumn{8}{c|}{$V_{133}$} & \multicolumn{6}{c}{$V_{56}$} \\
& & $1$ & $5_{a}$ & $5_{b}$ & $10_{a}$ & $10_{b}$ & $10_{c}$ & $24$ & $45$ & $1$ & $5_{a}$ & $5_{b}$ & $10_{a}$ & $10_{b}$ & $10_{c}$ \\ \hline
\endfirsthead
& & \multicolumn{8}{c|}{$V_{133}$} & \multicolumn{6}{c}{$V_{56}$} \\
& & $1$ & $5_{a}$ & $5_{b}$ & $10_{a}$ & $10_{b}$ & $10_{c}$ & $24$ & $45$ & $1$ & $5_{a}$ & $5_{b}$ & $10_{a}$ & $10_{b}$ & $10_{c}$ \\ \hline
\endhead
1) & \nongcr & 8 & 2 & 2 & 2 & 2 & 2 & 0 & 1 & 6 & 1 & 1 & 2 & 1 & 1 \\
2) & \nongcr & 8 & 4 & 4 & 0 & 2 & 2 & 0 & 1 & 6 & 3 & 3 & 0 & 1 & 1 \\
3) & \nongcr & 9 & 4 & 4 & 0 & 3 & 3 & 1 & 0 & 6 & 3 & 3 & 0 & 1 & 1 \\
\end{longtable}
\noindent Modules tested: V133, V912, L856, V56.\\

\setcounter{table}{143}
\begin{longtable}{rc|ccccc|cccc}
\caption{$M_{11} < E_7$, $p = 2$}\\
& & \multicolumn{5}{c|}{$V_{133}$} & \multicolumn{4}{c}{$V_{56}$} \\
& & $1$ & $10$ & $16_{a}$ & $16_{b}$ & $44$ & $1$ & $10$ & $16_{a}$ & $16_{b}$ \\ \hline
\endfirsthead
& & \multicolumn{5}{c|}{$V_{133}$} & \multicolumn{4}{c}{$V_{56}$} \\
& & $1$ & $10$ & $16_{a}$ & $16_{b}$ & $44$ & $1$ & $10$ & $16_{a}$ & $16_{b}$ \\ \hline
\endhead
1) & \nongcr & 5 & 2 & 2 & 2 & 1 & 4 & 2 & 1 & 1 \\
2) & \nongcr & 7 & 5 & 1 & 1 & 1 & 4 & 2 & 1 & 1 \\
\end{longtable}
\noindent Modules tested: V133, L132, V912, V56.\\

\setcounter{table}{144}
\begin{longtable}{rc|cccc|cc}
\caption{$2\cdot M_{12} < E_7$, $p = 0$ or $p \nmid |H|$}\\
& & \multicolumn{4}{c|}{$V_{133}$} & \multicolumn{2}{c}{$V_{56}$} \\
& & $1$ & $16_{a}$ & $16_{b}$ & $66$ & $12$ & $32$ \\ \hline
\endfirsthead
& & \multicolumn{4}{c|}{$V_{133}$} & \multicolumn{2}{c}{$V_{56}$} \\
& & $1$ & $16_{a}$ & $16_{b}$ & $66$ & $12$ & $32$ \\ \hline
\endhead
1) &  & 3 & 2 & 2 & 1 & 2 & 1 \\
\end{longtable}
\noindent Modules tested: V133, V912, V56.\\

\setcounter{table}{145}
\begin{longtable}{rc|ccc|cc}
\caption{$2\cdot M_{12} < E_7$, $p = 11$}\\
& & \multicolumn{3}{c|}{$V_{133}$} & \multicolumn{2}{c}{$V_{56}$} \\
& & $1$ & $16$ & $66$ & $12$ & $32$ \\ \hline
\endfirsthead
& & \multicolumn{3}{c|}{$V_{133}$} & \multicolumn{2}{c}{$V_{56}$} \\
& & $1$ & $16$ & $66$ & $12$ & $32$ \\ \hline
\endhead
1) &  & 3 & 4 & 1 & 2 & 1 \\
\end{longtable}
\noindent Modules tested: V133, V912, V56.\\

\setcounter{table}{146}
\begin{longtable}{rc|ccccc|cccc}
\caption{$M_{12} < E_7$, $p = 5$}\\
& & \multicolumn{5}{c|}{$V_{133}$} & \multicolumn{4}{c}{$V_{56}$} \\
& & $1$ & $11_{a}$ & $16_{a}$ & $16_{b}$ & $78$ & $1$ & $11_{a}$ & $16_{a}$ & $16_{b}$ \\ \hline
\endfirsthead
& & \multicolumn{5}{c|}{$V_{133}$} & \multicolumn{4}{c}{$V_{56}$} \\
& & $1$ & $11_{a}$ & $16_{a}$ & $16_{b}$ & $78$ & $1$ & $11_{a}$ & $16_{a}$ & $16_{b}$ \\ \hline
\endhead
1) &  & 1 & 2 & 1 & 1 & 1 & 2 & 2 & 1 & 1 \\
\end{longtable}
\noindent Modules tested: V133, V912, V56.\\

\setcounter{table}{147}
\begin{longtable}{rc|cccccc|cc}
\caption{$2\cdot M_{12} < E_7$, $p = 5$}\\
& & \multicolumn{6}{c|}{$V_{133}$} & \multicolumn{2}{c}{$V_{56}$} \\
& & $1$ & $16_{a}$ & $16_{b}$ & $55_{a}$ & $66$ & $78$ & $12$ & $32$ \\ \hline
\endfirsthead
& & \multicolumn{6}{c|}{$V_{133}$} & \multicolumn{2}{c}{$V_{56}$} \\
& & $1$ & $16_{a}$ & $16_{b}$ & $55_{a}$ & $66$ & $78$ & $12$ & $32$ \\ \hline
\endhead
1) & \possprim & 0 & 0 & 0 & 1 & 0 & 1 & 2 & 1 \\
2) &  & 3 & 2 & 2 & 0 & 1 & 0 & 2 & 1 \\
\end{longtable}
\noindent Modules tested: V133, V912, V56.\\

\setcounter{table}{148}
\begin{longtable}{rc|cccc|cccc}
\caption{$2\cdot M_{12} < E_7$, $p = 3$}\\
& & \multicolumn{4}{c|}{$V_{133}$} & \multicolumn{4}{c}{$V_{56}$} \\
& & $1$ & $15_{a}$ & $15_{b}$ & $34$ & $6_{a}$ & $6_{b}$ & $10_{c}$ & $10_{d}$ \\ \hline
\endfirsthead
& & \multicolumn{4}{c|}{$V_{133}$} & \multicolumn{4}{c}{$V_{56}$} \\
& & $1$ & $15_{a}$ & $15_{b}$ & $34$ & $6_{a}$ & $6_{b}$ & $10_{c}$ & $10_{d}$ \\ \hline
\endhead
1) & \nongcr & 9 & 3 & 3 & 1 & 3 & 3 & 1 & 1 \\
\end{longtable}
\noindent Modules tested: V133, V912, L856, V56.\\

\setcounter{table}{149}
\begin{longtable}{rc|ccccc|cccc}
\caption{$M_{12} < E_7$, $p = 2$}\\
& & \multicolumn{5}{c|}{$V_{133}$} & \multicolumn{4}{c}{$V_{56}$} \\
& & $1$ & $10$ & $16_{a}$ & $16_{b}$ & $44$ & $1$ & $10$ & $16_{a}$ & $16_{b}$ \\ \hline
\endfirsthead
& & \multicolumn{5}{c|}{$V_{133}$} & \multicolumn{4}{c}{$V_{56}$} \\
& & $1$ & $10$ & $16_{a}$ & $16_{b}$ & $44$ & $1$ & $10$ & $16_{a}$ & $16_{b}$ \\ \hline
\endhead
1) & \nongcr & 5 & 2 & 2 & 2 & 1 & 4 & 2 & 1 & 1 \\
\end{longtable}
\noindent Modules tested: V133, L132, V912, V56.\\

\setcounter{table}{150}
\begin{longtable}{rc|c|cc}
\caption{$2\cdot M_{22} < E_7$, $p = 5$}\\
& & \multicolumn{1}{c|}{$V_{133}$} & \multicolumn{2}{c}{$V_{56}$} \\
& & $133$ & $28_{a}$ & $28_{b}$ \\ \hline
\endfirsthead
& & \multicolumn{1}{c|}{$V_{133}$} & \multicolumn{2}{c}{$V_{56}$} \\
& & $133$ & $28_{a}$ & $28_{b}$ \\ \hline
\endhead
1) & \possprim & 1 & 1 & 1 \\
\end{longtable}
\noindent Modules tested: V133, V912, V56.\\

\setcounter{table}{151}
\begin{longtable}{rc|ccccc|cccc}
\caption{$J_{1} < E_7$, $p = 11$}\\
& & \multicolumn{5}{c|}{$V_{133}$} & \multicolumn{4}{c}{$V_{56}$} \\
& & $1$ & $7$ & $14$ & $27$ & $64$ & $1$ & $7$ & $14$ & $27$ \\ \hline
\endfirsthead
& & \multicolumn{5}{c|}{$V_{133}$} & \multicolumn{4}{c}{$V_{56}$} \\
& & $1$ & $7$ & $14$ & $27$ & $64$ & $1$ & $7$ & $14$ & $27$ \\ \hline
\endhead
1) &  & 3 & 5 & 1 & 3 & 0 & 0 & 4 & 2 & 0 \\
2) &  & 1 & 0 & 1 & 2 & 1 & 2 & 0 & 0 & 2 \\
3) &  & 21 & 14 & 1 & 0 & 0 & 14 & 6 & 0 & 0 \\
\end{longtable}
\noindent Modules tested: V133, V912, V56.\\

\setcounter{table}{152}
\begin{longtable}{rc|ccc|cc}
\caption{$2 \cdot J_{2} < E_7$, $p = 0$ or $p \nmid |H|$}\\
& & \multicolumn{3}{c|}{$V_{133}$} & \multicolumn{2}{c}{$V_{56}$} \\
& & $1$ & $14_{a}$ & $21_{b}$ & $6_{a}$ & $14_{c}$ \\ \hline
\endfirsthead
& & \multicolumn{3}{c|}{$V_{133}$} & \multicolumn{2}{c}{$V_{56}$} \\
& & $1$ & $14_{a}$ & $21_{b}$ & $6_{a}$ & $14_{c}$ \\ \hline
\endhead
1) &  & 14 & 7 & 1 & 7 & 1 \\
\end{longtable}
\noindent Modules tested: V133, V912, V56.\\

\setcounter{table}{153}
\begin{longtable}{rc|ccc|cc}
\caption{$2 \cdot J_{2} < E_7$, $p = 7$}\\
& & \multicolumn{3}{c|}{$V_{133}$} & \multicolumn{2}{c}{$V_{56}$} \\
& & $1$ & $14_{a}$ & $21_{b}$ & $6_{a}$ & $14_{c}$ \\ \hline
\endfirsthead
& & \multicolumn{3}{c|}{$V_{133}$} & \multicolumn{2}{c}{$V_{56}$} \\
& & $1$ & $14_{a}$ & $21_{b}$ & $6_{a}$ & $14_{c}$ \\ \hline
\endhead
1) &  & 14 & 7 & 1 & 7 & 1 \\
\end{longtable}
\noindent Modules tested: V133, V912, V56.\\

\setcounter{table}{154}
\begin{longtable}{rc|ccc|cc}
\caption{$2 \cdot J_{2} < E_7$, $p = 5$}\\
& & \multicolumn{3}{c|}{$V_{133}$} & \multicolumn{2}{c}{$V_{56}$} \\
& & $1$ & $14_{a}$ & $21$ & $6$ & $14_{b}$ \\ \hline
\endfirsthead
& & \multicolumn{3}{c|}{$V_{133}$} & \multicolumn{2}{c}{$V_{56}$} \\
& & $1$ & $14_{a}$ & $21$ & $6$ & $14_{b}$ \\ \hline
\endhead
1) &  & 14 & 7 & 1 & 7 & 1 \\
\end{longtable}
\noindent Modules tested: V133, V912, V56.\\

\setcounter{table}{155}
\begin{longtable}{rc|ccc|cc}
\caption{$2 \cdot J_{2} < E_7$, $p = 3$}\\
& & \multicolumn{3}{c|}{$V_{133}$} & \multicolumn{2}{c}{$V_{56}$} \\
& & $1$ & $13_{a}$ & $21_{a}$ & $6_{b}$ & $14$ \\ \hline
\endfirsthead
& & \multicolumn{3}{c|}{$V_{133}$} & \multicolumn{2}{c}{$V_{56}$} \\
& & $1$ & $13_{a}$ & $21_{a}$ & $6_{b}$ & $14$ \\ \hline
\endhead
1) & \nongcr & 21 & 7 & 1 & 7 & 1 \\
\end{longtable}
\noindent Modules tested: V133, V912, L856, V56.\\

\setcounter{table}{156}
\begin{longtable}{rc|ccccccc|ccccc}
\caption{$J_{2} < E_7$, $p = 2$}\\
& & \multicolumn{7}{c|}{$V_{133}$} & \multicolumn{5}{c}{$V_{56}$} \\
& & $1$ & $6_{a}$ & $6_{b}$ & $14_{a}$ & $14_{b}$ & $64_{b}$ & $84$ & $1$ & $6_{a}$ & $6_{b}$ & $14_{a}$ & $36$ \\ \hline
\endfirsthead
& & \multicolumn{7}{c|}{$V_{133}$} & \multicolumn{5}{c}{$V_{56}$} \\
& & $1$ & $6_{a}$ & $6_{b}$ & $14_{a}$ & $14_{b}$ & $64_{b}$ & $84$ & $1$ & $6_{a}$ & $6_{b}$ & $14_{a}$ & $36$ \\ \hline
\endhead
1) & \possprim & 1 & 1 & 0 & 1 & 2 & 0 & 1 & 2 & 1 & 2 & 0 & 1 \\
2) &  & 15 & 1 & 0 & 0 & 8 & 0 & 0 & 2 & 1 & 8 & 0 & 0 \\
3) & \nongcr & 11 & 8 & 3 & 4 & 0 & 0 & 0 & 4 & 4 & 0 & 2 & 0 \\
4) & \nongcr & 3 & 2 & 2 & 3 & 0 & 1 & 0 & 4 & 2 & 2 & 2 & 0 \\
5) & \nongcr & 35 & 14 & 0 & 1 & 0 & 0 & 0 & 20 & 6 & 0 & 0 & 0 \\
\end{longtable}
\noindent Modules tested: V133, L132, V912, V56.\\

\setcounter{table}{157}
\begin{longtable}{rc|c|cc}
\caption{$2 \cdot Ru < E_7$, $p = 5$}\\
& & \multicolumn{1}{c|}{$V_{133}$} & \multicolumn{2}{c}{$V_{56}$} \\
& & $133$ & $28_{a}$ & $28_{b}$ \\ \hline
\endfirsthead
& & \multicolumn{1}{c|}{$V_{133}$} & \multicolumn{2}{c}{$V_{56}$} \\
& & $133$ & $28_{a}$ & $28_{b}$ \\ \hline
\endhead
1) & \possprim & 1 & 1 & 1 \\
\end{longtable}
\noindent Modules not used (elements of large order): V912.\\
Modules tested: V133, V56.\\

\setcounter{table}{158}
\begin{longtable}{rc|c|cc}
\caption{$2 \cdot HS < E_7$, $p = 5$}\\
& & \multicolumn{1}{c|}{$V_{133}$} & \multicolumn{2}{c}{$V_{56}$} \\
& & $133_{a}$ & $28_{a}$ & $28_{b}$ \\ \hline
\endfirsthead
& & \multicolumn{1}{c|}{$V_{133}$} & \multicolumn{2}{c}{$V_{56}$} \\
& & $133_{a}$ & $28_{a}$ & $28_{b}$ \\ \hline
\endhead
1) & \possprim & 1 & 1 & 1 \\
\end{longtable}
\noindent Modules tested: V133, V912, V56.\\

\setcounter{table}{159}
\begin{longtable}{rc|cccccc|cccccc}
\caption{$L_2(7) < E_7$, $p = 0$ or $p \nmid |H|$}\\
& & \multicolumn{6}{c|}{$V_{133}$} & \multicolumn{6}{c}{$V_{56}$} \\
& & $1$ & $3_{a}$ & $3_{b}$ & $6$ & $7$ & $8$ & $1$ & $3_{a}$ & $3_{b}$ & $6$ & $7$ & $8$ \\ \hline
\endfirsthead
& & \multicolumn{6}{c|}{$V_{133}$} & \multicolumn{6}{c}{$V_{56}$} \\
& & $1$ & $3_{a}$ & $3_{b}$ & $6$ & $7$ & $8$ & $1$ & $3_{a}$ & $3_{b}$ & $6$ & $7$ & $8$ \\ \hline
\endhead
1) & \possprim & 0 & 1 & 1 & 6 & 5 & 7 & 0 & 2 & 2 & 0 & 4 & 2 \\
2) &  & 1 & 2 & 2 & 6 & 4 & 7 & 2 & 0 & 0 & 4 & 2 & 2 \\
3) &  & 2 & 3 & 3 & 6 & 3 & 7 & 0 & 2 & 2 & 0 & 4 & 2 \\
4) &  & 3 & 1 & 1 & 6 & 8 & 4 & 0 & 2 & 2 & 0 & 4 & 2 \\
5) &  & 3 & 4 & 4 & 6 & 2 & 7 & 4 & 2 & 2 & 4 & 0 & 2 \\
6) &  & 4 & 2 & 2 & 6 & 7 & 4 & 2 & 0 & 0 & 4 & 2 & 2 \\
7) &  & 5 & 3 & 3 & 6 & 6 & 4 & 0 & 2 & 2 & 0 & 4 & 2 \\
8) &  & 6 & 4 & 4 & 6 & 5 & 4 & 4 & 2 & 2 & 4 & 0 & 2 \\
9) &  & 8 & 0 & 0 & 2 & 7 & 8 & 0 & 1 & 1 & 6 & 2 & 0 \\
10) &  & 9 & 1 & 1 & 0 & 2 & 13 & 8 & 0 & 0 & 0 & 0 & 6 \\
11) &  & 10 & 2 & 2 & 2 & 5 & 8 & 2 & 3 & 3 & 6 & 0 & 0 \\
12) &  & 11 & 9 & 9 & 6 & 0 & 4 & 4 & 6 & 6 & 0 & 0 & 2 \\
13) &  & 12 & 1 & 1 & 0 & 5 & 10 & 8 & 0 & 0 & 0 & 0 & 6 \\
14) &  & 15 & 7 & 7 & 2 & 0 & 8 & 2 & 7 & 7 & 2 & 0 & 0 \\
15) &  & 21 & 1 & 1 & 0 & 14 & 1 & 14 & 0 & 0 & 0 & 6 & 0 \\
16) &  & 35 & 15 & 15 & 0 & 0 & 1 & 20 & 6 & 6 & 0 & 0 & 0 \\
\end{longtable}
\noindent Memoir rows not reproduced: 7) (not compatible with V912); 10) (not compatible with V912).\\
Modules tested: V133, V912, V56.\\

\setcounter{table}{160}
\begin{longtable}{rc|cccccc|ccccc}
\caption{$SL_2(7) < E_7$, $p = 0$ or $p \nmid |H|$}\\
& & \multicolumn{6}{c|}{$V_{133}$} & \multicolumn{5}{c}{$V_{56}$} \\
& & $1$ & $3_{a}$ & $3_{b}$ & $6_{a}$ & $7$ & $8_{a}$ & $4_{a}$ & $4_{b}$ & $6_{b}$ & $6_{c}$ & $8_{b}$ \\ \hline
\endfirsthead
& & \multicolumn{6}{c|}{$V_{133}$} & \multicolumn{5}{c}{$V_{56}$} \\
& & $1$ & $3_{a}$ & $3_{b}$ & $6_{a}$ & $7$ & $8_{a}$ & $4_{a}$ & $4_{b}$ & $6_{b}$ & $6_{c}$ & $8_{b}$ \\ \hline
\endhead
1) & \possprim & 0 & 4 & 4 & 3 & 5 & 7 & 2 & 2 & 3 & 1 & 2 \\
2) &  & 3 & 4 & 4 & 3 & 8 & 4 & 2 & 2 & 3 & 1 & 2 \\
3) &  & 14 & 1 & 1 & 7 & 1 & 8 & 1 & 1 & 7 & 1 & 0 \\
\end{longtable}
\noindent Modules tested: V133, V912, V56.\\

\setcounter{table}{161}
\begin{longtable}{rc|ccccc|ccccc}
\caption{$L_2(7) < E_7$, $p = 3$}\\
& & \multicolumn{5}{c|}{$V_{133}$} & \multicolumn{5}{c}{$V_{56}$} \\
& & $1$ & $3_{a}$ & $3_{b}$ & $6$ & $7$ & $1$ & $3_{a}$ & $3_{b}$ & $6$ & $7$ \\ \hline
\endfirsthead
& & \multicolumn{5}{c|}{$V_{133}$} & \multicolumn{5}{c}{$V_{56}$} \\
& & $1$ & $3_{a}$ & $3_{b}$ & $6$ & $7$ & $1$ & $3_{a}$ & $3_{b}$ & $6$ & $7$ \\ \hline
\endhead
1) & \possprim, \nongcr & 7 & 1 & 1 & 6 & 12 & 2 & 2 & 2 & 0 & 6 \\
2) & \nongcr & 8 & 2 & 2 & 6 & 11 & 4 & 0 & 0 & 4 & 4 \\
3) & \possprim, \nongcr & 9 & 3 & 3 & 6 & 10 & 2 & 2 & 2 & 0 & 6 \\
4) & \nongcr & 10 & 4 & 4 & 6 & 9 & 6 & 2 & 2 & 4 & 2 \\
5) & \nongcr & 15 & 9 & 9 & 6 & 4 & 6 & 6 & 6 & 0 & 2 \\
6) & \nongcr & 16 & 0 & 0 & 2 & 15 & 0 & 1 & 1 & 6 & 2 \\
7) & \nongcr & 18 & 2 & 2 & 2 & 13 & 2 & 3 & 3 & 6 & 0 \\
8) & \nongcr & 22 & 1 & 1 & 0 & 15 & 14 & 0 & 0 & 0 & 6 \\
9) & \nongcr & 23 & 7 & 7 & 2 & 8 & 2 & 7 & 7 & 2 & 0 \\
10) & \nongcr & 36 & 15 & 15 & 0 & 1 & 20 & 6 & 6 & 0 & 0 \\
\end{longtable}
\noindent Modules tested: V133, V912, L856, V56.\\

\setcounter{table}{162}
\begin{longtable}{rc|ccccc|cccc}
\caption{$SL_2(7) < E_7$, $p = 3$}\\
& & \multicolumn{5}{c|}{$V_{133}$} & \multicolumn{4}{c}{$V_{56}$} \\
& & $1$ & $3_{a}$ & $3_{b}$ & $6_{a}$ & $7$ & $4_{a}$ & $4_{b}$ & $6_{b}$ & $6_{c}$ \\ \hline
\endfirsthead
& & \multicolumn{5}{c|}{$V_{133}$} & \multicolumn{4}{c}{$V_{56}$} \\
& & $1$ & $3_{a}$ & $3_{b}$ & $6_{a}$ & $7$ & $4_{a}$ & $4_{b}$ & $6_{b}$ & $6_{c}$ \\ \hline
\endhead
1) & \possprim, \nongcr & 7 & 4 & 4 & 3 & 12 & 4 & 4 & 3 & 1 \\
2) & \possprim, \nongcr & 22 & 1 & 1 & 7 & 9 & 1 & 1 & 7 & 1 \\
\end{longtable}
\noindent Modules tested: V133, V912, L856, V56.\\

\setcounter{table}{163}
\begin{longtable}{rc|ccccccccc|ccccccccc}
\caption{$L_2(8) < E_7$, $p = 0$ or $p \nmid |H|$}\\
& & \multicolumn{9}{c|}{$V_{133}$} & \multicolumn{9}{c}{$V_{56}$} \\
& & $1$ & $7_{a}$ & $7_{b}$ & $7_{c}$ & $7_{d}$ & $8$ & $9_{a}$ & $9_{b}$ & $9_{c}$ & $1$ & $7_{a}$ & $7_{b}$ & $7_{c}$ & $7_{d}$ & $8$ & $9_{a}$ & $9_{b}$ & $9_{c}$ \\ \hline
\endfirsthead
& & \multicolumn{9}{c|}{$V_{133}$} & \multicolumn{9}{c}{$V_{56}$} \\
& & $1$ & $7_{a}$ & $7_{b}$ & $7_{c}$ & $7_{d}$ & $8$ & $9_{a}$ & $9_{b}$ & $9_{c}$ & $1$ & $7_{a}$ & $7_{b}$ & $7_{c}$ & $7_{d}$ & $8$ & $9_{a}$ & $9_{b}$ & $9_{c}$ \\ \hline
\endhead
1) &  & 1 & 2 & 1 & 1 & 1 & 2 & 3 & 3 & 3 & 0 & 2 & 2 & 2 & 2 & 0 & 0 & 0 & 0 \\
2) &  & 1 & 2 & 1 & 1 & 1 & 2 & 3 & 3 & 3 & 2 & 0 & 0 & 0 & 0 & 0 & 2 & 2 & 2 \\
3) &  & 3 & 1 & 3 & 2 & 1 & 0 & 3 & 3 & 3 & 0 & 2 & 3 & 1 & 2 & 0 & 0 & 0 & 0 \\
4) &  & 3 & 1 & 5 & 1 & 0 & 0 & 3 & 3 & 3 & 0 & 2 & 4 & 2 & 0 & 0 & 0 & 0 & 0 \\
5) &  & 3 & 2 & 1 & 1 & 1 & 4 & 4 & 0 & 3 & 4 & 0 & 0 & 0 & 0 & 2 & 2 & 0 & 2 \\
6) &  & 5 & 2 & 1 & 1 & 1 & 6 & 4 & 0 & 1 & 3 & 0 & 0 & 0 & 0 & 1 & 5 & 0 & 0 \\
7) &  & 6 & 5 & 1 & 1 & 1 & 1 & 4 & 0 & 3 & 6 & 2 & 0 & 0 & 0 & 0 & 2 & 0 & 2 \\
8) &  & 10 & 2 & 1 & 1 & 1 & 11 & 0 & 0 & 0 & 8 & 0 & 0 & 0 & 0 & 6 & 0 & 0 & 0 \\
9) &  & 13 & 5 & 1 & 1 & 1 & 8 & 0 & 0 & 0 & 10 & 2 & 0 & 0 & 0 & 4 & 0 & 0 & 0 \\
10) &  & 21 & 1 & 14 & 1 & 0 & 0 & 0 & 0 & 0 & 14 & 0 & 6 & 0 & 0 & 0 & 0 & 0 & 0 \\
\end{longtable}
\noindent Memoir rows not reproduced: 3) (not compatible with V912).\\
Modules tested: V133, V912, V56.\\

\setcounter{table}{164}
\begin{longtable}{rc|cccccc|cccccc}
\caption{$L_2(8) < E_7$, $p = 7$}\\
& & \multicolumn{6}{c|}{$V_{133}$} & \multicolumn{6}{c}{$V_{56}$} \\
& & $1$ & $7_{a}$ & $7_{b}$ & $7_{c}$ & $7_{d}$ & $8$ & $1$ & $7_{a}$ & $7_{b}$ & $7_{c}$ & $7_{d}$ & $8$ \\ \hline
\endfirsthead
& & \multicolumn{6}{c|}{$V_{133}$} & \multicolumn{6}{c}{$V_{56}$} \\
& & $1$ & $7_{a}$ & $7_{b}$ & $7_{c}$ & $7_{d}$ & $8$ & $1$ & $7_{a}$ & $7_{b}$ & $7_{c}$ & $7_{d}$ & $8$ \\ \hline
\endhead
1) & \possprim, \nongcr & 10 & 2 & 1 & 1 & 1 & 11 & 0 & 2 & 2 & 2 & 2 & 0 \\
2) & \nongcr & 10 & 2 & 1 & 1 & 1 & 11 & 8 & 0 & 0 & 0 & 0 & 6 \\
3) & \nongcr & 12 & 1 & 5 & 1 & 0 & 9 & 0 & 2 & 4 & 2 & 0 & 0 \\
4) & \nongcr & 12 & 1 & 3 & 2 & 1 & 9 & 0 & 2 & 3 & 1 & 2 & 0 \\
5) & \nongcr & 13 & 2 & 5 & 1 & 0 & 8 & 0 & 5 & 3 & 0 & 0 & 0 \\
6) & \nongcr & 13 & 5 & 1 & 1 & 1 & 8 & 10 & 2 & 0 & 0 & 0 & 4 \\
7) &  & 21 & 1 & 14 & 1 & 0 & 0 & 14 & 0 & 6 & 0 & 0 & 0 \\
\end{longtable}
\noindent Modules tested: V133, V912, V56.\\

\setcounter{table}{165}
\begin{longtable}{rc|ccccc|ccccc}
\caption{$L_2(8) < E_7$, $p = 3$}\\
& & \multicolumn{5}{c|}{$V_{133}$} & \multicolumn{5}{c}{$V_{56}$} \\
& & $1$ & $7$ & $9_{a}$ & $9_{b}$ & $9_{c}$ & $1$ & $7$ & $9_{a}$ & $9_{b}$ & $9_{c}$ \\ \hline
\endfirsthead
& & \multicolumn{5}{c|}{$V_{133}$} & \multicolumn{5}{c}{$V_{56}$} \\
& & $1$ & $7$ & $9_{a}$ & $9_{b}$ & $9_{c}$ & $1$ & $7$ & $9_{a}$ & $9_{b}$ & $9_{c}$ \\ \hline
\endhead
1) & \possprim, \nongcr & 3 & 7 & 3 & 3 & 3 & 0 & 8 & 0 & 0 & 0 \\
2) & \nongcr & 3 & 7 & 3 & 3 & 3 & 2 & 0 & 2 & 2 & 2 \\
3) & \nongcr & 5 & 8 & 3 & 3 & 2 & 2 & 0 & 1 & 3 & 2 \\
4) & \nongcr & 7 & 9 & 4 & 0 & 3 & 6 & 2 & 2 & 0 & 2 \\
5) & \nongcr & 11 & 11 & 4 & 0 & 1 & 4 & 1 & 5 & 0 & 0 \\
6) & \nongcr & 21 & 16 & 0 & 0 & 0 & 14 & 6 & 0 & 0 & 0 \\
\end{longtable}
\noindent Modules tested: V133, V912, L856, V56.\\

\setcounter{table}{166}
\begin{longtable}{rc|cccccccc|ccccccc}
\caption{$L_2(11) < E_7$, $p = 0$ or $p \nmid |H|$}\\
& & \multicolumn{8}{c|}{$V_{133}$} & \multicolumn{7}{c}{$V_{56}$} \\
& & $1$ & $5_{a}$ & $5_{b}$ & $10_{a}$ & $10_{b}$ & $11$ & $12_{a}$ & $12_{b}$ & $1$ & $5_{a}$ & $5_{b}$ & $10_{a}$ & $10_{b}$ & $11$ & $12_{a}$ \\ \hline
\endfirsthead
& & \multicolumn{8}{c|}{$V_{133}$} & \multicolumn{7}{c}{$V_{56}$} \\
& & $1$ & $5_{a}$ & $5_{b}$ & $10_{a}$ & $10_{b}$ & $11$ & $12_{a}$ & $12_{b}$ & $1$ & $5_{a}$ & $5_{b}$ & $10_{a}$ & $10_{b}$ & $11$ & $12_{a}$ \\ \hline
\endhead
1) &  & 1 & 2 & 2 & 1 & 2 & 2 & 3 & 2 & 2 & 1 & 1 & 0 & 2 & 0 & 2 \\
2) &  & 2 & 0 & 0 & 2 & 4 & 1 & 3 & 2 & 2 & 1 & 1 & 0 & 2 & 0 & 2 \\
3) &  & 4 & 2 & 2 & 1 & 2 & 5 & 1 & 1 & 4 & 1 & 1 & 0 & 2 & 2 & 0 \\
4) &  & 5 & 0 & 0 & 2 & 4 & 4 & 1 & 1 & 4 & 1 & 1 & 0 & 2 & 2 & 0 \\
5) &  & 9 & 4 & 4 & 6 & 0 & 0 & 1 & 1 & 6 & 3 & 3 & 2 & 0 & 0 & 0 \\
\end{longtable}
\noindent Modules tested: V133, V912, V56.\\

\setcounter{table}{167}
\begin{longtable}{rc|cccccccc|cccccc}
\caption{$SL_2(11) < E_7$, $p = 0$ or $p \nmid |H|$}\\
& & \multicolumn{8}{c|}{$V_{133}$} & \multicolumn{6}{c}{$V_{56}$} \\
& & $1$ & $5_{a}$ & $5_{b}$ & $10_{a}$ & $10_{b}$ & $11$ & $12_{a}$ & $12_{b}$ & $6_{a}$ & $6_{b}$ & $10_{d}$ & $10_{e}$ & $12_{c}$ & $12_{d}$ \\ \hline
\endfirsthead
& & \multicolumn{8}{c|}{$V_{133}$} & \multicolumn{6}{c}{$V_{56}$} \\
& & $1$ & $5_{a}$ & $5_{b}$ & $10_{a}$ & $10_{b}$ & $11$ & $12_{a}$ & $12_{b}$ & $6_{a}$ & $6_{b}$ & $10_{d}$ & $10_{e}$ & $12_{c}$ & $12_{d}$ \\ \hline
\endhead
1) & \possprim & 0 & 1 & 1 & 3 & 0 & 3 & 3 & 2 & 1 & 1 & 1 & 1 & 2 & 0 \\
2) &  & 1 & 1 & 1 & 3 & 0 & 4 & 3 & 1 & 0 & 0 & 1 & 1 & 1 & 2 \\
3) &  & 3 & 1 & 1 & 3 & 0 & 6 & 1 & 1 & 3 & 3 & 1 & 1 & 0 & 0 \\
4) &  & 8 & 3 & 3 & 0 & 6 & 1 & 1 & 1 & 3 & 3 & 1 & 1 & 0 & 0 \\
\end{longtable}
\noindent Modules tested: V133, V912, V56.\\

\setcounter{table}{168}
\begin{longtable}{rc|cccccc|cccccc}
\caption{$L_2(11) < E_7$, $p = 5$}\\
& & \multicolumn{6}{c|}{$V_{133}$} & \multicolumn{6}{c}{$V_{56}$} \\
& & $1$ & $5_{a}$ & $5_{b}$ & $10_{a}$ & $10_{b}$ & $11$ & $1$ & $5_{a}$ & $5_{b}$ & $10_{a}$ & $10_{b}$ & $11$ \\ \hline
\endfirsthead
& & \multicolumn{6}{c|}{$V_{133}$} & \multicolumn{6}{c}{$V_{56}$} \\
& & $1$ & $5_{a}$ & $5_{b}$ & $10_{a}$ & $10_{b}$ & $11$ & $1$ & $5_{a}$ & $5_{b}$ & $10_{a}$ & $10_{b}$ & $11$ \\ \hline
\endhead
1) & \nongcr & 6 & 2 & 2 & 1 & 2 & 7 & 4 & 1 & 1 & 0 & 2 & 2 \\
2) & \nongcr & 7 & 0 & 0 & 2 & 4 & 6 & 4 & 1 & 1 & 0 & 2 & 2 \\
3) & \nongcr & 11 & 4 & 4 & 6 & 0 & 2 & 6 & 3 & 3 & 2 & 0 & 0 \\
\end{longtable}
\noindent Modules tested: V133, V912, V56.\\

\setcounter{table}{169}
\begin{longtable}{rc|cccccc|cccc}
\caption{$SL_2(11) < E_7$, $p = 5$}\\
& & \multicolumn{6}{c|}{$V_{133}$} & \multicolumn{4}{c}{$V_{56}$} \\
& & $1$ & $5_{a}$ & $5_{b}$ & $10_{a}$ & $10_{b}$ & $11$ & $6_{a}$ & $6_{b}$ & $10_{d}$ & $10_{e}$ \\ \hline
\endfirsthead
& & \multicolumn{6}{c|}{$V_{133}$} & \multicolumn{4}{c}{$V_{56}$} \\
& & $1$ & $5_{a}$ & $5_{b}$ & $10_{a}$ & $10_{b}$ & $11$ & $6_{a}$ & $6_{b}$ & $10_{d}$ & $10_{e}$ \\ \hline
\endhead
1) & \possprim, \nongcr & 5 & 1 & 1 & 3 & 0 & 8 & 3 & 3 & 1 & 1 \\
2) & \nongcr & 10 & 3 & 3 & 0 & 6 & 3 & 3 & 3 & 1 & 1 \\
\end{longtable}
\noindent Modules tested: V133, V912, V56.\\

\setcounter{table}{170}
\begin{longtable}{rc|cccccc|ccccc}
\caption{$L_2(11) < E_7$, $p = 3$}\\
& & \multicolumn{6}{c|}{$V_{133}$} & \multicolumn{5}{c}{$V_{56}$} \\
& & $1$ & $5_{a}$ & $5_{b}$ & $10$ & $12_{a}$ & $12_{b}$ & $1$ & $5_{a}$ & $5_{b}$ & $10$ & $12_{a}$ \\ \hline
\endfirsthead
& & \multicolumn{6}{c|}{$V_{133}$} & \multicolumn{5}{c}{$V_{56}$} \\
& & $1$ & $5_{a}$ & $5_{b}$ & $10$ & $12_{a}$ & $12_{b}$ & $1$ & $5_{a}$ & $5_{b}$ & $10$ & $12_{a}$ \\ \hline
\endhead
1) & \nongcr & 3 & 4 & 4 & 3 & 3 & 2 & 2 & 3 & 3 & 0 & 2 \\
2) & \nongcr & 9 & 4 & 4 & 6 & 1 & 1 & 6 & 3 & 3 & 2 & 0 \\
\end{longtable}
\noindent Modules tested: V133, V912, L856, V56.\\

\setcounter{table}{171}
\begin{longtable}{rc|cccccc|ccccc}
\caption{$SL_2(11) < E_7$, $p = 3$}\\
& & \multicolumn{6}{c|}{$V_{133}$} & \multicolumn{5}{c}{$V_{56}$} \\
& & $1$ & $5_{a}$ & $5_{b}$ & $10_{a}$ & $12_{a}$ & $12_{b}$ & $6_{a}$ & $6_{b}$ & $10_{b}$ & $12_{c}$ & $12_{d}$ \\ \hline
\endfirsthead
& & \multicolumn{6}{c|}{$V_{133}$} & \multicolumn{5}{c}{$V_{56}$} \\
& & $1$ & $5_{a}$ & $5_{b}$ & $10_{a}$ & $12_{a}$ & $12_{b}$ & $6_{a}$ & $6_{b}$ & $10_{b}$ & $12_{c}$ & $12_{d}$ \\ \hline
\endhead
1) & \possprim, \nongcr & 3 & 1 & 1 & 6 & 3 & 2 & 1 & 1 & 2 & 2 & 0 \\
2) & \possprim, \nongcr & 5 & 1 & 1 & 7 & 3 & 1 & 0 & 0 & 2 & 1 & 2 \\
3) & \nongcr & 9 & 1 & 1 & 9 & 1 & 1 & 3 & 3 & 2 & 0 & 0 \\
4) &  & 9 & 9 & 9 & 1 & 1 & 1 & 3 & 3 & 2 & 0 & 0 \\
\end{longtable}
\noindent Modules tested: V133, V912, L856, V56.\\

\setcounter{table}{172}
\begin{longtable}{rc|cccccc|cccccc}
\caption{$L_2(11) < E_7$, $p = 2$}\\
& & \multicolumn{6}{c|}{$V_{133}$} & \multicolumn{6}{c}{$V_{56}$} \\
& & $1$ & $5_{a}$ & $5_{b}$ & $10$ & $12_{a}$ & $12_{b}$ & $1$ & $5_{a}$ & $5_{b}$ & $10$ & $12_{a}$ & $12_{b}$ \\ \hline
\endfirsthead
& & \multicolumn{6}{c|}{$V_{133}$} & \multicolumn{6}{c}{$V_{56}$} \\
& & $1$ & $5_{a}$ & $5_{b}$ & $10$ & $12_{a}$ & $12_{b}$ & $1$ & $5_{a}$ & $5_{b}$ & $10$ & $12_{a}$ & $12_{b}$ \\ \hline
\endhead
1) & \nongcr & 3 & 1 & 1 & 6 & 3 & 2 & 2 & 1 & 1 & 2 & 2 & 0 \\
2) & \nongcr & 3 & 4 & 4 & 3 & 3 & 2 & 2 & 1 & 1 & 2 & 2 & 0 \\
3) & \nongcr & 5 & 2 & 2 & 6 & 3 & 1 & 0 & 0 & 0 & 2 & 1 & 2 \\
4) & \possprim, \nongcr & 5 & 5 & 5 & 3 & 3 & 1 & 0 & 0 & 0 & 2 & 1 & 2 \\
5) & \nongcr & 9 & 4 & 4 & 6 & 1 & 1 & 6 & 3 & 3 & 2 & 0 & 0 \\
6) & \nongcr & 9 & 7 & 7 & 3 & 1 & 1 & 6 & 3 & 3 & 2 & 0 & 0 \\
\end{longtable}
\noindent Modules tested: V133, L132, V912, V56.\\

\setcounter{table}{173}
\begin{longtable}{rc|ccccccccc|ccccccc}
\caption{$L_2(13) < E_7$, $p = 0$ or $p \nmid |H|$}\\
& & \multicolumn{9}{c|}{$V_{133}$} & \multicolumn{7}{c}{$V_{56}$} \\
& & $1$ & $7_{a}$ & $7_{b}$ & $12_{a}$ & $12_{b}$ & $12_{c}$ & $13$ & $14_{a}$ & $14_{b}$ & $1$ & $7_{a}$ & $7_{b}$ & $12_{a}$ & $13$ & $14_{a}$ & $14_{b}$ \\ \hline
\endfirsthead
& & \multicolumn{9}{c|}{$V_{133}$} & \multicolumn{7}{c}{$V_{56}$} \\
& & $1$ & $7_{a}$ & $7_{b}$ & $12_{a}$ & $12_{b}$ & $12_{c}$ & $13$ & $14_{a}$ & $14_{b}$ & $1$ & $7_{a}$ & $7_{b}$ & $12_{a}$ & $13$ & $14_{a}$ & $14_{b}$ \\ \hline
\endhead
1) &  & 1 & 0 & 0 & 1 & 1 & 1 & 2 & 3 & 2 & 0 & 2 & 2 & 0 & 0 & 0 & 2 \\
2) &  & 1 & 0 & 0 & 1 & 1 & 1 & 2 & 3 & 2 & 2 & 0 & 0 & 0 & 2 & 2 & 0 \\
3) &  & 2 & 2 & 2 & 1 & 1 & 1 & 3 & 2 & 0 & 0 & 2 & 2 & 0 & 0 & 0 & 2 \\
4) &  & 2 & 2 & 2 & 1 & 1 & 1 & 3 & 2 & 0 & 2 & 0 & 0 & 0 & 2 & 2 & 0 \\
5) &  & 3 & 0 & 0 & 4 & 1 & 0 & 0 & 3 & 2 & 4 & 0 & 0 & 2 & 0 & 2 & 0 \\
6) &  & 3 & 5 & 0 & 0 & 0 & 0 & 3 & 3 & 1 & 0 & 4 & 0 & 0 & 0 & 0 & 2 \\
7) &  & 4 & 2 & 2 & 4 & 1 & 0 & 1 & 2 & 0 & 4 & 0 & 0 & 2 & 0 & 2 & 0 \\
8) &  & 21 & 14 & 0 & 0 & 0 & 0 & 0 & 0 & 1 & 14 & 6 & 0 & 0 & 0 & 0 & 0 \\
\end{longtable}
\noindent Modules tested: V133, V912, V56.\\

\setcounter{table}{174}
\begin{longtable}{rc|ccccccccc|ccccccc}
\caption{$SL_2(13) < E_7$, $p = 0$ or $p \nmid |H|$}\\
& & \multicolumn{9}{c|}{$V_{133}$} & \multicolumn{7}{c}{$V_{56}$} \\
& & $1$ & $7_{a}$ & $7_{b}$ & $12_{a}$ & $12_{b}$ & $12_{c}$ & $13$ & $14_{a}$ & $14_{b}$ & $6_{a}$ & $12_{d}$ & $12_{e}$ & $12_{f}$ & $14_{c}$ & $14_{d}$ & $14_{e}$ \\ \hline
\endfirsthead
& & \multicolumn{9}{c|}{$V_{133}$} & \multicolumn{7}{c}{$V_{56}$} \\
& & $1$ & $7_{a}$ & $7_{b}$ & $12_{a}$ & $12_{b}$ & $12_{c}$ & $13$ & $14_{a}$ & $14_{b}$ & $6_{a}$ & $12_{d}$ & $12_{e}$ & $12_{f}$ & $14_{c}$ & $14_{d}$ & $14_{e}$ \\ \hline
\endhead
1) & \possprim & 0 & 2 & 0 & 1 & 1 & 1 & 1 & 1 & 4 & 1 & 1 & 1 & 1 & 1 & 0 & 0 \\
2) & \possprim & 0 & 1 & 1 & 1 & 1 & 1 & 1 & 1 & 4 & 0 & 0 & 0 & 0 & 2 & 1 & 1 \\
3) &  & 1 & 2 & 0 & 2 & 1 & 1 & 0 & 1 & 4 & 1 & 1 & 2 & 0 & 1 & 0 & 0 \\
4) &  & 1 & 4 & 2 & 1 & 1 & 1 & 2 & 0 & 2 & 1 & 1 & 1 & 1 & 1 & 0 & 0 \\
5) &  & 1 & 3 & 3 & 1 & 1 & 1 & 2 & 0 & 2 & 0 & 0 & 0 & 0 & 2 & 1 & 1 \\
6) &  & 2 & 4 & 2 & 2 & 1 & 1 & 1 & 0 & 2 & 1 & 1 & 2 & 0 & 1 & 0 & 0 \\
7) &  & 14 & 1 & 0 & 0 & 0 & 0 & 0 & 7 & 1 & 7 & 0 & 0 & 0 & 1 & 0 & 0 \\
\end{longtable}
\noindent Modules tested: V133, V912, V56.\\

\setcounter{table}{175}
\begin{longtable}{rc|cccccc|cccccc}
\caption{$L_2(13) < E_7$, $p = 7$}\\
& & \multicolumn{6}{c|}{$V_{133}$} & \multicolumn{6}{c}{$V_{56}$} \\
& & $1$ & $7_{a}$ & $7_{b}$ & $12$ & $14_{a}$ & $14_{b}$ & $1$ & $7_{a}$ & $7_{b}$ & $12$ & $14_{a}$ & $14_{b}$ \\ \hline
\endfirsthead
& & \multicolumn{6}{c|}{$V_{133}$} & \multicolumn{6}{c}{$V_{56}$} \\
& & $1$ & $7_{a}$ & $7_{b}$ & $12$ & $14_{a}$ & $14_{b}$ & $1$ & $7_{a}$ & $7_{b}$ & $12$ & $14_{a}$ & $14_{b}$ \\ \hline
\endhead
1) & \possprim, \nongcr & 3 & 0 & 0 & 5 & 3 & 2 & 0 & 2 & 2 & 0 & 0 & 2 \\
2) & \nongcr & 3 & 0 & 0 & 5 & 3 & 2 & 4 & 0 & 0 & 2 & 2 & 0 \\
3) & \possprim, \nongcr & 5 & 2 & 2 & 6 & 2 & 0 & 0 & 2 & 2 & 0 & 0 & 2 \\
4) & \nongcr & 5 & 2 & 2 & 6 & 2 & 0 & 4 & 0 & 0 & 2 & 2 & 0 \\
5) & \nongcr & 6 & 5 & 0 & 3 & 3 & 1 & 0 & 4 & 0 & 0 & 0 & 2 \\
6) &  & 21 & 14 & 0 & 0 & 0 & 1 & 14 & 6 & 0 & 0 & 0 & 0 \\
\end{longtable}
\noindent Modules tested: V133, V912, V56.\\

\setcounter{table}{176}
\begin{longtable}{rc|cccccc|ccccc}
\caption{$SL_2(13) < E_7$, $p = 7$}\\
& & \multicolumn{6}{c|}{$V_{133}$} & \multicolumn{5}{c}{$V_{56}$} \\
& & $1$ & $7_{a}$ & $7_{b}$ & $12$ & $14_{a}$ & $14_{b}$ & $6_{a}$ & $6_{b}$ & $14_{c}$ & $14_{d}$ & $14_{e}$ \\ \hline
\endfirsthead
& & \multicolumn{6}{c|}{$V_{133}$} & \multicolumn{5}{c}{$V_{56}$} \\
& & $1$ & $7_{a}$ & $7_{b}$ & $12$ & $14_{a}$ & $14_{b}$ & $6_{a}$ & $6_{b}$ & $14_{c}$ & $14_{d}$ & $14_{e}$ \\ \hline
\endhead
1) & \possprim, \nongcr & 1 & 2 & 0 & 4 & 1 & 4 & 4 & 3 & 1 & 0 & 0 \\
2) & \possprim, \nongcr & 1 & 1 & 1 & 4 & 1 & 4 & 0 & 0 & 2 & 1 & 1 \\
3) & \possprim, \nongcr & 3 & 4 & 2 & 5 & 0 & 2 & 4 & 3 & 1 & 0 & 0 \\
4) & \possprim, \nongcr & 3 & 3 & 3 & 5 & 0 & 2 & 0 & 0 & 2 & 1 & 1 \\
5) &  & 14 & 1 & 0 & 0 & 7 & 1 & 7 & 0 & 1 & 0 & 0 \\
\end{longtable}
\noindent Modules tested: V133, V912, V56.\\

\setcounter{table}{177}
\begin{longtable}{rc|ccccccc|ccccc}
\caption{$L_2(13) < E_7$, $p = 3$}\\
& & \multicolumn{7}{c|}{$V_{133}$} & \multicolumn{5}{c}{$V_{56}$} \\
& & $1$ & $7_{a}$ & $7_{b}$ & $12_{a}$ & $12_{b}$ & $12_{c}$ & $13$ & $1$ & $7_{a}$ & $7_{b}$ & $12_{a}$ & $13$ \\ \hline
\endfirsthead
& & \multicolumn{7}{c|}{$V_{133}$} & \multicolumn{5}{c}{$V_{56}$} \\
& & $1$ & $7_{a}$ & $7_{b}$ & $12_{a}$ & $12_{b}$ & $12_{c}$ & $13$ & $1$ & $7_{a}$ & $7_{b}$ & $12_{a}$ & $13$ \\ \hline
\endhead
1) & \possprim, \nongcr & 4 & 2 & 2 & 1 & 1 & 1 & 5 & 0 & 4 & 4 & 0 & 0 \\
2) & \nongcr & 4 & 2 & 2 & 1 & 1 & 1 & 5 & 4 & 0 & 0 & 0 & 4 \\
3) & \nongcr & 6 & 6 & 1 & 0 & 0 & 0 & 6 & 0 & 6 & 2 & 0 & 0 \\
4) & \nongcr & 6 & 2 & 2 & 4 & 1 & 0 & 3 & 6 & 0 & 0 & 2 & 2 \\
5) &  & 21 & 15 & 1 & 0 & 0 & 0 & 0 & 14 & 6 & 0 & 0 & 0 \\
\end{longtable}
\noindent Modules tested: V133, V912, L856, V56.\\

\setcounter{table}{178}
\begin{longtable}{rc|ccccccc|ccccc}
\caption{$SL_2(13) < E_7$, $p = 3$}\\
& & \multicolumn{7}{c|}{$V_{133}$} & \multicolumn{5}{c}{$V_{56}$} \\
& & $1$ & $7_{a}$ & $7_{b}$ & $12_{a}$ & $12_{b}$ & $12_{c}$ & $13$ & $6_{a}$ & $12_{d}$ & $12_{e}$ & $12_{f}$ & $14$ \\ \hline
\endfirsthead
& & \multicolumn{7}{c|}{$V_{133}$} & \multicolumn{5}{c}{$V_{56}$} \\
& & $1$ & $7_{a}$ & $7_{b}$ & $12_{a}$ & $12_{b}$ & $12_{c}$ & $13$ & $6_{a}$ & $12_{d}$ & $12_{e}$ & $12_{f}$ & $14$ \\ \hline
\endhead
1) & \possprim, \nongcr & 1 & 6 & 4 & 1 & 1 & 1 & 2 & 1 & 1 & 1 & 1 & 1 \\
2) & \possprim, \nongcr & 1 & 5 & 5 & 1 & 1 & 1 & 2 & 0 & 0 & 0 & 0 & 4 \\
3) &  & 2 & 6 & 4 & 2 & 1 & 1 & 1 & 1 & 1 & 2 & 0 & 1 \\
4) & \nongcr & 21 & 2 & 1 & 0 & 0 & 0 & 7 & 7 & 0 & 0 & 0 & 1 \\
\end{longtable}
\noindent Modules tested: V133, V912, L856, V56.\\

\setcounter{table}{179}
\begin{longtable}{rc|ccccccc|ccccccc}
\caption{$L_2(13) < E_7$, $p = 2$}\\
& & \multicolumn{7}{c|}{$V_{133}$} & \multicolumn{7}{c}{$V_{56}$} \\
& & $1$ & $6_{a}$ & $6_{b}$ & $12_{a}$ & $12_{b}$ & $12_{c}$ & $14$ & $1$ & $6_{a}$ & $6_{b}$ & $12_{a}$ & $12_{b}$ & $12_{c}$ & $14$ \\ \hline
\endfirsthead
& & \multicolumn{7}{c|}{$V_{133}$} & \multicolumn{7}{c}{$V_{56}$} \\
& & $1$ & $6_{a}$ & $6_{b}$ & $12_{a}$ & $12_{b}$ & $12_{c}$ & $14$ & $1$ & $6_{a}$ & $6_{b}$ & $12_{a}$ & $12_{b}$ & $12_{c}$ & $14$ \\ \hline
\endhead
1) &  & 3 & 0 & 0 & 4 & 1 & 0 & 5 & 4 & 0 & 0 & 2 & 0 & 0 & 2 \\
2) & \nongcr & 3 & 2 & 0 & 2 & 1 & 1 & 5 & 2 & 1 & 2 & 1 & 2 & 0 & 0 \\
3) & \possprim, \nongcr & 3 & 3 & 1 & 1 & 1 & 1 & 5 & 2 & 1 & 2 & 1 & 1 & 1 & 0 \\
4) & \nongcr & 3 & 2 & 2 & 1 & 1 & 1 & 5 & 4 & 2 & 2 & 0 & 0 & 0 & 2 \\
5) & \nongcr & 5 & 5 & 0 & 0 & 0 & 0 & 7 & 4 & 4 & 0 & 0 & 0 & 0 & 2 \\
6) & \nongcr & 9 & 3 & 3 & 4 & 1 & 0 & 2 & 4 & 0 & 0 & 2 & 0 & 0 & 2 \\
7) & \nongcr & 9 & 5 & 3 & 2 & 1 & 1 & 2 & 2 & 1 & 2 & 1 & 2 & 0 & 0 \\
8) & \possprim, \nongcr & 9 & 6 & 4 & 1 & 1 & 1 & 2 & 2 & 1 & 2 & 1 & 1 & 1 & 0 \\
9) & \nongcr & 9 & 5 & 5 & 1 & 1 & 1 & 2 & 4 & 2 & 2 & 0 & 0 & 0 & 2 \\
10) & \nongcr & 11 & 8 & 3 & 0 & 0 & 0 & 4 & 4 & 4 & 0 & 0 & 0 & 0 & 2 \\
11) & \nongcr & 15 & 1 & 0 & 0 & 0 & 0 & 8 & 2 & 1 & 8 & 0 & 0 & 0 & 0 \\
12) & \nongcr & 35 & 14 & 0 & 0 & 0 & 0 & 1 & 20 & 6 & 0 & 0 & 0 & 0 & 0 \\
\end{longtable}
\noindent Modules tested: V133, L132, V912, V56.\\

\setcounter{table}{180}
\begin{longtable}{rc|ccccccc|cccccc}
\caption{$L_2(17) < E_7$, $p = 0$ or $p \nmid |H|$}\\
& & \multicolumn{7}{c|}{$V_{133}$} & \multicolumn{6}{c}{$V_{56}$} \\
& & $1$ & $9_{a}$ & $16_{a}$ & $17$ & $18_{a}$ & $18_{b}$ & $18_{c}$ & $1$ & $9_{a}$ & $9_{b}$ & $16_{a}$ & $17$ & $18_{a}$ \\ \hline
\endfirsthead
& & \multicolumn{7}{c|}{$V_{133}$} & \multicolumn{6}{c}{$V_{56}$} \\
& & $1$ & $9_{a}$ & $16_{a}$ & $17$ & $18_{a}$ & $18_{b}$ & $18_{c}$ & $1$ & $9_{a}$ & $9_{b}$ & $16_{a}$ & $17$ & $18_{a}$ \\ \hline
\endhead
1) &  & 1 & 3 & 1 & 1 & 2 & 1 & 1 & 2 & 2 & 0 & 0 & 0 & 2 \\
2) &  & 3 & 3 & 1 & 3 & 0 & 1 & 1 & 2 & 4 & 2 & 0 & 0 & 0 \\
3) &  & 3 & 3 & 1 & 3 & 0 & 1 & 1 & 4 & 2 & 0 & 0 & 2 & 0 \\
4) &  & 6 & 3 & 4 & 0 & 0 & 1 & 1 & 6 & 2 & 0 & 2 & 0 & 0 \\
\end{longtable}
\noindent Modules tested: V133, V912, V56.\\

\setcounter{table}{181}
\begin{longtable}{rc|cccccc|ccccc}
\caption{$L_2(17) < E_7$, $p = 3$}\\
& & \multicolumn{6}{c|}{$V_{133}$} & \multicolumn{5}{c}{$V_{56}$} \\
& & $1$ & $9_{a}$ & $16$ & $18_{a}$ & $18_{b}$ & $18_{c}$ & $1$ & $9_{a}$ & $9_{b}$ & $16$ & $18_{a}$ \\ \hline
\endfirsthead
& & \multicolumn{6}{c|}{$V_{133}$} & \multicolumn{5}{c}{$V_{56}$} \\
& & $1$ & $9_{a}$ & $16$ & $18_{a}$ & $18_{b}$ & $18_{c}$ & $1$ & $9_{a}$ & $9_{b}$ & $16$ & $18_{a}$ \\ \hline
\endhead
1) & \nongcr & 2 & 3 & 2 & 2 & 1 & 1 & 2 & 2 & 0 & 0 & 2 \\
2) & \nongcr & 6 & 3 & 4 & 0 & 1 & 1 & 2 & 4 & 2 & 0 & 0 \\
3) & \nongcr & 6 & 3 & 4 & 0 & 1 & 1 & 6 & 2 & 0 & 2 & 0 \\
\end{longtable}
\noindent Modules tested: V133, V912, L856, V56.\\

\setcounter{table}{182}
\begin{longtable}{rc|cccc|cccc}
\caption{$L_2(17) < E_7$, $p = 2$}\\
& & \multicolumn{4}{c|}{$V_{133}$} & \multicolumn{4}{c}{$V_{56}$} \\
& & $1$ & $8_{a}$ & $8_{b}$ & $16_{a}$ & $1$ & $8_{a}$ & $8_{b}$ & $16_{a}$ \\ \hline
\endfirsthead
& & \multicolumn{4}{c|}{$V_{133}$} & \multicolumn{4}{c}{$V_{56}$} \\
& & $1$ & $8_{a}$ & $8_{b}$ & $16_{a}$ & $1$ & $8_{a}$ & $8_{b}$ & $16_{a}$ \\ \hline
\endhead
1) & \nongcr & 13 & 5 & 2 & 4 & 8 & 2 & 0 & 2 \\
2) & \nongcr & 13 & 8 & 5 & 1 & 8 & 4 & 2 & 0 \\
\end{longtable}
\noindent Modules tested: V133, L132, V912, V56.\\

\setcounter{table}{183}
\begin{longtable}{rc|cccccccc|cccc}
\caption{$L_2(19) < E_7$, $p = 0$ or $p \nmid |H|$}\\
& & \multicolumn{8}{c|}{$V_{133}$} & \multicolumn{4}{c}{$V_{56}$} \\
& & $1$ & $9_{a}$ & $9_{b}$ & $18_{a}$ & $18_{d}$ & $20_{b}$ & $20_{c}$ & $20_{d}$ & $1$ & $9_{a}$ & $9_{b}$ & $18_{d}$ \\ \hline
\endfirsthead
& & \multicolumn{8}{c|}{$V_{133}$} & \multicolumn{4}{c}{$V_{56}$} \\
& & $1$ & $9_{a}$ & $9_{b}$ & $18_{a}$ & $18_{d}$ & $20_{b}$ & $20_{c}$ & $20_{d}$ & $1$ & $9_{a}$ & $9_{b}$ & $18_{d}$ \\ \hline
\endhead
1) &  & 1 & 1 & 1 & 1 & 2 & 1 & 1 & 1 & 2 & 1 & 1 & 2 \\
\end{longtable}
\noindent Modules tested: V133, V912, V56.\\

\setcounter{table}{184}
\begin{longtable}{rc|cccccc|ccccc}
\caption{$SL_2(19) < E_7$, $p = 0$ or $p \nmid |H|$}\\
& & \multicolumn{6}{c|}{$V_{133}$} & \multicolumn{5}{c}{$V_{56}$} \\
& & $18_{a}$ & $18_{b}$ & $19$ & $20_{b}$ & $20_{c}$ & $20_{d}$ & $10_{a}$ & $10_{b}$ & $18_{f}$ & $18_{h}$ & $20_{e}$ \\ \hline
\endfirsthead
& & \multicolumn{6}{c|}{$V_{133}$} & \multicolumn{5}{c}{$V_{56}$} \\
& & $18_{a}$ & $18_{b}$ & $19$ & $20_{b}$ & $20_{c}$ & $20_{d}$ & $10_{a}$ & $10_{b}$ & $18_{f}$ & $18_{h}$ & $20_{e}$ \\ \hline
\endhead
1) & \possprim & 2 & 1 & 1 & 1 & 1 & 1 & 1 & 1 & 1 & 1 & 0 \\
2) & \possprim & 2 & 1 & 1 & 1 & 1 & 1 & 0 & 0 & 1 & 1 & 1 \\
\end{longtable}
\noindent Modules not used (elements of large order): V912.\\
Modules tested: V133, V56.\\

\setcounter{table}{185}
\begin{longtable}{rc|ccccccc|ccc}
\caption{$L_2(19) < E_7$, $p = 5$}\\
& & \multicolumn{7}{c|}{$V_{133}$} & \multicolumn{3}{c}{$V_{56}$} \\
& & $1$ & $9_{a}$ & $9_{b}$ & $18$ & $20_{b}$ & $20_{c}$ & $20_{d}$ & $1$ & $9_{a}$ & $9_{b}$ \\ \hline
\endfirsthead
& & \multicolumn{7}{c|}{$V_{133}$} & \multicolumn{3}{c}{$V_{56}$} \\
& & $1$ & $9_{a}$ & $9_{b}$ & $18$ & $20_{b}$ & $20_{c}$ & $20_{d}$ & $1$ & $9_{a}$ & $9_{b}$ \\ \hline
\endhead
1) &  & 1 & 3 & 3 & 1 & 1 & 1 & 1 & 2 & 3 & 3 \\
\end{longtable}
\noindent Modules tested: V133, V912, V56.\\

\setcounter{table}{186}
\begin{longtable}{rc|ccccc|cccc}
\caption{$SL_2(19) < E_7$, $p = 5$}\\
& & \multicolumn{5}{c|}{$V_{133}$} & \multicolumn{4}{c}{$V_{56}$} \\
& & $1$ & $18_{a}$ & $20_{b}$ & $20_{c}$ & $20_{d}$ & $10_{a}$ & $10_{b}$ & $18_{b}$ & $20_{e}$ \\ \hline
\endfirsthead
& & \multicolumn{5}{c|}{$V_{133}$} & \multicolumn{4}{c}{$V_{56}$} \\
& & $1$ & $18_{a}$ & $20_{b}$ & $20_{c}$ & $20_{d}$ & $10_{a}$ & $10_{b}$ & $18_{b}$ & $20_{e}$ \\ \hline
\endhead
1) & \possprim, \nongcr & 1 & 4 & 1 & 1 & 1 & 0 & 0 & 2 & 1 \\
2) & \possprim, \nongcr & 1 & 4 & 1 & 1 & 1 & 1 & 1 & 2 & 0 \\
\end{longtable}
\noindent Modules not used (elements of large order): V912.\\
Modules tested: V133, V56.\\

\setcounter{table}{187}
\begin{longtable}{rc|cccccc|cccc}
\caption{$L_2(19) < E_7$, $p = 3$}\\
& & \multicolumn{6}{c|}{$V_{133}$} & \multicolumn{4}{c}{$V_{56}$} \\
& & $1$ & $9_{a}$ & $9_{b}$ & $18_{a}$ & $18_{d}$ & $19$ & $1$ & $9_{a}$ & $9_{b}$ & $18_{d}$ \\ \hline
\endfirsthead
& & \multicolumn{6}{c|}{$V_{133}$} & \multicolumn{4}{c}{$V_{56}$} \\
& & $1$ & $9_{a}$ & $9_{b}$ & $18_{a}$ & $18_{d}$ & $19$ & $1$ & $9_{a}$ & $9_{b}$ & $18_{d}$ \\ \hline
\endhead
1) & \nongcr & 4 & 1 & 1 & 1 & 2 & 3 & 2 & 1 & 1 & 2 \\
\end{longtable}
\noindent Modules tested: V133, V912, L856, V56.\\

\setcounter{table}{188}
\begin{longtable}{rc|cccc|cccc}
\caption{$SL_2(19) < E_7$, $p = 3$}\\
& & \multicolumn{4}{c|}{$V_{133}$} & \multicolumn{4}{c}{$V_{56}$} \\
& & $1$ & $18_{a}$ & $18_{b}$ & $19$ & $10_{a}$ & $10_{b}$ & $18_{f}$ & $18_{h}$ \\ \hline
\endfirsthead
& & \multicolumn{4}{c|}{$V_{133}$} & \multicolumn{4}{c}{$V_{56}$} \\
& & $1$ & $18_{a}$ & $18_{b}$ & $19$ & $10_{a}$ & $10_{b}$ & $18_{f}$ & $18_{h}$ \\ \hline
\endhead
1) & \possprim, \nongcr & 3 & 2 & 1 & 4 & 1 & 1 & 1 & 1 \\
\end{longtable}
\noindent Modules not used (elements of large order): V912, L856.\\
Modules tested: V133, V56.\\

\setcounter{table}{189}
\begin{longtable}{rc|cccccccc|cccccc}
\caption{$L_2(19) < E_7$, $p = 2$}\\
& & \multicolumn{8}{c|}{$V_{133}$} & \multicolumn{6}{c}{$V_{56}$} \\
& & $1$ & $9_{a}$ & $9_{b}$ & $18_{a}$ & $18_{b}$ & $20_{b}$ & $20_{c}$ & $20_{d}$ & $1$ & $9_{a}$ & $9_{b}$ & $18_{a}$ & $18_{b}$ & $20_{a}$ \\ \hline
\endfirsthead
& & \multicolumn{8}{c|}{$V_{133}$} & \multicolumn{6}{c}{$V_{56}$} \\
& & $1$ & $9_{a}$ & $9_{b}$ & $18_{a}$ & $18_{b}$ & $20_{b}$ & $20_{c}$ & $20_{d}$ & $1$ & $9_{a}$ & $9_{b}$ & $18_{a}$ & $18_{b}$ & $20_{a}$ \\ \hline
\endhead
1) &  & 1 & 0 & 0 & 3 & 1 & 1 & 1 & 1 & 0 & 1 & 1 & 0 & 1 & 1 \\
2) &  & 1 & 0 & 0 & 3 & 1 & 1 & 1 & 1 & 2 & 2 & 2 & 0 & 1 & 0 \\
3) & \possprim, \nongcr & 1 & 1 & 1 & 2 & 1 & 1 & 1 & 1 & 0 & 0 & 0 & 2 & 0 & 1 \\
4) & \nongcr & 1 & 1 & 1 & 2 & 1 & 1 & 1 & 1 & 2 & 1 & 1 & 2 & 0 & 0 \\
\end{longtable}
\noindent Modules tested: V133, L132, V912, V56.\\

\setcounter{table}{190}
\begin{longtable}{rc|ccccccc|ccc}
\caption{$L_2(25) < E_7$, $p = 0$ or $p \nmid |H|$}\\
& & \multicolumn{7}{c|}{$V_{133}$} & \multicolumn{3}{c}{$V_{56}$} \\
& & $1$ & $25$ & $26_{a}$ & $26_{b}$ & $26_{c}$ & $26_{d}$ & $26_{e}$ & $1$ & $26_{a}$ & $26_{c}$ \\ \hline
\endfirsthead
& & \multicolumn{7}{c|}{$V_{133}$} & \multicolumn{3}{c}{$V_{56}$} \\
& & $1$ & $25$ & $26_{a}$ & $26_{b}$ & $26_{c}$ & $26_{d}$ & $26_{e}$ & $1$ & $26_{a}$ & $26_{c}$ \\ \hline
\endhead
1) &  & 3 & 0 & 3 & 0 & 0 & 1 & 1 & 4 & 2 & 0 \\
2) &  & 4 & 1 & 2 & 2 & 0 & 0 & 0 & 4 & 2 & 0 \\
3) &  & 4 & 1 & 0 & 2 & 2 & 0 & 0 & 4 & 0 & 2 \\
\end{longtable}
\noindent Modules tested: V133, V912, V56.\\

\setcounter{table}{191}
\begin{longtable}{rc|ccccccc|ccc}
\caption{$L_2(25) < E_7$, $p = 13$}\\
& & \multicolumn{7}{c|}{$V_{133}$} & \multicolumn{3}{c}{$V_{56}$} \\
& & $1$ & $24$ & $26_{a}$ & $26_{b}$ & $26_{c}$ & $26_{d}$ & $26_{e}$ & $1$ & $26_{a}$ & $26_{c}$ \\ \hline
\endfirsthead
& & \multicolumn{7}{c|}{$V_{133}$} & \multicolumn{3}{c}{$V_{56}$} \\
& & $1$ & $24$ & $26_{a}$ & $26_{b}$ & $26_{c}$ & $26_{d}$ & $26_{e}$ & $1$ & $26_{a}$ & $26_{c}$ \\ \hline
\endhead
1) &  & 3 & 0 & 3 & 0 & 0 & 1 & 1 & 4 & 2 & 0 \\
2) &  & 5 & 1 & 2 & 2 & 0 & 0 & 0 & 4 & 2 & 0 \\
3) &  & 5 & 1 & 0 & 2 & 2 & 0 & 0 & 4 & 0 & 2 \\
\end{longtable}
\noindent Modules tested: V133, V912, V56.\\

\setcounter{table}{192}
\begin{longtable}{rc|ccccc|cccc}
\caption{$L_2(25) < E_7$, $p = 3$}\\
& & \multicolumn{5}{c|}{$V_{133}$} & \multicolumn{4}{c}{$V_{56}$} \\
& & $1$ & $13_{a}$ & $13_{b}$ & $25$ & $26$ & $1$ & $13_{a}$ & $13_{b}$ & $25$ \\ \hline
\endfirsthead
& & \multicolumn{5}{c|}{$V_{133}$} & \multicolumn{4}{c}{$V_{56}$} \\
& & $1$ & $13_{a}$ & $13_{b}$ & $25$ & $26$ & $1$ & $13_{a}$ & $13_{b}$ & $25$ \\ \hline
\endhead
1) &  & 4 & 2 & 2 & 1 & 2 & 4 & 2 & 2 & 0 \\
2) & \nongcr & 6 & 0 & 0 & 3 & 2 & 6 & 0 & 0 & 2 \\
\end{longtable}
\noindent Modules tested: V133, V912, L856, V56.\\

\setcounter{table}{193}
\begin{longtable}{rc|cccc|cc}
\caption{$L_2(25) < E_7$, $p = 2$}\\
& & \multicolumn{4}{c|}{$V_{133}$} & \multicolumn{2}{c}{$V_{56}$} \\
& & $1$ & $12_{a}$ & $12_{b}$ & $26$ & $1$ & $26$ \\ \hline
\endfirsthead
& & \multicolumn{4}{c|}{$V_{133}$} & \multicolumn{2}{c}{$V_{56}$} \\
& & $1$ & $12_{a}$ & $12_{b}$ & $26$ & $1$ & $26$ \\ \hline
\endhead
1) &  & 3 & 0 & 0 & 5 & 4 & 2 \\
2) & \nongcr & 9 & 3 & 3 & 2 & 4 & 2 \\
\end{longtable}
\noindent Modules tested: V133, L132, V912, V56.\\

\setcounter{table}{194}
\begin{longtable}{rc|cccc|cc}
\caption{$L_2(27) < E_7$, $p = 0$ or $p \nmid |H|$}\\
& & \multicolumn{4}{c|}{$V_{133}$} & \multicolumn{2}{c}{$V_{56}$} \\
& & $1$ & $26_{a}$ & $26_{b}$ & $26_{e}$ & $1$ & $26_{e}$ \\ \hline
\endfirsthead
& & \multicolumn{4}{c|}{$V_{133}$} & \multicolumn{2}{c}{$V_{56}$} \\
& & $1$ & $26_{a}$ & $26_{b}$ & $26_{e}$ & $1$ & $26_{e}$ \\ \hline
\endhead
1) &  & 3 & 1 & 1 & 3 & 4 & 2 \\
\end{longtable}
\noindent Modules tested: V133, V912, V56.\\

\setcounter{table}{195}
\begin{longtable}{rc|cccccc|cccc}
\caption{$SL_2(27) < E_7$, $p = 0$ or $p \nmid |H|$}\\
& & \multicolumn{6}{c|}{$V_{133}$} & \multicolumn{4}{c}{$V_{56}$} \\
& & $1$ & $26_{a}$ & $26_{b}$ & $26_{c}$ & $27$ & $28_{a}$ & $14_{a}$ & $14_{b}$ & $28_{j}$ & $28_{k}$ \\ \hline
\endfirsthead
& & \multicolumn{6}{c|}{$V_{133}$} & \multicolumn{4}{c}{$V_{56}$} \\
& & $1$ & $26_{a}$ & $26_{b}$ & $26_{c}$ & $27$ & $28_{a}$ & $14_{a}$ & $14_{b}$ & $28_{j}$ & $28_{k}$ \\ \hline
\endhead
1) & \possprim & 0 & 1 & 1 & 1 & 1 & 1 & 0 & 0 & 1 & 1 \\
2) &  & 1 & 1 & 1 & 1 & 2 & 0 & 2 & 2 & 0 & 0 \\
\end{longtable}
\noindent Modules tested: V133, V912, V56.\\

\setcounter{table}{196}
\begin{longtable}{rc|cccc|cc}
\caption{$L_2(27) < E_7$, $p = 13$}\\
& & \multicolumn{4}{c|}{$V_{133}$} & \multicolumn{2}{c}{$V_{56}$} \\
& & $1$ & $26_{a}$ & $26_{b}$ & $26_{e}$ & $1$ & $26_{e}$ \\ \hline
\endfirsthead
& & \multicolumn{4}{c|}{$V_{133}$} & \multicolumn{2}{c}{$V_{56}$} \\
& & $1$ & $26_{a}$ & $26_{b}$ & $26_{e}$ & $1$ & $26_{e}$ \\ \hline
\endhead
1) &  & 3 & 1 & 1 & 3 & 4 & 2 \\
\end{longtable}
\noindent Modules tested: V133, V912, V56.\\

\setcounter{table}{197}
\begin{longtable}{rc|ccccc|cc}
\caption{$SL_2(27) < E_7$, $p = 13$}\\
& & \multicolumn{5}{c|}{$V_{133}$} & \multicolumn{2}{c}{$V_{56}$} \\
& & $1$ & $26_{a}$ & $26_{b}$ & $26_{c}$ & $27$ & $14_{a}$ & $14_{b}$ \\ \hline
\endfirsthead
& & \multicolumn{5}{c|}{$V_{133}$} & \multicolumn{2}{c}{$V_{56}$} \\
& & $1$ & $26_{a}$ & $26_{b}$ & $26_{c}$ & $27$ & $14_{a}$ & $14_{b}$ \\ \hline
\endhead
1) & \possprim, \nongcr & 1 & 1 & 1 & 1 & 2 & 2 & 2 \\
\end{longtable}
\noindent Modules tested: V133, V912, V56.\\

\setcounter{table}{198}
\begin{longtable}{rc|cccc|ccc}
\caption{$L_2(27) < E_7$, $p = 7$}\\
& & \multicolumn{4}{c|}{$V_{133}$} & \multicolumn{3}{c}{$V_{56}$} \\
& & $1$ & $13_{a}$ & $13_{b}$ & $26$ & $1$ & $13_{a}$ & $13_{b}$ \\ \hline
\endfirsthead
& & \multicolumn{4}{c|}{$V_{133}$} & \multicolumn{3}{c}{$V_{56}$} \\
& & $1$ & $13_{a}$ & $13_{b}$ & $26$ & $1$ & $13_{a}$ & $13_{b}$ \\ \hline
\endhead
1) & \nongcr & 3 & 3 & 3 & 2 & 4 & 2 & 2 \\
\end{longtable}
\noindent Modules tested: V133, V912, V56.\\

\setcounter{table}{199}
\begin{longtable}{rc|ccc|cccc}
\caption{$SL_2(27) < E_7$, $p = 7$}\\
& & \multicolumn{3}{c|}{$V_{133}$} & \multicolumn{4}{c}{$V_{56}$} \\
& & $1$ & $26_{a}$ & $28_{a}$ & $14_{a}$ & $14_{b}$ & $28_{j}$ & $28_{k}$ \\ \hline
\endfirsthead
& & \multicolumn{3}{c|}{$V_{133}$} & \multicolumn{4}{c}{$V_{56}$} \\
& & $1$ & $26_{a}$ & $28_{a}$ & $14_{a}$ & $14_{b}$ & $28_{j}$ & $28_{k}$ \\ \hline
\endhead
1) & \possprim, \nongcr & 1 & 4 & 1 & 0 & 0 & 1 & 1 \\
2) & \possprim, \nongcr & 3 & 5 & 0 & 2 & 2 & 0 & 0 \\
\end{longtable}
\noindent Modules tested: V133, V912, V56.\\

\setcounter{table}{200}
\begin{longtable}{rc|ccccccc|cccccc}
\caption{$L_2(27) < E_7$, $p = 2$}\\
& & \multicolumn{7}{c|}{$V_{133}$} & \multicolumn{6}{c}{$V_{56}$} \\
& & $1$ & $13_{a}$ & $13_{b}$ & $26_{a}$ & $26_{b}$ & $26_{c}$ & $28_{a}$ & $1$ & $13_{a}$ & $13_{b}$ & $26_{a}$ & $28_{d}$ & $28_{e}$ \\ \hline
\endfirsthead
& & \multicolumn{7}{c|}{$V_{133}$} & \multicolumn{6}{c}{$V_{56}$} \\
& & $1$ & $13_{a}$ & $13_{b}$ & $26_{a}$ & $26_{b}$ & $26_{c}$ & $28_{a}$ & $1$ & $13_{a}$ & $13_{b}$ & $26_{a}$ & $28_{d}$ & $28_{e}$ \\ \hline
\endhead
1) & \possprim, \nongcr & 1 & 1 & 1 & 1 & 1 & 1 & 1 & 0 & 0 & 0 & 0 & 1 & 1 \\
2) &  & 3 & 0 & 0 & 4 & 0 & 1 & 0 & 4 & 0 & 0 & 2 & 0 & 0 \\
3) & \nongcr & 3 & 2 & 2 & 1 & 1 & 1 & 0 & 4 & 2 & 2 & 0 & 0 & 0 \\
\end{longtable}
\noindent Modules tested: V133, L132, V912, V56.\\

\setcounter{table}{201}
\begin{longtable}{rc|ccccc|cc}
\caption{$SL_2(29) < E_7$, $p = 0$ or $p \nmid |H|$}\\
& & \multicolumn{5}{c|}{$V_{133}$} & \multicolumn{2}{c}{$V_{56}$} \\
& & $15_{a}$ & $28_{b}$ & $30_{d}$ & $30_{e}$ & $30_{f}$ & $28_{k}$ & $28_{m}$ \\ \hline
\endfirsthead
& & \multicolumn{5}{c|}{$V_{133}$} & \multicolumn{2}{c}{$V_{56}$} \\
& & $15_{a}$ & $28_{b}$ & $30_{d}$ & $30_{e}$ & $30_{f}$ & $28_{k}$ & $28_{m}$ \\ \hline
\endhead
1) & \possprim & 1 & 1 & 1 & 1 & 1 & 1 & 1 \\
\end{longtable}
\noindent Modules not used (elements of large order): V912.\\
Modules tested: V133, V56.\\

\setcounter{table}{202}
\begin{longtable}{rc|ccc|cc}
\caption{$SL_2(29) < E_7$, $p = 7$}\\
& & \multicolumn{3}{c|}{$V_{133}$} & \multicolumn{2}{c}{$V_{56}$} \\
& & $15_{a}$ & $15_{b}$ & $28_{b}$ & $28_{k}$ & $28_{m}$ \\ \hline
\endfirsthead
& & \multicolumn{3}{c|}{$V_{133}$} & \multicolumn{2}{c}{$V_{56}$} \\
& & $15_{a}$ & $15_{b}$ & $28_{b}$ & $28_{k}$ & $28_{m}$ \\ \hline
\endhead
1) & \possprim & 4 & 3 & 1 & 1 & 1 \\
\end{longtable}
\noindent Modules not used (elements of large order): V912.\\
Modules tested: V133, V56.\\

\setcounter{table}{203}
\begin{longtable}{rc|ccccc|c}
\caption{$SL_2(29) < E_7$, $p = 5$}\\
& & \multicolumn{5}{c|}{$V_{133}$} & \multicolumn{1}{c}{$V_{56}$} \\
& & $15_{a}$ & $28_{b}$ & $30_{d}$ & $30_{e}$ & $30_{f}$ & $28_{c}$ \\ \hline
\endfirsthead
& & \multicolumn{5}{c|}{$V_{133}$} & \multicolumn{1}{c}{$V_{56}$} \\
& & $15_{a}$ & $28_{b}$ & $30_{d}$ & $30_{e}$ & $30_{f}$ & $28_{c}$ \\ \hline
\endhead
1) & \possprim & 1 & 1 & 1 & 1 & 1 & 2 \\
\end{longtable}
\noindent Modules not used (elements of large order): V912.\\
Modules tested: V133, V56.\\

\setcounter{table}{204}
\begin{longtable}{rc|ccccc|c}
\caption{$SL_2(29) < E_7$, $p = 3$}\\
& & \multicolumn{5}{c|}{$V_{133}$} & \multicolumn{1}{c}{$V_{56}$} \\
& & $15_{a}$ & $28_{b}$ & $30_{d}$ & $30_{e}$ & $30_{f}$ & $28_{d}$ \\ \hline
\endfirsthead
& & \multicolumn{5}{c|}{$V_{133}$} & \multicolumn{1}{c}{$V_{56}$} \\
& & $15_{a}$ & $28_{b}$ & $30_{d}$ & $30_{e}$ & $30_{f}$ & $28_{d}$ \\ \hline
\endhead
1) & \possprim & 1 & 1 & 1 & 1 & 1 & 2 \\
\end{longtable}
\noindent Modules not used (elements of large order): V912, L856.\\
Modules tested: V133, V56.\\

\setcounter{table}{205}
\begin{longtable}{rc|cccccc|cc}
\caption{$L_2(29) < E_7$, $p = 2$}\\
& & \multicolumn{6}{c|}{$V_{133}$} & \multicolumn{2}{c}{$V_{56}$} \\
& & $1$ & $14_{a}$ & $28_{b}$ & $30_{a}$ & $30_{b}$ & $30_{c}$ & $28_{d}$ & $28_{f}$ \\ \hline
\endfirsthead
& & \multicolumn{6}{c|}{$V_{133}$} & \multicolumn{2}{c}{$V_{56}$} \\
& & $1$ & $14_{a}$ & $28_{b}$ & $30_{a}$ & $30_{b}$ & $30_{c}$ & $28_{d}$ & $28_{f}$ \\ \hline
\endhead
1) & \possprim & 1 & 1 & 1 & 1 & 1 & 1 & 1 & 1 \\
\end{longtable}
\noindent Modules tested: V133, L132, V912, V56.\\

\setcounter{table}{206}
\begin{longtable}{l}
\caption{$SL_2(37) < E_7$, $p = 0$ or $p \nmid |H|$}\\
Not recomputed: no character table available in GAP. See the Memoir.
\end{longtable}

\setcounter{table}{207}
\begin{longtable}{l}
\caption{$SL_2(37) < E_7$, $p = 19$}\\
Not recomputed: no character table available in GAP. See the Memoir.
\end{longtable}

\setcounter{table}{208}
\begin{longtable}{l}
\caption{$SL_2(37) < E_7$, $p = 3$}\\
Not recomputed: no character table available in GAP. See the Memoir.
\end{longtable}

\setcounter{table}{209}
\begin{longtable}{l}
\caption{$L_2(37) < E_7$, $p = 2$}\\
Not recomputed: no Brauer table available in GAP. See the Memoir.
\end{longtable}

\setcounter{table}{210}
\begin{longtable}{rc|ccccccccccc|cccccc}
\caption{$L_3(3) < E_7$, $p = 0$ or $p \nmid |H|$}\\
& & \multicolumn{11}{c|}{$V_{133}$} & \multicolumn{6}{c}{$V_{56}$} \\
& & $1$ & $13$ & $16_{a}$ & $16_{b}$ & $16_{c}$ & $16_{d}$ & $26_{a}$ & $26_{b}$ & $26_{c}$ & $27$ & $39$ & $1$ & $12$ & $16_{c}$ & $16_{d}$ & $26_{a}$ & $27$ \\ \hline
\endfirsthead
& & \multicolumn{11}{c|}{$V_{133}$} & \multicolumn{6}{c}{$V_{56}$} \\
& & $1$ & $13$ & $16_{a}$ & $16_{b}$ & $16_{c}$ & $16_{d}$ & $26_{a}$ & $26_{b}$ & $26_{c}$ & $27$ & $39$ & $1$ & $12$ & $16_{c}$ & $16_{d}$ & $26_{a}$ & $27$ \\ \hline
\endhead
1) &  & 1 & 0 & 0 & 0 & 0 & 0 & 1 & 1 & 1 & 2 & 0 & 2 & 0 & 0 & 0 & 0 & 2 \\
2) &  & 2 & 1 & 2 & 2 & 0 & 0 & 0 & 0 & 0 & 2 & 0 & 0 & 2 & 1 & 1 & 0 & 0 \\
3) &  & 2 & 1 & 1 & 1 & 1 & 1 & 0 & 0 & 0 & 2 & 0 & 2 & 0 & 0 & 0 & 0 & 2 \\
4) &  & 3 & 0 & 0 & 0 & 0 & 0 & 3 & 1 & 1 & 0 & 0 & 4 & 0 & 0 & 0 & 2 & 0 \\
5) &  & 3 & 0 & 2 & 2 & 0 & 0 & 0 & 0 & 0 & 1 & 1 & 0 & 2 & 1 & 1 & 0 & 0 \\
6) &  & 3 & 0 & 1 & 1 & 1 & 1 & 0 & 0 & 0 & 1 & 1 & 2 & 0 & 0 & 0 & 0 & 2 \\
7) &  & 4 & 1 & 1 & 1 & 1 & 1 & 2 & 0 & 0 & 0 & 0 & 4 & 0 & 0 & 0 & 2 & 0 \\
\end{longtable}
\noindent Modules tested: V133, V912, V56.\\

\setcounter{table}{211}
\begin{longtable}{rc|cccccccc|cccc}
\caption{$L_3(3) < E_7$, $p = 13$}\\
& & \multicolumn{8}{c|}{$V_{133}$} & \multicolumn{4}{c}{$V_{56}$} \\
& & $1$ & $11$ & $13$ & $16$ & $26_{a}$ & $26_{b}$ & $26_{c}$ & $39$ & $1$ & $11$ & $16$ & $26_{a}$ \\ \hline
\endfirsthead
& & \multicolumn{8}{c|}{$V_{133}$} & \multicolumn{4}{c}{$V_{56}$} \\
& & $1$ & $11$ & $13$ & $16$ & $26_{a}$ & $26_{b}$ & $26_{c}$ & $39$ & $1$ & $11$ & $16$ & $26_{a}$ \\ \hline
\endhead
1) & \nongcr & 1 & 2 & 0 & 2 & 1 & 1 & 1 & 0 & 2 & 2 & 2 & 0 \\
2) & \nongcr & 2 & 2 & 1 & 6 & 0 & 0 & 0 & 0 & 2 & 2 & 2 & 0 \\
3) &  & 3 & 0 & 0 & 0 & 3 & 1 & 1 & 0 & 4 & 0 & 0 & 2 \\
4) &  & 3 & 1 & 0 & 5 & 0 & 0 & 0 & 1 & 2 & 2 & 2 & 0 \\
5) &  & 4 & 0 & 1 & 4 & 2 & 0 & 0 & 0 & 4 & 0 & 0 & 2 \\
\end{longtable}
\noindent Modules tested: V133, V912, V56.\\

\setcounter{table}{212}
\begin{longtable}{rc|ccccccc|ccccc}
\caption{$L_3(3) < E_7$, $p = 2$}\\
& & \multicolumn{7}{c|}{$V_{133}$} & \multicolumn{5}{c}{$V_{56}$} \\
& & $1$ & $12$ & $16_{a}$ & $16_{b}$ & $16_{c}$ & $16_{d}$ & $26$ & $1$ & $12$ & $16_{c}$ & $16_{d}$ & $26$ \\ \hline
\endfirsthead
& & \multicolumn{7}{c|}{$V_{133}$} & \multicolumn{5}{c}{$V_{56}$} \\
& & $1$ & $12$ & $16_{a}$ & $16_{b}$ & $16_{c}$ & $16_{d}$ & $26$ & $1$ & $12$ & $16_{c}$ & $16_{d}$ & $26$ \\ \hline
\endhead
1) & \nongcr & 3 & 0 & 0 & 0 & 0 & 0 & 5 & 4 & 0 & 0 & 0 & 2 \\
2) & \nongcr & 5 & 1 & 2 & 2 & 0 & 0 & 2 & 0 & 2 & 1 & 1 & 0 \\
3) & \nongcr & 5 & 1 & 1 & 1 & 1 & 1 & 2 & 4 & 0 & 0 & 0 & 2 \\
\end{longtable}
\noindent Modules tested: V133, L132, V912, V56.\\

\setcounter{table}{213}
\begin{longtable}{rc|ccc|c}
\caption{$2 \cdot L_3(4) < E_7$, $p = 0$ or $p \nmid |H|$}\\
& & \multicolumn{3}{c|}{$V_{133}$} & \multicolumn{1}{c}{$V_{56}$} \\
& & $35_{b}$ & $35_{c}$ & $63_{a}$ & $28_{a}$ \\ \hline
\endfirsthead
& & \multicolumn{3}{c|}{$V_{133}$} & \multicolumn{1}{c}{$V_{56}$} \\
& & $35_{b}$ & $35_{c}$ & $63_{a}$ & $28_{a}$ \\ \hline
\endhead
1) & \possprim & 1 & 1 & 1 & 2 \\
\end{longtable}
\noindent Modules tested: V133, V912, V56.\\

\setcounter{table}{214}
\begin{longtable}{rc|ccc|c}
\caption{$2 \cdot L_3(4) < E_7$, $p = 7$}\\
& & \multicolumn{3}{c|}{$V_{133}$} & \multicolumn{1}{c}{$V_{56}$} \\
& & $35_{b}$ & $35_{c}$ & $63_{a}$ & $28_{a}$ \\ \hline
\endfirsthead
& & \multicolumn{3}{c|}{$V_{133}$} & \multicolumn{1}{c}{$V_{56}$} \\
& & $35_{b}$ & $35_{c}$ & $63_{a}$ & $28_{a}$ \\ \hline
\endhead
1) & \possprim & 1 & 1 & 1 & 2 \\
\end{longtable}
\noindent Modules tested: V133, V912, V56.\\

\setcounter{table}{215}
\begin{longtable}{rc|ccc|c}
\caption{$2 \cdot L_3(4) < E_7$, $p = 5$}\\
& & \multicolumn{3}{c|}{$V_{133}$} & \multicolumn{1}{c}{$V_{56}$} \\
& & $35_{b}$ & $35_{c}$ & $63$ & $28$ \\ \hline
\endfirsthead
& & \multicolumn{3}{c|}{$V_{133}$} & \multicolumn{1}{c}{$V_{56}$} \\
& & $35_{b}$ & $35_{c}$ & $63$ & $28$ \\ \hline
\endhead
1) & \possprim & 1 & 1 & 1 & 2 \\
\end{longtable}
\noindent Modules tested: V133, V912, V56.\\

\setcounter{table}{216}
\begin{longtable}{rc|cccccc|cccc}
\caption{$2 \cdot L_3(4) < E_7$, $p = 3$}\\
& & \multicolumn{6}{c|}{$V_{133}$} & \multicolumn{4}{c}{$V_{56}$} \\
& & $1$ & $15_{a}$ & $15_{b}$ & $15_{c}$ & $19$ & $63_{a}$ & $6$ & $10_{a}$ & $10_{b}$ & $22_{b}$ \\ \hline
\endfirsthead
& & \multicolumn{6}{c|}{$V_{133}$} & \multicolumn{4}{c}{$V_{56}$} \\
& & $1$ & $15_{a}$ & $15_{b}$ & $15_{c}$ & $19$ & $63_{a}$ & $6$ & $10_{a}$ & $10_{b}$ & $22_{b}$ \\ \hline
\endhead
1) & \possprim, \nongcr & 2 & 0 & 1 & 1 & 2 & 1 & 2 & 0 & 0 & 2 \\
2) &  & 9 & 7 & 0 & 0 & 1 & 0 & 6 & 1 & 1 & 0 \\
\end{longtable}
\noindent Modules tested: V133, V912, L856, V56.\\

\setcounter{table}{217}
\begin{longtable}{rc|ccc|cc}
\caption{$L_4(3) < E_7$, $p = 2$}\\
& & \multicolumn{3}{c|}{$V_{133}$} & \multicolumn{2}{c}{$V_{56}$} \\
& & $1$ & $26_{a}$ & $26_{b}$ & $1$ & $26_{a}$ \\ \hline
\endfirsthead
& & \multicolumn{3}{c|}{$V_{133}$} & \multicolumn{2}{c}{$V_{56}$} \\
& & $1$ & $26_{a}$ & $26_{b}$ & $1$ & $26_{a}$ \\ \hline
\endhead
1) &  & 3 & 4 & 1 & 4 & 2 \\
\end{longtable}
\noindent Modules tested: V133, L132, V912, V56.\\

\setcounter{table}{218}
\begin{longtable}{rc|cccccccccccccc|cccccccccccc}
\caption{$U_3(3) < E_7$, $p = 0$ or $p \nmid |H|$}\\
& & \multicolumn{14}{c|}{$V_{133}$} & \multicolumn{12}{c}{$V_{56}$} \\
& & $1$ & $6$ & $7_{a}$ & $7_{b}$ & $7_{c}$ & $14$ & $21_{a}$ & $21_{b}$ & $21_{c}$ & $27$ & $28_{a}$ & $28_{b}$ & $32_{a}$ & $32_{b}$ & $1$ & $6$ & $7_{a}$ & $7_{b}$ & $7_{c}$ & $14$ & $21_{a}$ & $21_{b}$ & $21_{c}$ & $27$ & $28_{a}$ & $28_{b}$ \\ \hline
\endfirsthead
& & \multicolumn{14}{c|}{$V_{133}$} & \multicolumn{12}{c}{$V_{56}$} \\
& & $1$ & $6$ & $7_{a}$ & $7_{b}$ & $7_{c}$ & $14$ & $21_{a}$ & $21_{b}$ & $21_{c}$ & $27$ & $28_{a}$ & $28_{b}$ & $32_{a}$ & $32_{b}$ & $1$ & $6$ & $7_{a}$ & $7_{b}$ & $7_{c}$ & $14$ & $21_{a}$ & $21_{b}$ & $21_{c}$ & $27$ & $28_{a}$ & $28_{b}$ \\ \hline
\endhead
1) & \possprim & 0 & 0 & 0 & 0 & 0 & 0 & 0 & 1 & 1 & 1 & 0 & 0 & 1 & 1 & 0 & 0 & 0 & 0 & 0 & 0 & 0 & 0 & 0 & 0 & 1 & 1 \\
2) &  & 21 & 0 & 14 & 0 & 0 & 1 & 0 & 0 & 0 & 0 & 0 & 0 & 0 & 0 & 14 & 0 & 6 & 0 & 0 & 0 & 0 & 0 & 0 & 0 & 0 & 0 \\
3) & \newrow & 14 & 0 & 0 & 0 & 0 & 7 & 1 & 0 & 0 & 0 & 0 & 0 & 0 & 0 & 0 & 7 & 0 & 1 & 1 & 0 & 0 & 0 & 0 & 0 & 0 & 0 \\
4) & \newrow & 6 & 6 & 0 & 2 & 2 & 3 & 1 & 0 & 0 & 0 & 0 & 0 & 0 & 0 & 4 & 4 & 0 & 0 & 0 & 2 & 0 & 0 & 0 & 0 & 0 & 0 \\
5) & \newrow & 3 & 0 & 5 & 0 & 0 & 1 & 0 & 0 & 0 & 3 & 0 & 0 & 0 & 0 & 0 & 0 & 4 & 0 & 0 & 2 & 0 & 0 & 0 & 0 & 0 & 0 \\
6) & \newrow & 2 & 0 & 1 & 1 & 1 & 0 & 0 & 0 & 0 & 2 & 1 & 1 & 0 & 0 & 2 & 0 & 0 & 0 & 0 & 0 & 0 & 0 & 0 & 2 & 0 & 0 \\
7) & \newrow & 2 & 0 & 1 & 1 & 1 & 0 & 0 & 0 & 0 & 2 & 1 & 1 & 0 & 0 & 0 & 0 & 2 & 0 & 0 & 0 & 2 & 0 & 0 & 0 & 0 & 0 \\
8) & \newrow & 1 & 0 & 2 & 1 & 1 & 0 & 1 & 0 & 0 & 1 & 1 & 1 & 0 & 0 & 0 & 0 & 0 & 1 & 1 & 0 & 0 & 1 & 1 & 0 & 0 & 0 \\
9) & \newrow & 1 & 0 & 0 & 0 & 0 & 1 & 0 & 0 & 0 & 2 & 0 & 0 & 1 & 1 & 2 & 0 & 0 & 0 & 0 & 0 & 0 & 0 & 0 & 2 & 0 & 0 \\
10) & \newrow & 1 & 0 & 0 & 0 & 0 & 1 & 0 & 0 & 0 & 2 & 0 & 0 & 1 & 1 & 0 & 0 & 2 & 0 & 0 & 0 & 2 & 0 & 0 & 0 & 0 & 0 \\
11) & \newrow & 0 & 0 & 1 & 0 & 0 & 1 & 1 & 0 & 0 & 1 & 0 & 0 & 1 & 1 & 0 & 0 & 0 & 1 & 1 & 0 & 0 & 1 & 1 & 0 & 0 & 0 \\
\end{longtable}
\noindent Memoir rows not reproduced: 2) (not feasible on the Memoir's modules (all element orders, full power maps)); 3) (not feasible on the Memoir's modules (all element orders, full power maps)); 4) (not feasible on the Memoir's modules (all element orders, full power maps)); 5) (not feasible on the Memoir's modules (all element orders, full power maps)); 6) (not feasible on the Memoir's modules (all element orders, full power maps)); 7) (not feasible on the Memoir's modules (all element orders, full power maps)); 8) (not feasible on the Memoir's modules (all element orders, full power maps)); 9) (not feasible on the Memoir's modules (all element orders, full power maps)); 10) (not feasible on the Memoir's modules (all element orders, full power maps)).\\
Rows marked \newrow\ do not appear in the Memoir.\\
Modules tested: V133, V912, V56.\\

\setcounter{table}{219}
\begin{longtable}{rc|cccccccccccc|cccccccccccc}
\caption{$U_3(3) < E_7$, $p = 7$}\\
& & \multicolumn{12}{c|}{$V_{133}$} & \multicolumn{12}{c}{$V_{56}$} \\
& & $1$ & $6$ & $7_{a}$ & $7_{b}$ & $7_{c}$ & $14$ & $21_{a}$ & $21_{b}$ & $21_{c}$ & $26$ & $28_{a}$ & $28_{b}$ & $1$ & $6$ & $7_{a}$ & $7_{b}$ & $7_{c}$ & $14$ & $21_{a}$ & $21_{b}$ & $21_{c}$ & $26$ & $28_{a}$ & $28_{b}$ \\ \hline
\endfirsthead
& & \multicolumn{12}{c|}{$V_{133}$} & \multicolumn{12}{c}{$V_{56}$} \\
& & $1$ & $6$ & $7_{a}$ & $7_{b}$ & $7_{c}$ & $14$ & $21_{a}$ & $21_{b}$ & $21_{c}$ & $26$ & $28_{a}$ & $28_{b}$ & $1$ & $6$ & $7_{a}$ & $7_{b}$ & $7_{c}$ & $14$ & $21_{a}$ & $21_{b}$ & $21_{c}$ & $26$ & $28_{a}$ & $28_{b}$ \\ \hline
\endhead
1) & \possprim, \nongcr & 1 & 2 & 0 & 0 & 0 & 0 & 0 & 1 & 1 & 3 & 0 & 0 & 0 & 0 & 0 & 0 & 0 & 0 & 0 & 0 & 0 & 0 & 1 & 1 \\
2) & \possprim, \nongcr & 1 & 2 & 1 & 0 & 0 & 1 & 1 & 0 & 0 & 3 & 0 & 0 & 0 & 0 & 0 & 1 & 1 & 0 & 0 & 1 & 1 & 0 & 0 & 0 \\
3) &  & 2 & 0 & 2 & 1 & 1 & 0 & 1 & 0 & 0 & 1 & 1 & 1 & 0 & 0 & 0 & 1 & 1 & 0 & 0 & 1 & 1 & 0 & 0 & 0 \\
4) & \possprim, \nongcr & 3 & 2 & 0 & 0 & 0 & 1 & 0 & 0 & 0 & 4 & 0 & 0 & 0 & 0 & 2 & 0 & 0 & 0 & 2 & 0 & 0 & 0 & 0 & 0 \\
5) & \nongcr & 3 & 2 & 0 & 0 & 0 & 1 & 0 & 0 & 0 & 4 & 0 & 0 & 4 & 0 & 0 & 0 & 0 & 0 & 0 & 0 & 0 & 2 & 0 & 0 \\
6) & \nongcr & 4 & 0 & 1 & 1 & 1 & 0 & 0 & 0 & 0 & 2 & 1 & 1 & 0 & 0 & 2 & 0 & 0 & 0 & 2 & 0 & 0 & 0 & 0 & 0 \\
7) & \nongcr & 4 & 0 & 1 & 1 & 1 & 0 & 0 & 0 & 0 & 2 & 1 & 1 & 4 & 0 & 0 & 0 & 0 & 0 & 0 & 0 & 0 & 2 & 0 & 0 \\
8) & \nongcr & 6 & 0 & 5 & 0 & 0 & 1 & 0 & 0 & 0 & 3 & 0 & 0 & 0 & 0 & 4 & 0 & 0 & 2 & 0 & 0 & 0 & 0 & 0 & 0 \\
9) &  & 6 & 6 & 0 & 2 & 2 & 3 & 1 & 0 & 0 & 0 & 0 & 0 & 4 & 4 & 0 & 0 & 0 & 2 & 0 & 0 & 0 & 0 & 0 & 0 \\
10) &  & 14 & 0 & 0 & 0 & 0 & 7 & 1 & 0 & 0 & 0 & 0 & 0 & 0 & 7 & 0 & 1 & 1 & 0 & 0 & 0 & 0 & 0 & 0 & 0 \\
11) &  & 21 & 0 & 14 & 0 & 0 & 1 & 0 & 0 & 0 & 0 & 0 & 0 & 14 & 0 & 6 & 0 & 0 & 0 & 0 & 0 & 0 & 0 & 0 & 0 \\
\end{longtable}
\noindent Modules tested: V133, V912, V56.\\

\setcounter{table}{220}
\begin{longtable}{rc|c|c}
\caption{$U_3(8) < E_7$, $p \neq 2$}\\
& & \multicolumn{1}{c|}{$V_{133}$} & \multicolumn{1}{c}{$V_{56}$} \\
& & $133_{a}$ & $56$ \\ \hline
\endfirsthead
& & \multicolumn{1}{c|}{$V_{133}$} & \multicolumn{1}{c}{$V_{56}$} \\
& & $133_{a}$ & $56$ \\ \hline
\endhead
1) & \possprim & 1 & 1 \\
\end{longtable}
\noindent Modules tested: V133, V912, V56.\\

\setcounter{table}{221}
\begin{longtable}{rc|ccccccccc|ccccccc}
\caption{$U_4(2) < E_7$, $p = 0$ or $p \nmid |H|$}\\
& & \multicolumn{9}{c|}{$V_{133}$} & \multicolumn{7}{c}{$V_{56}$} \\
& & $1$ & $5_{a}$ & $5_{b}$ & $6$ & $10_{a}$ & $10_{b}$ & $15_{a}$ & $20$ & $24$ & $1$ & $5_{a}$ & $5_{b}$ & $6$ & $10_{a}$ & $10_{b}$ & $15_{a}$ \\ \hline
\endfirsthead
& & \multicolumn{9}{c|}{$V_{133}$} & \multicolumn{7}{c}{$V_{56}$} \\
& & $1$ & $5_{a}$ & $5_{b}$ & $6$ & $10_{a}$ & $10_{b}$ & $15_{a}$ & $20$ & $24$ & $1$ & $5_{a}$ & $5_{b}$ & $6$ & $10_{a}$ & $10_{b}$ & $15_{a}$ \\ \hline
\endhead
1) &  & 4 & 0 & 0 & 4 & 2 & 2 & 3 & 1 & 0 & 2 & 0 & 0 & 4 & 0 & 0 & 2 \\
2) &  & 8 & 0 & 0 & 0 & 0 & 0 & 7 & 1 & 0 & 0 & 0 & 0 & 6 & 1 & 1 & 0 \\
3) &  & 9 & 4 & 4 & 0 & 3 & 3 & 0 & 0 & 1 & 6 & 3 & 3 & 0 & 1 & 1 & 0 \\
\end{longtable}
\noindent Modules tested: V133, V912, V56.\\

\setcounter{table}{222}
\begin{longtable}{rc|ccccccccc|ccccccc}
\caption{$U_4(2) < E_7$, $p = 5$}\\
& & \multicolumn{9}{c|}{$V_{133}$} & \multicolumn{7}{c}{$V_{56}$} \\
& & $1$ & $5_{a}$ & $5_{b}$ & $6$ & $10_{a}$ & $10_{b}$ & $15_{a}$ & $20$ & $23$ & $1$ & $5_{a}$ & $5_{b}$ & $6$ & $10_{a}$ & $10_{b}$ & $15_{a}$ \\ \hline
\endfirsthead
& & \multicolumn{9}{c|}{$V_{133}$} & \multicolumn{7}{c}{$V_{56}$} \\
& & $1$ & $5_{a}$ & $5_{b}$ & $6$ & $10_{a}$ & $10_{b}$ & $15_{a}$ & $20$ & $23$ & $1$ & $5_{a}$ & $5_{b}$ & $6$ & $10_{a}$ & $10_{b}$ & $15_{a}$ \\ \hline
\endhead
1) &  & 4 & 0 & 0 & 4 & 2 & 2 & 3 & 1 & 0 & 2 & 0 & 0 & 4 & 0 & 0 & 2 \\
2) &  & 8 & 0 & 0 & 0 & 0 & 0 & 7 & 1 & 0 & 0 & 0 & 0 & 6 & 1 & 1 & 0 \\
3) &  & 10 & 4 & 4 & 0 & 3 & 3 & 0 & 0 & 1 & 6 & 3 & 3 & 0 & 1 & 1 & 0 \\
\end{longtable}
\noindent Modules tested: V133, V912, V56.\\

\setcounter{table}{223}
\begin{longtable}{rc|ccccc|cc}
\caption{$Sp_6(2) < E_7$, $p = 0$ or $p \nmid |H|$}\\
& & \multicolumn{5}{c|}{$V_{133}$} & \multicolumn{2}{c}{$V_{56}$} \\
& & $1$ & $7$ & $21_{a}$ & $27$ & $35_{a}$ & $7$ & $21_{a}$ \\ \hline
\endfirsthead
& & \multicolumn{5}{c|}{$V_{133}$} & \multicolumn{2}{c}{$V_{56}$} \\
& & $1$ & $7$ & $21_{a}$ & $27$ & $35_{a}$ & $7$ & $21_{a}$ \\ \hline
\endhead
1) &  & 1 & 2 & 1 & 1 & 2 & 2 & 2 \\
\end{longtable}
\noindent Modules tested: V133, V912, V56.\\

\setcounter{table}{224}
\begin{longtable}{rc|ccccc|cc}
\caption{$Sp_6(2) < E_7$, $p = 7$}\\
& & \multicolumn{5}{c|}{$V_{133}$} & \multicolumn{2}{c}{$V_{56}$} \\
& & $1$ & $7$ & $21_{a}$ & $26$ & $35_{a}$ & $7$ & $21_{a}$ \\ \hline
\endfirsthead
& & \multicolumn{5}{c|}{$V_{133}$} & \multicolumn{2}{c}{$V_{56}$} \\
& & $1$ & $7$ & $21_{a}$ & $26$ & $35_{a}$ & $7$ & $21_{a}$ \\ \hline
\endhead
1) &  & 2 & 2 & 1 & 1 & 2 & 2 & 2 \\
\end{longtable}
\noindent Modules tested: V133, V912, V56.\\

\setcounter{table}{225}
\begin{longtable}{rc|ccccc|cc}
\caption{$Sp_6(2) < E_7$, $p = 5$}\\
& & \multicolumn{5}{c|}{$V_{133}$} & \multicolumn{2}{c}{$V_{56}$} \\
& & $1$ & $7$ & $21_{a}$ & $27$ & $35_{a}$ & $7$ & $21_{a}$ \\ \hline
\endfirsthead
& & \multicolumn{5}{c|}{$V_{133}$} & \multicolumn{2}{c}{$V_{56}$} \\
& & $1$ & $7$ & $21_{a}$ & $27$ & $35_{a}$ & $7$ & $21_{a}$ \\ \hline
\endhead
1) &  & 1 & 2 & 1 & 1 & 2 & 2 & 2 \\
\end{longtable}
\noindent Modules tested: V133, V912, V56.\\

\setcounter{table}{226}
\begin{longtable}{rc|ccccc|cc}
\caption{$Sp_6(2) < E_7$, $p = 3$}\\
& & \multicolumn{5}{c|}{$V_{133}$} & \multicolumn{2}{c}{$V_{56}$} \\
& & $1$ & $7$ & $21$ & $27$ & $35$ & $7$ & $21$ \\ \hline
\endfirsthead
& & \multicolumn{5}{c|}{$V_{133}$} & \multicolumn{2}{c}{$V_{56}$} \\
& & $1$ & $7$ & $21$ & $27$ & $35$ & $7$ & $21$ \\ \hline
\endhead
1) &  & 1 & 2 & 1 & 1 & 2 & 2 & 2 \\
\end{longtable}
\noindent Modules tested: V133, V912, L856, V56.\\

\setcounter{table}{227}
\begin{longtable}{rc|cccc|c}
\caption{$\Omega_8^{+}(2) < E_7$, $p \neq 2$}\\
& & \multicolumn{4}{c|}{$V_{133}$} & \multicolumn{1}{c}{$V_{56}$} \\
& & $28$ & $35_{a}$ & $35_{b}$ & $35_{c}$ & $28$ \\ \hline
\endfirsthead
& & \multicolumn{4}{c|}{$V_{133}$} & \multicolumn{1}{c}{$V_{56}$} \\
& & $28$ & $35_{a}$ & $35_{b}$ & $35_{c}$ & $28$ \\ \hline
\endhead
1) & \possprim & 1 & 1 & 1 & 1 & 2 \\
\end{longtable}
\noindent Modules tested: V133, V912, V56.\\

\setcounter{table}{228}
\begin{longtable}{rc|ccc|cc}
\caption{$^{3}D_4(2) < E_7$, $p \neq 2, 3$}\\
& & \multicolumn{3}{c|}{$V_{133}$} & \multicolumn{2}{c}{$V_{56}$} \\
& & $1$ & $26$ & $52$ & $1$ & $26$ \\ \hline
\endfirsthead
& & \multicolumn{3}{c|}{$V_{133}$} & \multicolumn{2}{c}{$V_{56}$} \\
& & $1$ & $26$ & $52$ & $1$ & $26$ \\ \hline
\endhead
1) &  & 3 & 3 & 1 & 4 & 2 \\
\end{longtable}
\noindent Modules tested: V133, V912, V56.\\

\setcounter{table}{229}
\begin{longtable}{rc|ccc|cc}
\caption{$^{3}D_4(2) < E_7$, $p = 3$}\\
& & \multicolumn{3}{c|}{$V_{133}$} & \multicolumn{2}{c}{$V_{56}$} \\
& & $1$ & $25$ & $52$ & $1$ & $25$ \\ \hline
\endfirsthead
& & \multicolumn{3}{c|}{$V_{133}$} & \multicolumn{2}{c}{$V_{56}$} \\
& & $1$ & $25$ & $52$ & $1$ & $25$ \\ \hline
\endhead
1) & \nongcr & 6 & 3 & 1 & 6 & 2 \\
\end{longtable}
\noindent Modules tested: V133, V912, L856, V56.\\

\setcounter{table}{230}
\begin{longtable}{rc|cccc|ccc}
\caption{$^{2}F_4(2)' < E_7$, $p \neq 2, 3, 5$}\\
& & \multicolumn{4}{c|}{$V_{133}$} & \multicolumn{3}{c}{$V_{56}$} \\
& & $1$ & $27_{a}$ & $27_{b}$ & $78$ & $1$ & $27_{a}$ & $27_{b}$ \\ \hline
\endfirsthead
& & \multicolumn{4}{c|}{$V_{133}$} & \multicolumn{3}{c}{$V_{56}$} \\
& & $1$ & $27_{a}$ & $27_{b}$ & $78$ & $1$ & $27_{a}$ & $27_{b}$ \\ \hline
\endhead
1) &  & 1 & 1 & 1 & 1 & 2 & 1 & 1 \\
\end{longtable}
\noindent Modules tested: V133, V912, V56.\\

\setcounter{table}{231}
\begin{longtable}{rc|cccc|ccc}
\caption{$^{2}F_4(2)' < E_7$, $p = 5$}\\
& & \multicolumn{4}{c|}{$V_{133}$} & \multicolumn{3}{c}{$V_{56}$} \\
& & $1$ & $27_{a}$ & $27_{b}$ & $78$ & $1$ & $27_{a}$ & $27_{b}$ \\ \hline
\endfirsthead
& & \multicolumn{4}{c|}{$V_{133}$} & \multicolumn{3}{c}{$V_{56}$} \\
& & $1$ & $27_{a}$ & $27_{b}$ & $78$ & $1$ & $27_{a}$ & $27_{b}$ \\ \hline
\endhead
1) & \nongcr & 1 & 1 & 1 & 1 & 2 & 1 & 1 \\
\end{longtable}
\noindent Modules tested: V133, V912, V56.\\

\setcounter{table}{232}
\begin{longtable}{rc|cccc|ccc}
\caption{$^{2}F_4(2)' < E_7$, $p = 3$}\\
& & \multicolumn{4}{c|}{$V_{133}$} & \multicolumn{3}{c}{$V_{56}$} \\
& & $1$ & $27_{a}$ & $27_{b}$ & $77$ & $1$ & $27_{a}$ & $27_{b}$ \\ \hline
\endfirsthead
& & \multicolumn{4}{c|}{$V_{133}$} & \multicolumn{3}{c}{$V_{56}$} \\
& & $1$ & $27_{a}$ & $27_{b}$ & $77$ & $1$ & $27_{a}$ & $27_{b}$ \\ \hline
\endhead
1) &  & 2 & 1 & 1 & 1 & 2 & 1 & 1 \\
\end{longtable}
\noindent Modules tested: V133, V912, L856, V56.\\

\setcounter{table}{233}
\begin{longtable}{rc|ccccc|cc}
\caption{$^{2}B_2(8) < E_7$, $p = 5$}\\
& & \multicolumn{5}{c|}{$V_{133}$} & \multicolumn{2}{c}{$V_{56}$} \\
& & $14_{a}$ & $14_{b}$ & $35_{a}$ & $35_{b}$ & $35_{c}$ & $14_{a}$ & $14_{b}$ \\ \hline
\endfirsthead
& & \multicolumn{5}{c|}{$V_{133}$} & \multicolumn{2}{c}{$V_{56}$} \\
& & $14_{a}$ & $14_{b}$ & $35_{a}$ & $35_{b}$ & $35_{c}$ & $14_{a}$ & $14_{b}$ \\ \hline
\endhead
1) & \newrow & 1 & 1 & 1 & 1 & 1 & 2 & 2 \\
\end{longtable}
\noindent Memoir rows not reproduced: 1) (not feasible on the Memoir's modules (all element orders, full power maps)).\\
Rows marked \newrow\ do not appear in the Memoir.\\
Modules tested: V133, V912, V56.\\


\newpage
% Updated tables for E8 (Memoir Chapter 6), generated by tools/fg_tables.g
\section*{$E_8$: feasible characters on $L(G) = V_{248}$}

\setcounter{table}{234}
\begin{longtable}{l}
\caption{Alt$_{n} < E_8$, $n \ge 10$, unique feasible character}\\
Not recomputed: no character table available in GAP. See the Memoir.
\end{longtable}

\setcounter{table}{235}
\begin{longtable}{rc|ccccccc}
\caption{Alt$_{10} < E_8$, $p = 5$}\\
& & \multicolumn{7}{c}{$V_{248}$} \\
& & $1$ & $8$ & $28$ & $35_{a}$ & $35_{b}$ & $35_{c}$ & $56$ \\ \hline
\endfirsthead
& & \multicolumn{7}{c}{$V_{248}$} \\
& & $1$ & $8$ & $28$ & $35_{a}$ & $35_{b}$ & $35_{c}$ & $56$ \\ \hline
\endhead
1) & \possprim, \nongcr & 1 & 2 & 3 & 1 & 0 & 0 & 2 \\
2) &  & 3 & 0 & 5 & 1 & 1 & 1 & 0 \\
\end{longtable}
\noindent Modules tested: V248, V3875, V147250.\\

\setcounter{table}{236}
\begin{longtable}{rc|cccccccc}
\caption{Alt$_{10} < E_8$, $p = 2$}\\
& & \multicolumn{8}{c}{$V_{248}$} \\
& & $1$ & $8$ & $16$ & $26$ & $48$ & $64_{a}$ & $64_{b}$ & $198$ \\ \hline
\endfirsthead
& & \multicolumn{8}{c}{$V_{248}$} \\
& & $1$ & $8$ & $16$ & $26$ & $48$ & $64_{a}$ & $64_{b}$ & $198$ \\ \hline
\endhead
1) &  & 2 & 0 & 0 & 0 & 1 & 0 & 0 & 1 \\
2) & \nongcr & 4 & 2 & 0 & 2 & 1 & 1 & 1 & 0 \\
3) & \nongcr & 4 & 2 & 1 & 2 & 2 & 1 & 0 & 0 \\
4) & \possprim, \nongcr & 8 & 5 & 0 & 4 & 2 & 0 & 0 & 0 \\
5) & \nongcr & 16 & 1 & 1 & 8 & 0 & 0 & 0 & 0 \\
6) & \nongcr & 30 & 8 & 8 & 1 & 0 & 0 & 0 & 0 \\
\end{longtable}
\noindent Modules tested: V248, V3875, L3626, V147250, L143376.\\

\setcounter{table}{237}
\begin{longtable}{rc|ccccccc}
\caption{Alt$_9 < E_8$, $p = 0$ or $p \nmid |H|$}\\
& & \multicolumn{7}{c}{$V_{248}$} \\
& & $1$ & $8$ & $27$ & $28$ & $35_{a}$ & $35_{b}$ & $56$ \\ \hline
\endfirsthead
& & \multicolumn{7}{c}{$V_{248}$} \\
& & $1$ & $8$ & $27$ & $28$ & $35_{a}$ & $35_{b}$ & $56$ \\ \hline
\endhead
1) &  & 1 & 3 & 1 & 3 & 0 & 0 & 2 \\
2) &  & 3 & 1 & 1 & 5 & 1 & 1 & 0 \\
\end{longtable}
\noindent Modules tested: V248, V3875, V147250.\\

\setcounter{table}{238}
\begin{longtable}{rc|ccccccc}
\caption{Alt$_9 < E_8$, $p = 7$}\\
& & \multicolumn{7}{c}{$V_{248}$} \\
& & $1$ & $8$ & $19$ & $28$ & $35_{a}$ & $35_{b}$ & $56$ \\ \hline
\endfirsthead
& & \multicolumn{7}{c}{$V_{248}$} \\
& & $1$ & $8$ & $19$ & $28$ & $35_{a}$ & $35_{b}$ & $56$ \\ \hline
\endhead
1) &  & 1 & 4 & 1 & 3 & 0 & 0 & 2 \\
2) &  & 3 & 2 & 1 & 5 & 1 & 1 & 0 \\
\end{longtable}
\noindent Modules tested: V248, V3875, V147250, L147002.\\

\setcounter{table}{239}
\begin{longtable}{rc|ccccccc}
\caption{Alt$_9 < E_8$, $p = 5$}\\
& & \multicolumn{7}{c}{$V_{248}$} \\
& & $1$ & $8$ & $27$ & $28$ & $35_{a}$ & $35_{b}$ & $56$ \\ \hline
\endfirsthead
& & \multicolumn{7}{c}{$V_{248}$} \\
& & $1$ & $8$ & $27$ & $28$ & $35_{a}$ & $35_{b}$ & $56$ \\ \hline
\endhead
1) &  & 1 & 3 & 1 & 3 & 0 & 0 & 2 \\
2) &  & 3 & 1 & 1 & 5 & 1 & 1 & 0 \\
\end{longtable}
\noindent Memoir rows not reproduced: 1) (not compatible with V3875).\\
Modules tested: V248, V3875, V147250.\\

\setcounter{table}{240}
\begin{longtable}{rc|ccccc}
\caption{Alt$_9 < E_8$, $p = 3$}\\
& & \multicolumn{5}{c}{$V_{248}$} \\
& & $1$ & $7$ & $21$ & $27$ & $35$ \\ \hline
\endfirsthead
& & \multicolumn{5}{c}{$V_{248}$} \\
& & $1$ & $7$ & $21$ & $27$ & $35$ \\ \hline
\endhead
1) & \possprim, \nongcr & 4 & 6 & 5 & 1 & 2 \\
\end{longtable}
\noindent Modules tested: V248, V3875, V147250, L113243.\\

\setcounter{table}{241}
\begin{longtable}{rc|ccccccccc}
\caption{Alt$_9 < E_8$, $p = 2$}\\
& & \multicolumn{9}{c}{$V_{248}$} \\
& & $1$ & $8_{a}$ & $8_{b}$ & $8_{c}$ & $20_{a}$ & $20_{b}$ & $26$ & $48$ & $78$ \\ \hline
\endfirsthead
& & \multicolumn{9}{c}{$V_{248}$} \\
& & $1$ & $8_{a}$ & $8_{b}$ & $8_{c}$ & $20_{a}$ & $20_{b}$ & $26$ & $48$ & $78$ \\ \hline
\endhead
1) & \possprim, \nongcr & 2 & 2 & 1 & 1 & 1 & 1 & 0 & 2 & 1 \\
2) & \nongcr & 4 & 0 & 0 & 0 & 1 & 1 & 0 & 1 & 2 \\
3) & \possprim, \nongcr & 4 & 0 & 1 & 0 & 2 & 2 & 0 & 0 & 2 \\
4) & \possprim, \nongcr & 4 & 2 & 3 & 2 & 1 & 1 & 2 & 2 & 0 \\
5) & \possprim, \nongcr & 4 & 2 & 3 & 3 & 2 & 2 & 2 & 1 & 0 \\
6) & \possprim, \nongcr & 6 & 0 & 3 & 0 & 1 & 1 & 2 & 1 & 1 \\
7) & \possprim, \nongcr & 6 & 0 & 2 & 2 & 2 & 2 & 2 & 0 & 1 \\
8) & \possprim, \nongcr & 8 & 0 & 4 & 2 & 1 & 1 & 4 & 1 & 0 \\
9) & \possprim, \nongcr & 8 & 0 & 5 & 2 & 2 & 2 & 4 & 0 & 0 \\
10) & \nongcr & 8 & 5 & 0 & 0 & 0 & 0 & 4 & 2 & 0 \\
11) & \nongcr & 16 & 1 & 1 & 1 & 0 & 0 & 8 & 0 & 0 \\
12) &  & 30 & 8 & 8 & 8 & 0 & 0 & 1 & 0 & 0 \\
\end{longtable}
\noindent Modules tested: V248, V3875, L3626, V147250, L143376.\\

\setcounter{table}{242}
\begin{longtable}{rc|cccccccc}
\caption{Alt$_8 < E_8$, $p = 0$ or $p \nmid |H|$}\\
& & \multicolumn{8}{c}{$V_{248}$} \\
& & $1$ & $7$ & $14$ & $20$ & $21_{a}$ & $28$ & $35$ & $70$ \\ \hline
\endfirsthead
& & \multicolumn{8}{c}{$V_{248}$} \\
& & $1$ & $7$ & $14$ & $20$ & $21_{a}$ & $28$ & $35$ & $70$ \\ \hline
\endhead
1) &  & 3 & 2 & 1 & 0 & 0 & 4 & 1 & 1 \\
2) &  & 4 & 7 & 0 & 1 & 5 & 0 & 2 & 0 \\
\end{longtable}
\noindent Memoir rows not reproduced: 1) (not compatible with V147250).\\
Modules tested: V248, V3875, V147250.\\

\setcounter{table}{243}
\begin{longtable}{rc|ccccccccc}
\caption{Alt$_8 < E_8$, $p = 7$}\\
& & \multicolumn{9}{c}{$V_{248}$} \\
& & $1$ & $7$ & $14$ & $19$ & $21_{a}$ & $28$ & $35$ & $45$ & $70$ \\ \hline
\endfirsthead
& & \multicolumn{9}{c}{$V_{248}$} \\
& & $1$ & $7$ & $14$ & $19$ & $21_{a}$ & $28$ & $35$ & $45$ & $70$ \\ \hline
\endhead
1) &  & 3 & 2 & 1 & 0 & 0 & 4 & 1 & 0 & 1 \\
2) & \nongcr & 4 & 1 & 0 & 2 & 0 & 3 & 0 & 1 & 1 \\
3) &  & 5 & 7 & 0 & 1 & 5 & 0 & 2 & 0 & 0 \\
\end{longtable}
\noindent Modules tested: V248, V3875, V147250, L147002.\\

\setcounter{table}{244}
\begin{longtable}{rc|ccccccccc}
\caption{Alt$_8 < E_8$, $p = 5$}\\
& & \multicolumn{9}{c}{$V_{248}$} \\
& & $1$ & $7$ & $13$ & $20$ & $21_{a}$ & $21_{b}$ & $35$ & $43$ & $70$ \\ \hline
\endfirsthead
& & \multicolumn{9}{c}{$V_{248}$} \\
& & $1$ & $7$ & $13$ & $20$ & $21_{a}$ & $21_{b}$ & $35$ & $43$ & $70$ \\ \hline
\endhead
1) &  & 3 & 4 & 0 & 1 & 0 & 4 & 0 & 1 & 1 \\
2) &  & 3 & 4 & 0 & 1 & 1 & 3 & 0 & 1 & 1 \\
3) &  & 4 & 6 & 1 & 0 & 0 & 4 & 1 & 0 & 1 \\
4) &  & 4 & 7 & 0 & 1 & 5 & 0 & 2 & 0 & 0 \\
\end{longtable}
\noindent Modules tested: V248, V3875, V147250.\\

\setcounter{table}{245}
\begin{longtable}{rc|cccccc}
\caption{Alt$_8 < E_8$, $p = 3$}\\
& & \multicolumn{6}{c}{$V_{248}$} \\
& & $1$ & $7$ & $13$ & $21$ & $28$ & $35$ \\ \hline
\endfirsthead
& & \multicolumn{6}{c}{$V_{248}$} \\
& & $1$ & $7$ & $13$ & $21$ & $28$ & $35$ \\ \hline
\endhead
1) & \nongcr & 4 & 3 & 1 & 0 & 5 & 2 \\
2) & \nongcr & 4 & 8 & 1 & 5 & 0 & 2 \\
\end{longtable}
\noindent Modules tested: V248, V3875, V147250, L113243.\\

\setcounter{table}{246}
\begin{longtable}{rc|ccccccccc}
\caption{Alt$_7 < E_8$, $p = 0$ or $p \nmid |H|$}\\
& & \multicolumn{9}{c}{$V_{248}$} \\
& & $1$ & $6$ & $10_{a}$ & $10_{b}$ & $14_{a}$ & $14_{b}$ & $15$ & $21$ & $35$ \\ \hline
\endfirsthead
& & \multicolumn{9}{c}{$V_{248}$} \\
& & $1$ & $6$ & $10_{a}$ & $10_{b}$ & $14_{a}$ & $14_{b}$ & $15$ & $21$ & $35$ \\ \hline
\endhead
1) & \possprim & 0 & 0 & 1 & 1 & 2 & 0 & 4 & 0 & 4 \\
2) &  & 1 & 0 & 2 & 2 & 1 & 2 & 4 & 0 & 3 \\
3) &  & 2 & 0 & 3 & 3 & 0 & 4 & 4 & 0 & 2 \\
4) &  & 2 & 1 & 0 & 0 & 5 & 4 & 2 & 4 & 0 \\
5) &  & 2 & 1 & 0 & 0 & 2 & 7 & 2 & 4 & 0 \\
6) &  & 2 & 2 & 0 & 0 & 5 & 4 & 3 & 3 & 0 \\
7) &  & 2 & 2 & 0 & 0 & 2 & 7 & 3 & 3 & 0 \\
8) &  & 4 & 1 & 0 & 0 & 5 & 4 & 0 & 2 & 2 \\
9) &  & 4 & 1 & 0 & 0 & 2 & 7 & 0 & 2 & 2 \\
10) &  & 4 & 2 & 0 & 0 & 5 & 4 & 1 & 1 & 2 \\
11) &  & 4 & 2 & 0 & 0 & 2 & 7 & 1 & 1 & 2 \\
12) &  & 5 & 1 & 1 & 1 & 4 & 6 & 0 & 2 & 1 \\
13) &  & 5 & 2 & 1 & 1 & 4 & 6 & 1 & 1 & 1 \\
14) &  & 11 & 13 & 2 & 2 & 1 & 0 & 7 & 0 & 0 \\
\end{longtable}
\noindent Modules tested: V248, V3875, V147250.\\

\setcounter{table}{247}
\begin{longtable}{rc|ccccccc}
\caption{Alt$_7 < E_8$, $p = 7$}\\
& & \multicolumn{7}{c}{$V_{248}$} \\
& & $1$ & $5$ & $10$ & $14_{a}$ & $14_{b}$ & $21$ & $35$ \\ \hline
\endfirsthead
& & \multicolumn{7}{c}{$V_{248}$} \\
& & $1$ & $5$ & $10$ & $14_{a}$ & $14_{b}$ & $21$ & $35$ \\ \hline
\endhead
1) & \possprim & 0 & 0 & 8 & 0 & 4 & 2 & 2 \\
2) & \possprim & 0 & 2 & 7 & 1 & 2 & 1 & 3 \\
3) & \possprim & 0 & 4 & 6 & 2 & 0 & 0 & 4 \\
4) & \possprim, \nongcr & 1 & 2 & 9 & 0 & 4 & 1 & 2 \\
5) & \possprim, \nongcr & 1 & 4 & 8 & 1 & 2 & 0 & 3 \\
6) & \possprim, \nongcr & 2 & 3 & 0 & 3 & 5 & 4 & 1 \\
7) & \possprim, \nongcr & 2 & 3 & 0 & 6 & 2 & 4 & 1 \\
8) & \possprim, \nongcr & 2 & 4 & 10 & 0 & 4 & 0 & 2 \\
9) & \nongcr & 3 & 3 & 2 & 2 & 7 & 4 & 0 \\
10) & \nongcr & 3 & 3 & 2 & 5 & 4 & 4 & 0 \\
11) & \possprim, \nongcr & 3 & 5 & 1 & 3 & 5 & 3 & 1 \\
12) & \possprim, \nongcr & 3 & 5 & 1 & 6 & 2 & 3 & 1 \\
13) & \possprim, \nongcr & 4 & 5 & 3 & 2 & 7 & 3 & 0 \\
14) & \possprim, \nongcr & 4 & 5 & 3 & 5 & 4 & 3 & 0 \\
15) &  & 5 & 1 & 0 & 2 & 7 & 2 & 2 \\
16) &  & 5 & 1 & 0 & 5 & 4 & 2 & 2 \\
17) &  & 6 & 1 & 2 & 4 & 6 & 2 & 1 \\
18) & \nongcr & 6 & 3 & 1 & 2 & 7 & 1 & 2 \\
19) & \nongcr & 6 & 3 & 1 & 5 & 4 & 1 & 2 \\
20) & \nongcr & 7 & 3 & 3 & 4 & 6 & 1 & 1 \\
21) & \nongcr & 15 & 2 & 2 & 1 & 0 & 9 & 0 \\
22) & \nongcr & 24 & 20 & 11 & 1 & 0 & 0 & 0 \\
\end{longtable}
\noindent Modules tested: V248, V3875, V147250, L147002.\\

\setcounter{table}{248}
\begin{longtable}{rc|cccccccc}
\caption{Alt$_7 < E_8$, $p = 5$}\\
& & \multicolumn{8}{c}{$V_{248}$} \\
& & $1$ & $6$ & $8$ & $10_{a}$ & $10_{b}$ & $13$ & $15$ & $35$ \\ \hline
\endfirsthead
& & \multicolumn{8}{c}{$V_{248}$} \\
& & $1$ & $6$ & $8$ & $10_{a}$ & $10_{b}$ & $13$ & $15$ & $35$ \\ \hline
\endhead
1) & \possprim & 0 & 0 & 4 & 1 & 1 & 2 & 2 & 4 \\
2) & \possprim & 0 & 1 & 3 & 1 & 1 & 1 & 3 & 4 \\
3) & \possprim & 0 & 2 & 2 & 1 & 1 & 0 & 4 & 4 \\
4) & \possprim & 0 & 4 & 12 & 4 & 4 & 1 & 0 & 1 \\
5) &  & 2 & 0 & 2 & 1 & 1 & 0 & 0 & 6 \\
6) & \nongcr & 3 & 0 & 2 & 2 & 2 & 3 & 3 & 3 \\
7) & \nongcr & 3 & 1 & 1 & 2 & 2 & 2 & 4 & 3 \\
8) & \nongcr & 3 & 3 & 11 & 5 & 5 & 3 & 0 & 0 \\
9) & \possprim, \nongcr & 3 & 7 & 14 & 0 & 0 & 7 & 0 & 0 \\
10) &  & 5 & 0 & 0 & 2 & 2 & 1 & 1 & 5 \\
11) & \nongcr & 6 & 0 & 0 & 3 & 3 & 4 & 4 & 2 \\
12) & \possprim, \nongcr & 6 & 6 & 9 & 0 & 0 & 8 & 2 & 0 \\
13) & \possprim, \nongcr & 6 & 7 & 8 & 0 & 0 & 7 & 3 & 0 \\
14) & \nongcr & 8 & 6 & 7 & 0 & 0 & 6 & 0 & 2 \\
15) & \nongcr & 8 & 7 & 6 & 0 & 0 & 5 & 1 & 2 \\
16) & \possprim, \nongcr & 9 & 3 & 6 & 0 & 0 & 11 & 2 & 0 \\
17) & \possprim, \nongcr & 9 & 4 & 5 & 0 & 0 & 10 & 3 & 0 \\
18) & \nongcr & 11 & 3 & 4 & 0 & 0 & 9 & 0 & 2 \\
19) & \nongcr & 11 & 4 & 3 & 0 & 0 & 8 & 1 & 2 \\
20) & \nongcr & 11 & 5 & 6 & 1 & 1 & 8 & 0 & 1 \\
21) & \nongcr & 11 & 6 & 5 & 1 & 1 & 7 & 1 & 1 \\
22) &  & 11 & 14 & 1 & 2 & 2 & 0 & 7 & 0 \\
23) &  & 28 & 0 & 25 & 1 & 1 & 0 & 0 & 0 \\
\end{longtable}
\noindent Modules tested: V248, V3875, V147250.\\

\setcounter{table}{249}
\begin{longtable}{rc|cccccc}
\caption{Alt$_7 < E_8$, $p = 3$}\\
& & \multicolumn{6}{c}{$V_{248}$} \\
& & $1$ & $6$ & $10_{a}$ & $10_{b}$ & $13$ & $15$ \\ \hline
\endfirsthead
& & \multicolumn{6}{c}{$V_{248}$} \\
& & $1$ & $6$ & $10_{a}$ & $10_{b}$ & $13$ & $15$ \\ \hline
\endhead
1) & \possprim, \nongcr & 10 & 0 & 5 & 5 & 6 & 4 \\
2) & \possprim, \nongcr & 11 & 5 & 0 & 0 & 9 & 6 \\
3) & \nongcr & 12 & 13 & 2 & 2 & 1 & 7 \\
4) & \possprim, \nongcr & 17 & 3 & 2 & 2 & 11 & 2 \\
\end{longtable}
\noindent Modules tested: V248, V3875, V147250, L113243.\\

\setcounter{table}{250}
\begin{longtable}{rc|cccccc}
\caption{Alt$_7 < E_8$, $p = 2$}\\
& & \multicolumn{6}{c}{$V_{248}$} \\
& & $1$ & $4_{a}$ & $4_{b}$ & $6$ & $14$ & $20$ \\ \hline
\endfirsthead
& & \multicolumn{6}{c}{$V_{248}$} \\
& & $1$ & $4_{a}$ & $4_{b}$ & $6$ & $14$ & $20$ \\ \hline
\endhead
1) & \possprim, \nongcr & 8 & 1 & 1 & 0 & 8 & 6 \\
2) & \possprim, \nongcr & 8 & 1 & 1 & 1 & 9 & 5 \\
3) & \possprim, \nongcr & 8 & 1 & 1 & 2 & 10 & 4 \\
4) & \possprim, \nongcr & 8 & 4 & 4 & 4 & 6 & 5 \\
5) & \possprim, \nongcr & 8 & 4 & 4 & 5 & 7 & 4 \\
6) & \possprim, \nongcr & 8 & 4 & 4 & 6 & 8 & 3 \\
7) & \nongcr & 8 & 7 & 7 & 8 & 4 & 4 \\
8) & \nongcr & 8 & 7 & 7 & 9 & 5 & 3 \\
9) & \nongcr & 8 & 7 & 7 & 10 & 6 & 2 \\
10) & \nongcr & 18 & 2 & 2 & 9 & 0 & 8 \\
11) & \nongcr & 18 & 2 & 2 & 10 & 1 & 7 \\
12) & \nongcr & 18 & 2 & 2 & 17 & 8 & 0 \\
13) &  & 46 & 16 & 16 & 9 & 0 & 1 \\
14) &  & 46 & 16 & 16 & 10 & 1 & 0 \\
\end{longtable}
\noindent Modules tested: V248, V3875, L3626, V147250, L143376.\\

\setcounter{table}{251}
\begin{longtable}{rc|ccccccc}
\caption{Alt$_6 < E_8$, $p = 0$ or $p \nmid |H|$}\\
& & \multicolumn{7}{c}{$V_{248}$} \\
& & $1$ & $5_{a}$ & $5_{b}$ & $8_{a}$ & $8_{b}$ & $9$ & $10$ \\ \hline
\endfirsthead
& & \multicolumn{7}{c}{$V_{248}$} \\
& & $1$ & $5_{a}$ & $5_{b}$ & $8_{a}$ & $8_{b}$ & $9$ & $10$ \\ \hline
\endhead
1) & \possprim & 0 & 0 & 0 & 6 & 6 & 8 & 8 \\
2) & \possprim & 0 & 6 & 0 & 4 & 4 & 6 & 10 \\
3) & \possprim & 0 & 2 & 2 & 7 & 7 & 4 & 8 \\
4) & \possprim & 0 & 3 & 3 & 4 & 4 & 6 & 10 \\
5) &  & 1 & 4 & 1 & 6 & 6 & 4 & 9 \\
6) &  & 1 & 2 & 2 & 10 & 5 & 3 & 8 \\
7) &  & 1 & 5 & 2 & 3 & 3 & 6 & 11 \\
8) &  & 1 & 7 & 7 & 6 & 6 & 9 & 0 \\
9) &  & 2 & 0 & 0 & 7 & 7 & 6 & 8 \\
10) &  & 2 & 6 & 0 & 5 & 5 & 4 & 10 \\
11) &  & 2 & 1 & 1 & 4 & 4 & 8 & 10 \\
12) &  & 2 & 4 & 1 & 9 & 4 & 3 & 9 \\
13) &  & 2 & 3 & 3 & 5 & 5 & 4 & 10 \\
14) &  & 2 & 4 & 4 & 2 & 2 & 6 & 12 \\
15) &  & 2 & 9 & 6 & 5 & 5 & 9 & 1 \\
16) &  & 3 & 0 & 0 & 10 & 5 & 5 & 8 \\
17) &  & 3 & 3 & 0 & 3 & 3 & 8 & 11 \\
18) &  & 3 & 6 & 0 & 8 & 3 & 3 & 10 \\
19) &  & 3 & 5 & 2 & 4 & 4 & 4 & 11 \\
20) &  & 3 & 3 & 3 & 8 & 3 & 3 & 10 \\
21) &  & 3 & 5 & 5 & 6 & 6 & 11 & 0 \\
22) &  & 3 & 11 & 5 & 4 & 4 & 9 & 2 \\
23) &  & 3 & 7 & 7 & 7 & 7 & 7 & 0 \\
24) &  & 3 & 8 & 8 & 4 & 4 & 9 & 2 \\
25) &  & 4 & 1 & 1 & 5 & 5 & 6 & 10 \\
26) &  & 4 & 2 & 2 & 2 & 2 & 8 & 12 \\
27) &  & 4 & 5 & 2 & 7 & 2 & 3 & 11 \\
28) &  & 4 & 4 & 4 & 3 & 3 & 4 & 12 \\
29) &  & 4 & 7 & 4 & 5 & 5 & 11 & 1 \\
30) &  & 4 & 9 & 6 & 6 & 6 & 7 & 1 \\
31) &  & 4 & 7 & 7 & 10 & 5 & 6 & 0 \\
32) &  & 4 & 10 & 7 & 3 & 3 & 9 & 3 \\
33) &  & 5 & 3 & 0 & 4 & 4 & 6 & 11 \\
34) &  & 5 & 1 & 1 & 8 & 3 & 5 & 10 \\
35) &  & 5 & 9 & 3 & 4 & 4 & 11 & 2 \\
36) &  & 5 & 4 & 4 & 6 & 1 & 3 & 12 \\
37) &  & 5 & 5 & 5 & 7 & 7 & 9 & 0 \\
38) &  & 5 & 11 & 5 & 5 & 5 & 7 & 2 \\
39) &  & 5 & 6 & 6 & 4 & 4 & 11 & 2 \\
40) &  & 5 & 9 & 6 & 9 & 4 & 6 & 1 \\
41) &  & 5 & 8 & 8 & 5 & 5 & 7 & 2 \\
42) &  & 5 & 9 & 9 & 2 & 2 & 9 & 4 \\
43) &  & 6 & 3 & 0 & 7 & 2 & 5 & 11 \\
44) &  & 6 & 2 & 2 & 3 & 3 & 6 & 12 \\
45) &  & 6 & 7 & 4 & 6 & 6 & 9 & 1 \\
46) &  & 6 & 5 & 5 & 10 & 5 & 8 & 0 \\
47) &  & 6 & 8 & 5 & 3 & 3 & 11 & 3 \\
48) &  & 6 & 11 & 5 & 8 & 3 & 6 & 2 \\
49) &  & 6 & 10 & 7 & 4 & 4 & 7 & 3 \\
50) &  & 6 & 8 & 8 & 8 & 3 & 6 & 2 \\
51) &  & 7 & 2 & 2 & 6 & 1 & 5 & 12 \\
52) &  & 7 & 9 & 3 & 5 & 5 & 9 & 2 \\
53) &  & 7 & 7 & 4 & 9 & 4 & 8 & 1 \\
54) &  & 7 & 6 & 6 & 5 & 5 & 9 & 2 \\
55) &  & 7 & 7 & 7 & 2 & 2 & 11 & 4 \\
56) &  & 7 & 10 & 7 & 7 & 2 & 6 & 3 \\
57) &  & 7 & 9 & 9 & 3 & 3 & 7 & 4 \\
58) &  & 8 & 0 & 0 & 20 & 0 & 0 & 8 \\
59) &  & 8 & 9 & 3 & 8 & 3 & 8 & 2 \\
60) &  & 8 & 8 & 5 & 4 & 4 & 9 & 3 \\
61) &  & 8 & 6 & 6 & 8 & 3 & 8 & 2 \\
62) &  & 8 & 9 & 9 & 6 & 1 & 6 & 4 \\
63) &  & 9 & 8 & 5 & 7 & 2 & 8 & 3 \\
64) &  & 9 & 7 & 7 & 20 & 0 & 1 & 0 \\
65) &  & 9 & 7 & 7 & 3 & 3 & 9 & 4 \\
66) &  & 10 & 7 & 7 & 6 & 1 & 8 & 4 \\
67) &  & 11 & 0 & 0 & 4 & 4 & 17 & 2 \\
68) &  & 11 & 5 & 5 & 20 & 0 & 3 & 0 \\
69) &  & 13 & 0 & 0 & 5 & 5 & 15 & 2 \\
70) &  & 13 & 1 & 1 & 2 & 2 & 17 & 4 \\
71) &  & 13 & 5 & 5 & 11 & 11 & 1 & 0 \\
72) &  & 14 & 0 & 0 & 8 & 3 & 14 & 2 \\
73) &  & 14 & 7 & 4 & 10 & 10 & 1 & 1 \\
74) &  & 15 & 1 & 1 & 3 & 3 & 15 & 4 \\
75) &  & 15 & 9 & 3 & 9 & 9 & 1 & 2 \\
76) &  & 15 & 6 & 6 & 9 & 9 & 1 & 2 \\
77) &  & 16 & 1 & 1 & 6 & 1 & 14 & 4 \\
78) &  & 16 & 8 & 5 & 8 & 8 & 1 & 3 \\
79) &  & 17 & 7 & 7 & 7 & 7 & 1 & 4 \\
80) &  & 21 & 0 & 0 & 9 & 9 & 7 & 2 \\
81) &  & 23 & 1 & 1 & 7 & 7 & 7 & 4 \\
82) &  & 24 & 21 & 0 & 0 & 0 & 1 & 11 \\
83) &  & 28 & 0 & 0 & 25 & 0 & 0 & 2 \\
\end{longtable}
\noindent Modules tested: V248, V3875, V147250.\\

\setcounter{table}{252}
\begin{longtable}{rc|ccccc}
\caption{Alt$_6 < E_8$, $p = 5$}\\
& & \multicolumn{5}{c}{$V_{248}$} \\
& & $1$ & $5_{a}$ & $5_{b}$ & $8$ & $10$ \\ \hline
\endfirsthead
& & \multicolumn{5}{c}{$V_{248}$} \\
& & $1$ & $5_{a}$ & $5_{b}$ & $8$ & $10$ \\ \hline
\endhead
1) & \possprim, \nongcr & 4 & 2 & 2 & 18 & 8 \\
2) & \possprim, \nongcr & 5 & 4 & 1 & 16 & 9 \\
3) & \possprim, \nongcr & 6 & 6 & 0 & 14 & 10 \\
4) & \possprim, \nongcr & 6 & 3 & 3 & 14 & 10 \\
5) & \possprim, \nongcr & 7 & 5 & 2 & 12 & 11 \\
6) & \possprim, \nongcr & 8 & 0 & 0 & 20 & 8 \\
7) & \possprim, \nongcr & 8 & 4 & 4 & 10 & 12 \\
8) & \possprim, \nongcr & 10 & 1 & 1 & 16 & 10 \\
9) & \possprim, \nongcr & 10 & 7 & 7 & 21 & 0 \\
10) & \possprim, \nongcr & 11 & 3 & 0 & 14 & 11 \\
11) & \possprim, \nongcr & 11 & 9 & 6 & 19 & 1 \\
12) & \nongcr & 12 & 2 & 2 & 12 & 12 \\
13) & \possprim, \nongcr & 12 & 11 & 5 & 17 & 2 \\
14) & \possprim, \nongcr & 12 & 8 & 8 & 17 & 2 \\
15) & \possprim, \nongcr & 13 & 10 & 7 & 15 & 3 \\
16) & \possprim, \nongcr & 14 & 5 & 5 & 23 & 0 \\
17) & \nongcr & 14 & 9 & 9 & 13 & 4 \\
18) & \possprim, \nongcr & 15 & 7 & 4 & 21 & 1 \\
19) & \possprim, \nongcr & 16 & 9 & 3 & 19 & 2 \\
20) & \possprim, \nongcr & 16 & 6 & 6 & 19 & 2 \\
21) & \nongcr & 17 & 8 & 5 & 17 & 3 \\
22) & \nongcr & 18 & 7 & 7 & 15 & 4 \\
23) &  & 25 & 21 & 0 & 1 & 11 \\
24) & \nongcr & 28 & 0 & 0 & 25 & 2 \\
25) & \nongcr & 30 & 1 & 1 & 21 & 4 \\
\end{longtable}
\noindent Memoir rows not reproduced: 23) (not compatible with V3875).\\
Modules tested: V248, V3875, V147250.\\

\setcounter{table}{253}
\begin{longtable}{rc|ccccc}
\caption{Alt$_5 < E_8$, $p = 0$ or $p \nmid |H|$}\\
& & \multicolumn{5}{c}{$V_{248}$} \\
& & $1$ & $3_{a}$ & $3_{b}$ & $4$ & $5$ \\ \hline
\endfirsthead
& & \multicolumn{5}{c}{$V_{248}$} \\
& & $1$ & $3_{a}$ & $3_{b}$ & $4$ & $5$ \\ \hline
\endhead
1) & \possprim & 0 & 14 & 14 & 16 & 20 \\
2) &  & 2 & 15 & 15 & 14 & 20 \\
3) &  & 3 & 18 & 13 & 13 & 20 \\
4) &  & 3 & 14 & 14 & 19 & 17 \\
5) &  & 5 & 15 & 15 & 17 & 17 \\
6) &  & 6 & 18 & 13 & 16 & 17 \\
7) &  & 6 & 14 & 14 & 22 & 14 \\
8) &  & 8 & 6 & 6 & 16 & 28 \\
9) &  & 8 & 28 & 8 & 8 & 20 \\
10) &  & 8 & 15 & 15 & 20 & 14 \\
11) &  & 9 & 18 & 13 & 19 & 14 \\
12) &  & 10 & 7 & 7 & 14 & 28 \\
13) &  & 10 & 19 & 19 & 6 & 20 \\
14) &  & 11 & 10 & 5 & 13 & 28 \\
15) &  & 11 & 6 & 6 & 19 & 25 \\
16) &  & 11 & 28 & 8 & 11 & 17 \\
17) &  & 13 & 7 & 7 & 17 & 25 \\
18) &  & 13 & 19 & 19 & 9 & 17 \\
19) &  & 14 & 10 & 5 & 16 & 25 \\
20) &  & 14 & 6 & 6 & 22 & 22 \\
21) &  & 14 & 28 & 8 & 14 & 14 \\
22) &  & 16 & 20 & 0 & 8 & 28 \\
23) &  & 16 & 7 & 7 & 20 & 22 \\
24) &  & 16 & 19 & 19 & 12 & 14 \\
25) &  & 17 & 10 & 5 & 19 & 22 \\
26) &  & 18 & 11 & 11 & 6 & 28 \\
27) &  & 19 & 20 & 0 & 11 & 25 \\
28) &  & 20 & 35 & 10 & 2 & 17 \\
29) &  & 21 & 11 & 11 & 9 & 25 \\
30) &  & 22 & 20 & 0 & 14 & 22 \\
31) &  & 23 & 35 & 10 & 5 & 14 \\
32) &  & 24 & 11 & 11 & 12 & 22 \\
33) &  & 28 & 50 & 0 & 0 & 14 \\
34) &  & 28 & 27 & 2 & 2 & 25 \\
35) &  & 31 & 27 & 2 & 5 & 22 \\
36) &  & 35 & 6 & 6 & 43 & 1 \\
37) &  & 37 & 7 & 7 & 41 & 1 \\
38) &  & 38 & 10 & 5 & 40 & 1 \\
39) &  & 43 & 20 & 0 & 35 & 1 \\
40) &  & 45 & 11 & 11 & 33 & 1 \\
41) &  & 52 & 27 & 2 & 26 & 1 \\
42) &  & 78 & 55 & 0 & 0 & 1 \\
\end{longtable}
\noindent Modules tested: V248, V3875, V147250.\\

\setcounter{table}{254}
\begin{longtable}{rc|cccc}
\caption{Alt$_5 < E_8$, $p = 3$}\\
& & \multicolumn{4}{c}{$V_{248}$} \\
& & $1$ & $3_{a}$ & $3_{b}$ & $4$ \\ \hline
\endfirsthead
& & \multicolumn{4}{c}{$V_{248}$} \\
& & $1$ & $3_{a}$ & $3_{b}$ & $4$ \\ \hline
\endhead
1) & \possprim, \nongcr & 20 & 14 & 14 & 36 \\
2) & \possprim, \nongcr & 22 & 15 & 15 & 34 \\
3) & \possprim, \nongcr & 23 & 18 & 13 & 33 \\
4) & \nongcr & 28 & 28 & 8 & 28 \\
5) & \nongcr & 30 & 19 & 19 & 26 \\
6) & \possprim, \nongcr & 36 & 6 & 6 & 44 \\
7) & \nongcr & 37 & 35 & 10 & 19 \\
8) & \possprim, \nongcr & 38 & 7 & 7 & 42 \\
9) & \possprim, \nongcr & 39 & 10 & 5 & 41 \\
10) & \nongcr & 42 & 50 & 0 & 14 \\
11) & \nongcr & 44 & 20 & 0 & 36 \\
12) & \nongcr & 46 & 11 & 11 & 34 \\
13) & \nongcr & 53 & 27 & 2 & 27 \\
14) &  & 79 & 55 & 0 & 1 \\
\end{longtable}
\noindent Modules tested: V248, V3875, V147250, L113243.\\

\setcounter{table}{255}
\begin{longtable}{rc|ccccccc}
\caption{$M_{11} < E_8$, $p = 0$ or $p \nmid |H|$}\\
& & \multicolumn{7}{c}{$V_{248}$} \\
& & $1$ & $10_{a}$ & $11$ & $16_{a}$ & $16_{b}$ & $45$ & $55$ \\ \hline
\endfirsthead
& & \multicolumn{7}{c}{$V_{248}$} \\
& & $1$ & $10_{a}$ & $11$ & $16_{a}$ & $16_{b}$ & $45$ & $55$ \\ \hline
\endhead
1) &  & 9 & 0 & 6 & 4 & 4 & 1 & 0 \\
2) &  & 10 & 0 & 5 & 4 & 4 & 0 & 1 \\
3) &  & 15 & 6 & 0 & 4 & 4 & 1 & 0 \\
\end{longtable}
\noindent Modules tested: V248, V3875, V147250.\\

\setcounter{table}{256}
\begin{longtable}{rc|cccccccc}
\caption{$M_{11} < E_8$, $p = 11$}\\
& & \multicolumn{8}{c}{$V_{248}$} \\
& & $1$ & $9$ & $10_{a}$ & $10_{b}$ & $11$ & $16$ & $44$ & $55$ \\ \hline
\endfirsthead
& & \multicolumn{8}{c}{$V_{248}$} \\
& & $1$ & $9$ & $10_{a}$ & $10_{b}$ & $11$ & $16$ & $44$ & $55$ \\ \hline
\endhead
1) & \possprim & 0 & 2 & 3 & 3 & 0 & 1 & 1 & 2 \\
2) & \possprim, \nongcr & 2 & 3 & 0 & 0 & 1 & 2 & 4 & 0 \\
3) & \possprim, \nongcr & 3 & 4 & 0 & 0 & 2 & 0 & 3 & 1 \\
4) & \possprim, \nongcr & 4 & 6 & 4 & 4 & 0 & 0 & 0 & 2 \\
5) &  & 9 & 1 & 1 & 1 & 6 & 9 & 0 & 0 \\
6) &  & 10 & 0 & 0 & 0 & 5 & 8 & 0 & 1 \\
7) & \nongcr & 21 & 7 & 1 & 1 & 0 & 9 & 0 & 0 \\
\end{longtable}
\noindent Modules tested: V248, V3875, V147250.\\

\setcounter{table}{257}
\begin{longtable}{rc|ccccccc}
\caption{$M_{11} < E_8$, $p = 5$}\\
& & \multicolumn{7}{c}{$V_{248}$} \\
& & $1$ & $10_{a}$ & $11$ & $16_{a}$ & $16_{b}$ & $45$ & $55$ \\ \hline
\endfirsthead
& & \multicolumn{7}{c}{$V_{248}$} \\
& & $1$ & $10_{a}$ & $11$ & $16_{a}$ & $16_{b}$ & $45$ & $55$ \\ \hline
\endhead
1) &  & 3 & 1 & 0 & 0 & 0 & 4 & 1 \\
2) & \nongcr & 4 & 0 & 0 & 2 & 2 & 4 & 0 \\
3) & \nongcr & 7 & 1 & 7 & 2 & 2 & 2 & 0 \\
4) & \nongcr & 8 & 1 & 6 & 2 & 2 & 1 & 1 \\
5) & \nongcr & 9 & 0 & 6 & 4 & 4 & 1 & 0 \\
6) & \nongcr & 9 & 1 & 5 & 2 & 2 & 0 & 2 \\
7) & \nongcr & 10 & 0 & 5 & 4 & 4 & 0 & 1 \\
8) & \nongcr & 14 & 7 & 0 & 2 & 2 & 1 & 1 \\
9) & \nongcr & 15 & 6 & 0 & 4 & 4 & 1 & 0 \\
\end{longtable}
\noindent Modules tested: V248, V3875, V147250.\\

\setcounter{table}{258}
\begin{longtable}{rc|cccccccc}
\caption{$M_{11} < E_8$, $p = 3$}\\
& & \multicolumn{8}{c}{$V_{248}$} \\
& & $1$ & $5_{a}$ & $5_{b}$ & $10_{a}$ & $10_{b}$ & $10_{c}$ & $24$ & $45$ \\ \hline
\endfirsthead
& & \multicolumn{8}{c}{$V_{248}$} \\
& & $1$ & $5_{a}$ & $5_{b}$ & $10_{a}$ & $10_{b}$ & $10_{c}$ & $24$ & $45$ \\ \hline
\endhead
1) & \possprim & 0 & 0 & 0 & 2 & 0 & 0 & 2 & 4 \\
2) & \possprim, \nongcr & 1 & 2 & 2 & 0 & 1 & 1 & 3 & 3 \\
3) & \possprim, \nongcr & 2 & 2 & 2 & 0 & 2 & 2 & 4 & 2 \\
4) & \possprim, \nongcr & 4 & 5 & 5 & 5 & 0 & 0 & 6 & 0 \\
5) & \possprim, \nongcr & 4 & 7 & 7 & 3 & 0 & 0 & 6 & 0 \\
6) & \nongcr & 23 & 4 & 4 & 6 & 4 & 4 & 0 & 1 \\
7) & \nongcr & 23 & 10 & 10 & 0 & 4 & 4 & 0 & 1 \\
8) & \nongcr & 24 & 10 & 10 & 0 & 5 & 5 & 1 & 0 \\
\end{longtable}
\noindent Modules tested: V248, V3875, V147250, L113243.\\

\setcounter{table}{259}
\begin{longtable}{rc|ccccc}
\caption{$M_{11} < E_8$, $p = 2$}\\
& & \multicolumn{5}{c}{$V_{248}$} \\
& & $1$ & $10$ & $16_{a}$ & $16_{b}$ & $44$ \\ \hline
\endfirsthead
& & \multicolumn{5}{c}{$V_{248}$} \\
& & $1$ & $10$ & $16_{a}$ & $16_{b}$ & $44$ \\ \hline
\endhead
1) & \nongcr & 14 & 3 & 5 & 5 & 1 \\
2) & \nongcr & 16 & 6 & 4 & 4 & 1 \\
3) & \nongcr & 18 & 9 & 3 & 3 & 1 \\
\end{longtable}
\noindent Modules tested: V248, V3875, L3626, V147250, L143376.\\

\setcounter{table}{260}
\begin{longtable}{rc|ccccc}
\caption{$M_{12} < E_8$, $p = 5$}\\
& & \multicolumn{5}{c}{$V_{248}$} \\
& & $1$ & $11_{a}$ & $16_{a}$ & $16_{b}$ & $78$ \\ \hline
\endfirsthead
& & \multicolumn{5}{c}{$V_{248}$} \\
& & $1$ & $11_{a}$ & $16_{a}$ & $16_{b}$ & $78$ \\ \hline
\endhead
1) &  & 8 & 6 & 3 & 3 & 1 \\
\end{longtable}
\noindent Modules tested: V248, V3875, V147250.\\

\setcounter{table}{261}
\begin{longtable}{rc|ccccc}
\caption{$M_{12} < E_8$, $p = 2$}\\
& & \multicolumn{5}{c}{$V_{248}$} \\
& & $1$ & $10$ & $16_{a}$ & $16_{b}$ & $44$ \\ \hline
\endfirsthead
& & \multicolumn{5}{c}{$V_{248}$} \\
& & $1$ & $10$ & $16_{a}$ & $16_{b}$ & $44$ \\ \hline
\endhead
1) & \nongcr & 16 & 6 & 4 & 4 & 1 \\
\end{longtable}
\noindent Modules tested: V248, V3875, L3626, V147250, L143376.\\

\setcounter{table}{262}
\begin{longtable}{rc|cccccc}
\caption{$J_{1} < E_8$, $p = 11$}\\
& & \multicolumn{6}{c}{$V_{248}$} \\
& & $1$ & $7$ & $14$ & $27$ & $64$ & $77_{c}$ \\ \hline
\endfirsthead
& & \multicolumn{6}{c}{$V_{248}$} \\
& & $1$ & $7$ & $14$ & $27$ & $64$ & $77_{c}$ \\ \hline
\endhead
1) &  & 1 & 0 & 3 & 0 & 2 & 1 \\
2) &  & 6 & 13 & 5 & 3 & 0 & 0 \\
3) &  & 8 & 0 & 1 & 6 & 1 & 0 \\
4) &  & 52 & 26 & 1 & 0 & 0 & 0 \\
\end{longtable}
\noindent Modules tested: V248, V3875, V147250.\\

\setcounter{table}{263}
\begin{longtable}{rc|cccccccc}
\caption{$J_{2} < E_8$, $p = 2$}\\
& & \multicolumn{8}{c}{$V_{248}$} \\
& & $1$ & $6_{a}$ & $6_{b}$ & $14_{a}$ & $14_{b}$ & $36$ & $64_{b}$ & $84$ \\ \hline
\endfirsthead
& & \multicolumn{8}{c}{$V_{248}$} \\
& & $1$ & $6_{a}$ & $6_{b}$ & $14_{a}$ & $14_{b}$ & $36$ & $64_{b}$ & $84$ \\ \hline
\endhead
1) &  & 2 & 1 & 0 & 1 & 3 & 1 & 1 & 1 \\
2) & \nongcr & 4 & 2 & 2 & 4 & 0 & 1 & 2 & 0 \\
3) & \nongcr & 8 & 4 & 3 & 2 & 1 & 2 & 0 & 1 \\
4) & \nongcr & 14 & 6 & 6 & 7 & 0 & 0 & 1 & 0 \\
5) & \nongcr & 16 & 8 & 8 & 1 & 1 & 3 & 0 & 0 \\
6) & \nongcr & 22 & 16 & 3 & 8 & 0 & 0 & 0 & 0 \\
7) & \nongcr & 78 & 26 & 0 & 1 & 0 & 0 & 0 & 0 \\
\end{longtable}
\noindent Modules tested: V248, V3875, L3626, V147250, L143376.\\

\setcounter{table}{264}
\begin{longtable}{rc|ccc}
\caption{$J_{3} < E_8$, $p = 2$}\\
& & \multicolumn{3}{c}{$V_{248}$} \\
& & $80$ & $84_{a}$ & $84_{b}$ \\ \hline
\endfirsthead
& & \multicolumn{3}{c}{$V_{248}$} \\
& & $80$ & $84_{a}$ & $84_{b}$ \\ \hline
\endhead
1) & \possprim & 1 & 1 & 1 \\
\end{longtable}
\noindent Modules tested: V248, V3875, L3626, V147250, L143376.\\

\setcounter{table}{265}
\begin{longtable}{l}
\caption{$Th < E_8$, $p = 3$ (irreducible on $V_{248}$)}\\
Not recomputed: no Brauer table available in GAP. See the Memoir.
\end{longtable}

\setcounter{table}{266}
\begin{longtable}{rc|cccccc}
\caption{$L_2(7) < E_8$, $p = 0$ or $p \nmid |H|$}\\
& & \multicolumn{6}{c}{$V_{248}$} \\
& & $1$ & $3_{a}$ & $3_{b}$ & $6$ & $7$ & $8$ \\ \hline
\endfirsthead
& & \multicolumn{6}{c}{$V_{248}$} \\
& & $1$ & $3_{a}$ & $3_{b}$ & $6$ & $7$ & $8$ \\ \hline
\endhead
1) & \possprim & 0 & 5 & 5 & 6 & 10 & 14 \\
2) &  & 2 & 7 & 7 & 6 & 8 & 14 \\
3) &  & 3 & 0 & 0 & 14 & 7 & 14 \\
4) &  & 3 & 5 & 5 & 6 & 13 & 11 \\
5) &  & 5 & 2 & 2 & 14 & 5 & 14 \\
6) &  & 5 & 7 & 7 & 6 & 11 & 11 \\
7) &  & 6 & 0 & 0 & 14 & 10 & 11 \\
8) &  & 6 & 5 & 5 & 6 & 16 & 8 \\
9) &  & 8 & 2 & 2 & 14 & 8 & 11 \\
10) &  & 8 & 7 & 7 & 6 & 14 & 8 \\
11) &  & 9 & 0 & 0 & 14 & 13 & 8 \\
12) &  & 11 & 2 & 2 & 14 & 11 & 8 \\
13) &  & 14 & 8 & 8 & 14 & 2 & 11 \\
14) &  & 17 & 8 & 8 & 14 & 5 & 8 \\
15) &  & 22 & 21 & 21 & 6 & 0 & 8 \\
16) &  & 28 & 1 & 1 & 0 & 2 & 25 \\
17) &  & 31 & 1 & 1 & 0 & 5 & 22 \\
18) &  & 52 & 1 & 1 & 0 & 26 & 1 \\
19) &  & 78 & 27 & 27 & 0 & 0 & 1 \\
\end{longtable}
\noindent Modules tested: V248, V3875, V147250.\\

\setcounter{table}{267}
\begin{longtable}{rc|ccccc}
\caption{$L_2(7) < E_8$, $p = 3$}\\
& & \multicolumn{5}{c}{$V_{248}$} \\
& & $1$ & $3_{a}$ & $3_{b}$ & $6$ & $7$ \\ \hline
\endfirsthead
& & \multicolumn{5}{c}{$V_{248}$} \\
& & $1$ & $3_{a}$ & $3_{b}$ & $6$ & $7$ \\ \hline
\endhead
1) & \possprim, \nongcr & 14 & 5 & 5 & 6 & 24 \\
2) & \possprim, \nongcr & 16 & 7 & 7 & 6 & 22 \\
3) & \possprim, \nongcr & 17 & 0 & 0 & 14 & 21 \\
4) & \nongcr & 19 & 2 & 2 & 14 & 19 \\
5) & \nongcr & 25 & 8 & 8 & 14 & 13 \\
6) & \nongcr & 30 & 21 & 21 & 6 & 8 \\
7) & \nongcr & 53 & 1 & 1 & 0 & 27 \\
8) &  & 79 & 27 & 27 & 0 & 1 \\
\end{longtable}
\noindent Modules tested: V248, V3875, V147250, L113243.\\

\setcounter{table}{268}
\begin{longtable}{rc|ccccccccc}
\caption{$L_2(8) < E_8$, $p = 0$ or $p \nmid |H|$}\\
& & \multicolumn{9}{c}{$V_{248}$} \\
& & $1$ & $7_{a}$ & $7_{b}$ & $7_{c}$ & $7_{d}$ & $8$ & $9_{a}$ & $9_{b}$ & $9_{c}$ \\ \hline
\endfirsthead
& & \multicolumn{9}{c}{$V_{248}$} \\
& & $1$ & $7_{a}$ & $7_{b}$ & $7_{c}$ & $7_{d}$ & $8$ & $9_{a}$ & $9_{b}$ & $9_{c}$ \\ \hline
\endhead
1) & \possprim & 0 & 3 & 5 & 5 & 5 & 4 & 4 & 3 & 3 \\
2) & \possprim & 0 & 6 & 5 & 4 & 3 & 4 & 4 & 3 & 3 \\
3) &  & 1 & 3 & 5 & 5 & 5 & 5 & 3 & 3 & 3 \\
4) &  & 1 & 4 & 5 & 5 & 5 & 3 & 4 & 3 & 3 \\
5) &  & 1 & 6 & 5 & 4 & 3 & 5 & 3 & 3 & 3 \\
6) &  & 2 & 2 & 8 & 6 & 4 & 2 & 4 & 3 & 3 \\
7) &  & 2 & 2 & 7 & 5 & 6 & 2 & 4 & 3 & 3 \\
8) &  & 2 & 3 & 5 & 5 & 5 & 6 & 5 & 2 & 1 \\
9) &  & 2 & 4 & 5 & 5 & 5 & 4 & 3 & 3 & 3 \\
10) &  & 2 & 5 & 8 & 5 & 2 & 2 & 4 & 3 & 3 \\
11) &  & 2 & 6 & 5 & 4 & 3 & 6 & 5 & 2 & 1 \\
12) &  & 3 & 2 & 8 & 6 & 4 & 3 & 3 & 3 & 3 \\
13) &  & 3 & 2 & 7 & 5 & 6 & 3 & 3 & 3 & 3 \\
14) &  & 3 & 4 & 5 & 5 & 5 & 5 & 5 & 2 & 1 \\
15) &  & 3 & 5 & 8 & 5 & 2 & 3 & 3 & 3 & 3 \\
16) &  & 3 & 6 & 5 & 5 & 5 & 1 & 4 & 3 & 3 \\
17) &  & 4 & 2 & 8 & 6 & 4 & 4 & 5 & 2 & 1 \\
18) &  & 4 & 2 & 7 & 5 & 6 & 4 & 5 & 2 & 1 \\
19) &  & 4 & 5 & 8 & 5 & 2 & 4 & 5 & 2 & 1 \\
20) &  & 4 & 6 & 5 & 5 & 5 & 2 & 3 & 3 & 3 \\
21) &  & 4 & 7 & 11 & 2 & 2 & 0 & 4 & 3 & 3 \\
22) &  & 5 & 0 & 1 & 1 & 1 & 3 & 8 & 7 & 7 \\
23) &  & 5 & 6 & 5 & 5 & 5 & 3 & 5 & 2 & 1 \\
24) &  & 5 & 7 & 11 & 2 & 2 & 1 & 3 & 3 & 3 \\
25) &  & 6 & 0 & 1 & 1 & 1 & 4 & 7 & 7 & 7 \\
26) &  & 6 & 5 & 13 & 5 & 0 & 0 & 3 & 3 & 3 \\
27) &  & 6 & 5 & 9 & 4 & 5 & 0 & 3 & 3 & 3 \\
28) &  & 6 & 7 & 11 & 2 & 2 & 2 & 5 & 2 & 1 \\
29) &  & 7 & 0 & 1 & 1 & 1 & 5 & 12 & 2 & 6 \\
30) &  & 7 & 0 & 1 & 1 & 1 & 5 & 9 & 6 & 5 \\
31) &  & 7 & 2 & 1 & 1 & 1 & 1 & 8 & 7 & 7 \\
32) &  & 7 & 5 & 13 & 5 & 0 & 1 & 5 & 2 & 1 \\
33) &  & 7 & 5 & 9 & 4 & 5 & 1 & 5 & 2 & 1 \\
34) &  & 8 & 2 & 1 & 1 & 1 & 2 & 7 & 7 & 7 \\
35) &  & 8 & 12 & 11 & 1 & 0 & 0 & 5 & 2 & 1 \\
36) &  & 9 & 2 & 1 & 1 & 1 & 3 & 12 & 2 & 6 \\
37) &  & 9 & 2 & 1 & 1 & 1 & 3 & 9 & 6 & 5 \\
38) &  & 10 & 1 & 5 & 0 & 1 & 0 & 7 & 7 & 7 \\
39) &  & 11 & 1 & 5 & 0 & 1 & 1 & 12 & 2 & 6 \\
40) &  & 11 & 1 & 5 & 0 & 1 & 1 & 9 & 6 & 5 \\
41) &  & 12 & 0 & 1 & 1 & 1 & 10 & 8 & 0 & 7 \\
42) &  & 12 & 0 & 1 & 1 & 1 & 10 & 14 & 0 & 1 \\
43) &  & 14 & 2 & 1 & 1 & 1 & 8 & 8 & 0 & 7 \\
44) &  & 14 & 2 & 1 & 1 & 1 & 8 & 14 & 0 & 1 \\
45) &  & 16 & 1 & 5 & 0 & 1 & 6 & 8 & 0 & 7 \\
46) &  & 16 & 1 & 5 & 0 & 1 & 6 & 14 & 0 & 1 \\
47) &  & 21 & 9 & 1 & 1 & 1 & 1 & 8 & 0 & 7 \\
48) &  & 21 & 9 & 1 & 1 & 1 & 1 & 14 & 0 & 1 \\
49) &  & 27 & 0 & 1 & 1 & 1 & 25 & 0 & 0 & 0 \\
50) &  & 29 & 2 & 1 & 1 & 1 & 23 & 0 & 0 & 0 \\
51) &  & 31 & 1 & 5 & 0 & 1 & 21 & 0 & 0 & 0 \\
52) &  & 36 & 9 & 1 & 1 & 1 & 16 & 0 & 0 & 0 \\
53) &  & 52 & 1 & 26 & 1 & 0 & 0 & 0 & 0 & 0 \\
\end{longtable}
\noindent Modules tested: V248, V3875, V147250.\\

\setcounter{table}{269}
\begin{longtable}{rc|cccccc}
\caption{$L_2(8) < E_8$, $p = 7$}\\
& & \multicolumn{6}{c}{$V_{248}$} \\
& & $1$ & $7_{a}$ & $7_{b}$ & $7_{c}$ & $7_{d}$ & $8$ \\ \hline
\endfirsthead
& & \multicolumn{6}{c}{$V_{248}$} \\
& & $1$ & $7_{a}$ & $7_{b}$ & $7_{c}$ & $7_{d}$ & $8$ \\ \hline
\endhead
1) & \possprim, \nongcr & 10 & 3 & 5 & 5 & 5 & 14 \\
2) & \possprim, \nongcr & 10 & 6 & 5 & 4 & 3 & 14 \\
3) & \possprim, \nongcr & 11 & 4 & 5 & 5 & 5 & 13 \\
4) & \nongcr & 12 & 2 & 8 & 6 & 4 & 12 \\
5) & \nongcr & 12 & 2 & 7 & 5 & 6 & 12 \\
6) & \nongcr & 12 & 5 & 8 & 5 & 2 & 12 \\
7) & \nongcr & 13 & 6 & 5 & 5 & 5 & 11 \\
8) & \nongcr & 14 & 7 & 11 & 2 & 2 & 10 \\
9) & \nongcr & 15 & 5 & 13 & 5 & 0 & 9 \\
10) & \nongcr & 15 & 5 & 9 & 4 & 5 & 9 \\
11) & \nongcr & 16 & 12 & 11 & 1 & 0 & 8 \\
12) & \nongcr & 20 & 13 & 5 & 5 & 5 & 4 \\
13) & \nongcr & 27 & 0 & 1 & 1 & 1 & 25 \\
14) & \nongcr & 29 & 2 & 1 & 1 & 1 & 23 \\
15) & \nongcr & 31 & 1 & 5 & 0 & 1 & 21 \\
16) & \nongcr & 36 & 9 & 1 & 1 & 1 & 16 \\
17) &  & 52 & 1 & 26 & 1 & 0 & 0 \\
\end{longtable}
\noindent Modules tested: V248, V3875, V147250, L147002.\\

\setcounter{table}{270}
\begin{longtable}{rc|ccccc}
\caption{$L_2(8) < E_8$, $p = 3$}\\
& & \multicolumn{5}{c}{$V_{248}$} \\
& & $1$ & $7$ & $9_{a}$ & $9_{b}$ & $9_{c}$ \\ \hline
\endfirsthead
& & \multicolumn{5}{c}{$V_{248}$} \\
& & $1$ & $7$ & $9_{a}$ & $9_{b}$ & $9_{c}$ \\ \hline
\endhead
1) & \possprim, \nongcr & 4 & 22 & 4 & 3 & 3 \\
2) & \possprim, \nongcr & 6 & 23 & 3 & 3 & 3 \\
3) & \nongcr & 8 & 6 & 8 & 7 & 7 \\
4) & \possprim, \nongcr & 8 & 24 & 5 & 2 & 1 \\
5) & \nongcr & 10 & 7 & 7 & 7 & 7 \\
6) & \nongcr & 12 & 8 & 12 & 2 & 6 \\
7) & \nongcr & 12 & 8 & 9 & 6 & 5 \\
8) & \nongcr & 22 & 13 & 14 & 0 & 1 \\
9) & \nongcr & 22 & 13 & 8 & 0 & 7 \\
10) & \nongcr & 52 & 28 & 0 & 0 & 0 \\
\end{longtable}
\noindent Modules tested: V248, V3875, V147250, L113243.\\

\setcounter{table}{271}
\begin{longtable}{rc|cccccccc}
\caption{$L_2(11) < E_8$, $p = 0$ or $p \nmid |H|$}\\
& & \multicolumn{8}{c}{$V_{248}$} \\
& & $1$ & $5_{a}$ & $5_{b}$ & $10_{a}$ & $10_{b}$ & $11$ & $12_{a}$ & $12_{b}$ \\ \hline
\endfirsthead
& & \multicolumn{8}{c}{$V_{248}$} \\
& & $1$ & $5_{a}$ & $5_{b}$ & $10_{a}$ & $10_{b}$ & $11$ & $12_{a}$ & $12_{b}$ \\ \hline
\endhead
1) & \possprim & 0 & 2 & 2 & 4 & 0 & 4 & 6 & 6 \\
2) &  & 1 & 0 & 0 & 5 & 2 & 3 & 6 & 6 \\
3) &  & 2 & 2 & 2 & 4 & 0 & 6 & 5 & 5 \\
4) &  & 3 & 0 & 0 & 5 & 2 & 5 & 5 & 5 \\
5) &  & 3 & 2 & 2 & 4 & 0 & 7 & 7 & 2 \\
6) &  & 4 & 0 & 0 & 5 & 2 & 6 & 7 & 2 \\
7) &  & 7 & 1 & 1 & 9 & 1 & 1 & 5 & 5 \\
8) &  & 7 & 1 & 1 & 1 & 9 & 1 & 5 & 5 \\
9) &  & 7 & 4 & 4 & 1 & 6 & 1 & 5 & 5 \\
10) &  & 8 & 1 & 1 & 9 & 1 & 2 & 7 & 2 \\
11) &  & 8 & 1 & 1 & 1 & 9 & 2 & 7 & 2 \\
12) &  & 8 & 2 & 2 & 2 & 8 & 0 & 5 & 5 \\
13) &  & 8 & 4 & 4 & 1 & 6 & 2 & 7 & 2 \\
14) &  & 8 & 8 & 8 & 2 & 2 & 0 & 5 & 5 \\
15) &  & 9 & 2 & 2 & 2 & 8 & 1 & 7 & 2 \\
16) &  & 9 & 8 & 8 & 2 & 2 & 1 & 7 & 2 \\
17) &  & 10 & 2 & 2 & 4 & 0 & 14 & 1 & 1 \\
18) &  & 11 & 0 & 0 & 5 & 2 & 13 & 1 & 1 \\
19) &  & 15 & 1 & 1 & 9 & 1 & 9 & 1 & 1 \\
20) &  & 15 & 1 & 1 & 1 & 9 & 9 & 1 & 1 \\
21) &  & 15 & 4 & 4 & 1 & 6 & 9 & 1 & 1 \\
22) &  & 16 & 2 & 2 & 2 & 8 & 8 & 1 & 1 \\
23) &  & 16 & 8 & 8 & 2 & 2 & 8 & 1 & 1 \\
24) &  & 24 & 10 & 10 & 10 & 0 & 0 & 1 & 1 \\
\end{longtable}
\noindent Modules tested: V248, V3875, V147250.\\

\setcounter{table}{272}
\begin{longtable}{rc|cccccc}
\caption{$L_2(11) < E_8$, $p = 5$}\\
& & \multicolumn{6}{c}{$V_{248}$} \\
& & $1$ & $5_{a}$ & $5_{b}$ & $10_{a}$ & $10_{b}$ & $11$ \\ \hline
\endfirsthead
& & \multicolumn{6}{c}{$V_{248}$} \\
& & $1$ & $5_{a}$ & $5_{b}$ & $10_{a}$ & $10_{b}$ & $11$ \\ \hline
\endhead
1) & \possprim, \nongcr & 12 & 2 & 2 & 4 & 0 & 16 \\
2) & \possprim, \nongcr & 13 & 0 & 0 & 5 & 2 & 15 \\
3) & \nongcr & 17 & 1 & 1 & 9 & 1 & 11 \\
4) & \nongcr & 17 & 1 & 1 & 1 & 9 & 11 \\
5) & \nongcr & 17 & 4 & 4 & 1 & 6 & 11 \\
6) & \nongcr & 18 & 2 & 2 & 2 & 8 & 10 \\
7) & \nongcr & 18 & 8 & 8 & 2 & 2 & 10 \\
8) & \nongcr & 26 & 10 & 10 & 10 & 0 & 2 \\
\end{longtable}
\noindent Modules tested: V248, V3875, V147250.\\

\setcounter{table}{273}
\begin{longtable}{rc|cccccc}
\caption{$L_2(11) < E_8$, $p = 3$}\\
& & \multicolumn{6}{c}{$V_{248}$} \\
& & $1$ & $5_{a}$ & $5_{b}$ & $10$ & $12_{a}$ & $12_{b}$ \\ \hline
\endfirsthead
& & \multicolumn{6}{c}{$V_{248}$} \\
& & $1$ & $5_{a}$ & $5_{b}$ & $10$ & $12_{a}$ & $12_{b}$ \\ \hline
\endhead
1) & \possprim, \nongcr & 4 & 2 & 2 & 8 & 6 & 6 \\
2) &  & 4 & 10 & 10 & 0 & 6 & 6 \\
3) & \possprim, \nongcr & 8 & 2 & 2 & 10 & 5 & 5 \\
4) & \nongcr & 8 & 10 & 10 & 2 & 5 & 5 \\
5) & \possprim, \nongcr & 10 & 2 & 2 & 11 & 7 & 2 \\
6) & \nongcr & 10 & 10 & 10 & 3 & 7 & 2 \\
7) & \nongcr & 24 & 2 & 2 & 18 & 1 & 1 \\
8) & \nongcr & 24 & 10 & 10 & 10 & 1 & 1 \\
\end{longtable}
\noindent Modules tested: V248, V3875, V147250, L113243.\\

\setcounter{table}{274}
\begin{longtable}{rc|cccccc}
\caption{$L_2(11) < E_8$, $p = 2$}\\
& & \multicolumn{6}{c}{$V_{248}$} \\
& & $1$ & $5_{a}$ & $5_{b}$ & $10$ & $12_{a}$ & $12_{b}$ \\ \hline
\endfirsthead
& & \multicolumn{6}{c}{$V_{248}$} \\
& & $1$ & $5_{a}$ & $5_{b}$ & $10$ & $12_{a}$ & $12_{b}$ \\ \hline
\endhead
1) &  & 4 & 0 & 0 & 10 & 6 & 6 \\
2) & \possprim, \nongcr & 4 & 3 & 3 & 7 & 6 & 6 \\
3) & \possprim, \nongcr & 4 & 6 & 6 & 4 & 6 & 6 \\
4) & \nongcr & 8 & 2 & 2 & 10 & 5 & 5 \\
5) & \possprim, \nongcr & 8 & 5 & 5 & 7 & 5 & 5 \\
6) & \possprim, \nongcr & 8 & 8 & 8 & 4 & 5 & 5 \\
7) & \nongcr & 10 & 3 & 3 & 10 & 7 & 2 \\
8) & \possprim, \nongcr & 10 & 6 & 6 & 7 & 7 & 2 \\
9) & \possprim, \nongcr & 10 & 9 & 9 & 4 & 7 & 2 \\
10) & \nongcr & 24 & 10 & 10 & 10 & 1 & 1 \\
11) & \possprim, \nongcr & 24 & 13 & 13 & 7 & 1 & 1 \\
12) & \possprim, \nongcr & 24 & 16 & 16 & 4 & 1 & 1 \\
\end{longtable}
\noindent Modules tested: V248, V3875, L3626, V147250, L143376.\\

\setcounter{table}{275}
\begin{longtable}{rc|ccccccccc}
\caption{$L_2(13) < E_8$, $p = 0$ or $p \nmid |H|$}\\
& & \multicolumn{9}{c}{$V_{248}$} \\
& & $1$ & $7_{a}$ & $7_{b}$ & $12_{a}$ & $12_{b}$ & $12_{c}$ & $13$ & $14_{a}$ & $14_{b}$ \\ \hline
\endfirsthead
& & \multicolumn{9}{c}{$V_{248}$} \\
& & $1$ & $7_{a}$ & $7_{b}$ & $12_{a}$ & $12_{b}$ & $12_{c}$ & $13$ & $14_{a}$ & $14_{b}$ \\ \hline
\endhead
1) & \possprim & 0 & 1 & 1 & 3 & 3 & 2 & 2 & 2 & 6 \\
2) &  & 1 & 1 & 1 & 3 & 3 & 3 & 1 & 2 & 6 \\
3) &  & 1 & 3 & 3 & 3 & 3 & 2 & 3 & 1 & 4 \\
4) &  & 2 & 1 & 1 & 5 & 1 & 4 & 0 & 2 & 6 \\
5) &  & 2 & 2 & 2 & 1 & 1 & 0 & 2 & 4 & 8 \\
6) &  & 2 & 3 & 3 & 3 & 3 & 3 & 2 & 1 & 4 \\
7) &  & 2 & 5 & 5 & 3 & 3 & 2 & 4 & 0 & 2 \\
8) &  & 3 & 2 & 2 & 1 & 1 & 1 & 1 & 4 & 8 \\
9) &  & 3 & 3 & 3 & 5 & 1 & 4 & 1 & 1 & 4 \\
10) &  & 3 & 4 & 4 & 1 & 1 & 0 & 3 & 3 & 6 \\
11) &  & 3 & 5 & 5 & 3 & 3 & 3 & 3 & 0 & 2 \\
12) &  & 4 & 0 & 0 & 1 & 1 & 0 & 4 & 10 & 2 \\
13) &  & 4 & 4 & 4 & 1 & 1 & 1 & 2 & 3 & 6 \\
14) &  & 4 & 5 & 5 & 5 & 1 & 4 & 2 & 0 & 2 \\
15) &  & 4 & 6 & 6 & 1 & 1 & 0 & 4 & 2 & 4 \\
16) &  & 5 & 0 & 0 & 1 & 1 & 1 & 3 & 10 & 2 \\
17) &  & 5 & 2 & 2 & 1 & 1 & 0 & 5 & 9 & 0 \\
18) &  & 5 & 6 & 6 & 1 & 1 & 1 & 3 & 2 & 4 \\
19) &  & 6 & 13 & 0 & 0 & 0 & 0 & 3 & 3 & 5 \\
20) &  & 6 & 2 & 2 & 1 & 1 & 1 & 4 & 9 & 0 \\
21) &  & 7 & 0 & 0 & 1 & 1 & 0 & 7 & 7 & 2 \\
22) &  & 8 & 0 & 0 & 1 & 1 & 1 & 6 & 7 & 2 \\
23) &  & 8 & 2 & 2 & 1 & 1 & 0 & 8 & 6 & 0 \\
24) &  & 9 & 2 & 2 & 1 & 1 & 1 & 7 & 6 & 0 \\
25) &  & 14 & 0 & 0 & 8 & 1 & 0 & 0 & 7 & 2 \\
26) &  & 15 & 2 & 2 & 8 & 1 & 0 & 1 & 6 & 0 \\
27) &  & 52 & 26 & 0 & 0 & 0 & 0 & 0 & 0 & 1 \\
\end{longtable}
\noindent Modules tested: V248, V3875, V147250.\\

\setcounter{table}{276}
\begin{longtable}{rc|cccccc}
\caption{$L_2(13) < E_8$, $p = 7$}\\
& & \multicolumn{6}{c}{$V_{248}$} \\
& & $1$ & $7_{a}$ & $7_{b}$ & $12$ & $14_{a}$ & $14_{b}$ \\ \hline
\endfirsthead
& & \multicolumn{6}{c}{$V_{248}$} \\
& & $1$ & $7_{a}$ & $7_{b}$ & $12$ & $14_{a}$ & $14_{b}$ \\ \hline
\endhead
1) & \possprim, \nongcr & 2 & 1 & 1 & 10 & 2 & 6 \\
2) & \nongcr & 4 & 2 & 2 & 4 & 4 & 8 \\
3) & \possprim, \nongcr & 4 & 3 & 3 & 11 & 1 & 4 \\
4) & \nongcr & 6 & 4 & 4 & 5 & 3 & 6 \\
5) & \possprim, \nongcr & 6 & 5 & 5 & 12 & 0 & 2 \\
6) & \nongcr & 8 & 0 & 0 & 6 & 10 & 2 \\
7) & \nongcr & 8 & 6 & 6 & 6 & 2 & 4 \\
8) & \nongcr & 9 & 13 & 0 & 3 & 3 & 5 \\
9) & \nongcr & 10 & 2 & 2 & 7 & 9 & 0 \\
10) & \nongcr & 14 & 0 & 0 & 9 & 7 & 2 \\
11) & \nongcr & 16 & 2 & 2 & 10 & 6 & 0 \\
12) &  & 52 & 26 & 0 & 0 & 0 & 1 \\
\end{longtable}
\noindent Modules tested: V248, V3875, V147250, L147002.\\

\setcounter{table}{277}
\begin{longtable}{rc|ccccccc}
\caption{$L_2(13) < E_8$, $p = 3$}\\
& & \multicolumn{7}{c}{$V_{248}$} \\
& & $1$ & $7_{a}$ & $7_{b}$ & $12_{a}$ & $12_{b}$ & $12_{c}$ & $13$ \\ \hline
\endfirsthead
& & \multicolumn{7}{c}{$V_{248}$} \\
& & $1$ & $7_{a}$ & $7_{b}$ & $12_{a}$ & $12_{b}$ & $12_{c}$ & $13$ \\ \hline
\endhead
1) & \possprim, \nongcr & 2 & 7 & 7 & 3 & 3 & 2 & 4 \\
2) & \nongcr & 3 & 7 & 7 & 3 & 3 & 3 & 3 \\
3) & \nongcr & 4 & 7 & 7 & 5 & 1 & 4 & 2 \\
4) & \nongcr & 6 & 10 & 10 & 1 & 1 & 0 & 6 \\
5) & \nongcr & 7 & 10 & 10 & 1 & 1 & 1 & 5 \\
6) & \nongcr & 9 & 18 & 5 & 0 & 0 & 0 & 6 \\
7) & \nongcr & 14 & 2 & 2 & 1 & 1 & 0 & 14 \\
8) & \nongcr & 15 & 2 & 2 & 1 & 1 & 1 & 13 \\
9) & \nongcr & 21 & 2 & 2 & 8 & 1 & 0 & 7 \\
10) &  & 52 & 27 & 1 & 0 & 0 & 0 & 0 \\
\end{longtable}
\noindent Modules tested: V248, V3875, V147250, L113243.\\

\setcounter{table}{278}
\begin{longtable}{rc|ccccccc}
\caption{$L_2(13) < E_8$, $p = 2$}\\
& & \multicolumn{7}{c}{$V_{248}$} \\
& & $1$ & $6_{a}$ & $6_{b}$ & $12_{a}$ & $12_{b}$ & $12_{c}$ & $14$ \\ \hline
\endfirsthead
& & \multicolumn{7}{c}{$V_{248}$} \\
& & $1$ & $6_{a}$ & $6_{b}$ & $12_{a}$ & $12_{b}$ & $12_{c}$ & $14$ \\ \hline
\endhead
1) &  & 4 & 1 & 1 & 5 & 1 & 4 & 8 \\
2) & \nongcr & 4 & 2 & 2 & 3 & 3 & 3 & 8 \\
3) & \possprim, \nongcr & 4 & 3 & 3 & 3 & 3 & 2 & 8 \\
4) & \nongcr & 8 & 3 & 3 & 1 & 1 & 1 & 12 \\
5) & \nongcr & 8 & 4 & 4 & 1 & 1 & 0 & 12 \\
6) & \nongcr & 10 & 4 & 4 & 5 & 1 & 4 & 5 \\
7) & \nongcr & 10 & 5 & 5 & 3 & 3 & 3 & 5 \\
8) & \possprim, \nongcr & 10 & 6 & 6 & 3 & 3 & 2 & 5 \\
9) &  & 14 & 0 & 0 & 8 & 1 & 0 & 9 \\
10) & \nongcr & 14 & 6 & 6 & 1 & 1 & 1 & 9 \\
11) & \nongcr & 14 & 7 & 7 & 1 & 1 & 0 & 9 \\
12) & \nongcr & 16 & 13 & 0 & 0 & 0 & 0 & 11 \\
13) & \nongcr & 16 & 2 & 2 & 10 & 3 & 2 & 2 \\
14) & \nongcr & 16 & 7 & 7 & 5 & 1 & 4 & 2 \\
15) & \nongcr & 16 & 8 & 8 & 3 & 3 & 3 & 2 \\
16) & \possprim, \nongcr & 16 & 9 & 9 & 3 & 3 & 2 & 2 \\
17) & \nongcr & 20 & 3 & 3 & 8 & 1 & 0 & 6 \\
18) & \nongcr & 20 & 9 & 9 & 1 & 1 & 1 & 6 \\
19) & \nongcr & 20 & 10 & 10 & 1 & 1 & 0 & 6 \\
20) & \nongcr & 22 & 16 & 3 & 0 & 0 & 0 & 8 \\
21) & \nongcr & 78 & 26 & 0 & 0 & 0 & 0 & 1 \\
\end{longtable}
\noindent Modules tested: V248, V3875, L3626, V147250, L143376.\\

\setcounter{table}{279}
\begin{longtable}{rc|ccccccccccccccccc}
\caption{$L_2(16) < E_8$, $p = 0$ or $p \nmid |H|$}\\
& & \multicolumn{17}{c}{$V_{248}$} \\
& & $1$ & $15_{a}$ & $15_{b}$ & $15_{c}$ & $15_{d}$ & $15_{e}$ & $15_{f}$ & $15_{g}$ & $15_{h}$ & $16$ & $17_{a}$ & $17_{b}$ & $17_{c}$ & $17_{d}$ & $17_{e}$ & $17_{f}$ & $17_{g}$ \\ \hline
\endfirsthead
& & \multicolumn{17}{c}{$V_{248}$} \\
& & $1$ & $15_{a}$ & $15_{b}$ & $15_{c}$ & $15_{d}$ & $15_{e}$ & $15_{f}$ & $15_{g}$ & $15_{h}$ & $16$ & $17_{a}$ & $17_{b}$ & $17_{c}$ & $17_{d}$ & $17_{e}$ & $17_{f}$ & $17_{g}$ \\ \hline
\endhead
1) & \possprim & 0 & 3 & 1 & 1 & 0 & 2 & 2 & 2 & 1 & 0 & 0 & 0 & 0 & 1 & 1 & 1 & 1 \\
2) & \possprim & 0 & 2 & 2 & 2 & 2 & 1 & 1 & 1 & 1 & 0 & 0 & 0 & 0 & 1 & 1 & 1 & 1 \\
3) & \possprim & 0 & 2 & 2 & 1 & 1 & 2 & 1 & 1 & 2 & 0 & 0 & 0 & 0 & 1 & 1 & 1 & 1 \\
4) &  & 1 & 3 & 1 & 1 & 0 & 2 & 2 & 2 & 1 & 1 & 1 & 1 & 1 & 0 & 0 & 0 & 0 \\
5) &  & 1 & 2 & 2 & 2 & 2 & 1 & 1 & 1 & 1 & 1 & 1 & 1 & 1 & 0 & 0 & 0 & 0 \\
6) &  & 1 & 2 & 2 & 1 & 1 & 2 & 1 & 1 & 2 & 1 & 1 & 1 & 1 & 0 & 0 & 0 & 0 \\
7) &  & 2 & 4 & 3 & 1 & 0 & 0 & 1 & 3 & 1 & 0 & 1 & 1 & 1 & 0 & 0 & 0 & 0 \\
8) &  & 2 & 3 & 1 & 1 & 2 & 2 & 0 & 3 & 1 & 0 & 1 & 1 & 1 & 0 & 0 & 0 & 0 \\
\end{longtable}
\noindent Modules tested: V248, V3875, V147250.\\

\setcounter{table}{280}
\begin{longtable}{rc|ccccccccc}
\caption{$L_2(16) < E_8$, $p = 17$}\\
& & \multicolumn{9}{c}{$V_{248}$} \\
& & $1$ & $15$ & $17_{a}$ & $17_{b}$ & $17_{c}$ & $17_{d}$ & $17_{e}$ & $17_{f}$ & $17_{g}$ \\ \hline
\endfirsthead
& & \multicolumn{9}{c}{$V_{248}$} \\
& & $1$ & $15$ & $17_{a}$ & $17_{b}$ & $17_{c}$ & $17_{d}$ & $17_{e}$ & $17_{f}$ & $17_{g}$ \\ \hline
\endhead
1) & \possprim & 0 & 12 & 0 & 0 & 0 & 1 & 1 & 1 & 1 \\
2) & \possprim & 2 & 13 & 1 & 1 & 1 & 0 & 0 & 0 & 0 \\
3) &  & 10 & 0 & 0 & 3 & 3 & 4 & 0 & 4 & 0 \\
4) &  & 10 & 0 & 1 & 6 & 0 & 4 & 0 & 0 & 3 \\
5) &  & 10 & 0 & 1 & 2 & 1 & 4 & 4 & 1 & 1 \\
6) &  & 10 & 0 & 8 & 0 & 0 & 3 & 0 & 2 & 1 \\
7) &  & 10 & 0 & 8 & 2 & 1 & 2 & 0 & 0 & 1 \\
8) &  & 10 & 0 & 8 & 3 & 3 & 0 & 0 & 0 & 0 \\
9) &  & 12 & 1 & 7 & 1 & 1 & 1 & 1 & 1 & 1 \\
\end{longtable}
\noindent Modules tested: V248, V3875, V147250.\\

\setcounter{table}{281}
\begin{longtable}{rc|ccccccccccc}
\caption{$L_2(16) < E_8$, $p = 5$}\\
& & \multicolumn{11}{c}{$V_{248}$} \\
& & $1$ & $15_{a}$ & $15_{b}$ & $15_{c}$ & $15_{d}$ & $15_{e}$ & $15_{f}$ & $15_{g}$ & $15_{h}$ & $16$ & $17$ \\ \hline
\endfirsthead
& & \multicolumn{11}{c}{$V_{248}$} \\
& & $1$ & $15_{a}$ & $15_{b}$ & $15_{c}$ & $15_{d}$ & $15_{e}$ & $15_{f}$ & $15_{g}$ & $15_{h}$ & $16$ & $17$ \\ \hline
\endhead
1) & \possprim & 0 & 2 & 2 & 2 & 2 & 1 & 1 & 1 & 1 & 0 & 4 \\
2) & \possprim & 0 & 2 & 2 & 1 & 1 & 2 & 1 & 1 & 2 & 0 & 4 \\
3) & \possprim & 0 & 3 & 1 & 1 & 0 & 2 & 2 & 2 & 1 & 0 & 4 \\
4) &  & 3 & 2 & 2 & 2 & 2 & 1 & 1 & 1 & 1 & 3 & 1 \\
5) &  & 3 & 2 & 2 & 1 & 1 & 2 & 1 & 1 & 2 & 3 & 1 \\
6) &  & 3 & 3 & 1 & 1 & 0 & 2 & 2 & 2 & 1 & 3 & 1 \\
7) &  & 4 & 3 & 1 & 1 & 2 & 2 & 0 & 3 & 1 & 2 & 1 \\
8) &  & 4 & 4 & 3 & 1 & 0 & 0 & 1 & 3 & 1 & 2 & 1 \\
9) &  & 5 & 4 & 4 & 0 & 0 & 2 & 2 & 1 & 1 & 1 & 1 \\
10) &  & 5 & 5 & 1 & 3 & 0 & 3 & 1 & 0 & 1 & 1 & 1 \\
11) &  & 5 & 3 & 2 & 2 & 0 & 2 & 2 & 3 & 0 & 1 & 1 \\
\end{longtable}
\noindent Modules tested: V248, V3875, V147250.\\

\setcounter{table}{282}
\begin{longtable}{rc|cccccccccccc}
\caption{$L_2(16) < E_8$, $p = 3$}\\
& & \multicolumn{12}{c}{$V_{248}$} \\
& & $1$ & $15_{a}$ & $15_{b}$ & $15_{c}$ & $15_{d}$ & $15_{e}$ & $15_{f}$ & $15_{g}$ & $15_{h}$ & $16$ & $17_{a}$ & $17_{b}$ \\ \hline
\endfirsthead
& & \multicolumn{12}{c}{$V_{248}$} \\
& & $1$ & $15_{a}$ & $15_{b}$ & $15_{c}$ & $15_{d}$ & $15_{e}$ & $15_{f}$ & $15_{g}$ & $15_{h}$ & $16$ & $17_{a}$ & $17_{b}$ \\ \hline
\endhead
1) & \possprim & 0 & 3 & 1 & 1 & 0 & 2 & 2 & 2 & 1 & 0 & 2 & 2 \\
2) & \possprim & 0 & 2 & 2 & 2 & 2 & 1 & 1 & 1 & 1 & 0 & 2 & 2 \\
3) & \possprim & 0 & 2 & 2 & 1 & 1 & 2 & 1 & 1 & 2 & 0 & 2 & 2 \\
4) &  & 2 & 3 & 1 & 1 & 0 & 2 & 2 & 2 & 1 & 2 & 1 & 1 \\
5) &  & 2 & 2 & 2 & 2 & 2 & 1 & 1 & 1 & 1 & 2 & 1 & 1 \\
6) &  & 2 & 2 & 2 & 1 & 1 & 2 & 1 & 1 & 2 & 2 & 1 & 1 \\
7) &  & 3 & 4 & 3 & 1 & 0 & 0 & 1 & 3 & 1 & 1 & 1 & 1 \\
8) &  & 3 & 3 & 1 & 1 & 2 & 2 & 0 & 3 & 1 & 1 & 1 & 1 \\
9) &  & 4 & 3 & 2 & 2 & 0 & 2 & 2 & 3 & 0 & 0 & 1 & 1 \\
10) &  & 4 & 5 & 1 & 3 & 0 & 3 & 1 & 0 & 1 & 0 & 1 & 1 \\
11) &  & 4 & 4 & 4 & 0 & 0 & 2 & 2 & 1 & 1 & 0 & 1 & 1 \\
\end{longtable}
\noindent Modules tested: V248, V3875, V147250, L113243.\\

\setcounter{table}{283}
\begin{longtable}{rc|ccccccccccc}
\caption{$L_2(17) < E_8$, $p = 0$ or $p \nmid |H|$}\\
& & \multicolumn{11}{c}{$V_{248}$} \\
& & $1$ & $9_{a}$ & $9_{b}$ & $16_{a}$ & $16_{b}$ & $16_{c}$ & $16_{d}$ & $17$ & $18_{a}$ & $18_{b}$ & $18_{c}$ \\ \hline
\endfirsthead
& & \multicolumn{11}{c}{$V_{248}$} \\
& & $1$ & $9_{a}$ & $9_{b}$ & $16_{a}$ & $16_{b}$ & $16_{c}$ & $16_{d}$ & $17$ & $18_{a}$ & $18_{b}$ & $18_{c}$ \\ \hline
\endhead
1) & \possprim & 0 & 1 & 0 & 0 & 2 & 2 & 2 & 1 & 1 & 3 & 3 \\
2) & \possprim & 0 & 1 & 0 & 3 & 2 & 1 & 0 & 1 & 1 & 3 & 3 \\
3) &  & 1 & 1 & 0 & 1 & 2 & 2 & 2 & 0 & 1 & 3 & 3 \\
4) &  & 7 & 8 & 1 & 1 & 0 & 0 & 0 & 0 & 6 & 1 & 1 \\
5) &  & 8 & 7 & 0 & 1 & 0 & 0 & 0 & 1 & 6 & 1 & 1 \\
6) &  & 10 & 11 & 4 & 1 & 0 & 0 & 0 & 3 & 0 & 1 & 1 \\
7) &  & 14 & 7 & 0 & 1 & 0 & 0 & 0 & 7 & 0 & 1 & 1 \\
8) &  & 21 & 7 & 0 & 8 & 0 & 0 & 0 & 0 & 0 & 1 & 1 \\
\end{longtable}
\noindent Modules tested: V248, V3875, V147250.\\

\setcounter{table}{284}
\begin{longtable}{rc|ccccccc}
\caption{$L_2(17) < E_8$, $p = 3$}\\
& & \multicolumn{7}{c}{$V_{248}$} \\
& & $1$ & $9_{a}$ & $9_{b}$ & $16$ & $18_{a}$ & $18_{b}$ & $18_{c}$ \\ \hline
\endfirsthead
& & \multicolumn{7}{c}{$V_{248}$} \\
& & $1$ & $9_{a}$ & $9_{b}$ & $16$ & $18_{a}$ & $18_{b}$ & $18_{c}$ \\ \hline
\endhead
1) & \possprim, \nongcr & 1 & 1 & 0 & 7 & 1 & 3 & 3 \\
2) &  & 7 & 8 & 1 & 1 & 6 & 1 & 1 \\
3) & \nongcr & 9 & 7 & 0 & 2 & 6 & 1 & 1 \\
4) & \nongcr & 13 & 11 & 4 & 4 & 0 & 1 & 1 \\
5) & \nongcr & 21 & 7 & 0 & 8 & 0 & 1 & 1 \\
\end{longtable}
\noindent Modules tested: V248, V3875, V147250, L113243.\\

\setcounter{table}{285}
\begin{longtable}{rc|ccccccc}
\caption{$L_2(17) < E_8$, $p = 2$}\\
& & \multicolumn{7}{c}{$V_{248}$} \\
& & $1$ & $8_{a}$ & $8_{b}$ & $16_{a}$ & $16_{b}$ & $16_{c}$ & $16_{d}$ \\ \hline
\endfirsthead
& & \multicolumn{7}{c}{$V_{248}$} \\
& & $1$ & $8_{a}$ & $8_{b}$ & $16_{a}$ & $16_{b}$ & $16_{c}$ & $16_{d}$ \\ \hline
\endhead
1) & \nongcr & 16 & 4 & 3 & 2 & 6 & 1 & 2 \\
2) & \nongcr & 16 & 6 & 5 & 3 & 2 & 2 & 2 \\
3) & \nongcr & 16 & 8 & 7 & 1 & 2 & 2 & 2 \\
4) & \possprim, \nongcr & 16 & 9 & 8 & 0 & 2 & 2 & 2 \\
5) & \possprim, \nongcr & 16 & 9 & 8 & 3 & 2 & 1 & 0 \\
6) & \nongcr & 32 & 9 & 2 & 8 & 0 & 0 & 0 \\
7) & \nongcr & 32 & 16 & 9 & 1 & 0 & 0 & 0 \\
\end{longtable}
\noindent Modules tested: V248, V3875, L3626, V147250, L143376.\\

\setcounter{table}{286}
\begin{longtable}{rc|ccccccccccc}
\caption{$L_2(19) < E_8$, $p = 0$ or $p \nmid |H|$}\\
& & \multicolumn{11}{c}{$V_{248}$} \\
& & $1$ & $9_{a}$ & $9_{b}$ & $18_{a}$ & $18_{b}$ & $18_{d}$ & $19$ & $20_{a}$ & $20_{b}$ & $20_{c}$ & $20_{d}$ \\ \hline
\endfirsthead
& & \multicolumn{11}{c}{$V_{248}$} \\
& & $1$ & $9_{a}$ & $9_{b}$ & $18_{a}$ & $18_{b}$ & $18_{d}$ & $19$ & $20_{a}$ & $20_{b}$ & $20_{c}$ & $20_{d}$ \\ \hline
\endhead
1) & \possprim & 0 & 1 & 1 & 2 & 2 & 0 & 2 & 0 & 3 & 1 & 2 \\
2) & \possprim & 0 & 1 & 1 & 2 & 2 & 0 & 2 & 3 & 1 & 1 & 1 \\
3) &  & 1 & 0 & 0 & 3 & 2 & 1 & 1 & 0 & 3 & 1 & 2 \\
4) &  & 1 & 0 & 0 & 3 & 2 & 1 & 1 & 3 & 1 & 1 & 1 \\
5) &  & 1 & 1 & 1 & 2 & 2 & 0 & 3 & 2 & 1 & 1 & 1 \\
6) &  & 2 & 0 & 0 & 3 & 2 & 1 & 2 & 2 & 1 & 1 & 1 \\
7) &  & 3 & 1 & 1 & 2 & 2 & 0 & 5 & 0 & 1 & 1 & 1 \\
8) &  & 4 & 0 & 0 & 3 & 2 & 1 & 4 & 0 & 1 & 1 & 1 \\
9) &  & 8 & 3 & 3 & 1 & 0 & 6 & 0 & 0 & 1 & 1 & 1 \\
\end{longtable}
\noindent Modules tested: V248, V3875, V147250.\\

\setcounter{table}{287}
\begin{longtable}{rc|cccccccc}
\caption{$L_2(19) < E_8$, $p = 5$}\\
& & \multicolumn{8}{c}{$V_{248}$} \\
& & $1$ & $9_{a}$ & $9_{b}$ & $18$ & $20_{a}$ & $20_{b}$ & $20_{c}$ & $20_{d}$ \\ \hline
\endfirsthead
& & \multicolumn{8}{c}{$V_{248}$} \\
& & $1$ & $9_{a}$ & $9_{b}$ & $18$ & $20_{a}$ & $20_{b}$ & $20_{c}$ & $20_{d}$ \\ \hline
\endhead
1) & \possprim, \nongcr & 2 & 1 & 1 & 6 & 0 & 3 & 1 & 2 \\
2) & \possprim, \nongcr & 2 & 1 & 1 & 6 & 3 & 1 & 1 & 1 \\
3) & \possprim, \nongcr & 4 & 1 & 1 & 7 & 2 & 1 & 1 & 1 \\
4) & \possprim, \nongcr & 8 & 1 & 1 & 9 & 0 & 1 & 1 & 1 \\
5) & \nongcr & 8 & 9 & 9 & 1 & 0 & 1 & 1 & 1 \\
\end{longtable}
\noindent Modules tested: V248, V3875, V147250.\\

\setcounter{table}{288}
\begin{longtable}{rc|ccccccc}
\caption{$L_2(19) < E_8$, $p = 3$}\\
& & \multicolumn{7}{c}{$V_{248}$} \\
& & $1$ & $9_{a}$ & $9_{b}$ & $18_{a}$ & $18_{b}$ & $18_{d}$ & $19$ \\ \hline
\endfirsthead
& & \multicolumn{7}{c}{$V_{248}$} \\
& & $1$ & $9_{a}$ & $9_{b}$ & $18_{a}$ & $18_{b}$ & $18_{d}$ & $19$ \\ \hline
\endhead
1) & \possprim, \nongcr & 6 & 1 & 1 & 2 & 2 & 0 & 8 \\
2) & \nongcr & 7 & 0 & 0 & 3 & 2 & 1 & 7 \\
3) & \nongcr & 11 & 3 & 3 & 1 & 0 & 6 & 3 \\
\end{longtable}
\noindent Modules tested: V248, V3875, V147250, L113243.\\

\setcounter{table}{289}
\begin{longtable}{rc|ccccccccc}
\caption{$L_2(19) < E_8$, $p = 2$}\\
& & \multicolumn{9}{c}{$V_{248}$} \\
& & $1$ & $9_{a}$ & $9_{b}$ & $18_{a}$ & $18_{b}$ & $20_{a}$ & $20_{b}$ & $20_{c}$ & $20_{d}$ \\ \hline
\endfirsthead
& & \multicolumn{9}{c}{$V_{248}$} \\
& & $1$ & $9_{a}$ & $9_{b}$ & $18_{a}$ & $18_{b}$ & $20_{a}$ & $20_{b}$ & $20_{c}$ & $20_{d}$ \\ \hline
\endhead
1) &  & 2 & 0 & 0 & 6 & 1 & 0 & 3 & 1 & 2 \\
2) &  & 2 & 0 & 0 & 6 & 1 & 3 & 1 & 1 & 1 \\
3) & \nongcr & 2 & 1 & 1 & 3 & 3 & 0 & 3 & 1 & 2 \\
4) & \nongcr & 2 & 1 & 1 & 3 & 3 & 3 & 1 & 1 & 1 \\
5) & \possprim, \nongcr & 2 & 3 & 3 & 2 & 2 & 0 & 3 & 1 & 2 \\
6) & \possprim, \nongcr & 2 & 3 & 3 & 2 & 2 & 3 & 1 & 1 & 1 \\
7) & \nongcr & 4 & 1 & 1 & 6 & 1 & 2 & 1 & 1 & 1 \\
8) & \nongcr & 4 & 2 & 2 & 3 & 3 & 2 & 1 & 1 & 1 \\
9) & \possprim, \nongcr & 4 & 4 & 4 & 2 & 2 & 2 & 1 & 1 & 1 \\
10) & \nongcr & 8 & 3 & 3 & 6 & 1 & 0 & 1 & 1 & 1 \\
11) & \nongcr & 8 & 4 & 4 & 3 & 3 & 0 & 1 & 1 & 1 \\
12) & \possprim, \nongcr & 8 & 6 & 6 & 2 & 2 & 0 & 1 & 1 & 1 \\
\end{longtable}
\noindent Modules tested: V248, V3875, L3626, V147250, L143376.\\

\setcounter{table}{290}
\begin{longtable}{rc|ccccccccccccc}
\caption{$L_2(25) < E_8$, $p = 0$ or $p \nmid |H|$}\\
& & \multicolumn{13}{c}{$V_{248}$} \\
& & $1$ & $24_{a}$ & $24_{b}$ & $24_{c}$ & $24_{d}$ & $24_{e}$ & $24_{f}$ & $25$ & $26_{a}$ & $26_{b}$ & $26_{c}$ & $26_{d}$ & $26_{e}$ \\ \hline
\endfirsthead
& & \multicolumn{13}{c}{$V_{248}$} \\
& & $1$ & $24_{a}$ & $24_{b}$ & $24_{c}$ & $24_{d}$ & $24_{e}$ & $24_{f}$ & $25$ & $26_{a}$ & $26_{b}$ & $26_{c}$ & $26_{d}$ & $26_{e}$ \\ \hline
\endhead
1) & \possprim & 0 & 2 & 1 & 1 & 0 & 1 & 1 & 0 & 0 & 0 & 0 & 2 & 2 \\
2) & \possprim & 0 & 1 & 1 & 1 & 1 & 1 & 1 & 0 & 0 & 0 & 0 & 2 & 2 \\
3) &  & 14 & 0 & 0 & 0 & 0 & 0 & 0 & 0 & 7 & 0 & 0 & 1 & 1 \\
4) &  & 15 & 0 & 0 & 0 & 0 & 0 & 0 & 1 & 6 & 2 & 0 & 0 & 0 \\
5) &  & 15 & 0 & 0 & 0 & 0 & 0 & 0 & 1 & 0 & 2 & 6 & 0 & 0 \\
\end{longtable}
\noindent Modules tested: V248, V3875, V147250.\\

\setcounter{table}{291}
\begin{longtable}{rc|ccccccc}
\caption{$L_2(25) < E_8$, $p = 13$}\\
& & \multicolumn{7}{c}{$V_{248}$} \\
& & $1$ & $24$ & $26_{a}$ & $26_{b}$ & $26_{c}$ & $26_{d}$ & $26_{e}$ \\ \hline
\endfirsthead
& & \multicolumn{7}{c}{$V_{248}$} \\
& & $1$ & $24$ & $26_{a}$ & $26_{b}$ & $26_{c}$ & $26_{d}$ & $26_{e}$ \\ \hline
\endhead
1) & \possprim & 0 & 6 & 0 & 0 & 0 & 2 & 2 \\
2) &  & 14 & 0 & 7 & 0 & 0 & 1 & 1 \\
3) &  & 16 & 1 & 6 & 2 & 0 & 0 & 0 \\
4) &  & 16 & 1 & 0 & 2 & 6 & 0 & 0 \\
\end{longtable}
\noindent Modules tested: V248, V3875, V147250.\\

\setcounter{table}{292}
\begin{longtable}{rc|ccccccccccc}
\caption{$L_2(25) < E_8$, $p = 3$}\\
& & \multicolumn{11}{c}{$V_{248}$} \\
& & $1$ & $13_{a}$ & $13_{b}$ & $24_{a}$ & $24_{b}$ & $24_{c}$ & $24_{d}$ & $24_{e}$ & $24_{f}$ & $25$ & $26$ \\ \hline
\endfirsthead
& & \multicolumn{11}{c}{$V_{248}$} \\
& & $1$ & $13_{a}$ & $13_{b}$ & $24_{a}$ & $24_{b}$ & $24_{c}$ & $24_{d}$ & $24_{e}$ & $24_{f}$ & $25$ & $26$ \\ \hline
\endhead
1) & \possprim & 0 & 0 & 0 & 2 & 1 & 1 & 0 & 1 & 1 & 0 & 4 \\
2) &  & 15 & 6 & 6 & 0 & 0 & 0 & 0 & 0 & 0 & 1 & 2 \\
3) & \nongcr & 21 & 0 & 0 & 0 & 0 & 0 & 0 & 0 & 0 & 7 & 2 \\
4) & \newrow & 0 & 0 & 0 & 1 & 1 & 1 & 1 & 1 & 1 & 0 & 4 \\
\end{longtable}
\noindent Rows marked \newrow\ do not appear in the Memoir.\\
Modules tested: V248, V3875, V147250, L113243.\\

\setcounter{table}{293}
\begin{longtable}{rc|cccccccccc}
\caption{$L_2(25) < E_8$, $p = 2$}\\
& & \multicolumn{10}{c}{$V_{248}$} \\
& & $1$ & $12_{a}$ & $12_{b}$ & $24_{a}$ & $24_{b}$ & $24_{c}$ & $24_{d}$ & $24_{e}$ & $24_{f}$ & $26$ \\ \hline
\endfirsthead
& & \multicolumn{10}{c}{$V_{248}$} \\
& & $1$ & $12_{a}$ & $12_{b}$ & $24_{a}$ & $24_{b}$ & $24_{c}$ & $24_{d}$ & $24_{e}$ & $24_{f}$ & $26$ \\ \hline
\endhead
1) & \possprim & 0 & 0 & 0 & 2 & 1 & 1 & 0 & 1 & 1 & 4 \\
2) & \possprim & 0 & 0 & 0 & 1 & 1 & 1 & 1 & 1 & 1 & 4 \\
3) & \nongcr & 6 & 2 & 2 & 2 & 2 & 2 & 0 & 1 & 0 & 1 \\
4) & \nongcr & 6 & 3 & 3 & 1 & 1 & 1 & 1 & 1 & 1 & 1 \\
5) & \nongcr & 6 & 3 & 3 & 2 & 1 & 1 & 0 & 1 & 1 & 1 \\
6) & \nongcr & 10 & 7 & 2 & 0 & 0 & 0 & 0 & 0 & 0 & 5 \\
7) &  & 14 & 0 & 0 & 0 & 0 & 0 & 0 & 0 & 0 & 9 \\
8) & \nongcr & 16 & 10 & 5 & 0 & 0 & 0 & 0 & 0 & 0 & 2 \\
9) & \nongcr & 20 & 3 & 3 & 0 & 0 & 0 & 0 & 0 & 0 & 6 \\
\end{longtable}
\noindent Modules tested: V248, V3875, L3626, V147250, L143376.\\

\setcounter{table}{294}
\begin{longtable}{rc|cccc}
\caption{$L_2(27) < E_8$, $p = 0$, $13$ or $p \nmid |H|$}\\
& & \multicolumn{4}{c}{$V_{248}$} \\
& & $1$ & $26_{a}$ & $26_{b}$ & $26_{e}$ \\ \hline
\endfirsthead
& & \multicolumn{4}{c}{$V_{248}$} \\
& & $1$ & $26_{a}$ & $26_{b}$ & $26_{e}$ \\ \hline
\endhead
1) &  & 14 & 1 & 1 & 7 \\
\end{longtable}
\noindent Modules tested: V248, V3875, V147250.\\

\setcounter{table}{295}
\begin{longtable}{rc|cccc}
\caption{$L_2(27) < E_8$, $p = 7$}\\
& & \multicolumn{4}{c}{$V_{248}$} \\
& & $1$ & $13_{a}$ & $13_{b}$ & $26$ \\ \hline
\endfirsthead
& & \multicolumn{4}{c}{$V_{248}$} \\
& & $1$ & $13_{a}$ & $13_{b}$ & $26$ \\ \hline
\endhead
1) & \nongcr & 14 & 7 & 7 & 2 \\
\end{longtable}
\noindent Modules tested: V248, V3875, V147250, L147002.\\

\setcounter{table}{296}
\begin{longtable}{rc|cccccccccccc}
\caption{$L_2(27) < E_8$, $p = 2$}\\
& & \multicolumn{12}{c}{$V_{248}$} \\
& & $1$ & $13_{a}$ & $13_{b}$ & $26_{a}$ & $26_{b}$ & $26_{c}$ & $28_{a}$ & $28_{b}$ & $28_{c}$ & $28_{d}$ & $28_{e}$ & $28_{f}$ \\ \hline
\endfirsthead
& & \multicolumn{12}{c}{$V_{248}$} \\
& & $1$ & $13_{a}$ & $13_{b}$ & $26_{a}$ & $26_{b}$ & $26_{c}$ & $28_{a}$ & $28_{b}$ & $28_{c}$ & $28_{d}$ & $28_{e}$ & $28_{f}$ \\ \hline
\endhead
1) &  & 2 & 0 & 0 & 1 & 1 & 1 & 2 & 1 & 1 & 1 & 0 & 1 \\
2) &  & 2 & 0 & 0 & 1 & 1 & 1 & 1 & 1 & 1 & 1 & 1 & 1 \\
3) & \nongcr & 2 & 1 & 1 & 1 & 1 & 0 & 2 & 1 & 1 & 1 & 0 & 1 \\
4) & \nongcr & 2 & 1 & 1 & 1 & 1 & 0 & 1 & 1 & 1 & 1 & 1 & 1 \\
5) & \nongcr & 4 & 1 & 1 & 1 & 1 & 1 & 4 & 0 & 0 & 0 & 0 & 1 \\
6) & \nongcr & 4 & 1 & 1 & 1 & 1 & 1 & 2 & 2 & 0 & 1 & 0 & 0 \\
7) & \nongcr & 4 & 2 & 2 & 1 & 1 & 0 & 4 & 0 & 0 & 0 & 0 & 1 \\
8) & \nongcr & 4 & 2 & 2 & 1 & 1 & 0 & 2 & 2 & 0 & 1 & 0 & 0 \\
9) &  & 14 & 0 & 0 & 8 & 0 & 1 & 0 & 0 & 0 & 0 & 0 & 0 \\
10) & \nongcr & 14 & 6 & 6 & 1 & 1 & 1 & 0 & 0 & 0 & 0 & 0 & 0 \\
11) & \nongcr & 14 & 7 & 7 & 1 & 1 & 0 & 0 & 0 & 0 & 0 & 0 & 0 \\
\end{longtable}
\noindent Modules tested: V248, V3875, L3626, V147250, L143376.\\

\setcounter{table}{297}
\begin{longtable}{rc|cccccccccccccc}
\caption{$L_2(29) < E_8$, $p = 0$ or $p \nmid |H|$}\\
& & \multicolumn{14}{c}{$V_{248}$} \\
& & $1$ & $15_{a}$ & $15_{b}$ & $28_{a}$ & $28_{b}$ & $28_{c}$ & $28_{d}$ & $28_{e}$ & $28_{f}$ & $28_{g}$ & $29$ & $30_{d}$ & $30_{e}$ & $30_{f}$ \\ \hline
\endfirsthead
& & \multicolumn{14}{c}{$V_{248}$} \\
& & $1$ & $15_{a}$ & $15_{b}$ & $28_{a}$ & $28_{b}$ & $28_{c}$ & $28_{d}$ & $28_{e}$ & $28_{f}$ & $28_{g}$ & $29$ & $30_{d}$ & $30_{e}$ & $30_{f}$ \\ \hline
\endhead
1) & \possprim & 0 & 1 & 0 & 1 & 1 & 1 & 0 & 0 & 0 & 0 & 1 & 2 & 1 & 1 \\
2) & \possprim & 0 & 2 & 1 & 1 & 1 & 1 & 0 & 0 & 0 & 0 & 1 & 1 & 1 & 1 \\
3) & \newrow & 1 & 2 & 1 & 0 & 0 & 0 & 1 & 1 & 1 & 1 & 0 & 1 & 1 & 1 \\
4) & \newrow & 1 & 1 & 0 & 0 & 0 & 0 & 1 & 1 & 1 & 1 & 0 & 2 & 1 & 1 \\
\end{longtable}
\noindent Memoir rows read with the published correction: one constituent 28 replaced by 29 in each row.\\
Rows marked \newrow\ do not appear in the Memoir.\\
Modules tested: V248, V3875, V147250.\\

\setcounter{table}{298}
\begin{longtable}{rc|ccccccccccc}
\caption{$L_2(29) < E_8$, $p = 7$}\\
& & \multicolumn{11}{c}{$V_{248}$} \\
& & $1$ & $15_{a}$ & $15_{b}$ & $28_{a}$ & $28_{b}$ & $28_{c}$ & $28_{d}$ & $28_{e}$ & $28_{f}$ & $28_{g}$ & $29$ \\ \hline
\endfirsthead
& & \multicolumn{11}{c}{$V_{248}$} \\
& & $1$ & $15_{a}$ & $15_{b}$ & $28_{a}$ & $28_{b}$ & $28_{c}$ & $28_{d}$ & $28_{e}$ & $28_{f}$ & $28_{g}$ & $29$ \\ \hline
\endhead
1) & \possprim & 0 & 5 & 4 & 1 & 1 & 1 & 0 & 0 & 0 & 0 & 1 \\
2) &  & 1 & 5 & 4 & 0 & 0 & 0 & 1 & 1 & 1 & 1 & 0 \\
\end{longtable}
\noindent Modules tested: V248, V3875, V147250, L147002.\\

\setcounter{table}{299}
\begin{longtable}{rc|cccccccc}
\caption{$L_2(29) < E_8$, $p = 5$}\\
& & \multicolumn{8}{c}{$V_{248}$} \\
& & $1$ & $15_{a}$ & $15_{b}$ & $28_{a}$ & $28_{b}$ & $30_{d}$ & $30_{e}$ & $30_{f}$ \\ \hline
\endfirsthead
& & \multicolumn{8}{c}{$V_{248}$} \\
& & $1$ & $15_{a}$ & $15_{b}$ & $28_{a}$ & $28_{b}$ & $30_{d}$ & $30_{e}$ & $30_{f}$ \\ \hline
\endhead
1) &  & 1 & 1 & 0 & 4 & 0 & 2 & 1 & 1 \\
2) & \possprim & 1 & 1 & 0 & 1 & 3 & 2 & 1 & 1 \\
3) &  & 1 & 2 & 1 & 4 & 0 & 1 & 1 & 1 \\
4) & \possprim & 1 & 2 & 1 & 1 & 3 & 1 & 1 & 1 \\
\end{longtable}
\noindent Memoir rows not reproduced: 3) (not feasible on the Memoir's modules (all element orders, full power maps)); 4) (not feasible on the Memoir's modules (all element orders, full power maps)).\\
Modules tested: V248, V3875, V147250.\\

\setcounter{table}{300}
\begin{longtable}{rc|ccccccccc}
\caption{$L_2(29) < E_8$, $p = 3$}\\
& & \multicolumn{9}{c}{$V_{248}$} \\
& & $1$ & $15_{a}$ & $15_{b}$ & $28_{a}$ & $28_{b}$ & $28_{c}$ & $30_{d}$ & $30_{e}$ & $30_{f}$ \\ \hline
\endfirsthead
& & \multicolumn{9}{c}{$V_{248}$} \\
& & $1$ & $15_{a}$ & $15_{b}$ & $28_{a}$ & $28_{b}$ & $28_{c}$ & $30_{d}$ & $30_{e}$ & $30_{f}$ \\ \hline
\endhead
1) &  & 1 & 1 & 0 & 0 & 2 & 2 & 2 & 1 & 1 \\
2) & \possprim, \nongcr & 1 & 1 & 0 & 2 & 1 & 1 & 2 & 1 & 1 \\
3) &  & 1 & 2 & 1 & 0 & 2 & 2 & 1 & 1 & 1 \\
4) & \possprim, \nongcr & 1 & 2 & 1 & 2 & 1 & 1 & 1 & 1 & 1 \\
\end{longtable}
\noindent Modules tested: V248, V3875, V147250, L113243.\\

\setcounter{table}{301}
\begin{longtable}{rc|ccccccccccccc}
\caption{$L_2(29) < E_8$, $p = 2$}\\
& & \multicolumn{13}{c}{$V_{248}$} \\
& & $1$ & $14_{a}$ & $14_{b}$ & $28_{a}$ & $28_{b}$ & $28_{c}$ & $28_{d}$ & $28_{e}$ & $28_{f}$ & $28_{g}$ & $30_{a}$ & $30_{b}$ & $30_{c}$ \\ \hline
\endfirsthead
& & \multicolumn{13}{c}{$V_{248}$} \\
& & $1$ & $14_{a}$ & $14_{b}$ & $28_{a}$ & $28_{b}$ & $28_{c}$ & $28_{d}$ & $28_{e}$ & $28_{f}$ & $28_{g}$ & $30_{a}$ & $30_{b}$ & $30_{c}$ \\ \hline
\endhead
1) &  & 2 & 1 & 0 & 0 & 0 & 0 & 1 & 1 & 1 & 1 & 2 & 1 & 1 \\
2) & \possprim, \nongcr & 2 & 2 & 1 & 1 & 1 & 1 & 0 & 0 & 0 & 0 & 2 & 1 & 1 \\
3) &  & 4 & 1 & 0 & 0 & 1 & 0 & 2 & 0 & 2 & 0 & 1 & 1 & 1 \\
4) &  & 4 & 1 & 0 & 1 & 1 & 0 & 1 & 2 & 0 & 0 & 1 & 1 & 1 \\
5) & \nongcr & 4 & 2 & 1 & 0 & 0 & 0 & 1 & 1 & 1 & 1 & 1 & 1 & 1 \\
6) & \possprim & 4 & 3 & 2 & 1 & 1 & 1 & 0 & 0 & 0 & 0 & 1 & 1 & 1 \\
\end{longtable}
\noindent Modules tested: V248, V3875, L3626, V147250, L143376.\\

\setcounter{table}{302}
\begin{longtable}{rc|cccccccccccccc}
\caption{$L_2(31) < E_8$, $p = 0$ or $p \nmid |H|$}\\
& & \multicolumn{14}{c}{$V_{248}$} \\
& & $1$ & $15_{a}$ & $15_{b}$ & $30_{a}$ & $30_{b}$ & $30_{c}$ & $31$ & $32_{a}$ & $32_{b}$ & $32_{c}$ & $32_{d}$ & $32_{e}$ & $32_{f}$ & $32_{g}$ \\ \hline
\endfirsthead
& & \multicolumn{14}{c}{$V_{248}$} \\
& & $1$ & $15_{a}$ & $15_{b}$ & $30_{a}$ & $30_{b}$ & $30_{c}$ & $31$ & $32_{a}$ & $32_{b}$ & $32_{c}$ & $32_{d}$ & $32_{e}$ & $32_{f}$ & $32_{g}$ \\ \hline
\endhead
1) & \possprim & 0 & 0 & 0 & 0 & 2 & 2 & 0 & 0 & 0 & 0 & 1 & 1 & 1 & 1 \\
2) & \possprim & 0 & 0 & 0 & 2 & 1 & 1 & 0 & 0 & 0 & 0 & 1 & 1 & 1 & 1 \\
3) & \possprim & 0 & 1 & 1 & 1 & 2 & 0 & 0 & 0 & 0 & 0 & 1 & 1 & 1 & 1 \\
4) & \possprim & 0 & 1 & 1 & 1 & 1 & 1 & 0 & 0 & 0 & 0 & 1 & 1 & 1 & 1 \\
5) &  & 1 & 0 & 0 & 0 & 2 & 2 & 1 & 1 & 1 & 1 & 0 & 0 & 0 & 0 \\
6) &  & 1 & 0 & 0 & 2 & 1 & 1 & 1 & 1 & 1 & 1 & 0 & 0 & 0 & 0 \\
7) &  & 1 & 1 & 1 & 1 & 2 & 0 & 1 & 1 & 1 & 1 & 0 & 0 & 0 & 0 \\
8) &  & 1 & 1 & 1 & 1 & 1 & 1 & 1 & 1 & 1 & 1 & 0 & 0 & 0 & 0 \\
\end{longtable}
\noindent Modules tested: V248, V3875, V147250.\\

\setcounter{table}{303}
\begin{longtable}{rc|cccccccc}
\caption{$L_2(31) < E_8$, $p = 5$}\\
& & \multicolumn{8}{c}{$V_{248}$} \\
& & $1$ & $15_{a}$ & $15_{b}$ & $30_{a}$ & $30_{b}$ & $30_{c}$ & $31$ & $32$ \\ \hline
\endfirsthead
& & \multicolumn{8}{c}{$V_{248}$} \\
& & $1$ & $15_{a}$ & $15_{b}$ & $30_{a}$ & $30_{b}$ & $30_{c}$ & $31$ & $32$ \\ \hline
\endhead
1) & \possprim & 0 & 0 & 0 & 0 & 2 & 2 & 0 & 4 \\
2) & \possprim & 0 & 0 & 0 & 2 & 1 & 1 & 0 & 4 \\
3) & \possprim & 0 & 1 & 1 & 1 & 2 & 0 & 0 & 4 \\
4) & \possprim & 0 & 1 & 1 & 1 & 1 & 1 & 0 & 4 \\
5) &  & 3 & 0 & 0 & 0 & 2 & 2 & 3 & 1 \\
6) &  & 3 & 0 & 0 & 2 & 1 & 1 & 3 & 1 \\
7) &  & 3 & 1 & 1 & 1 & 2 & 0 & 3 & 1 \\
8) &  & 3 & 1 & 1 & 1 & 1 & 1 & 3 & 1 \\
9) &  & 5 & 2 & 2 & 2 & 1 & 1 & 1 & 1 \\
10) &  & 5 & 3 & 3 & 1 & 1 & 1 & 1 & 1 \\
\end{longtable}
\noindent Memoir rows not reproduced: 13) (not compatible with V3875).\\
Modules tested: V248, V3875, V147250.\\

\setcounter{table}{304}
\begin{longtable}{rc|ccccccccc}
\caption{$L_2(31) < E_8$, $p = 3$}\\
& & \multicolumn{9}{c}{$V_{248}$} \\
& & $1$ & $15_{a}$ & $15_{b}$ & $30_{a}$ & $30_{b}$ & $30_{c}$ & $31$ & $32_{a}$ & $32_{b}$ \\ \hline
\endfirsthead
& & \multicolumn{9}{c}{$V_{248}$} \\
& & $1$ & $15_{a}$ & $15_{b}$ & $30_{a}$ & $30_{b}$ & $30_{c}$ & $31$ & $32_{a}$ & $32_{b}$ \\ \hline
\endhead
1) & \possprim & 0 & 0 & 0 & 0 & 2 & 2 & 0 & 2 & 2 \\
2) & \possprim & 0 & 0 & 0 & 2 & 1 & 1 & 0 & 2 & 2 \\
3) & \possprim & 0 & 1 & 1 & 1 & 2 & 0 & 0 & 2 & 2 \\
4) & \possprim & 0 & 1 & 1 & 1 & 1 & 1 & 0 & 2 & 2 \\
5) &  & 2 & 0 & 0 & 0 & 2 & 2 & 2 & 1 & 1 \\
6) &  & 2 & 0 & 0 & 2 & 1 & 1 & 2 & 1 & 1 \\
7) &  & 2 & 1 & 1 & 1 & 2 & 0 & 2 & 1 & 1 \\
8) &  & 2 & 1 & 1 & 1 & 1 & 1 & 2 & 1 & 1 \\
9) &  & 4 & 2 & 2 & 2 & 1 & 1 & 0 & 1 & 1 \\
10) &  & 4 & 3 & 3 & 1 & 1 & 1 & 0 & 1 & 1 \\
\end{longtable}
\noindent Modules tested: V248, V3875, V147250, L113243.\\

\setcounter{table}{305}
\begin{longtable}{rc|cccccccccc}
\caption{$L_2(31) < E_8$, $p = 2$}\\
& & \multicolumn{10}{c}{$V_{248}$} \\
& & $1$ & $15_{a}$ & $15_{b}$ & $32_{a}$ & $32_{b}$ & $32_{c}$ & $32_{d}$ & $32_{e}$ & $32_{f}$ & $32_{g}$ \\ \hline
\endfirsthead
& & \multicolumn{10}{c}{$V_{248}$} \\
& & $1$ & $15_{a}$ & $15_{b}$ & $32_{a}$ & $32_{b}$ & $32_{c}$ & $32_{d}$ & $32_{e}$ & $32_{f}$ & $32_{g}$ \\ \hline
\endhead
1) & \possprim & 0 & 4 & 4 & 0 & 0 & 0 & 1 & 1 & 1 & 1 \\
2) & \possprim & 2 & 5 & 5 & 1 & 1 & 1 & 0 & 0 & 0 & 0 \\
\end{longtable}
\noindent Modules tested: V248, V3875, L3626, V147250, L143376.\\

\setcounter{table}{306}
\begin{longtable}{rc|cccccccc}
\caption{$L_2(32) < E_8$, $p = 0$, $31$ or $p \nmid |H|$}\\
& & \multicolumn{8}{c}{$V_{248}$} \\
& & $31_{b}$ & $31_{c}$ & $31_{d}$ & $31_{e}$ & $31_{g}$ & $31_{i}$ & $31_{l}$ & $31_{n}$ \\ \hline
\endfirsthead
& & \multicolumn{8}{c}{$V_{248}$} \\
& & $31_{b}$ & $31_{c}$ & $31_{d}$ & $31_{e}$ & $31_{g}$ & $31_{i}$ & $31_{l}$ & $31_{n}$ \\ \hline
\endhead
1) & \possprim & 1 & 1 & 1 & 1 & 1 & 1 & 1 & 1 \\
\end{longtable}
\noindent Modules tested: V248, V3875, V147250.\\

\setcounter{table}{307}
\begin{longtable}{rc|cc}
\caption{$L_2(32) < E_8$, $p = 11$}\\
& & \multicolumn{2}{c}{$V_{248}$} \\
& & $31_{a}$ & $31_{b}$ \\ \hline
\endfirsthead
& & \multicolumn{2}{c}{$V_{248}$} \\
& & $31_{a}$ & $31_{b}$ \\ \hline
\endhead
1) & \possprim & 7 & 1 \\
2) & \possprim & 4 & 4 \\
\end{longtable}
\noindent Modules tested: V248, V3875, V147250.\\

\setcounter{table}{308}
\begin{longtable}{rc|ccccc}
\caption{$L_2(32) < E_8$, $p = 3$}\\
& & \multicolumn{5}{c}{$V_{248}$} \\
& & $31_{b}$ & $31_{c}$ & $31_{d}$ & $31_{e}$ & $31_{f}$ \\ \hline
\endfirsthead
& & \multicolumn{5}{c}{$V_{248}$} \\
& & $31_{b}$ & $31_{c}$ & $31_{d}$ & $31_{e}$ & $31_{f}$ \\ \hline
\endhead
1) & \possprim & 3 & 1 & 3 & 1 & 0 \\
2) & \possprim & 2 & 2 & 2 & 1 & 1 \\
\end{longtable}
\noindent Modules tested: V248, V3875, V147250, L113243.\\

\setcounter{table}{309}
\begin{longtable}{l}
\caption{$L_2(37) < E_8$, $p = 2$}\\
Not recomputed: no Brauer table available in GAP. See the Memoir.
\end{longtable}

\setcounter{table}{310}
\begin{longtable}{rc|cccccc}
\caption{$L_2(41) < E_8$, $p = 0$, $3$ or $p \nmid |H|$}\\
& & \multicolumn{6}{c}{$V_{248}$} \\
& & $40_{c}$ & $40_{f}$ & $42_{a}$ & $42_{c}$ & $42_{g}$ & $42_{i}$ \\ \hline
\endfirsthead
& & \multicolumn{6}{c}{$V_{248}$} \\
& & $40_{c}$ & $40_{f}$ & $42_{a}$ & $42_{c}$ & $42_{g}$ & $42_{i}$ \\ \hline
\endhead
1) & \possprim & 1 & 1 & 1 & 1 & 1 & 1 \\
\end{longtable}
\noindent Modules not used (elements of large order): V3875, V147250.\\
Modules tested: V248.\\

\setcounter{table}{311}
\begin{longtable}{rc|ccccc}
\caption{$L_2(41) < E_8$, $p = 7$}\\
& & \multicolumn{5}{c}{$V_{248}$} \\
& & $40_{b}$ & $42_{a}$ & $42_{c}$ & $42_{g}$ & $42_{i}$ \\ \hline
\endfirsthead
& & \multicolumn{5}{c}{$V_{248}$} \\
& & $40_{b}$ & $42_{a}$ & $42_{c}$ & $42_{g}$ & $42_{i}$ \\ \hline
\endhead
1) & \possprim & 2 & 1 & 1 & 1 & 1 \\
\end{longtable}
\noindent Modules not used (elements of large order): V3875, V147250, L147002.\\
Modules tested: V248.\\

\setcounter{table}{312}
\begin{longtable}{rc|ccc}
\caption{$L_2(41) < E_8$, $p = 5$}\\
& & \multicolumn{3}{c}{$V_{248}$} \\
& & $40_{c}$ & $40_{f}$ & $42$ \\ \hline
\endfirsthead
& & \multicolumn{3}{c}{$V_{248}$} \\
& & $40_{c}$ & $40_{f}$ & $42$ \\ \hline
\endhead
1) & \possprim & 1 & 1 & 4 \\
\end{longtable}
\noindent Modules not used (elements of large order): V3875, V147250.\\
Modules tested: V248.\\

\setcounter{table}{313}
\begin{longtable}{l}
\caption{$L_2(41) < E_8$, $p = 2$}\\
Not recomputed: no Brauer table available in GAP. See the Memoir.
\end{longtable}

\setcounter{table}{314}
\begin{longtable}{rc|ccccc}
\caption{$L_2(49) < E_8$, $p = 0$ or $p \nmid |H|$}\\
& & \multicolumn{5}{c}{$V_{248}$} \\
& & $48_{a}$ & $50_{h}$ & $50_{i}$ & $50_{j}$ & $50_{k}$ \\ \hline
\endfirsthead
& & \multicolumn{5}{c}{$V_{248}$} \\
& & $48_{a}$ & $50_{h}$ & $50_{i}$ & $50_{j}$ & $50_{k}$ \\ \hline
\endhead
1) & \possprim & 1 & 1 & 1 & 1 & 1 \\
\end{longtable}
\noindent Modules tested: V248, V3875, V147250.\\

\setcounter{table}{315}
\begin{longtable}{rc|ccccc}
\caption{$L_2(49) < E_8$, $p = 5$}\\
& & \multicolumn{5}{c}{$V_{248}$} \\
& & $48$ & $50_{h}$ & $50_{i}$ & $50_{j}$ & $50_{k}$ \\ \hline
\endfirsthead
& & \multicolumn{5}{c}{$V_{248}$} \\
& & $48$ & $50_{h}$ & $50_{i}$ & $50_{j}$ & $50_{k}$ \\ \hline
\endhead
1) & \possprim & 1 & 1 & 1 & 1 & 1 \\
\end{longtable}
\noindent Modules tested: V248, V3875, V147250.\\

\setcounter{table}{316}
\begin{longtable}{rc|ccc}
\caption{$L_2(49) < E_8$, $p = 3$}\\
& & \multicolumn{3}{c}{$V_{248}$} \\
& & $48_{a}$ & $50_{b}$ & $50_{c}$ \\ \hline
\endfirsthead
& & \multicolumn{3}{c}{$V_{248}$} \\
& & $48_{a}$ & $50_{b}$ & $50_{c}$ \\ \hline
\endhead
1) & \possprim & 1 & 2 & 2 \\
\end{longtable}
\noindent Modules tested: V248, V3875, V147250, L113243.\\

\setcounter{table}{317}
\begin{longtable}{rc|cccccc}
\caption{$L_2(49) < E_8$, $p = 2$}\\
& & \multicolumn{6}{c}{$V_{248}$} \\
& & $1$ & $24_{a}$ & $24_{b}$ & $48_{a}$ & $48_{b}$ & $50$ \\ \hline
\endfirsthead
& & \multicolumn{6}{c}{$V_{248}$} \\
& & $1$ & $24_{a}$ & $24_{b}$ & $48_{a}$ & $48_{b}$ & $50$ \\ \hline
\endhead
1) & \possprim & 0 & 0 & 0 & 1 & 0 & 4 \\
2) &  & 6 & 1 & 1 & 2 & 1 & 1 \\
3) &  & 6 & 3 & 3 & 1 & 0 & 1 \\
\end{longtable}
\noindent Modules tested: V248, V3875, L3626, V147250, L143376.\\

\setcounter{table}{318}
\begin{longtable}{rc|cccc}
\caption{$L_2(61) < E_8$, $p = 0$, $31$ or $p \nmid |H|$}\\
& & \multicolumn{4}{c}{$V_{248}$} \\
& & $62_{a}$ & $62_{g}$ & $62_{k}$ & $62_{m}$ \\ \hline
\endfirsthead
& & \multicolumn{4}{c}{$V_{248}$} \\
& & $62_{a}$ & $62_{g}$ & $62_{k}$ & $62_{m}$ \\ \hline
\endhead
1) & \possprim & 1 & 1 & 1 & 1 \\
\end{longtable}
\noindent Modules not used (elements of large order): V3875, V147250.\\
Modules tested: V248.\\

\setcounter{table}{319}
\begin{longtable}{rc|c}
\caption{$L_2(61) < E_8$, $p = 5$}\\
& & \multicolumn{1}{c}{$V_{248}$} \\
& & $62_{a}$ \\ \hline
\endfirsthead
& & \multicolumn{1}{c}{$V_{248}$} \\
& & $62_{a}$ \\ \hline
\endhead
1) & \possprim & 4 \\
\end{longtable}
\noindent Modules not used (elements of large order): V3875, V147250.\\
Modules tested: V248.\\

\setcounter{table}{320}
\begin{longtable}{rc|cc}
\caption{$L_2(61) < E_8$, $p = 3$}\\
& & \multicolumn{2}{c}{$V_{248}$} \\
& & $62_{a}$ & $62_{c}$ \\ \hline
\endfirsthead
& & \multicolumn{2}{c}{$V_{248}$} \\
& & $62_{a}$ & $62_{c}$ \\ \hline
\endhead
1) & \possprim & 2 & 2 \\
\end{longtable}
\noindent Modules not used (elements of large order): V3875, V147250, L113243.\\
Modules tested: V248.\\

\setcounter{table}{321}
\begin{longtable}{l}
\caption{$L_2(61) < E_8$, $p = 2$}\\
Not recomputed: no Brauer table available in GAP. See the Memoir.
\end{longtable}

\setcounter{table}{322}
\begin{longtable}{rc|cccccccccccc}
\caption{$L_3(3) < E_8$, $p = 0$ or $p \nmid |H|$}\\
& & \multicolumn{12}{c}{$V_{248}$} \\
& & $1$ & $12$ & $13$ & $16_{a}$ & $16_{b}$ & $16_{c}$ & $16_{d}$ & $26_{a}$ & $26_{b}$ & $26_{c}$ & $27$ & $39$ \\ \hline
\endfirsthead
& & \multicolumn{12}{c}{$V_{248}$} \\
& & $1$ & $12$ & $13$ & $16_{a}$ & $16_{b}$ & $16_{c}$ & $16_{d}$ & $26_{a}$ & $26_{b}$ & $26_{c}$ & $27$ & $39$ \\ \hline
\endhead
1) & \possprim & 0 & 0 & 1 & 0 & 0 & 0 & 0 & 1 & 2 & 2 & 1 & 2 \\
2) & \possprim & 0 & 2 & 0 & 0 & 0 & 0 & 0 & 4 & 0 & 0 & 3 & 1 \\
3) &  & 1 & 0 & 0 & 0 & 0 & 0 & 0 & 1 & 2 & 2 & 0 & 3 \\
4) &  & 1 & 0 & 2 & 0 & 0 & 0 & 0 & 3 & 2 & 2 & 0 & 1 \\
5) &  & 2 & 0 & 1 & 1 & 1 & 1 & 1 & 0 & 1 & 1 & 0 & 3 \\
6) &  & 2 & 0 & 3 & 1 & 1 & 1 & 1 & 2 & 1 & 1 & 0 & 1 \\
7) &  & 3 & 0 & 4 & 2 & 2 & 2 & 2 & 1 & 0 & 0 & 0 & 1 \\
8) &  & 4 & 4 & 0 & 1 & 1 & 1 & 1 & 1 & 1 & 1 & 2 & 0 \\
9) &  & 5 & 4 & 1 & 2 & 2 & 2 & 2 & 0 & 0 & 0 & 2 & 0 \\
10) &  & 6 & 4 & 0 & 2 & 2 & 2 & 2 & 0 & 0 & 0 & 1 & 1 \\
11) &  & 8 & 0 & 0 & 0 & 0 & 0 & 0 & 1 & 1 & 1 & 6 & 0 \\
12) &  & 9 & 0 & 1 & 1 & 1 & 1 & 1 & 0 & 0 & 0 & 6 & 0 \\
13) &  & 10 & 0 & 0 & 1 & 1 & 1 & 1 & 0 & 0 & 0 & 5 & 1 \\
14) &  & 14 & 0 & 0 & 0 & 0 & 0 & 0 & 7 & 1 & 1 & 0 & 0 \\
15) &  & 15 & 0 & 1 & 1 & 1 & 1 & 1 & 6 & 0 & 0 & 0 & 0 \\
\end{longtable}
\noindent Memoir rows not reproduced: 5) (not compatible with V3875).\\
Modules tested: V248, V3875, V147250.\\

\setcounter{table}{323}
\begin{longtable}{rc|cccccccc}
\caption{$L_3(3) < E_8$, $p = 13$}\\
& & \multicolumn{8}{c}{$V_{248}$} \\
& & $1$ & $11$ & $13$ & $16$ & $26_{a}$ & $26_{b}$ & $26_{c}$ & $39$ \\ \hline
\endfirsthead
& & \multicolumn{8}{c}{$V_{248}$} \\
& & $1$ & $11$ & $13$ & $16$ & $26_{a}$ & $26_{b}$ & $26_{c}$ & $39$ \\ \hline
\endhead
1) & \possprim & 0 & 1 & 1 & 1 & 1 & 2 & 2 & 2 \\
2) &  & 1 & 0 & 0 & 0 & 1 & 2 & 2 & 3 \\
3) &  & 1 & 0 & 2 & 0 & 3 & 2 & 2 & 1 \\
4) &  & 1 & 1 & 2 & 5 & 0 & 1 & 1 & 2 \\
5) &  & 2 & 0 & 1 & 4 & 0 & 1 & 1 & 3 \\
6) &  & 2 & 0 & 3 & 4 & 2 & 1 & 1 & 1 \\
7) & \possprim, \nongcr & 2 & 5 & 0 & 3 & 4 & 0 & 0 & 1 \\
8) &  & 3 & 0 & 4 & 8 & 1 & 0 & 0 & 1 \\
9) & \nongcr & 8 & 6 & 0 & 6 & 1 & 1 & 1 & 0 \\
10) & \nongcr & 9 & 6 & 1 & 10 & 0 & 0 & 0 & 0 \\
11) & \nongcr & 10 & 5 & 0 & 9 & 0 & 0 & 0 & 1 \\
12) &  & 14 & 0 & 0 & 0 & 7 & 1 & 1 & 0 \\
13) &  & 15 & 0 & 1 & 4 & 6 & 0 & 0 & 0 \\
\end{longtable}
\noindent Modules tested: V248, V3875, V147250.\\

\setcounter{table}{324}
\begin{longtable}{rc|ccccccc}
\caption{$L_3(3) < E_8$, $p = 2$}\\
& & \multicolumn{7}{c}{$V_{248}$} \\
& & $1$ & $12$ & $16_{a}$ & $16_{b}$ & $16_{c}$ & $16_{d}$ & $26$ \\ \hline
\endfirsthead
& & \multicolumn{7}{c}{$V_{248}$} \\
& & $1$ & $12$ & $16_{a}$ & $16_{b}$ & $16_{c}$ & $16_{d}$ & $26$ \\ \hline
\endhead
1) & \possprim, \nongcr & 4 & 3 & 0 & 0 & 0 & 0 & 8 \\
2) & \possprim, \nongcr & 6 & 4 & 1 & 1 & 1 & 1 & 5 \\
3) & \nongcr & 8 & 5 & 2 & 2 & 2 & 2 & 2 \\
4) & \nongcr & 14 & 0 & 0 & 0 & 0 & 0 & 9 \\
5) & \nongcr & 16 & 1 & 1 & 1 & 1 & 1 & 6 \\
\end{longtable}
\noindent Modules tested: V248, V3875, L3626, V147250, L143376.\\

\setcounter{table}{325}
\begin{longtable}{rc|cc}
\caption{$L_3(5) < E_8$, $p \neq 2$, $5$}\\
& & \multicolumn{2}{c}{$V_{248}$} \\
& & $124_{e}$ & $124_{f}$ \\ \hline
\endfirsthead
& & \multicolumn{2}{c}{$V_{248}$} \\
& & $124_{e}$ & $124_{f}$ \\ \hline
\endhead
1) & \possprim & 1 & 1 \\
\end{longtable}
\noindent Modules tested: V248, V3875, V147250.\\

\setcounter{table}{326}
\begin{longtable}{rc|c}
\caption{$L_3(5) < E_8$, $p = 2$}\\
& & \multicolumn{1}{c}{$V_{248}$} \\
& & $124_{b}$ \\ \hline
\endfirsthead
& & \multicolumn{1}{c}{$V_{248}$} \\
& & $124_{b}$ \\ \hline
\endhead
1) & \possprim & 2 \\
\end{longtable}
\noindent Memoir rows not reproduced: 2) (not compatible with V3875; nor with L3626; nor with V147250; nor with L143376).\\
Modules tested: V248, V3875, L3626, V147250, L143376.\\

\setcounter{table}{327}
\begin{longtable}{rc|ccc}
\caption{$L_4(3) < E_8$, $p = 2$}\\
& & \multicolumn{3}{c}{$V_{248}$} \\
& & $1$ & $26_{a}$ & $26_{b}$ \\ \hline
\endfirsthead
& & \multicolumn{3}{c}{$V_{248}$} \\
& & $1$ & $26_{a}$ & $26_{b}$ \\ \hline
\endhead
1) &  & 14 & 8 & 1 \\
\end{longtable}
\noindent Memoir rows not reproduced: 1) (not compatible with L3626); 2) (not compatible with L3626).\\
Modules tested: V248, V3875, L3626, V147250, L143376.\\

\setcounter{table}{328}
\begin{longtable}{l}
\caption{$L_4(5) < E_8$, $p = 2$ (irreducible on $V_{248}$)}\\
Not recomputed: no Brauer table available in GAP. See the Memoir.
\end{longtable}

\setcounter{table}{329}
\begin{longtable}{rc|cccccccccccccc}
\caption{$U_3(3) < E_8$, $p = 0$ or $p \nmid |H|$}\\
& & \multicolumn{14}{c}{$V_{248}$} \\
& & $1$ & $6$ & $7_{a}$ & $7_{b}$ & $7_{c}$ & $14$ & $21_{a}$ & $21_{b}$ & $21_{c}$ & $27$ & $28_{a}$ & $28_{b}$ & $32_{a}$ & $32_{b}$ \\ \hline
\endfirsthead
& & \multicolumn{14}{c}{$V_{248}$} \\
& & $1$ & $6$ & $7_{a}$ & $7_{b}$ & $7_{c}$ & $14$ & $21_{a}$ & $21_{b}$ & $21_{c}$ & $27$ & $28_{a}$ & $28_{b}$ & $32_{a}$ & $32_{b}$ \\ \hline
\endhead
1) &  & 1 & 0 & 0 & 0 & 0 & 3 & 1 & 0 & 0 & 0 & 1 & 1 & 2 & 2 \\
2) &  & 2 & 0 & 1 & 0 & 0 & 0 & 1 & 1 & 1 & 0 & 2 & 2 & 1 & 1 \\
3) &  & 2 & 0 & 1 & 1 & 1 & 2 & 1 & 0 & 0 & 0 & 2 & 2 & 1 & 1 \\
4) &  & 3 & 0 & 0 & 0 & 0 & 0 & 0 & 1 & 1 & 1 & 2 & 2 & 1 & 1 \\
5) &  & 3 & 0 & 1 & 2 & 2 & 1 & 1 & 2 & 2 & 1 & 0 & 0 & 1 & 1 \\
6) &  & 3 & 0 & 2 & 2 & 2 & 1 & 1 & 0 & 0 & 0 & 3 & 3 & 0 & 0 \\
7) &  & 4 & 0 & 4 & 0 & 0 & 1 & 4 & 0 & 0 & 2 & 0 & 0 & 1 & 1 \\
8) &  & 4 & 0 & 2 & 3 & 3 & 0 & 1 & 2 & 2 & 1 & 1 & 1 & 0 & 0 \\
9) &  & 5 & 0 & 5 & 1 & 1 & 0 & 4 & 0 & 0 & 2 & 1 & 1 & 0 & 0 \\
10) &  & 6 & 0 & 13 & 0 & 0 & 5 & 0 & 0 & 0 & 3 & 0 & 0 & 0 & 0 \\
11) &  & 8 & 0 & 0 & 0 & 0 & 1 & 0 & 0 & 0 & 6 & 0 & 0 & 1 & 1 \\
12) &  & 9 & 0 & 1 & 1 & 1 & 0 & 0 & 0 & 0 & 6 & 1 & 1 & 0 & 0 \\
13) &  & 17 & 14 & 0 & 2 & 2 & 7 & 1 & 0 & 0 & 0 & 0 & 0 & 0 & 0 \\
14) &  & 52 & 0 & 26 & 0 & 0 & 1 & 0 & 0 & 0 & 0 & 0 & 0 & 0 & 0 \\
\end{longtable}
\noindent Modules tested: V248, V3875, V147250.\\

\setcounter{table}{330}
\begin{longtable}{rc|cccccccccccc}
\caption{$U_3(3) < E_8$, $p = 7$}\\
& & \multicolumn{12}{c}{$V_{248}$} \\
& & $1$ & $6$ & $7_{a}$ & $7_{b}$ & $7_{c}$ & $14$ & $21_{a}$ & $21_{b}$ & $21_{c}$ & $26$ & $28_{a}$ & $28_{b}$ \\ \hline
\endfirsthead
& & \multicolumn{12}{c}{$V_{248}$} \\
& & $1$ & $6$ & $7_{a}$ & $7_{b}$ & $7_{c}$ & $14$ & $21_{a}$ & $21_{b}$ & $21_{c}$ & $26$ & $28_{a}$ & $28_{b}$ \\ \hline
\endhead
1) & \possprim, \nongcr & 1 & 2 & 0 & 0 & 0 & 0 & 3 & 1 & 1 & 5 & 0 & 0 \\
2) & \possprim, \nongcr & 1 & 4 & 0 & 0 & 0 & 3 & 1 & 0 & 0 & 4 & 1 & 1 \\
3) & \nongcr & 2 & 2 & 1 & 0 & 0 & 0 & 1 & 1 & 1 & 2 & 2 & 2 \\
4) & \nongcr & 2 & 2 & 1 & 1 & 1 & 2 & 1 & 0 & 0 & 2 & 2 & 2 \\
5) &  & 3 & 0 & 2 & 2 & 2 & 1 & 1 & 0 & 0 & 0 & 3 & 3 \\
6) & \nongcr & 4 & 2 & 0 & 0 & 0 & 0 & 0 & 1 & 1 & 3 & 2 & 2 \\
7) & \nongcr & 4 & 2 & 1 & 2 & 2 & 1 & 1 & 2 & 2 & 3 & 0 & 0 \\
8) &  & 5 & 0 & 2 & 3 & 3 & 0 & 1 & 2 & 2 & 1 & 1 & 1 \\
9) & \nongcr & 6 & 2 & 4 & 0 & 0 & 1 & 4 & 0 & 0 & 4 & 0 & 0 \\
10) & \nongcr & 7 & 0 & 5 & 1 & 1 & 0 & 4 & 0 & 0 & 2 & 1 & 1 \\
11) & \nongcr & 9 & 0 & 13 & 0 & 0 & 5 & 0 & 0 & 0 & 3 & 0 & 0 \\
12) & \nongcr & 14 & 2 & 0 & 0 & 0 & 1 & 0 & 0 & 0 & 8 & 0 & 0 \\
13) & \nongcr & 15 & 0 & 1 & 1 & 1 & 0 & 0 & 0 & 0 & 6 & 1 & 1 \\
14) &  & 17 & 14 & 0 & 2 & 2 & 7 & 1 & 0 & 0 & 0 & 0 & 0 \\
15) &  & 52 & 0 & 26 & 0 & 0 & 1 & 0 & 0 & 0 & 0 & 0 & 0 \\
\end{longtable}
\noindent Modules tested: V248, V3875, V147250, L147002.\\

\setcounter{table}{331}
\begin{longtable}{rc|ccccc}
\caption{$U_3(8) < E_8$, $p \neq 2, 3$}\\
& & \multicolumn{5}{c}{$V_{248}$} \\
& & $1$ & $56$ & $57_{a}$ & $57_{b}$ & $133_{a}$ \\ \hline
\endfirsthead
& & \multicolumn{5}{c}{$V_{248}$} \\
& & $1$ & $56$ & $57_{a}$ & $57_{b}$ & $133_{a}$ \\ \hline
\endhead
1) &  & 1 & 0 & 1 & 1 & 1 \\
2) &  & 3 & 2 & 0 & 0 & 1 \\
\end{longtable}
\noindent Modules tested: V248, V3875, V147250.\\

\setcounter{table}{332}
\begin{longtable}{rc|ccc}
\caption{$U_3(8) < E_8$, $p = 3$}\\
& & \multicolumn{3}{c}{$V_{248}$} \\
& & $1$ & $56$ & $133_{a}$ \\ \hline
\endfirsthead
& & \multicolumn{3}{c}{$V_{248}$} \\
& & $1$ & $56$ & $133_{a}$ \\ \hline
\endhead
1) & \nongcr & 3 & 2 & 1 \\
\end{longtable}
\noindent Modules tested: V248, V3875, V147250, L113243.\\

\setcounter{table}{333}
\begin{longtable}{rc|cccccccccccccc}
\caption{$U_4(2) < E_8$, $p = 0$ or $p \nmid |H|$}\\
& & \multicolumn{14}{c}{$V_{248}$} \\
& & $1$ & $5_{a}$ & $5_{b}$ & $6$ & $10_{a}$ & $10_{b}$ & $15_{a}$ & $20$ & $24$ & $40_{a}$ & $40_{b}$ & $45_{a}$ & $45_{b}$ & $64$ \\ \hline
\endfirsthead
& & \multicolumn{14}{c}{$V_{248}$} \\
& & $1$ & $5_{a}$ & $5_{b}$ & $6$ & $10_{a}$ & $10_{b}$ & $15_{a}$ & $20$ & $24$ & $40_{a}$ & $40_{b}$ & $45_{a}$ & $45_{b}$ & $64$ \\ \hline
\endhead
1) & \possprim & 0 & 0 & 0 & 0 & 0 & 0 & 2 & 0 & 0 & 0 & 0 & 1 & 1 & 2 \\
2) & \possprim & 0 & 1 & 1 & 0 & 1 & 1 & 0 & 0 & 2 & 1 & 1 & 1 & 1 & 0 \\
3) &  & 11 & 0 & 0 & 12 & 2 & 2 & 7 & 1 & 0 & 0 & 0 & 0 & 0 & 0 \\
4) &  & 24 & 10 & 10 & 0 & 5 & 5 & 0 & 0 & 1 & 0 & 0 & 0 & 0 & 0 \\
\end{longtable}
\noindent Modules tested: V248, V3875, V147250.\\

\setcounter{table}{334}
\begin{longtable}{rc|cccccccccccccc}
\caption{$U_4(2) < E_8$, $p = 5$}\\
& & \multicolumn{14}{c}{$V_{248}$} \\
& & $1$ & $5_{a}$ & $5_{b}$ & $6$ & $10_{a}$ & $10_{b}$ & $15_{a}$ & $20$ & $23$ & $40_{a}$ & $40_{b}$ & $45_{a}$ & $45_{b}$ & $58$ \\ \hline
\endfirsthead
& & \multicolumn{14}{c}{$V_{248}$} \\
& & $1$ & $5_{a}$ & $5_{b}$ & $6$ & $10_{a}$ & $10_{b}$ & $15_{a}$ & $20$ & $23$ & $40_{a}$ & $40_{b}$ & $45_{a}$ & $45_{b}$ & $58$ \\ \hline
\endhead
1) & \possprim & 0 & 0 & 0 & 2 & 0 & 0 & 2 & 0 & 0 & 0 & 0 & 1 & 1 & 2 \\
2) & \nongcr & 2 & 1 & 1 & 0 & 1 & 1 & 0 & 0 & 2 & 1 & 1 & 1 & 1 & 0 \\
3) &  & 11 & 0 & 0 & 12 & 2 & 2 & 7 & 1 & 0 & 0 & 0 & 0 & 0 & 0 \\
4) &  & 25 & 10 & 10 & 0 & 5 & 5 & 0 & 0 & 1 & 0 & 0 & 0 & 0 & 0 \\
\end{longtable}
\noindent Modules tested: V248, V3875, V147250.\\

\setcounter{table}{335}
\begin{longtable}{rc|cccccc}
\caption{$PSp_4(5) < E_8$, $p = 2$}\\
& & \multicolumn{6}{c}{$V_{248}$} \\
& & $1$ & $12_{a}$ & $12_{b}$ & $40$ & $64$ & $104_{b}$ \\ \hline
\endfirsthead
& & \multicolumn{6}{c}{$V_{248}$} \\
& & $1$ & $12_{a}$ & $12_{b}$ & $40$ & $64$ & $104_{b}$ \\ \hline
\endhead
1) & \possprim & 0 & 0 & 0 & 1 & 0 & 2 \\
2) & \possprim, \nongcr & 8 & 4 & 4 & 2 & 1 & 0 \\
\end{longtable}
\noindent Modules tested: V248, V3875, L3626, V147250, L143376.\\

\setcounter{table}{336}
\begin{longtable}{rc|ccccc}
\caption{$Sp_6(2) < E_8$, $p = 0$, $5$ or $p \nmid |H|$}\\
& & \multicolumn{5}{c}{$V_{248}$} \\
& & $1$ & $7$ & $21_{a}$ & $27$ & $35_{a}$ \\ \hline
\endfirsthead
& & \multicolumn{5}{c}{$V_{248}$} \\
& & $1$ & $7$ & $21_{a}$ & $27$ & $35_{a}$ \\ \hline
\endhead
1) &  & 4 & 6 & 5 & 1 & 2 \\
\end{longtable}
\noindent Modules tested: V248, V3875, V147250.\\

\setcounter{table}{337}
\begin{longtable}{rc|ccccc}
\caption{$Sp_6(2) < E_8$, $p = 7$}\\
& & \multicolumn{5}{c}{$V_{248}$} \\
& & $1$ & $7$ & $21_{a}$ & $26$ & $35_{a}$ \\ \hline
\endfirsthead
& & \multicolumn{5}{c}{$V_{248}$} \\
& & $1$ & $7$ & $21_{a}$ & $26$ & $35_{a}$ \\ \hline
\endhead
1) &  & 5 & 6 & 5 & 1 & 2 \\
\end{longtable}
\noindent Modules tested: V248, V3875, V147250, L147002.\\

\setcounter{table}{338}
\begin{longtable}{rc|ccccc}
\caption{$Sp_6(2) < E_8$, $p = 3$}\\
& & \multicolumn{5}{c}{$V_{248}$} \\
& & $1$ & $7$ & $21$ & $27$ & $35$ \\ \hline
\endfirsthead
& & \multicolumn{5}{c}{$V_{248}$} \\
& & $1$ & $7$ & $21$ & $27$ & $35$ \\ \hline
\endhead
1) &  & 4 & 6 & 5 & 1 & 2 \\
\end{longtable}
\noindent Modules tested: V248, V3875, V147250, L113243.\\

\setcounter{table}{339}
\begin{longtable}{rc|ccccc}
\caption{$\Omega_8^{+}(2) < E_8$, $p \neq 2$}\\
& & \multicolumn{5}{c}{$V_{248}$} \\
& & $1$ & $28$ & $35_{a}$ & $35_{b}$ & $35_{c}$ \\ \hline
\endfirsthead
& & \multicolumn{5}{c}{$V_{248}$} \\
& & $1$ & $28$ & $35_{a}$ & $35_{b}$ & $35_{c}$ \\ \hline
\endhead
1) &  & 3 & 5 & 1 & 1 & 1 \\
\end{longtable}
\noindent Modules tested: V248, V3875, V147250.\\

\setcounter{table}{340}
\begin{longtable}{rc|cccccc}
\caption{$G_2(3) < E_8$, $p \neq 2, 3$}\\
& & \multicolumn{6}{c}{$V_{248}$} \\
& & $1$ & $14$ & $64_{a}$ & $64_{b}$ & $78$ & $91_{a}$ \\ \hline
\endfirsthead
& & \multicolumn{6}{c}{$V_{248}$} \\
& & $1$ & $14$ & $64_{a}$ & $64_{b}$ & $78$ & $91_{a}$ \\ \hline
\endhead
1) &  & 1 & 0 & 0 & 0 & 2 & 1 \\
2) &  & 1 & 2 & 1 & 1 & 0 & 1 \\
\end{longtable}
\noindent Modules tested: V248, V3875, V147250.\\

\setcounter{table}{341}
\begin{longtable}{rc|cccccc}
\caption{$G_2(3) < E_8$, $p = 2$}\\
& & \multicolumn{6}{c}{$V_{248}$} \\
& & $1$ & $14$ & $64_{a}$ & $64_{b}$ & $78$ & $90_{a}$ \\ \hline
\endfirsthead
& & \multicolumn{6}{c}{$V_{248}$} \\
& & $1$ & $14$ & $64_{a}$ & $64_{b}$ & $78$ & $90_{a}$ \\ \hline
\endhead
1) &  & 2 & 0 & 0 & 0 & 2 & 1 \\
2) &  & 2 & 2 & 1 & 1 & 0 & 1 \\
\end{longtable}
\noindent Modules tested: V248, V3875, L3626, V147250, L143376.\\

\setcounter{table}{342}
\begin{longtable}{rc|ccc}
\caption{$^{3}D_4(2) < E_8$, $p \neq 2, 3$}\\
& & \multicolumn{3}{c}{$V_{248}$} \\
& & $1$ & $26$ & $52$ \\ \hline
\endfirsthead
& & \multicolumn{3}{c}{$V_{248}$} \\
& & $1$ & $26$ & $52$ \\ \hline
\endhead
1) &  & 14 & 7 & 1 \\
\end{longtable}
\noindent Memoir rows not reproduced: 1) (not feasible on the Memoir's modules (all element orders, full power maps)).\\
Modules tested: V248, V3875, V147250.\\

\setcounter{table}{343}
\begin{longtable}{rc|cccc}
\caption{$^{3}D_4(2) < E_8$, $p = 3$}\\
& & \multicolumn{4}{c}{$V_{248}$} \\
& & $1$ & $25$ & $52$ & $196$ \\ \hline
\endfirsthead
& & \multicolumn{4}{c}{$V_{248}$} \\
& & $1$ & $25$ & $52$ & $196$ \\ \hline
\endhead
1) & \possprim & 0 & 0 & 1 & 1 \\
2) & \nongcr & 21 & 7 & 1 & 0 \\
\end{longtable}
\noindent Modules tested: V248, V3875, V147250, L113243.\\

\setcounter{table}{344}
\begin{longtable}{rc|cccc}
\caption{$^{2}F_4(2)' < E_8$, $p \neq 2, 3, 5$}\\
& & \multicolumn{4}{c}{$V_{248}$} \\
& & $1$ & $27_{a}$ & $27_{b}$ & $78$ \\ \hline
\endfirsthead
& & \multicolumn{4}{c}{$V_{248}$} \\
& & $1$ & $27_{a}$ & $27_{b}$ & $78$ \\ \hline
\endhead
1) &  & 8 & 3 & 3 & 1 \\
\end{longtable}
\noindent Modules tested: V248, V3875, V147250.\\

\setcounter{table}{345}
\begin{longtable}{rc|cccc}
\caption{$^{2}F_4(2)' < E_8$, $p = 5$}\\
& & \multicolumn{4}{c}{$V_{248}$} \\
& & $1$ & $27_{a}$ & $27_{b}$ & $78$ \\ \hline
\endfirsthead
& & \multicolumn{4}{c}{$V_{248}$} \\
& & $1$ & $27_{a}$ & $27_{b}$ & $78$ \\ \hline
\endhead
1) & \nongcr & 8 & 3 & 3 & 1 \\
\end{longtable}
\noindent Modules tested: V248, V3875, V147250.\\

\setcounter{table}{346}
\begin{longtable}{rc|cccccc}
\caption{$^{2}F_4(2)' < E_8$, $p = 3$}\\
& & \multicolumn{6}{c}{$V_{248}$} \\
& & $1$ & $27_{a}$ & $27_{b}$ & $77$ & $124_{a}$ & $124_{b}$ \\ \hline
\endfirsthead
& & \multicolumn{6}{c}{$V_{248}$} \\
& & $1$ & $27_{a}$ & $27_{b}$ & $77$ & $124_{a}$ & $124_{b}$ \\ \hline
\endhead
1) &  & 0 & 0 & 0 & 0 & 1 & 1 \\
2) &  & 9 & 3 & 3 & 1 & 0 & 0 \\
\end{longtable}
\noindent Modules tested: V248, V3875, V147250, L113243.\\

\setcounter{table}{347}
\begin{longtable}{rc|cccccc}
\caption{$^{2}B_2(8) < E_8$, $p = 0$ or $p \nmid |H|$}\\
& & \multicolumn{6}{c}{$V_{248}$} \\
& & $1$ & $14_{a}$ & $14_{b}$ & $64$ & $65_{a}$ & $91$ \\ \hline
\endfirsthead
& & \multicolumn{6}{c}{$V_{248}$} \\
& & $1$ & $14_{a}$ & $14_{b}$ & $64$ & $65_{a}$ & $91$ \\ \hline
\endhead
1) & \possprim & 0 & 1 & 1 & 1 & 1 & 1 \\
2) &  & 1 & 1 & 1 & 2 & 0 & 1 \\
\end{longtable}
\noindent Memoir rows not reproduced: 2) (not compatible with V3875; nor with V147250); 4) (not compatible with V3875; nor with V147250); 5) (not compatible with V3875; nor with V147250).\\
Modules tested: V248, V3875, V147250.\\

\setcounter{table}{348}
\begin{longtable}{rc|cccccc}
\caption{$^{2}B_2(8) < E_8$, $p = 13$}\\
& & \multicolumn{6}{c}{$V_{248}$} \\
& & $1$ & $14_{a}$ & $14_{b}$ & $35$ & $65_{a}$ & $91$ \\ \hline
\endfirsthead
& & \multicolumn{6}{c}{$V_{248}$} \\
& & $1$ & $14_{a}$ & $14_{b}$ & $35$ & $65_{a}$ & $91$ \\ \hline
\endhead
1) & \possprim & 1 & 2 & 2 & 1 & 1 & 1 \\
2) & \possprim & 3 & 3 & 3 & 2 & 0 & 1 \\
\end{longtable}
\noindent Memoir rows not reproduced: 1) (not compatible with V3875; nor with V147250); 3) (not compatible with V3875; nor with V147250); 4) (not compatible with V3875; nor with V147250).\\
Modules tested: V248, V3875, V147250.\\

\setcounter{table}{349}
\begin{longtable}{rc|ccccc}
\caption{$^{2}B_2(8) < E_8$, $p = 7$}\\
& & \multicolumn{5}{c}{$V_{248}$} \\
& & $1$ & $14_{a}$ & $14_{b}$ & $64$ & $91$ \\ \hline
\endfirsthead
& & \multicolumn{5}{c}{$V_{248}$} \\
& & $1$ & $14_{a}$ & $14_{b}$ & $64$ & $91$ \\ \hline
\endhead
1) & \possprim & 1 & 1 & 1 & 2 & 1 \\
2) &  & 8 & 4 & 4 & 2 & 0 \\
\end{longtable}
\noindent Memoir rows not reproduced: 2) (not compatible with V3875; nor with V147250; nor with L147002); 3) (not compatible with V3875; nor with V147250; nor with L147002); 5) (not compatible with V3875; nor with V147250; nor with L147002).\\
Modules tested: V248, V3875, V147250, L147002.\\

\setcounter{table}{350}
\begin{longtable}{rc|cccccccc}
\caption{$^{2}B_2(8) < E_8$, $p = 5$}\\
& & \multicolumn{8}{c}{$V_{248}$} \\
& & $1$ & $14_{a}$ & $14_{b}$ & $35_{a}$ & $35_{b}$ & $35_{c}$ & $63$ & $65_{a}$ \\ \hline
\endfirsthead
& & \multicolumn{8}{c}{$V_{248}$} \\
& & $1$ & $14_{a}$ & $14_{b}$ & $35_{a}$ & $35_{b}$ & $35_{c}$ & $63$ & $65_{a}$ \\ \hline
\endhead
1) & \possprim & 1 & 2 & 2 & 0 & 0 & 0 & 2 & 1 \\
2) &  & 3 & 0 & 0 & 4 & 0 & 3 & 0 & 0 \\
3) &  & 3 & 2 & 2 & 0 & 0 & 0 & 3 & 0 \\
4) &  & 3 & 5 & 5 & 1 & 1 & 1 & 0 & 0 \\
\end{longtable}
\noindent Modules tested: V248, V3875, V147250.\\

\setcounter{table}{351}
\begin{longtable}{rc|cc}
\caption{$^{2}B_2(32) < E_8$, $p = 5$}\\
& & \multicolumn{2}{c}{$V_{248}$} \\
& & $124_{a}$ & $124_{b}$ \\ \hline
\endfirsthead
& & \multicolumn{2}{c}{$V_{248}$} \\
& & $124_{a}$ & $124_{b}$ \\ \hline
\endhead
1) & \possprim & 1 & 1 \\
\end{longtable}
\noindent Modules not used (elements of large order): V3875, V147250.\\
Modules tested: V248.\\


\newpage
\end{document}
